\chapter{III. Sach-Register}

zu den beiden ersten Bänden der Sämmtlichen Schriften Luthers.

Abel. I, 295. ff. 307. ff. 330. ff. — Bedeutung des Namens I, 383. A. lebt, ob er wohl erschlagen wurde I, 403. — S. \gesperrt{Kain} 2.

Abeloniten. Ketzer, warum so genannt II, 1266.

Abendmahl, das heilige. Darin ist noch etwas Anderes als in der Taufe, nämlich: Leib und Blut Christi I, 279. f., ob er wohl zur Rechten Gottes ist I, 1487. f.; doch nicht, wenn man es nicht genießt I, 727. Dazu soll man offenbare Sünder nicht lassen II, 385. Die Einsetzungsworte stärken die Hoffnung von Vergebung I, 1066.; es macht den Einzelnen seiner Seligkeit gewiß I, 1066.; Gott speist uns, nicht der Diener II, 1842. f.; man soll es in der gemeinen Versammlung gebrauchen II, 779, mit Kniebeugen ehren II, 1842. Der \gesperrt{Pabst} hat den Kelch geraubt, aus dem A. ein Opfer gemacht I, 844. 1063.; wie er den Irrthum von der Einen Gestalt bestätigen will I, 727. 1175. f.; Christus hat Brod und Wein nicht geopfert I, 902.; die Papisten haben das A. in Abgötterei verwandelt I, 1664.; machen es zu einem äußerlichen Werk, damit man der Kirche Gehorsam beweist II, 38.; die Einsetzungsworte durften nur die Priester, und zwar nicht mit lauter Stimme lesen I, 1664. — S. \gesperrt{Gnadenmittel}.

Aberglaube. Entsteht aus dem rechten Gottesdienst II, 13.; bleibt eine Zeit wie die andere I, 727. f.; ist gar ein schädlicher Regent I, 731.

Abfall. A. von Gott und seiner Gnade ist möglich I, 520. 676.; erfolgt gemeiniglich nach 20 Jahren I, 1636.; die Nachkommen fallen gemeiniglich ab I, 1453. f. 1455. 1638. f. 1662. f.; II, 1825. f.; geschieht zuerst im Herzen, dann folgen die äußerlichen Sünden I, 635., und vielfältige Abgötterei II, 912. f.; wird nicht ungestraft hingehen I, 553.; zog die Sprachverwirrung nach sich I, 692; wohin er endlich führt I, 480. — S. \gesperrt{Abtrünnige}.

Abgötterei. A. hat gewöhnlich einen herrlichen Ursprung I, 1667. f.; entsteht aus der rechten Religion I, 663.; II, 13.; folgt aus dem Abfall vom wahren Gott II, 912. f.; schon von Kain aufgebracht I, 304. f.; ist aller Gottesdienst ohne Gottes Gebot I, 778.; ist der Ursprung aller Sünde II, 910.; ist Jammer und Noth II, 1684.; wird allein erkannt aus Gottes Wort I, 736.; man strafe ihn I, 730. — S. \gesperrt{Gestirndienst}, \gesperrt{Götzendienst}, \gesperrt{Kälberdienst}, \gesperrt{Moloch}.

Abimelech. 1. A. war Königsname, nicht Name einer einzelnen Person II, 185.

2. \gesperrt{Abimelech}, Abrahams Zeitgenosse. A. war fromm I, 1463. 1478. f.; gläubig I, 1322. ff.; gottselig I, 1313. 1318.; keusch I, 1322. f.; liebte das Predigtamt I, 1356. f.; war ein Bischof in seinem Reich I, 1335. 1340. — Er sündigte dadurch, daß er als Obrigkeit nicht wußte, daß Abraham Unrecht geschehen I, 1466. ff. Er kam zur Buße und einer klaren Erkenntniß Gottes I, 1335. 1345. — S. \gesperrt{Abraham} 20.

3. \gesperrt{Abimelech}, Isaaks Zeitgenosse. Wohl nicht derselbe A., der Abrahams Zeitgenosse war II, 192., sondern dessen Sohn II, 225. f. War listiger als sein Vater II, 225. f.; fromm II, 192.; gottesfürchtig und haßte den Ehebruch II, 185. ff. — S. \gesperrt{Isaak} 12.

Ablaß. S. \gesperrt{Fegfeuer}.

Abortion. Ist Sünde I, 1748.

Abraham. 1. \gesperrt{Im Allgemeinen} I, 717. ff. Seine Geschichte wird verachtet von der Welt I, 1300. ff. und den Papisten I, 1302. An die Stelle des Geschlechts A.’s ist das Geschlecht Christi getreten I, 1056.

2. \gesperrt{Ehre und Bedeutung} A.’s. Mit A. beginnt eine neue Kirche I, 718. f.; ein großer Mann, mit Verheißungen geziert I, 1064.; steht über den andern Patriarchen I, 1054. f. 1332.; das Haupt aller Heiligen I, 918. 921.; 1686.; 1155. 1355. f.; übertrifft Aeneas, Achilles, Agamemnon I, 919. Die Juden wissen ihn nicht genug zu rühmen I, 733., wird aber in der Schrift nicht wie Noah gelobt I, 729. f. Seine sonderliche Ehre ist, daß Gott mit ihm redet I, 718.

3. \gesperrt{Geburt} A.’s. Ob er der Erstgeborene war I, 298. 435. 720. f. 769. 1599. f.

4. \gesperrt{Namen} A.’s. Hebräer I, 715. f. 885.; Gottes Freund und Geliebter I, 1117.; Fürst Gottes I, 1628. 1631.; Gesegneter Gottes I, 908. f. — Von der Veränderung seines Namens aus Abram in Abraham I, 1036. 1043. f. 1052. f.

5. \gesperrt{Armuth und Reichthum} A.’s. I, 826. ff. — An seinem Reichthum stoßen sich die Mönche und Andere I, 829. f.; er hing sein Herz nicht daran I, 918.; war recht arm, nicht wie die Mönche I, 1164., die mit ihm nicht zu vergleichen I, 738. f., insonderheit nicht in dem, wie er Alles verlassen I, 739. 761.

\quellseite{541}
6. \gesperrt{Reisen} A.’s. I, 770. ff. 783. ff. 826. ff. 1300. ff. — A.’s Reisen geschehen nach Gottes Willen I, 1352. f.; wie die Propheten seine Wanderschaft angesehen I, 1305.; sie waren beschwerlich und gefährlich I, 771. f.; 1302.

7. \gesperrt{Berufung} A.’s. I, 728. ff. Wie alt A. war I, 1600.; wohl durch Sem berufen I, 711. f. 734. f.; berufen aus Gnaden I, 730. f. 732. f.; dadurch ein neuer Mensch und Patriarch I, 731.

8. \gesperrt{Erscheinungen und Offenbarungen}. Mit keinem Andern hat Gott so oft geredet I, 1054. f.; dies geschah wohl achtmal I, 1544.; Gott redet mit ihm wie mit einem Freunde I, 938. f.; dies ist das Vornehmste in A.’s Geschichte I, 861.; geschah wohl durch das Predigtamt, durch Sem I, 869. f. — Die Erscheinung der \gesperrt{drei Männer} I, 1130. ff. 1154. ff. 1181. ff.; des \gesperrt{Engels} bei Isaaks Opferung I, 1522. ff.; der \gesperrt{HErr} erschien ihm wohl, ehe er zu Sem kam I, 777.

9. \gesperrt{Verheißungen}, A. gegeben. I, 740. ff. 860. ff. 936. ff. 951. ff. 1035. ff. 1059. ff. 1091. ff. 1547. ff. 1559. ff. — Worin A. uns etwas voraus hatte, wir dem A. I, 1385. Die leiblichen Verheißungen schließen Christum und das ewige Leben in sich I, 1096. ff.; A. verstand die leiblichen Verheißungen zugleich geistlich I, 1107.; A. ließ sich an den geistlichen genügen I, 908. 916. ff. A. hat Gott damit nicht versucht, daß er ein Zeichen zur Bestätigung der Verheißung vom Besitz Canaans begehrte I, 953. ff.; Gott gab diese Verheißung, daß man wissen könnte, woher der verheißene Same kommen sollte II, 141.; A. war ein Fremdling in Canaan in seinem Leben, nach seinem Tode besaß er es I, 1060. f.; mußte trotz der Verheißung wie Kain hin und her ziehen I, 782. Die Verheißung, daß Christus von ihm sollte geboren werden, empfing er zum Lohn, da er schon gerecht war II, 143.

10. \gesperrt{Glaube und Religion} A.’s. A. hatte zum \gesperrt{Lehrer} Noah I, 655. f., Sem I, 898.; ward ein \gesperrt{Götzendiener} I, 722. 711.; wohl ein „frommer“ Mann, doch des Teufels Leibeigener I, 731. — A. hatte die \gesperrt{Erkenntniß} von der heiligen Dreieinigkeit I, 1153. 1183. 1673. f., von Christo I, 748. 1627., daß der verheißene Same der Heiland der Welt sei I, 1657. und Gottes Sohn I, 1674.; von dem Dienst der heiligen Engel I, 1684. ff. A. nahm zu an Erkenntniß II, 77. f.; wie er die Verheißung, daß der Größere dem Kleineren dienen sollte, verstanden II, 80. f. A. \gesperrt{glaubte} eine Auferstehung der Todten u. ein ewiges Leben I, 964. f. 1060. ff. 1064. 1488. f. — A. glaubte an den zukünftigen, wir an den gekommenen Christus I, 949. f. Die Schrift rühmt seinen Glauben I, 739. 773. 791. 804.; je größer die Gefahr, desto stärker sein Glaube I, 809., der größer als Gideons I, 891., mit dem der unsrige nicht zu vergleichen I, 743. f. 772. f. Er glaubte, obwohl die Verheißung der Vernunft dunkel war I, 1098., und unmöglich schien I, 740. ff.; glaubte dem von Gott gegebenen, erbetenen Zeichen I, 958.; hielt sich an die Verheißung von Isaak I, 1384. Worin er seinen Glauben bewies I, 758. 771.; 775.; 1487. f.; 1692. Sein Glaube begann zu straucheln I, 789. f. A. hatte auch Schwachheiten I, 1354. f.

11. \gesperrt{Bund} A.’s mit Gott I, 1054. ff. War ein doppelter: der Beschneidung und der Verheißung I, 1114. f.

12. \gesperrt{Beschneidung}. I, 1036. ff. 1108. ff. Aehnlichkeit der B. A.’s und der B. Christi I, 1024. f.

13. \gesperrt{Gerechtigkeit} A.’s \gesperrt{vor Gott}. A. war gerecht vor Gott durch den Glauben I, 940. ff. 952. f. 1017. f. 1580. 1583. f., nicht durch Werke I, 944. 1536. f. 1579. ff., nicht durch die Beschneidung I, 1017. f.; ehe er im Gehorsam aus Ur auszog I, 760., oder sich beschneiden ließ I, 1010. f. 1017. f. 1038. Die Papisten behaupten, er ward durch Werke gerecht I, 1579. ff.

14. \gesperrt{Tugenden} A.’s. Im \gesperrt{Allgemeinen} I, 915. ff.; an A. findet sich eine große Anzahl geistlicher und weltlicher Tugenden I, 1640. f. — A. war bescheiden I, 889. f.; beständig I, 781. f. 826. ff.; billig und rechtschaffen I, 919.; demüthig und ehrerbietig I, 915. ff. 1142. ff. 1148. 1182. 1476. f. 1596.; fromm I, 899.; gastfrei I, 1137. ff. 1142. ff. 1154. ff., woher er diese Tugend gelernt I, 1134. ff. 1144. ff. 1154. ff.; geduldig I, 840. ff. 908. 916. 1304. f.; II, 237.; gehorsam I, 736. ff. 755. ff. 869. f. 1074. f. 1123. ff. 1262. ff. 1375. 1412. ff. 1496. ff., mit A.’s Gehorsam ist der der Mönche nicht zu vergleichen I, 758. f. 1412. ff. 1417. ff.; nicht geizig I, 915. ff.; höflich I, 1475. 1628. ff.; keusch I, 1170. 1744. f.; liebreich I, 840. ff. 1197. ff.; mitleidig I, 1286. f.; wahrhaftig I, 788. ff. 1313.

15. \gesperrt{Werke} A.’s. Im Allgemeinen I, 826. ff. 886. ff. — A. ist ein Vater auch der guten Werke I, 1141. f.; verrichtet sie alle im Glauben I, 1305. 1682. ff.; braucht dabei seine Vernunft I, 786. 890. f.; handelt wie ein geistlicher, verständiger I, 1675. ff. und weltkluger Mann I, 1474. ff. 1638. f.; thut keine sonderliche Werke, beweist aber damit Glauben und Liebe I, 919. f.; wird seiner geringen Werke wegen von den Mönchen verachtet I, 1521.; gegen A.’s sind des Pabstthums Werke stinkender Dreck I, 774.; es folgt nicht: Was Gott an A. gefiel, gefällt ihm auch an uns I, 779. — A.’s Werke als \gesperrt{Hausvater} I, 918. f.; seine Pilgrimschaft I, 1305. f.; pflanzt Bäume I, 1479. f.; kauft Begräbnißstätte I, 1616. f.; II, 103. f.; sorgt für Isaak um ein Weib I, 1682.; als \gesperrt{Ehemann} I, 1742. ff. 1749. f.; im \gesperrt{Verkehr} mit Anderen I, 1675. f.; als \gesperrt{Bischof und Priester}, der er war I, 778. 1521.; wie er sein Amt verwaltet I, 781. f.; war Andern ein Segen I, 746. 767. 906. f.; baute Altäre und predigte I, 778. 781. 839. 874.; 826. ff. 838. 1478. f. 1631.; lehrte I, 805. 819. 837. f. 1332. f. 1352. 1361.; II, 80.; I, 773. 957. 1661. f. 1667. f.; hat wohl eine kleine Historie zusammengetragen I, 1753.; betete I, 1197. ff. 1321. 1360. ff.; 1107. ff. 1288. f. 1362. f.; als \gesperrt{Kriegsheld} I, 886. ff.; ist ein größerer Held als Alexander etc. I, 889. ff. und Simson I, 890.; uns ein Beispiel I, 889. f.

16. \gesperrt{Sünden} A.’s. Daß er bei den Gerariten Sarah für seine Schwester ausgab, war Sünde I, 1325. f., die jedoch zum Besten diente I, 1328. S. oben 10.

17. \gesperrt{Anfechtungen und Leiden} A.’s. A. mußte vor anderen Patriarchen schwere \gesperrt{Anfechtungen} erdulden I, 1482.; eine folgte der anderen I, 1482. ff.; ist ein Märtyrer über alle Märtyrer I, 1417.; so ward sein Fleisch immer getödtet I, 1350., A. vor Vermessenheit bewahrt 923. ff. 928. f. An ihm sehen wir zweierlei Tödtung sein selbst I, 1421.; wie er sein Fleisch getödtet I, 1505. f.; gegen die Tödtung seines Fleisches ist die Kasteiung der Mönche nichts I, 1414. f. A. ward \gesperrt{angefochten} I, 783. ff.; 802. ff. 810.; 874.; 920. ff. 928. f. 932. f. 934. f.; 1070. ff.; 1284. f.; 1286. 1313. f.; 1298.; 1318.; 1482. ff.; 1503. f.; 1528. f.; 790. In diesen Anfechtungen hielt er sich an die Verheißungen I, 1487.; ward er getröstet I, 920. ff. 929. 1060. 1323. 1340. ff. 1342.; siegte er durch den Glauben I, 785. 814. 874. 1318. f. 1421. f. 1487. ff. 1529. ff.; II, 784. — A. mußte unter Feinden wohnen I, 776., wurde \gesperrt{verfolgt} I, 780. f. 839. f., \gesperrt{verspottet} I, 1131.

18. \gesperrt{Gesinde} A.’s. War groß I, 772. 887. f. 1618.; eine heilige Kirche I, 773.; fromm und gehorsam I, 919. 1129. 1131.; frömmer als das jetzige I, 739. 773.; zu loben, daß es A. ins Elend folgte I, 773.; dies war A. sehr tröstlich I, 774.

19. \gesperrt{Kirche}. Empfängt zu A.’s Zeit ein größer Licht I, 719.; eine neue fängt mit A. an I, 718.; A.’s Haus war die rechte I, 1193. f. 1360. f.; ihm folgten viele Verfolgte I, 1213.; Kirche ward gemehrt und getröstet I, 907.; A. bekehrte Viele I, 1303.; seine Kirche stand in Ansehen I, 1631. ff.

20. \gesperrt{A. und Abimelech}. I, 1300. ff. 1330. ff. 1345. ff. 1356. ff. 1360. ff. 1460. ff. A.’s Argwohn gegen Abimelech I, 1465. ff.

21. \gesperrt{A. und die Egypter}. I, 802. ff. 811. ff. 818. ff. A.’s Argwohn gegen die Egypter I, 794. ff. 798. 801. f. Ob A. in Egypten die Sternkunde gelehrt I, 805.

22. \gesperrt{A. und Elieser}. I, 1640. ff. 1656. ff. 1681. ff. 1706. ff. 1722. ff.

23. \gesperrt{A. und Hagar}. I, 974. ff. 983. ff. 1410. ff. A. legte sich nicht aus Unkeuschheit zu Hagar I, 969. f. 972. f. 987.; dadurch kein Ehebrecher I, 971. ff.; liebte sie I, 1389. f.; Ausstoßung ward ihm schwer I, 1394. 1411. f. 1420. ff. 1422. f.

24. \gesperrt{A. und die Hebroniten}. I, 1628. ff.

25. \gesperrt{A. und Isaak}. I, 1091. ff. 1368. ff. 1496. ff. 1522. ff. 1640. ff.

26. \gesperrt{A. und Ismael}. I, 1108. ff. 1386. ff. 1410. ff. — A. hielt Ismael für den verheißenen Samen I, 1007. 1035.; liebte ihn I, 1108. 1389. f.; seine Ausstoßung ward A. schwer I, 1394. f. 1411. f. 1420. ff.

27. \gesperrt{A. und Ketura}. I, 1742. ff. — Heirathet sie nicht aus Geilheit I, 1744. f. 1746. ff.; vielleicht, weil Isaaks Ehe so lange unfruchtbar blieb II, 3.

28. \gesperrt{A. und Lot}. I, 840. ff. — Warum Lot die Gegend am Jordan wählt, A. in Canaan bleibt I, 853.; wie A. die Nachricht von Lots Gefängniß aufnimmt I, 886.; Lot um A.’s Gebetes willen errettet I, 1288. f.; Lot und seine Töchter von A. aufgenommen I, 1297. 1342. f., über die Blutschande getröstet I, 1300. f. 1342. f.

29. \gesperrt{A. und Mamre}. A. findet Aufnahme bei Mamre I, 874. f. 1141.; bringt ihn zur Erkenntniß I, 875. f.

30. \gesperrt{A. und Melchisedek}. I, 897. ff.

31. \gesperrt{A. und Nahor}. I, 1592. ff.

32. \gesperrt{A. und die Philister}. I, 1457. ff.

33. \gesperrt{A. und Sarah}. I, 968. ff. 974. ff. 983. ff. 1386. 1602. ff. 1615. ff. Hatten viel Anfechtungen und Trübsale I, 968. ff.; ihr Verhalten darin I, 976. ff.; hofften auf Gott I, 813.; beteten I, 991.; liebten Ismael I, 1007. 1035.; trösteten Hagar und Ismael I, 1412. f.; warum A. Sarah in Canaan begraben wollte I, 1627. Als \gesperrt{Eheleute} zu bewundern I, 793.; ein Exempel für Eheleute I, 1160. 1163. f. 1353. f.; 794. 986. f. 1253.; 793.; hatten langen heftigen Streit I, 1386. ff. 1390. ff. 1410.; II, 1873. f. A. theilte Sarah den Befehl, Isaak zu opfern, nicht mit I, 1490. f.

34. \gesperrt{A. und Sem}. A. hat wohl Sem auf seinen Reisen aufgesucht I, 774. f.; bekam durch ihn Befehl, Sarah zu gehorchen I, 1395.

35. \gesperrt{A. u. die Sodomiter}. I, 895. 913. ff. 1181. ff.

36. \gesperrt{Tod} A.’s. I, 1752. ff. Wie Gott A.’s Lohn sein kann, da er doch stirbt I, 964.

37. \gesperrt{Kinder oder Same} A.’s. Die Juden sind es nicht I, 1669.; leiblicher und geistlicher Same I, 1396. ff.; heiliger und gottloser I, 1585. Zu A.’s Kindern macht die Verheißung und der Glaube I, 1053. 1101. 1132. 1396. ff. 1438.; II, 1872.; A.’s Geschlecht ward von allen Völkern gehaßt II, 1133.

Abrahams Same. Die Lehre davon ist die wichtigste in der Schrift II, 143.; in der Verheißung von ihm ist fast die ganze christliche Lehre begriffen I, 1578. 1583.; wahrer Gott, wahrer Mensch, ohne Sünde I, 1570. f. 1674.; II, 143.; ist Christus I, 1560. 1569. f.; ist selbst der Segen I, 1566. f. Sein \gesperrt{Segen} ist ein geistlicher I, 1564. f.; ohne denselben ist kein Volk gesegnet I, 1591. f.; nimmt den Tod und allen Fluch weg I, 1571.; macht voll guter Früchte I, 1578.; schließt allen Ruhm und eigene Gerechtigkeit aus I, 1572. f.; geht alle Menschen an I, 1571., auch die Heiden I, 1560.; desselben wird man durch den Glauben theilhaftig I, 1367. f.; Werkheilige haben kein Theil daran I, 1368. f.; darnach haben sich die Heiligen des A.\,T. herzlich gesehnt I, 1591. Ob Gott ihn hätte erwecken können, wenn Jakob von Esau erwürgt worden wäre II, 762. Jüdische Auslegung von A.’s Samen I, 1564.

Absolution. Jeder Christ kann sie sprechen II, 1003., auch ein Weib oder Kind II, 2057. ff.; beim Pfarrherrn oder einem christlichen Bruder zu suchen II, 779.; dabei nicht auf den absolvirenden Menschen zu sehen II, 287.; wird aus göttlichem Befehl gesprochen II, 300. f. 1899.; Gott selbst absolviret II, 435. 437.; sie gibt Vergebung der Sünde II, 298. f. 369., erlöst von Sünde, Teufel, Hölle II, 1891. 2009., verkündigt und gibt das Reich Gottes II, 411.; Absolutionsformel zu Luthers Zeit II, 299.

Abtrünnige. Verfolger ihres Ordens I, 689.; die größten Feinde der Kirche II, 629.

Accursius. Jurist, bekennt seine Unwissenheit II, 1927.

Achilles. Mit Abraham nicht zu vergleichen I, 919.; ein Held und doch freundlich gegen Briseis II, 603.; konnte sein Herz nicht bezwingen II, 1387.

Achsib. 1. Stadt in Asser II, 1163.

2. Stadt in Dan II, 1163. 1186.

Ackerbau. Von Gott geboten I, 1479.

Ackerbauer. A. haben den glücklichsten Stand I, 259.; II, 553.

Ada. II, 990. — Bedeutet Schmuck, Halsband I, 386.; II, 989. — S. \gesperrt{Esau} 3.

Adam. 1. \gesperrt{Adam}. I, 66. ff. 101. ff. 124. ff. Bedeutung des \gesperrt{Namens} und Anlaß dazu I, 383. 410.; II, 68. f.; ob nur ein Verheiratheter A. heißen könne I, 412. A. hatte trefflichen \gesperrt{Verstand und Weisheit} I, 145. f.; wußte wohl aus Eingebung des Heiligen Geistes, daß Abel erschlagen war I, 346.; prophezeite im Paradies I, 167. 170.; hatte die Hoffnung auf ein ewiges Leben I, 269.; seine Erkenntniß natürlicher Dinge selbst nach dem Fall viel vortrefflicher als unsere mühsam erlangte Erkenntniß I, 146. f. Mußte den \gesperrt{Schweiß} des Haus-, Regierungs- und Kirchenstandes zugleich schwitzen I, 261. Ein größerer \gesperrt{Lehrer} denn Mose I, 328.; lehrte vom Glauben I, 328. \gesperrt{Starb} zur Zeit Lamechs I, 391.; wo \gesperrt{begraben} I, 1610. f. Zeitgenossen A.’s I, 407. 416. War eine \gesperrt{Figur Christi} I, 268. — S. \gesperrt{Schooß}.

2. \gesperrt{Adam und Eva}. Im Allgemeinen I, 160. ff.; 172. ff.; 196. ff.; 207. ff.; 223. ff.; 242. ff.; 268. ff.;
\quellseite{542}
290. ff.; 302. ff.; 317. ff.; 330. ff.; 347. ff.; 410. ff. — Worin A. und E. gleich waren I, 81. 141. 226., worin ungleich I, 84. 185. 268.; die Scholastiker halten A. für den obersten, E. für den untersten Theil der Vernunft I, 195. — Sie hatten die \gesperrt{Hoffnung} auf ein ewiges Leben I, 269. f.; hielten sich im \gesperrt{Glauben} an die erste Verheißung I, 235. ff.; womit sie sich nach dem Fall getröstet I, 244. 304.; II, 919. f.; zeugten Kinder, um den verheißenen Samen zu erlangen I, 290. f. — Aus Gottesfurcht und \gesperrt{Gehorsam} stießen sie Kain aus I, 375. f. — Ihre \gesperrt{Anfechtungen} über Kains Brudermord I, 342. ff.; hatten den Trost nicht, daß es Andern auch schon so ergangen I, 344.; ihr Herzeleid über Kains Ausstoßung I, 347. 375. f.; sie waren Eines Gemüths und Eines Willens I, 82.; A. hat E. im Paradiese unterrichtet I, 176.; A. gelobte nicht ewige Keuschheit nach Abels Tod I, 413. — Hatten wohl vor Seths Geburt eine Verheißung I, 395.; wie sie ihre Kinder unterrichtet und erzogen I, 301. ff. 408. — S. \gesperrt{Stand der Unschuld}, \gesperrt{Sündenfall}.

Adel. Soll sich der Tugend befleißigen und den Seinen dienen I, 879.

Adoption. Ist in der Schrift sehr gemein II, 1000.; A. Ephraims und Manasses II, 1848.

Affecten. Die natürlichen A. sind nicht gar zu verwerfen I, 1614.; sind an sich gut, doch durch die Erbsünde verderbt II, 1501. 1672.; bleiben auch in den Gottseligen II, 1669. f.; wann zu unterdrücken II, 1679. f. — S. \gesperrt{Natur} 3., \gesperrt{Neigungen}, \gesperrt{Väter} II, 5., \gesperrt{Zorn} 2.

Agamemnon. Nicht mit Abraham zu vergleichen I, 919.

Agatha. Märtyrerin, getrost II, 1989. 1994.

Agnes. Märtyrerin, keusch II, 1268. f.; todesmuthig II, 1903. — S. \gesperrt{Kloster}.

Agricola. Toller, unerfahrener Geist I, 1212.; wollte für heilig und gerecht gelten II, 708.

Ahab. Gottlos, hatte ein herrliches Regiment II, 2043.

Ahalibama. Heißt Schloßfrau II, 989. — S. \gesperrt{Esau} 3.

Ahas. Abgöttisch II, 2014.; wollte, Abraham nach, seinen Sohn opfern I, 1496.; handelte damit wider Gottes Wort I, 1524., ward ein Mörder I, 1179. f.

Albrecht, Bischof von Mainz, verwarf die reine Lehre wissentlich I, 1646. — S. \gesperrt{Bischof} 3.

Alexander, der Große. Führte selbst bei Allem die Aufsicht II, 1241.; mit ihm ist kein jüdischer oder christlicher König zu vergleichen II, 462.; war kein solcher Kriegsheld wie Abraham I, 889., oder Joseph II, 1376.; konnte sich nicht bezwingen II, 1387.; seine Geschichte ist mit den Historien der heiligen Schrift nicht zu vergleichen II, 462.; hatte Glück um der Kirche willen I, 1409.; ward hoffährtig und fiel II, 1634.; sein Leichnam dem Augustus gezeigt II, 2073.

Alkoran. Warum Mahomet seinem Buche diesen Namen gegeben II, 223.; der Teufel ist der Meister des A.’s II, 401.; ist ein Ungeheuer, vermischt Gesetz und Evangelium I, 1009.; was er sagt: vom Fall des Teufels II, 401., vom Leviathan II, 1948.; bei Auslegung des A.’s befleißigen sich die Türken der Allegorien II, 559. — S. \gesperrt{Mahomet}.

Allegorien. I, 610. ff. — Von den närrischen Mönchen hergekommen I, 559.; eine Zeitlang dem historischen Verstande vorgezogen I, 284. 610.; sind wie eine schöne Hure II, 559.; welche Väter sich der A. beflissen II, 557. \gesperrt{Schaden und Nutzen} der A.: gefährlich, verdunkeln die Lehre I, 284.; II, 557. f. 802.; beweisen nichts I, 230. f. 286. 903. 952. 1150. f.; II, 235. 802.; zieren nur I, 285. f. 610. 1150. f.; II, 774. 802. \gesperrt{Gebrauch der} A.: man soll nicht bloße A. suchen II, 774.; sind nur der Schaum der heiligen Schrift I, 285.; man soll sie nicht ganz verwerfen I, 611.; man soll sie mit Vorsicht brauchen I, 611.; II, 559. f.; welche und welche nicht zu verwerfen I, 611.; um zwischen den A. der Väter zu unterscheiden, bedarf man eines guten Verstandes II, 811.; mit A. hat nicht Jeder Glück II, 774. \gesperrt{Erfordernisse} guter A.: der buchstäbliche Sinn muß die Grundlage der A. sein I, 286. f.; II, 802.; müssen auf den Glauben gehen I, 612. 616. f. 626. f.; II, 402. 560. \gesperrt{Beispiele von} A.: alberne I, 225. f.; gute I, 286. f. 514. f. 603–627.; 903. 952. 1204. 1495. 1541.; II, 235. 402. 475. 811. ff.

Allmacht. Was unter Gottes ordentlicher und unter seiner unmittelbaren A. zu verstehen ist I, 1250.; A. offenbart sich in der Schöpfung I, 31. 60.; kann aus Nichts Alles machen II, 1677. 1700. 1702., aus Bösem Gutes II, 893. f., das Erschaffene erhalten und vermehren I, 1540. f.; nach Gottes ordentlicher A. sollen wir uns richten I, 1250.

Allwissenheit. Alles ist Gott gegenwärtig I, 960. 963.

Altar. Was einen „A. bauen“ heißt I, 778.; II, 222. 854. 1751.; wozu wir Altäre in den Kirchen haben I, 115.

Alter und Kindheit. Sind nicht zu vergleichen II, 249.; Einfluß des A.’s auf Sinne und Kräfte II, 1846. 1858.

Amalek. I, 994.

Amalekiter. Nachkommen Esaus von Amalek, mütterlicherseits Hethiter II, 994.; sehr stolz ib.; wohnten zwischen Egypten und Canaan I, 881.

Ambrosius. Ehelos, wartete seines Berufs I, 1320. f.; hatte mit der Brunst des Fleisches zu kämpfen II, 530.; hatte sonderliche Freude an dem Werk der Menschwerdung Christi II, 400.; legte die Schrift nicht so gut aus wie Augustin II, 79.; verstand die active und passive Bedeutung der Wörter nicht II, 1488.; A.’s Meinung über Hiob II, 1014. f.

Amiris. Flucht aus Sibaris I, 1258. f.

Ammi. Bedeutung des Namens und Anlaß I, 1298. ff.

Ammoniter. Vor andern Völkern herrlich gesegnet I, 1298. f.; nicht zu Kirchenämtern in Israel zugelassen I, 368. 1298. f.; viele selig I, 368.

Amos. Lehrte im Reich Ephraim II, 2048.

Amram. I, 959.

Amraphel. Ob A. Nimrod war I, 876. f.

Amt. Das A. entdeckt den Mann I, 834. Im Amt dient man Gott II, 1411.; wie man sein A. führen soll II, 1394. f.; man fliehe die Beschwerden nicht II, 1408. f.; thöricht, auf ein A. stolz sein II, 1395.; Aemter sind Strafen II, 1395.; machen uns zu Dienern II, 1397.; einst ist Rechenschaft abzulegen II, 1394. Man soll sich von \gesperrt{weltlichen Aemtern} nicht entziehen I, 833. f. — S. \gesperrt{Beruf}, \gesperrt{Stände}.

Ana. Ob durch Blutschande gezeugt II, 995.; ob er die Maulpferde erfunden II, 997. 1007. f.

Anachoreten. Sind Einsiedler I, 997.

Analogie des Glaubens. Was die A. d. G. ist I, 1531.

Anastasia. Märtyrerin II, 1982. f.

Anbetung. Verschiedene Arten der A. I, 1635. II, 306. 820. f.; was „anbeten“ heißt I, 1058.; II, 1046.

Aeneas. Ist mit Abraham nicht zu vergleichen I, 919.

Aner. Leutseliger, frommer Mann I, 886. 888.; hielt sich zu Abrahams Kirche II, 1069.

Anfang. Was Mose „Anfang“ nennt I, 12.

Anfechtungen. 1. Die \gesperrt{Lehre} von den A. muß man um der künftigen Prediger willen treiben II, 783. f.

2. \gesperrt{Art und Natur} der A. \gesperrt{Mancherlei} A. fallen den Gottseligen vor I, 1072. 1532.; II, 417. f.; eine folgt der anderen I, 1482.; je frömmer der Christ, desto mehr A. I, 1486.; wenn sie da sind, scheinen sie mit der Verheißung zu streiten II, 823.; darin ist es schwer sich der Gnade Gottes zu trösten II, 968.; der Christ schwankt oft zwischen Ungeduld und Hoffnung II, 1317.; sind der rechte stetige Tod II, 1094. Das \gesperrt{Hauptstück} aller A. ist, daß man zweifelt, ob Gott wirklich dies oder jenes gesagt habe I, 187. Zu den \gesperrt{leiblichen} A. fügt Satan auch \gesperrt{geistliche} I, 1424. ff. 1484.; die leiblichen werden heftiger, wenn der Schrecken im Gewissen dazu kommt II, 967. f. Erst versucht Satan uns vom Wort abzuführen, dann folgen die Sünden wider die zweite Tafel I, 199. \gesperrt{Es sind} A., wenn Gott sich zu widersprechen scheint I, 1487. f.; wenn man Gottes Gnade nicht fühlt I, 927.; wenn man meint, man sei von Gott verlassen I, 556. ff.; II, 1092. und ganz in Satans Gewalt gegeben II, 733. f. Wir werden angefochten \gesperrt{über} die Größe der Sünde I, 1177.; darüber, als ob wir nicht glaubten II, 179.; ob wir selig werden I, 1177. Die \gesperrt{schwersten} A. sind: wenn man darüber disputirt, warum Gott dies oder jenes also gemacht I, 186.; Unglaube und Verzweiflung II, 782. f. 1319., welche A. nicht Jeder fassen kann ib.; Vermessenheit, Ehrgeiz, Selbstvertrauen I, 924. f.; die \gesperrt{allerhöchste} A. ist, wenn man Gott im Herzen feind wird I, 1484.; in den \gesperrt{hohen} A. muß man mit Gott kämpfen II, 786. Die \gesperrt{Jugend} ficht Brunst des Fleisches, das \gesperrt{männliche} Alter Ungeduld und Rachgier an I, 558. f.; II, 20. ff. 1227. f. 1252. ff.; Eheleute haben schwere A. I, 984. 1392.; ebenso die allerhöchsten Heiligen, von Sünden wider die erste Tafel I, 923., Unglauben und Verzweiflung II, 782. f. In A. ist alles schwarz und traurig II, 797. f. 1538. f., werden die Heiligen furchtsam I, 927., zweifeln an der Verheißung I, 934., meinen, Gott nehme sich ihrer nicht an I, 1003.; die Sinne werden irre II, 787. f. 792.; der Leib entkräftet I, 1451. Hohe A. kommen, wenn wir allein sind I, 939., in der Nacht II, 771. f.

3. \gesperrt{Urheber} der A. Satan richtet die A. an I, 1140.; wie er in den A. zu uns redet II, 755.; er ficht uns ohne Aufhören an, um uns matt zu machen I, 938.; II, 1269.

4. \gesperrt{Gefährlichkeit} der A. Die A. sind dem Fleisch ohne den Geist unerträglich I, 560.; geistliche A. tödten Viele II, 1320.; wo nicht Trost folgte, würde der Glaube sinken I, 828. und kein Mensch selig I, 777.

5. \gesperrt{Trost in} A. Gott ist gewiß in den A. bei uns I, 777. und uns günstig II, 968.; Gott läßt die A. nicht größer werden, als wir sie ertragen können I, 784. f.; Gott kann in seinen Verheißungen nicht lügen I, 1489.; fühlt man auch die Gnade nicht, so ist sie doch unwandelbar I, 927. \gesperrt{Von dem Trost}, wenn Satan die Sünde groß macht I, 1177.; wenn es scheint, als sei man ganz in Satan Gewalt gegeben II, 733. f.; bei der A. über Beharrlichkeit im Glauben II, 180., darüber, ob man selig werden wird I, 1177. 1484. f.

6. \gesperrt{Kampf in} A. \gesperrt{und Sieg}. Ohne den Heiligen Geist ist Sieg nicht möglich II, 787.; Gott hilft I, 807.; Gottselige müssen doch endlich siegen II, 800.; der Glaube kämpft II, 1098.; wie es dabei zugeht II, 762.; der Glaube siegt II, 897. f. 1098. f., durch das Wort I, 786.; II, 417. f.; wie gegen A. zur Unzucht zu kämpfen II, 20. ff.; auf A. folgt Freude I, 811.; II, 800. u. Erkenntniß d. Liebe Gottes II, 808. f. — S. \gesperrt{Glaube} 11.

7. \gesperrt{Rechtes Verhalten in} A. In A. soll man die Einsamkeit fliehen I, 840. f., die Gedanken der A. nicht ins Urtheil und Gewissen aufnehmen, sondern durch Gottes Wort zurücktreiben II, 785. f., geduldig sein II, 1296., in der Gottseligkeit verharren II, 1255., beständig sein im Glauben II, 193. 1384., auf Gott vertrauen II, 744. ff., nicht verzagen II, 57. 1569.; wie wir den Glauben stärken sollen I, 771.; sollen uns an die Verheißungen halten I, 1484. ff. 1489.; II, 754., ans Wort II, 52. ff. 56. 197. 201. f. 305. 763. 1090. f. 1385. f., an die Sacramente II, 1092. f., in der heimlichen Beichte Trost suchen II, 1002. f., die A. Anderer betrachten I, 1490.; II, 806. f., beten I, 1485. In Anfechtungen wegen hoher Gaben soll man auf die Vergebung der Sünde als die höchste Gabe sehen I, 926.; wenn man denkt, Gott habe Einen verworfen, soll man dafür halten, er habe ihn in den Armen und herze ihn II, 806.; in der A., ob Gott uns selig machen wolle, sollen wir uns an die Verheißung halten I, 1484.

8. \gesperrt{Zweck und Nutzen} der A. Außer den A. dünken wir uns stark II, 741. f.; A. sind uns nöthig II, 193. 201. f.; werden geschickt, daß wir uns der Gaben nicht überheben I, 926. ff.; schaffen mancherlei Frucht I, 826. f. 1457.; II, 306. 1365. f. 1384. f.; lehren aufs Wort merken I, 1391.; II, 2038.; durch A. wird das Fleisch getödtet, der Geist lebendig gemacht II, 797. 1094.; fördern in der Erkenntniß und Glauben II, 1096. f.; lehren beten I, 939.; II, 1646. f.; lehren Geduld I, 559.

9. \gesperrt{Behandlung Angefochtener}. Angefochtenen Herzen ist nicht so bald geholfen I, 574.; II, 1685. ff.; wie zu trösten II, 784. 1686. f.; denselben muß man gewissen, starken Trost einreden I, 574.; wie Luther eine angefochtene Person getröstet, welche meinte, sie glaube nicht II, 179. — S. \gesperrt{Gewissen} 2, 3.

Vgl. \gesperrt{Abraham} 17, \gesperrt{David}, \gesperrt{Isaak} 11, \gesperrt{Jakob} 15, \gesperrt{Joseph} 11, \gesperrt{Lot} 1, \gesperrt{Noah} 11, \gesperrt{Rebekka}, \gesperrt{Sarah}, \gesperrt{Väter} II, 10., \gesperrt{Versuchung}.

Angesicht. 1. \gesperrt{Gottes}. Bedeutet seine Gnade, der Rücken aber seinen Zorn I, 1004.; was es heißt „Gottes A. leuchtet“ II, 837.; Gottes A. ist das, dadurch er sich offenbart I, 13. f. 376. f.; wann er uns sein freundliches A. zeigt und wann er es von uns wendet II, 837.; woran man sein freundliches A. erkennt II, 838.

2. \gesperrt{A. der Menschen}. Der Heilige Geist macht ein frei, fröhlich Angesicht I, 1250.

Anrufen. „Den Namen des HErrn anrufen“ I, 398. f.; Unterschied zwischen „den Namen“ und „im Namen“ des HErrn anrufen I, 837. f.; II, 223.

Anthropomorphismus. Die Schrift redet oft von Gott wie von einem Menschen I, 18. f. 487.; Beispiele: I, 701. 1195. f. 1535. f.; warum sie dies thut I, 487. f. 1535. f.; wie solche Stellen zu erklären I, 487. ff.; solcher A. ist sehr lieblich I, 1535. f. — S. \gesperrt{Menschenworte}.

Anthropomorphiten. Was sie gelehrt I, 17. ff. 487.; mit Unrecht vom Pabst verdammt I, 17. ff.; in welchem Falle mit Recht verdammt I, 487.

Antichrist. Von der Meinung, der A. komme aus dem Stamme Dan II, 2023. ff.; wie er in der Schrift beschrieben wird I, 554.; Pabst und Türke sind der A. I, 649. 1062.; Pabst der eigentliche A. I, 506. 1062. ff. 1404.; II, 1679. 2024. ff.; der A. hat trotz aller Sacramente und Gnadenzeichen lange regiert I, 1668.; dem Pabst, dem A., ist man billig feind II, 1895.

Anticipatio. Redefigur II, 548. f.; muß in der Schrift beachtet werden II, 1510. f. — S. \gesperrt{Hysteron}.

\quellseite{543}
Antilogisten. Wer und warum so genannt II, 1023.; beschuldigen die Schrift der Dunkelheit II, 1023. f.

Antinomer. Was sie wollten und lehrten Gesetz und Evangelium betreffend I, 1188. 1205. ff. 1244. 1301. 1429. 1669. f.; ihre Lehre falsch und gefährlich I, 1189. f. 1207. f. 1212. 1259. 1426. f.; sind tolle unerfahrene Geister I, 1212.; genießen die Güter dieser Welt, verschieben die Buße II, 116.; wollten sich mit Luthers Beispiel schützen I, 1428. f. Gott ist kein A. I, 1330.

Antiphrasis. Redefigur I, 1379.; II, 911. f. 1206.; Beispiele: II, 1206. 1532. f. 1738. f.

Antonius. 1. \gesperrt{Sanct} A. Man mag ihn unter die Heiligen rechnen I, 1320.; ward zu einem Gerber im Traum gewiesen I, 1181. 1436. f.; man soll seinem Exempel nicht nachfolgen I, 1075.

2. \gesperrt{Dr.} A. Englischer Märtyrer I, 352.

Apis. Abgott der Egypter I, 455.; II, 1615. 1827.

Aposiopesis. Redefigur, Beispiele I, 370.; II, 1933.

Apostel. Geistlich trunken nach Christi Auferstehung II, 1981. ff.; hatten ihre sonderlichen Offenbarungen I, 536.; erklärten die Verheißung: Ich will dein Gott sein I, 1067. — S. \gesperrt{Jünger}.

Araber. Ismaeliten II, 14.; Räuber I, 997.; streitbare Leute I, 1454.; feindeten Perser und Römer an I, 998.; jetzt unter den Türken II, 14.; Gott hatte unter ihnen eine zahlreiche Kirche I, 997.; hießen Loban II, 1307.

Aram. Ländername I, 680.; früher: Haran I, 772.

Ararat. Gebirge I, 562.

Arbaces. Wohl der Arphachsad der Schrift I, 679. f.

Arbeit. Zur A., nicht zum Müßiggang ist der Mensch geschaffen I, 126.; Gott will A. haben II, 1239., segnet sie II, 1239.; beschwerlich durch die Sünde I, 124. f. 257. f.; jetzt eine Strafe I, 125., Erinnerung an die Sünde I, 100.; durch A. sollen wir das Fleisch kreuzigen II, 1241.; die Menschen fliehen sie II, 1789.; wer zur A. geschickt und willig, wird werth gehalten II, 1244.; ob alle Menschen Ackerbauer sein oder mit den Händen arbeiten sollen I, 258.; wir sollen arbeiten, das Andere Gott befehlen II, 1443. f. 1446.; A. soll den geistlichen Schlaf des Glaubens nicht hindern II, 1239.; in der Kirche heißt A.: lehren u.\,s.\,w. I, 261. — S. \gesperrt{Gebet} 12, \gesperrt{Hausarbeit}.

Arche. I, 512. ff. — Spitzfindige Fragen I, 516.; ob man Stücke von der A. gefunden I, 562.; warum Gott grade durch die A. Noah retten wollte I, 526. 541.; hatte die Proportionen des menschlichen Leibes I, 514.; allegorische Deutung derselben I, 514. f. 616.; päbstische Allegorien I, 612. — S. \gesperrt{Thiere} 2.

Aergerniß. Zank und A. findet sich in allen Ständen I, 984.; A., von Rechtgläubigen gegeben, ist schädlicher, als wenn von Verfolgern der Wahrheit gegeben II, 1090.; wie solch A. zu überwinden I, 636.; über bloß genommenes A. mache man sich kein Gewissen I, 793.

Argumente. Vom Nöthigen, Nützlichen, Wohlanständigen II, 352. 1565. — S. \gesperrt{Beweise}.

Argwohn. Ist Verleumdung I, 794.; schädlich I, 795.; soll nicht unter Eheleuten sein I, 800.; was die heidnischen Philosophen davon gehalten I, 795.; wie doch die Schrift ihn gebieten kann I, 795. f.; A. gegen offenbare Bösewichter ist nicht unrecht II, 1474. f. — S. \gesperrt{Eheleute} 4, \gesperrt{Freundschaft}.

Arianer. Lehrten, vor dem „Anfang“ seien die Engel und der Sohn erschaffen I, 13.

Arioch. König, wohl in Assyrien I, 878.

Aristoteles. Uebertrifft Cicero I, 681.; sein 5. Buch der Ethik ist das allerherrlichste II, 1876.; \gesperrt{sagt}: Wahrheit ist die Mutter der Liebe I, 795., das Amt entdeckt den Mann I, 834. A.’s \gesperrt{Meinungen} über: Gott I, 1551.; Materie I, 7.; Ewigkeit der Welt I, 5.; des Himmels Natur I, 32.; Ewigkeit der Sonne und des Menschen I, 151.; den ersten und letzten Menschen I, 150.; Ursprung des Menschen I, 681. ff.; Zeugung des Menschen I, 151.; Definition der Seele II, 1670. f.; Tugend, Definition II, 1877. 1880., sie zu üben sei der Mensch geschaffen I, 159.; epieikeia oder Billigkeit II, 1876. 1880.; Vernunft des Menschen, sie strebe nach dem Besten I, 577.; natürliche Ursachen des Regenbogens I, 605. f.; Entstehung der Mäuse I, 63.

Arius. Mußte in Irrthum gerathen, denn er glaubte dem Wort nicht I, 17.; lehrte, Christus sei geschaffen, doch vor den andern Creaturen I, 15. f.; hielt diese Lehre für Weisheit I, 196. f.; fand Beifall I, 1283. f.

Arme. Man soll ihnen helfen II, 84. f.

Armuth. Rechtes Verhalten I, 836. — S. \gesperrt{Abraham} 5.

Arnold. Hatte den Wahrsagergeist II, 1050.

Arphachsad. Wann geboren I, 713. f. — S. \gesperrt{Arbaces}.

Arzenei. Entdeckung neuer A. bei neuen Krankheiten geschieht durch der Engel Dienst I, 1246.; A. soll man gebrauchen I, 525.

Arzt. Junger, muß neuen Kirchhof haben I, 847.

Askenas. Held und frommer Mann I, 663.; welcher Völker Vater, Vorfahre der Deutschen I, 660.

Asnath. Josephs Weib, heilig, gottselig II, 1417.

Asser. 1. \gesperrt{Jakobs Sohn}. Bedeutung des Namens I, 564.; Jakobs Segen über ihn II, 2033. f. — S. \gesperrt{Jakob} I, 19.; II; \gesperrt{Joseph} 12.

2. \gesperrt{Der Stamm}. Wohnsitz II, 2018.; fruchtbares Land II, 2034.

Assur. 1. \gesperrt{Sems Sohn}. Warum von Babel ausgezogen I, 674. f. 679. 713.; baut Ninive I, 673. 713.

2. Sind die Assyrer I, 680.

Assyrien. In der Schrift gelobt I, 675. f.

Astharoth Karnaim. Königreich I, 880.

Astrologie. Hat in 1 Mos. 1, 14. keinen Grund I, 54. f.; ist keine Wissenschaft I, 54. f.; nicht gar zu verachten I, 38.; man mag sich wohl damit beschäftigen I, 54.; es ist Thorheit, sich um die Zukunft ängstigen I, 55.

Astronomie. Lobenswerth I, 35.; darin Egypter und Chaldäer berühmt I, 805.; Lehren der A. I, 33. ff. 49. f.

Astyanax. Dem Hector nicht gleich II, 1368.

Atlanten. Isaak, Jakob, Joseph waren A. II, 1155.

Auferstehung. Werden einst trotz Verwesung auferstehen I, 965. \gesperrt{Beweise} dafür aus der Natur I, 63.; II, 16. f. aus der Schrift A.\,T.’s I, 243. 348. f. 406. 420. 428. 1517. f. 1519. f. 1755. f.; II, 16. 141. 214. ff. 981. \gesperrt{Zeugniß} davon hatte die Welt zu allen Zeiten I, 426. f., auch im A.\,T. II, 1305. Diese Lehre ist eine Kunst der Heiligen, geht über die Vernunft I, 1518.; II, 16. 1962. f.; um der Verheißung willen zu glauben I, 1384. f.; ist tröstlich I, 1488. f. 1509. ff.; wer sie leugnet, glaubt gar nichts I, 1520. f. 1764.; II, 1305. — S. \gesperrt{Tod} 2, \gesperrt{Unsterblichkeit}. \gesperrt{Christi} A. s. \gesperrt{Christus} 2.

Aufrichtig sollen wir vor Gott wandeln II, 259.

Aufrührer. Christen sind nicht Aufrührer I, 1465.

Aufsicht. Bei dem Gesinde wichtiger als eigene Arbeit II, 1240.; nöthig im Regiment II, 1423. und Kirchenamt II, 1424.

Augen. Verschiedene Arten II, 2007.; schöne, liebliche A. eine seltene Gabe II, 491. f.; Falkenaugen besonders gelobt II, 492.; kleine, blöde A. entstellen II, 492.; rothe A. bekommen, die immer toll und voll sind II, 1986. f. 1995. f. A. bedeuten in der Schrift Regenten und Lehrer II, 2006.

Augustinus. Hatte \gesperrt{vor seiner Bekehrung} eine Concubine I, 1750.; belachte, als \gesperrt{Manichäer}, die Patriarchen als die geilsten Menschen II, 556. 570.; was ihm besonders an den Manichäern gefiel I, 451. \gesperrt{Nach seiner Bekehrung} hat er nicht in Ruhe gelebt und sich gepflegt I, 260.; hat an der Brunst des Fleisches gelitten II, 530., einmal angefangen, an der Glaubwürdigkeit der Schrift zu zweifeln II, 1023.; wollte keine Engelerscheinung I, 1527.; was er für ein Mönchsleben angerichtet I, 847.; II, 1879.; war weise und verständig II, 1879. Ist ein besserer \gesperrt{Schriftausleger} als Ambrosius II, 79.; verstand nicht: die active und passive Bedeutung der Worte II, 1488., noch die Figur hysteron proteron II, 1067.; allegorisirt I, 285. 610.; Beispiele von Allegorien I, 6. 612.; II, 774. 810. f.; sie sind fehlerhaft I, 612. 626.; wodurch vom Allegorisiren abgezogen II, 557. — \gesperrt{Lehren und Meinungen} über: 1 Mos. 46, 8. ff. II, 1762. 1764. 1 Mos. 47, 31. II, 1838. f.; Ps. 85, 11. II, 1810.; Abeloniten II, 1266.; Abenderkenntniß und Morgenerkenntniß I, 1535.; Abrahams Gastmahl I, 1382.; Allmacht Gottes I, 525.; Arche I, 514.; Dreieinigkeit I, 60. f. 1151.; Ehemann, der sein Weib liebt II, 168.; Fegfeuer II, 2068.; Gaben Gottes II, 1390.; Gefahr II, 764.; Gewissen, böses II, 1590.; Gleichheit: Non aequaliter omnes etc. II, 512.; Gnade Gottes I, 1346.; II, 1961. 2084. f.; Gott, was er vor dem Anfang gemacht I, 13.; Güte Gottes II, 685. 1700. 2083.; Hiob II, 1014. f.; hysteron proteron II, 1510. f.; Isaak II, 1067.; Jakob II, 490. f. 836.; Joseph II, 1065. 1130. f. 1483.; Kinder, von Herodes getödtet II, 1150.; Kirchenväter I, 1316. 1382.; Lot I, 1230.; Lügen I, 788.; Materie I, 10.; Monica II, 1308.; Name, guter II, 1303. f.; Potiphar II, 1414.; Regulus I, 581. f.; Rom II, 2025.; Sacrament I, 1045.; Scheol II, 1542.; Seele II, 1770.; die Siebenzig II, 1766.; Unsterblichkeit I, 53.; Unzucht und Versuchung dazu II, 1253.; Werke Gottes nach außen II, 1868.; Wille, freier I, 480.; Deus coronat intus voluntatem etc. II, 1115.; Wunderwerke II, 1245. f. — S. \gesperrt{Monica}.

Augustus. Weiser, glückseliger Kaiser II, 1409.; ihm wurde das Regiment schwer II, 1409.

Auserwählte. Ihnen dient Alles zum Besten I, 894.; um ihretwillen thut Gott Anderen Gutes II, 1234. f., wie Kains Beispiel lehrt I, 368. f. 370.; um ihretwillen erhielt er das jüdische Reich II, 2046. f.

Ausnahmen von der Regel muß man oft machen II, 511. f.

Aussätzige. Heftig zur Unzucht gereizt I, 75.

Averrois. Hält die Sphären für vernünftig I, 35.

\section{B}

Baal. Bedeutung des Worts II, 1610. f.; warum man die Götzen B. genannt I, 1322.; II, 1109.

Baal Hanan. Bedeutung des Namens II, 1017.; wird durch Umstellung „Hannibal“ II, 1017.

Baal Peor. Abgott I, 1665. — S. \gesperrt{Priapus}.

Baalsdiener. Nannten sich Bräute Gottes II, 1611.

Baalsdienst. Bei den Juden II, 1109.

Babel. Stadt und Reich I, 672. ff. 694. ff. 877.; wovon seine Fruchtbarkeit abhing II, 1701.

Babylon. S. Babel.

Balsam. Aller Harze König II, 1124.; verschiedene Sorten II, 1124. f.; wo er wächst I, 851.; Gebrauch II, 1582.

Bann. Eine Strafe der Kirche, wird verachtet II, 1922.; wen man bannen soll I, 1034. f.; nimmt Alles weg I, 1450.; Kain gebannt I, 347. f. 365.; Hagar absolviret I, 1448.; der Papisten B., ein Mißbrauch des Namens Gottes I, 1403.

Barfüßermönche. Kein unleidlicher und giftiger Ding ist auf Erden denn ein B. II, 1628.; reinigten den Kirchhof, wenn eine Frauensperson darüber gegangen war II, 164. f.; betrogen und sogen das Volk aus II, 1955.

Barmherzigkeit. 1. \gesperrt{Gottes}. Gott ist barmherzig gegen die, die auf seinen Wegen gehen I, 512. f.; gepriesen durch Abrahams Beruf I, 730.; man muß sie erfahren II, 1383.

2. \gesperrt{Der Menschen}. Welches die Werke der B. sind I, 1136.; wo sie aufhören muß I, 1208. — S. \gesperrt{Gastfreiheit}.

3. \gesperrt{B. und Treue}. Unterschied II, 757. f.

Basmath. Bedeutung des Namens II, 989. — S. \gesperrt{Esau} 3.

Bauernaufruhr. I, 823.

Baum. 1. \gesperrt{B. des Erkenntnisses} des Guten und Bösen I, 115. f. Stand wohl da, wo Christus gekreuzigt wurde I, 678. 851.; II, 428.; wäre ein Haupttempel gewesen I, 116. 129. 188.; daran sollte Adam seinen Gottesdienst und Gehorsam beweisen I, 114. f. 118. 189. 278.; woher der Name I, 114. 118.; an sich ein guter Baum I, 118.; warum er tödtete I, 280.

2. \gesperrt{B. des Lebens} I, 112. ff. Stand wohl, wo später Gethsemane war II, 1554.; hatte die Kraft, die Jugendfrische zu erhalten, durch Gottes Wort I, 113. f. 277. ff. 280.; ob Adam und Eva den Tod überwunden hätten, wenn sie nach dem Fall davon gegessen hätten I, 277. f.; selbst dann hätte sie Satan wieder versuchen können I, 281.; warum Adam nach dem Fall nicht davon essen durfte I, 280. f.

Bayern. Galten für närrisch II, 851.

Bazra. Name zweier Städte in Idumäa II, 1014.

Beamte. Ausschreitungen derselben II, 1782.

Beelzebub. Urheber des Falls der Engel I, 138.

Begräbniß. Abraham kaufte ein Begräbniß I, 1616. ff.; Sinn christlicher Gebräuche beim B. II, 1760. — S. \gesperrt{Leichenbegängniß}, \gesperrt{Todtenbestattung}.

Begrüßungen. Dabei wünscht man Glück II, 1604.

Beichte. Vornehmste Kraft derselben besteht im Trost II, 1003.; die heimliche B. soll man behalten II, 1002.

Beispiel. Ist nicht alsbald eine Regel I, 1233. f.

Bekehrung. Geschieht nicht ohne Schrecken über die Sünde I, 1324.; die zuvor sonderlich böse waren, werden trefflich fromme Leute II, 1776.

Bekenntniß. 1. \gesperrt{B. der Sünde}. Fällt uns schwer II, 1514. ff., 1518. f., 1526.; ohne Hoffnung auf Vergebung will der Mensch nicht bekennen I, 217. f., 221. f.; ist keine Schande II, 871.; Gott will es I, 346.; nöthig, um Vergebung zu erlangen II, 1515. f., 1518.; überwindet Gottes Zorn I, 336. f., 339. — S. \gesperrt{Glaube} 8.

2. \gesperrt{B. des Glaubens}. Man soll Christum und sein Wort bekennen I, 1553.; darauf folgt Verfolgung II, 1297. f. — S. \gesperrt{Geduld} 3, \gesperrt{Glaube} 8.

Bekümmerniß. „Gott bekümmerte sich“ I, 492. f. — S. \gesperrt{Betrübniß}, \gesperrt{Menschenworte}.

Benjamin. 1. \gesperrt{Jakobs Sohn} II, 550. ff. Name aus Benoni in B. verwandelt II, 969. f.; warum er der Sohn der Rechten heißt II, 969. f.; Jakobs Segen über ihn II, 2053. f. und Erfüllung desselben ib. — S. \gesperrt{Jakob} I, 19.; II.; \gesperrt{Joseph} 12.

2. \gesperrt{Der Stamm}. Wohnsitz II, 2054.

Benno. Meißnischer Abgott I, 1531.

\quellseite{544}
Bernhardus. Verstand die Religion wohl, steht über Franciscus und Anderen I, 759.; sonderliche Freude an der Menschwerdung Christi II, 400.; marterte seinen Leib I, 1711.; II, 211. f.; trank in tiefen Gedanken Oel für Wein II, 291. Seine \gesperrt{Auslegekunst} II, 1844. \gesperrt{Schrift}: Brief an Pabst Eugenius I, 1470. — \gesperrt{Meinungen und Aussprüche} über: das Gebet I, 1704.; Marias Glauben II, 417.; des Teufels Fall II, 400.; Undankbarkeit II, 1262. 1315.; bei Weibern wohnen und sie nicht erkennen, ist ein eben so großes Wunder, wie Todte auferwecken II, 1264. 1266.

Bernstein. Als Arznei gebrauchtes Harz II, 1124. f.

Bersaba. Lage I, 1423.; II, 199.; woher der Name II, 1352. f.; Viehzucht daselbst I, 1423. f.; zu verwundern, daß man da keine Abgötterei trieb I, 1478.

Beruf. I. 1. Wie man den Beruf \gesperrt{erhält} I, 954.

2. \gesperrt{Heiligkeit} des B.’s. Heilig ist der Beruf, der Gottes Wort hat I, 567.; Jeder soll sich freuen, daß er einen göttlichen B. hat I, 1310.; wir können darin des göttlichen Wortes gewiß sein I, 1500.

3. \gesperrt{Verschiedenheit} des B.’s ist nöthig I, 1073. f.

4. \gesperrt{Verhalten} im B. Die Wenigsten sind mit ihrem Beruf zufrieden I, 1071. f.; man suche nicht sonderliche Werke I, 1525.; folge nicht den Werken, sondern dem Glauben der heiligen Väter nach I, 568. ff.; sehe darauf, daß man auch leibliche Verheißungen darin habe II, 1732.; bleibe beim Sittengesetz und dem gemeinen Recht II, 514.; bei Gefahr, die man in jedem Beruf erwarten muß, rufen Christen Gott an und werden errettet I, 1320.; man rufe Gott um Segen und Verstand an II, 277.; handle klüglich I, 1638. f.; sei Gott gehorsam I, 1073.; warte des B.’s fleißig I, 1500.; II, 460. f., 467. f., 1626. f., wenn auch Thorheiten nicht zu vermeiden sind II, 276., und befehle das Uebrige Gott II, 468. 748. 1302. f.; verzage nicht, wenn auch Sünden vorkommen und nicht Alles nach Wunsch geht II, 1567.; Jeder thue das Seine in Demuth, um Gottes willen II, 1344. 1391. f., frage nach den Anderen nichts I, 1071. 1073. f. 1270. 1281. 1525. Trotz aller Undankbarkeit diene man dem Andern II, 1237. 1281. 1315.; trotz Haß lehre und strafe man in seinem B. II, 1040.; bei Kreuz tröste man sich mit Gottes Wort I, 1502. f.; II, 139. f. 753. Wir sollen des B.’s warten, ob wir wohl auf Gott vertrauen sollen II, 1446. Gläubig in seinem B. einhergehen ist Gottseligkeit I, 867.; des B.’s warten ist rechter Gottesdienst I, 1058. 1071. f. 1478.; Christ soll sich der weltlichen Dinge nicht entschlagen I, 1477. f.; bleibt in gutem Stande, auch wenn er seines B.’s wartet I, 858.; des B.’s warten und sich um Andere nicht kümmern ist große Weisheit I, 1071.

5. \gesperrt{Werke} des B.’s. Uebertreffen alle Werke der Mönche und Nonnen I, 858. f. 1320. f.; sind mit Gottes Wort geziert I, 867.; Gott fordert nach dem mannigfaltigen B. auch mancherlei Werke I, 568. f.; sind die rechten guten Werke I, 1181. Alle im Glauben gethanen sind gut I, 858. f. 866., ein Gottesdienst I, 1164., gefallen Gott I, 1309.; II, 459. f.; selbst die gemeinsten I, 1310. f.; II, 468. ff. 472. f. 474. f. Die Papisten verachten die Werke des B.’s I, 1180. 1302. 1367. f., ebenso die Welt, die etwas Sonderliches will I, 1300. ff. 1305. ff.

6. \gesperrt{Verlassen und Wechsel} des B.’s. Jeder bleibe in seinem Beruf I, 262. 1075. 1281.; wir sollen den B. nicht eigenwillig verlassen I, 993. 996.; Mönche und Nonnen verlassen den B. I, 1071.

S. \gesperrt{Amt}, \gesperrt{Stände}.

II. \gesperrt{B. der Kinder Gottes}, der Juden im A.\,T. I, 1075., der Christen I, 1075.

III. \gesperrt{B. zum Lehren}. Als Gottes B. anzunehmen I, 763.

IV. \gesperrt{B. und Erwählung} soll man fest machen I, 1123. 1535.

Bescheidenheit. Bescheidene muß man schier zwingen, Wohlthaten anzunehmen I, 1216.; man soll B. üben I, 1216.; B. übten die Engel I, 1215. f.

Beschneidung. 1. \gesperrt{Im Allgemeinen} I, 1006. ff. 1036. ff. 1067. ff. Es ist wichtig, von der Beschneidung zu handeln I, 1009. f.; ein Sacrament I, 1019. 1021. 1049. ff.; ein Leiden, nicht ein Werk I, 1038.

2. \gesperrt{Wer} beschnitten werden sollte. Nur die \gesperrt{Juden} I, 1014. f. 1022. 1028. f.; warum I, 1014. f. 1026. f.; so sollten die Heiden zum Glauben gereizt werden I, 1028. f. 1031. 1043. f. 1055.; nicht alle Völker I, 1012. ff. 1022. ff. 1037. 1052. 1055. 1075.; II, 228. f. 1351. f. 1367.; Proselyten hätten nicht sollen beschnitten werden II, 229.; wer sie den Heiden aufnöthigen will, sündigt I, 1024.; Christen frei von der B. I, 1027.; nur das \gesperrt{männliche Geschlecht} unter den Juden I, 1041.; warum I, 1077. 1081. f.; gläubige Weiber ohne B. selig I, 1041. 1078. — S. \gesperrt{Joseph} 15.

3. \gesperrt{Glied}, an dem die B. vollzogen werden sollte. Warum an der Vorhaut I, 1040. 1080. f. 1658.

4. \gesperrt{Zeit}, da die B. vollzogen werden sollte. Warum am 8. Tage I, 1082. f. 1086. f.; geistliche Deutung des 8. Tages I, 1087. f.; Kinder, die vorher starben, nicht zu verdammen I, 1040. f. 1078. 1083. f. 1086. 1090. f.

5. \gesperrt{Wie lange} die B. währen sollte. Nur auf eine bestimmte Zeit eingesetzt I, 1025. f. 1068.; nur so lange Abrahams Same ein Volk wäre I, 1056. f. 1100.; bis auf Christum I, 1014. f. 1024. f. 1037. 1043. 1050.; bis auf Einsetzung der Taufe I, 1015.; jetzt aufgehoben I, 1069. ff. durch Christum I, 1092. 1129. f.

6. \gesperrt{Zweck, Nutzen und Wirkungen} der B. Die B. diente: zur Absonderung Israels von andern Völkern I, 1044.; dazu, daß man wüßte, woher der Messias kommen sollte II, 229. f.; zu einem Kennzeichen der Kirche I, 1045. 1048. 1068.; zur Uebung des Glaubens und Gehorsams I, 1012. 1020. f. War ein Zeichen, das Sünde und Gnade darstellte I, 1081. f.; ein Zeichen der Gerechtigkeit I, 1012. 1019. 1038.; ein Zeichen für Abrahams Nachkommen, daß sie Erben der Verheißung sein sollten I, 1020. f.; das bloße Werk der Beschneidung machte nicht gerecht I, 1010. ff. 1015. 1017. f. 1019. 1037. 1039., nicht zu Gottes Volk I, 1012. 1015., nicht selig I, 1049.; sie reichte die Gerechtigkeit dar I, 1038. f., machte die Gläubigen gerecht I, 1021. 1050. und selig I, 1044. 1116.; die Knäblein kamen durch sie zum Glauben I, 1039. 1041. f., wurden selig I, 1039. 1049.

7. \gesperrt{Nothwendigkeit der} B. \gesperrt{zur Seligkeit}. Im \gesperrt{Alten Testament} war die Strafe der Verachtung: geistliche Ausrottung I, 1090., Verdammniß I, 1041.; siehe oben 4. Im \gesperrt{Neuen Testament}: Juden behaupten die Nothwendigkeit der B. I, 1018. ff.; noch keine christliche Secte hat sie als nothwendig gefordert I, 1008. f. außer die Sabbather II, 1351.; ohne gute Erkenntniß glaubt man solche Nothwendigkeit leicht I, 1010.; sie ließe sich mit scheinbaren Gründen beweisen I, 1009.; gegen den Irrthum von der Nothwendigkeit mußten die Apostel kämpfen I, 1009. 1016. f.; B. ist nicht nothwendig I, 1010. 1022.

8. \gesperrt{Geistliche} B. ist Tödtung des Fleisches I, 1129. f. im ewigen Leben erfüllt I, 1088. — S. \gesperrt{Kreuzigung}.

Beständigkeit. Wir sollen nicht nach einem andern Weg oder Evangelium sehen I, 1282.; B. im Glauben, nothwendig und schwer I, 559.

Beste. Daraus wird oft das Aergste II, 996.

Bestürzung. In der B. kommt man von sich selbst I, 1450. f., sieht und fühlt man nicht II, 500. 1122. 1193., begehen Gläubige oft Sünden I, 1292. f. — S. \gesperrt{Erregung}, \gesperrt{Gedanken}.

Besuchen. Anständig, gebräuchlich II, 844. f. 853.

Beter. Muß glauben I, 1365.; man soll ihn ehren I, 1363.; Gott erhört ihn I, 1542. f. — S. \gesperrt{Gebet}.

Bethel. Lage I, 780.; II, 425. ff. 903.; Bedeutung des Namens I, 781.; II, 426. f.; hieß früher Lus I, 781.; berühmter Ort II, 958.; viel Götzendienst daselbst getrieben I, 779. 781.; II, 425. f. 903. f. 958. f.; ob das B., da Jakob die Himmelsleiter sah, Jerusalem war II, 427. ff. 438.

Bethlehem. Bedeutung des Namens II, 971.; Lage II, 1856. Den Kindern zu B. konnte Herodes keinen größeren Nutzen bereiten II, 1150.

Bethuel. Jüngerer Bruder Labans I, 1720.

Betrogen. Wie wollen wir es so gut haben, wenn wir Christo glauben und von ihm betrogen sind II, 1549.

Betrübniß. Gläubige werden betrübt, wenn Gott mit seinem Gerichte verzieht I, 506. 508.; wenn sie Gottes Zorn über die Welt um der Sünden willen fühlen I, 462. f., welche B. der Heilige Geist wirkt I, 493.; in den Heiligen wird Gott betrübt I, 454. f. — S. \gesperrt{Bekümmerniß}.

Betrübte. Satan mehrt ihre Betrübniß II, 392. f.; schlägt sie gar darnieder I, 1425.; nur die über ihre Sünden Betrübten geht der Segen Abrahams an I, 1569. — S. \gesperrt{Anfechtungen} 2.

Betrug und List. Wann erlaubt II, 265. 311. f.

Betrügen. Wie man es anfangen muß, Einen zu betrügen II, 274.; betrügen kann man Leichtgläubige leicht, Ungläubige nur schwer II, 273. f.

Bettler. Der Teufel hat auch seine B. I, 1136.; welchen B. man nichts geben soll I, 1140.; II, 1602.; welches die rechten B. sind, denen man helfen soll I, 1136.; die rechten B. sind bescheiden II, 1602., die falschen nehmen, was die rechten haben sollten II, 1602.

Beweise. Nicht rhetorische, sondern dialektische soll man brauchen I, 1175., sonderlich in theologischen Sachen I, 1177.; B. a posteriori können auch etwas gelten I, 570. — S. \gesperrt{Argumente}, \gesperrt{Fallacia}.

Bias. Meinung über: Freundschaft I, 795. f., Regieramt II, 1405.

Bibel. Was die B. im Stande der Unschuld gewesen wäre I, 128.; Abrahams B. I, 1684., Josephs II, 1823. — S. \gesperrt{Schrift}.

Bildersprache. In der alten Welt gebraucht I, 1476.

Bileam. Wer B. war I, 1601. ff.; II, 383. f. 1015.; großer Mann und Prophet I, 1601.; II, 385.; tief gefallen I, 1601. ff.; II, 1016.

Bilha. S. \gesperrt{Jakob} I, 18, \gesperrt{Ruben} 2.

Billigkeit fordert, daß wir Andere nicht zwingen, unserm Beispiel zu folgen I, 919.; ziemt sich bei Handhabung der Gesetze II, 1876. ff.

Bischof. 1. Was sein \gesperrt{Amt} sei I, 628.

2. „\gesperrt{B. in unserer Nachbarschaft}“. Ließ M. Georg ermorden II, 1108.

3. \gesperrt{B. von Mainz}. Wollte für heilig und gerecht gelten II, 708.; frug nichts nach Christi Ehre und der göttl. Wahrheit II, 671. f.; ging mit List um II, 1099.

Bischofsstab. Ursprung II, 1843.

Bitumen. Was es sei I, 513.

Blindheit der Sodomiter I, 1236. f.; andere Beispiele solcher Blindheit I, 1227.; II, 1471.

Blut. Blutessen im A.\,T. verboten I, 595. f.

Blutschande. Gott ist der B. feind II, 1935.; B. mit des Vaters Weibe, auch den Heiden ein Greuel II, 978.

Blutschänder. Ruben ein warnendes Exempel II, 1922. ff.

Bogen. Früher im Kriege gebraucht I, 1454.

Bonaventura. Hat die Menschwerdung Christi lieb gehabt II, 400., sich an dem ehelichen Leben der hl. Väter geärgert II, 556.; Meinung über den Ehestand II, 166.

Brandopfer. Ganz angezündet I, 571.

Braut. Jemanden die B. stehlen, ist die größte Gewalt und Schmach II, 495. ff.

Brautgeschenke. Nützlich I, 1683.; nicht zu verwerfen I, 1738. — S. \gesperrt{Geschenke}.

Brautliebe. Ein Geschöpf Gottes II, 477.; in der Schrift nicht verdammt II, 493.; ungeduldig II, 490. f.; wann am heftigsten II, 493.; zart II, 496.; Jakobs B. regierte der Heilige Geist II, 491.

Brautwerber. Ist recht, wenn B. den Bräutigam loben I, 1728. f.; Elieser wirbt um Rebekka I, 1722. ff.

Brechen. „Brechen“ heißt „Speise“ II, 1486.; „Brod brechen“, II, 1439. f. 1489. 1509. 1532.

Breitenbach. Ueber die Araber und Troglodyten I, 997.

Bria. Woher der Name II, 1369.

Brief und Siegel. I, 1638. — S. \gesperrt{Bund}.

Brigitta. Von ihren Offenbarungen II, 559.

Brod. B. heißt oft Speise und Trank II, 1613.; „Brod auflegen“ II, 1613. \gesperrt{Ungesäuertes} B. essen stand bis zu Moses Zeit frei I, 1221. f.; zu Lots Zeit I, 1220. f.; Moses Gesetz gilt nicht mehr I, 1221. — S. \gesperrt{Brechen}.

Bruder. Bezeichnet oft: Bruderskind II, 307. 478. f., Verwandte II, 695. f. 714.

Brügge in Flandern. Hat keinen Hafen mehr I, 595.

Brunnen. Unterschied zwischen puteus und fons II, 203.; B. \gesperrt{der großen Tiefe}, wie sie bei der Sintfluth aufgebrochen I, 546.; B. \gesperrt{des Lebendigen und Sehenden}, woher der Name I, 1005. f.; II, 9.

Bücher und Commentare haben die Bibel in den Hintergrund gedrängt I, 1299. f.; was die Kirchenväter gethan, damit der B. nicht zu viele würden I, 1289.; welche man allein lesen soll I, 1290.

Buhlteufel. Was davon zu halten I, 447. f.

Bund. B. und Verträge werden treulos gebrochen I, 1473.; so der B. zwischen Maximilian, Ludwig und dem Pabst I, 967.; Gebräuche zu Abrahams Zeit I, 1472. ff.

Burgensis. Vom Judenthum zum Christenthum bekehrt I, 1566.; redet von: einem Gesetz der Juden, daß sie die Heiden umbringen dürfen II, 889., Jakobs Liebe zu Joseph II, 1034. ff.

Bus. Wohnsitz I, 1601.

Buße. Was es heißt, Buße thun II, 316. 663.; \gesperrt{besteht} nicht in einem äußerlichen Werke, als, Fasten II, 1108., auch nicht in bloßer Reue II, 1495., sondern in Reue und Glauben II, 1457. f. 1467. 2089.; dabei sucht man Gottes Gnade II, 314.; Unterschied zwischen rechter B. und Galgenreue II, 313. f. Aufrichtige B. findet allezeit Raum bei Gott II, 313. f.; tilgt die Sünde II, 866. 872.; ist die einzige Errettung vom Zorn Gottes II, 1632. f.; man soll B. nicht aufschieben II, 116.; die \gesperrt{Papisten} wissen nicht, was B. ist II, 1464. — S. \gesperrt{Esau} 1., \gesperrt{Galgenreue}, \gesperrt{Heuchler}.

Bußfertige. S. \gesperrt{Gesetz} 2.

\quellseite{545}
\section{C}

Caiphas. Vom Stamme Simeon II, 1117.

Calah. Ein Theil Ninives I, 674.

Camarim. Jüdische Secte II, 1610.

Canaan. 1. \gesperrt{Hams Sohn}. Bedeutung des Namens I, 665. f.; berühmtester Sohn Hams I, 665. C.’s \gesperrt{Nachkommen} I, 677. f.; hatten trotz Noahs Fluch Glück I, 677. f.; endlich ausgetrieben und ausgerottet I, 678.

2. \gesperrt{Das Land}. Ein schönes Land I, 678.; reich an Würze II, 1124. f.; für Sem bestimmt, von Nimrod eingenommen I, 776.; auf wie lange Abrahams Samen verheißen I, 872. f.; seine vornehmste Ehre war, daß an dem leiblichen Orte das ewige Leben, Sieg über Sünde und Tod war II, 1907.; jetzt unfruchtbar I, 121.

Cananiter. Bedeutung des Namens II, 1162.; wohl etliche C. fromm zur Zeit Abrahams I, 787. 837. und Jakobs II, 1164.

Canonisten. Haben die Welt mit Büchern über Menschensatzungen erfüllt, dachten nicht an Auslegung der Zehen Gebote II, 1896.; setzen ihre Canones über die Schrift I, 1642.; Definition von Ehestand I, 1645.; Meinung über: Simonie I, 1618. ff., heimliche Verlöbnisse I, 1641. ff.

Caphthorim. I, 676. f.

Carl V. Versäumte die Gelegenheit wider die Türken II, 98.; Gott nahm ihm den Muth II, 932.; hatte, was er hatte, um der Kirche willen I, 1409.

Carlstadt. Nach wem genannt I, 1610.; ein Schwarmgeist II, 1385.; machte Doctoren der Theologie, obwohl er es für Sünde hielt I, 697.

Carmel. Fruchtbarer Berg II, 2034.

Carthäuser. Sind Baaliten II, 1611.; halten streng auf ihre Regeln: kein Fleisch zu essen I, 845. 1316., nicht zu schlafen I, 845.; brachten die Sodomie nach Deutschland I, 1222.

Cäsar, Julius. Ein ungläubiger Heide II, 1305. f.; inwiefern kein jüdischer oder christlicher König mit ihm zu vergleichen II, 462.; mehrere seiner Träume II, 1310. f.; Meinung über einen schnellen Tod II, 1305.

Casluhim. I, 676.

Catachresis. Redefigur II, 1014. 1844.

Cato. Uebertraf an Rath Cicero, war aber nicht so beredt I, 998.; Ausspruch über Träume II, 420. 1051.

Cerberus. Was die Poeten von ihm fabeln II, 701.

Ceremonialgesetz. Sollte nur bis auf Christum bestehen I, 1026.; jetzt aufgehoben I, 873. f.; II, 925. ff.; man muß darin Christum suchen II, 1179.

Ceremonien. Nützlich und nöthig im Gottesdienst II, 6. f. 293. 914.; machen nicht gerecht II, 914. — S. \gesperrt{Gottesdienst} 3.

Chaldäer. Beteten das Feuer an I, 1665.; in der Astronomie berühmt I, 805.

Cherubim. Engel mit junger lieblicher Gestalt I, 288.

Chesib. Stadt, sonst Achsib genannt II, 1163.

Christ. \gesperrt{Wer} ein Christ sei I, 1239.; II, 1997. 2063.; Christen \gesperrt{sind} Brüder, darum geht sie auch die Gefahr ihrer Mitchristen an I, 337.; haben denselben Geist, wenn auch nicht in gleichem Maße I, 1367.; sind Abrahams Same aus Juden und Heiden I, 1101.; Gesegnete I, 1567.; II, 142.; Kinder des Segens, passivisch und activisch II, 1892.; Pilgrime I, 427. 1304. f. 1617. 1626.; II, 1801. f.; Kriegsleute II, 408. \gesperrt{Kennzeichen} eines C. II, 464. \gesperrt{Herrlichkeit}: ein C. sein ist große Herrlichkeit II, 1611. f.; haben großen Vorzug vor Abraham I, 863. 868. f.; mehr Ursache zur Freude als er I, 1105.; ihr Ruhm II, 66.; sind ein Schauspiel Gottes und der Engel I, 1519. 1522.; II, 1098. f.; an ihnen nichts gering in Gottes Augen II, 471.; Gott liebt sie I, 1342.; an ihnen gefällt Gott Alles II, 479. 617. f.; sind alle Propheten I, 1367. und Gottes Volk I, 963. 1401.; selige Leute I, 862. f.; heilige Esel, trunken von dem Heiligen Geist II, 1985. ff.; können rühmen und trotzen wider Welt und Teufel I, 1589.; Alles muß ihnen zum Besten dienen II, 848., auch ihre Fehler II, 685. 705.; sind Ueberwinder Gottes II, 796. 811. f. 952.; zu allen guten Werken geschickt I, 1581.; gehorchen Gott freiwillig, furchtlos II, 703.; Herren des ewigen Lebens II, 358. f. \gesperrt{Leben}: ihr Beruf ist glauben, was unsichtbar ist II, 1917. f. 1920.; geistlich trunken sein II, 1983. 1989.; sich nicht fürchten II, 836.; glauben, sich taufen lassen, irdischen Berufs warten I, 1075.; ihr Leben besteht in Glauben und Anrufung I, 1321.; II, 706.; sollen Israeliten und Jakobiten sein II, 957.; sollen kämpfen II, 408., heilsame Zungen haben II, 349. ff., in einander Gott ehren I, 1144. f.; lassen Alles über sich ergehen um des Bruders willen I, 1147.; werden nicht immer vom Heiligen Geist getrieben, der sie in manchen Dingen ihre Vernunft gebrauchen läßt I, 1308. f.; sehen allein auf Gottes Ehre und der Kirche Nutzen I, 695.; thun alles, weil sie ein ander Leben suchen II, 1237. Gesetze machen keine Christen, sondern das \gesperrt{Wort} II, 2001. — S. \gesperrt{Fleisch} 2, \gesperrt{Fromme}, \gesperrt{Gläubige}, \gesperrt{Gottselige}, \gesperrt{Heilige} 1, \gesperrt{Kinder} 2, 4, \gesperrt{Leben} 9, \gesperrt{Obrigkeit} 13, \gesperrt{Prophet} 1, \gesperrt{Vereinigung}, \gesperrt{Welt} 3, \gesperrt{Wiedergeborene}.

Christophorus. Schutzheiliger der Bauern II, 1919.

Christus. 1. \gesperrt{Person und Namen Christi}. Ist der dem Abraham verheißene Segen I, 754., Abrahams Same I, 1569.; ist und mußte sein wahrer Gott und wahrer Mensch I, 1570.; II, 404. und zwar ohne Sünde I, 1570. f.; ist wichtig und gewiß, daß er wahrer Mensch ist II, 409. f.; seine Menschheit war ein Köder, dadurch Satan betrogen ward II, 312.; die Vereinigung der beiden Naturen ist auch den Engeln ein Geheimniß II, 396. f., dessen wir uns in der Ewigkeit freuen werden II, 397. f. 400.; heißt „die Furcht Isaaks“ II, 703. — S. \gesperrt{Gemeinschaft}, \gesperrt{Menschwerdung}, \gesperrt{Mittheilung}.

2. \gesperrt{Stände Christi}. \gesperrt{Empfängniß und Geburt}: ist von einer Jungfrau geboren I, 237., ohne Zuthun eines Mannes I, 1167. 1570. f., die ihn bis in den 10ten Monden getragen hat I, 1167. Ist aus sündlichem Fleisch geboren, das der Heilige Geist im Augenblick der Empfängniß gereinigt hat II, 1170. ff. 1195. 1201. f.; wollte unter seinen Vorfahren große Sünder haben I, 733.; hat heidnisches und jüdisches Geblüt II, 1172. f. 1177. 1200. f. 1418. — \gesperrt{Leben und Leiden}: warum er \gesperrt{getauft} worden I, 1021. 1024. Hatte Eigenthum I, 835. f. Christi \gesperrt{Leiden} mit dem Leiden der Patriarchen verglichen II, 345. 1589.; nicht so faßlich wie der Patriarchen Leiden II, 1078.; wir könnens nicht fassen II, 345. 1553.; Leiden am Kreuz nicht so groß wie im Garten II, 1555.; litt im Garten die Schmerzen der Hölle II, 1554. f.; war damit angefochten, daß er Gottes Gnade nicht fühlte I, 927. — S. \gesperrt{Kreuz} 10, \gesperrt{Rock}. — \gesperrt{Tod}: C. ist von den Stämmen Simeon und Levi getödtet II, 1117. f. 1936. f.; sein Tod ein wirklicher, ob C. wohl nicht verwest ist I, 1609.; Unterschied zwischen seinem Tod und dem Tode anderer Heiligen I, 1609.; sein Tod ist ein Sacrament, das uns versichert, daß unser Tod nichts sei, und macht uns lebendig I, 1607. ff. — S. \gesperrt{Grab}. — \gesperrt{Höllenfahrt}: wie es damit zugegangen II, 1544.; was er in der Hölle gethan II, 1555.; was Petrus lehrt I, 535. f. — \gesperrt{Auferstehung}: wer mit C. auferstanden I, 1606.; II, 2064. f. 2070. 2090.; dadurch sollte C. Auferstehung geehrt werden II, 2070. — \gesperrt{Himmelfahrt} C.: wichtiger als die Henochs oder Elias I, 427.; beweist, daß der Schlange Kopf zertreten I, 427.

3. \gesperrt{Werk und Amt Christi}: C. war von Ewigkeit zu unserem \gesperrt{Heilande} bestimmt I, 1541.; ists auch nach der menschlichen Natur I, 280.; er ists allein, nicht auch Maria I, 235.; er war nöthig, die menschliche Natur konnte sich selbst nicht erlösen I, 241.; außer ihm ist Alles verflucht I, 1573., kein Heil und Segen I, 1591. f.; in ihm allein Gnade und Gerechtigkeit I, 1025. Die \gesperrt{Erlösung} ist eine wunderbare Gotteswohlthat I, 749.; wovon uns Christus erlöst I, 240. f.; inniger Zusammenhang der Erlösung von Teufel, Tod, Sünde, Gesetz I, 240. f.; C. ist des Gesetzes Ende I, 1010. 1024. f.; II, 2025.; was er für gute Werke gethan I, 1367.; ist der wahre Weg zur Gerechtigkeit I, 760.; macht durch seine Keuschheit uns keusch und selig I, 1661.; verwandelt Strafe der Sünde in seliges Kreuz I, 244. f. Christus ist ein \gesperrt{König} I, 1558. 1585. f. 1589., hat dem Teufel die Macht genommen I, 233. f., hat immer etwas Gutes mit uns vor II, 1490.; ein \gesperrt{Priester} I, 1558. 1585. f. 1589. ewiglich I, 901. 1589. f. nach der Ordnung Melchisedeks I, 904. f., sein Opfer ist vollkommen I, 904.; hat den Segen über alle Völker ausgebreitet I, 748.; kommt aus Abrahams Lenden, heiligt den Ehestand I, 1659. f.; ein \gesperrt{Prophet}: der große Prophet II, 2060.; hat die Beschneidung und das ganze Gesetz aufgehoben I, 1092., eine neue Lehre und Ceremonien gebracht I, 1026. C. will die Heiden nicht verwerfen II, 1177.

4. \gesperrt{Christus im A.\,T.} Die \gesperrt{Verheißung} von C. wird immer klarer gemacht I, 519.; Propheten haben aus Mose geschöpft I, 749.; C. der Edelstein auch in den leiblichen Verheißungen II, 704.; warum Canaan den Vätern verheißen war I, 719.; II, 141.; Verheißung von Christo von Abraham auf David geleitet I, 1586. ff. \gesperrt{Vorbilder}: Adam I, 268.; Melchisedek I, 900. ff.; Abrahams Widder I, 1541.; Joseph s. \gesperrt{Joseph} 2.; C. muß man in den äußerlichen Ceremonien des A.\,T.’s suchen II, 1179. Ist den Vätern oft erschienen II, 800.

5. \gesperrt{Erkenntniß und Verachtung Christi}. Wo C. nicht erkannt wird, da ist Uneinigkeit und Blindheit I, 692., da werden aus Sünden Tugenden I, 1671; Verachtung Christi kommt v. Teufel, Welt u. Fleisch I, 432.

6. \gesperrt{Christus und die heil. Schrift}. Um seinetwillen ist dieselbe gegeben II, 1177.; ihn sollen wir darin suchen I, 1047.; C. fast an allen Orten im A.\,T. unter dem Namen Gottes offenbart II, 853.

Vgl. \gesperrt{Abrahams Same}, \gesperrt{Messias}, \gesperrt{Myrrhe}, \gesperrt{Schiloh}, \gesperrt{Schooß}, \gesperrt{Weibessame}, \gesperrt{Wort} 1.

Churfürsten zu Sachsen. Nahmen sich der Kirche und Schule an I, 1357. 1632.

Chus. Hams Sohn I, 665. f.; Vater der Aethiopier I, 665.; sein Wohnsitz I, 666.

Cicero. Ist Aristoteles nicht gleich I, 681. \gesperrt{Aussprüche}: O me falso etc. II, 1405.; Meine Thränen hindern etc. II, 1652. \gesperrt{Meinungen} über: Begabte II, 1726., Freundschaft I, 795., Glück der Gottlosen, Unglück der Frommen I, 808., Gott I, 683. 815., die Seele I, 683., das Unbegrenzte I, 682.

Clemens, Pabst. Soll nicht getauft gewesen sein I, 1091.; scheußliches Ungeheuer II, 1660.; beraubte den Tempel zu Toledo II, 587.

Cochläus. Lästerte Luthers Lehre I, 1151.; II, 1273. f.

Cölibat. S. \gesperrt{Ehelosigkeit}.

Conversen. Wer so genannt II, 554.; Conversus Dominus II, 554.

Crates. Warf Geld weg, lebte vom Bettel I, 830. 856.; besiegte scheinbar den Geldgeiz, unterlag dem Ehrgeiz I, 856.

Creaturen. Sind nur Worte in Gottes Grammatik I, 60.; in ihnen Merkmale der Gottheit I, 1613.; Wesen, Gestalt und Güte der Creaturen I, 1613. f.; ihren Nutzen erkennen nur die Frommen I, 1613. f.; man soll sie gebrauchen I, 832. f. 836.; Mißbrauch derselben beweist die Bosheit des menschlichen Herzens I, 855. f.

Creutziger, Dr. Caspar I, 1.

Cybele. Heidnische Göttin II, 1222.

Cyklopen. Fabeln der Poeten I, 475. 596.

Cyprian. Sein Traum von Christo II, 1308.; braucht „Genugthuung“ unvorsichtig II, 1455. \gesperrt{Aussprüche} über: Märtyrer, durch Unzucht versucht II, 1259.; von Kirchendienern, die bei nackten Weibern zu schlafen pflegten II, 1266. f.

Cyrus. Erkannte den wahren Gott I, 1023.; II, 1773.; richtet den Gottesdienst wieder an I, 585.

\section{D}

Dan. 1. \gesperrt{Jakobs Sohn}. II, 544. ff.; Bedeutung des Namens II, 545. 2024.; Segen Jakobs II, 2021. ff., in Simson erfüllt ib. — S. \gesperrt{Jakob} I, 19.; II; \gesperrt{Joseph} 12.

2. \gesperrt{Stamm}. Wohnsitz I, 891. 2021. f. — S. \gesperrt{Antichrist}.

3. \gesperrt{Stadt}. Hieß früher Lais II, 2022.

Daniel. Schöpfte aus Mose II, 430.; ein Beter I, 1705.; bittet für Jerusalem und den Tempel II, 1652. f.; bekehrte Viele II, 1773. — S. \gesperrt{Propheten} 2.

Danken soll man Gott für Großes und Kleines I, 1718. 1738., für leibliche Güter II, 296. f.

Dankopfer. Sind alle Werke und Leiden der Christen II, 446. f.; dazu soll uns Gottes Gnade treiben II, 451.

Darius. Erkannte den wahren Gott I, 1023.

David. \gesperrt{Glaube}: Er erkannte die Wohlthat des Wortes I, 862. f., Christus werde König und Priester sein I, 1558. f. 1585. f. 1589., Christi Priesterthum werde ewig währen, das Levitische aufhören I, 1589. ff.; voll Glaubens II, 1963. ff.; unterschied sich durch den Glauben von Cäsar etc. II, 1566.; eignet sich den Eid, den Gott Abraham geschworen, zu I, 1557. \gesperrt{Leben und Wandel}: führte einen heiligen Wandel in allen drei Ständen II, 460.; war ein Pilgrim I, 738. f. 774. \gesperrt{Werke}: diente in seinem Amte dem Willen Gottes II, 1411.; hat durch den Glauben Großes gethan I, 1501.; II, 1964., Goliath überwunden II, 287. 1963. f.; gefiel Gott auch in seinen Kriegen wohl II, 459. f.; hatte ein gottseliges Reich II, 1823. f.; in ihm erfüllt sich theilweise Jakobs Segen über Juda II, 1961. ff. \gesperrt{Eheliches Leben}: Lebte nach Absaloms Blutschande ohne Weib II, 977. \gesperrt{Sünden}: fiel, da er aufgeblasen ward I, 1558.; wie die Sünde mit Bathseba gestraft worden II, 1116. \gesperrt{Anfechtungen und Leiden}: Hatte mit Fleisch und Unglauben zu kämpfen II, 287. f. 2087., mit Verzweiflung II, 443.; zweifelte an der Verheißung I, 934.; wie er sich getröstet II, 304. \gesperrt{Schriften}: Hat Mose fleißig betrachtet I, 1618.; hat in seinen Psalmen aus Mose geschöpft I, 1559. 1585. ff. 1685.; II, 1133.; hat viele Psalmen um seiner Verfolgungen willen gemacht II, 1962.; Ps. 1. II, 2041.; Ps. 107. I, 1319.; Ps. 110. und 132. I, 1557. ff.

\quellseite{546}
Debora. II, 942. ff. War heilig II, 942.; gottselig und weise II, 946.; Prophetin II, 949.; Großmutter im Hause Jakobs II, 946. f.; geehrt und geliebt II, 942. 947.; erlebte Josephs Verkauf nicht II, 1102.

Delphin. König im Meer I, 62.

Demosthenes. Ausspruch über das Regieramt II, 1405. 1409.

Demuth. Eine seltene Gabe II, 1393. f.; größere Tugend denn Geduld II, 1395. f.; schöne Tugend, besonders an Hochgestellten II, 1393.; Gott ein angenehmes Rauchopfer I, 1453.; ist auch in weltlichen Dingen ein gut Werk I, 1629. f.; darum soll bitten, wer ein Amt hat I, 635.; was uns dazu treiben soll II, 1396., unsere Thorheiten II, 276. f., unser Elend und Jammer II, 1388. 1390.; zur Uebung darin dient Kreuz I, 823. f.

Demüthige. Gott ist ein Gott derselben I, 1446.; tröstet sie I, 1430.; Christen sollen gegen einander demüthig sein I, 1144. f.

Demüthigung. Gott demüthigt Hohe und Begabte I, 437. f.; D. thut wehe I, 988.; dient zur Läuterung, nicht Genugthuung II, 1660.; wie wir darin beten sollen I, 1434.; Gott tröstet darin I, 1444.

Deukalion. Fabel von D. I, 687.

Deutsche. Ihr Stammvater vielleicht Askenas I, 660.; ihre Vorfahren große, treffliche Leute I, 663.; galten allezeit für die löblichste Nation I, 660.; sind Bierbäuche II, 1360., Schlemmer II, 1809. 1819. — S. \gesperrt{Askenas}.

Deutschland. Hat herrliche Gaben I, 663.; nicht einig gegen den Türken I, 706. — S. \gesperrt{Prophezeiungen} 2.

Dialektik. 1. Was dazu gehört II, 1357.; fordert: daß man unterscheide und theile I, 1580., die Worte nur in Einer Bedeutung gebrauche I, 1407.; II, 1373.

2. \gesperrt{D. und Rhetorik} I, 1149. f. 1153. 1175. f.; II, 203.

Diana. Hieß πολύμαστος I, 1016.; II, 1585.

Dichten des menschlichen Herzens. Was das sei I, 483. f. 579.; ist böse I, 579.

Dina. Wann geboren II, 551. 571. f.; war vorwitzig, ohne der Eltern Wissen ausgegangen, da sie geschwächt ward II, 862.; ward Hausmutter in Jakobs Hause II, 1565. — S. \gesperrt{Sichem} 2.

Diocletian. Weise und klug, verfolgte die Christen II, 1658. f.; 1984.

Dionysius. 1. \gesperrt{D. der Aeltere, aus Syrakus}. Plünderte das Bild Jupiters II, 923.

2. \gesperrt{St. Dionysius}. War nicht St. Pauli Jünger I, 287. f.; soll Erscheinungen gehabt haben I, 1120.; Meinung über die Hierarchien I, 287. f.

Dodanim. Untergegangene Völker I, 662.

Donat. Schulbuch II, 58. f.

Doppelzüngige. Reden anders, denn sie es meinen II, 583.; Fromme sind nicht doppelzüngig II, 255.

Dornen und Disteln. Erst seit dem Fall I, 46. 255. f., oder vorher unschädlich I, 94.

Dreieinigkeit. 1. Die \gesperrt{Lehre von der} D. I, 1148. ff.; Unterschied der Personen bezieht sich auf die Werke und Offenbarungen Gottes I, 71.; der Vater gibt, der Sohn erklärt, der Heilige Geist erfüllt die Träume II, 1053.; ist lieblich, Fußstapfen der D. in allen Creaturen zu suchen I, 1151. Jetzt wie ehedem glaubt man diese Lehre nicht I, 25. f.; II, 940.; Juden und Türken spotten derselben I, 71. f. 1149. f. 1278. ff.; II, 213.

2. \gesperrt{Offenbarung der Lehre von der} D. Diese Lehre ist in der Schrift immer klarer geworden I, 274.; die deutliche Offenbarung gehörte ins N.\,T. I, 22. 707. ff.; II, 940.; warum I, 72. \gesperrt{Offenbarung} im A.\,T.: man soll sich befleißigen, diese Lehre aus dem A.\,T. zu beweisen II, 915.; \gesperrt{Mehrzahl der Personen} angezeigt, wo von Elohim die Rede ist II, 939. f.; was die Juden auf diesen Beweis erwiedern II, 1529.; wiederholt von Mose und sonst im A.\,T. gelehrt I, 24. ff.; von Mose I, 11. f. 14. f. 24. ff. 70. f. 707. 1278. f. 1352. 1584. 1718.; II, 213. ff. 703. f. 853. 915.; \gesperrt{Unterschied} der Personen angezeigt I, 20.; \gesperrt{Einheit} des Wesens und \gesperrt{Mehrzahl} der Personen angezeigt I, 273. ff.; \gesperrt{Dreizahl} der Personen I, 60. f.; ob sich diese Lehre aus 1 Mos. 18, 2–5. erweisen läßt I, 1148. ff. 1152. f. 1183. \gesperrt{Offenbarung} im N.\,T. bei der Taufe Christi I, 1152.

Drohungen. 1. \gesperrt{Drohungen Gottes}. D. wider die Gottlosen scheinen oft an den Frommen in Erfüllung zu gehen I, 360.; Gott verzieht mit der Erfüllung I, 641.; II, 1917.; Welt verachtet und verlacht sie I, 550. f. 641. 1243.

2. \gesperrt{D. und Verheißungen Gottes}. Stets bei einander I, 1076.; streiten nicht wider einander I, 550. ff.; Gott verzieht mit der Erfüllung II, 1917.; erfüllt sie dann desto herrlicher oder schrecklicher ib.; Erfüllung scheint oft widersinnig II, 1813.; fordern Furcht und Glauben II, 1917. f. 1920.; D. soll man einschränken, die V. erweitern I, 1083.; das Fleisch verachtet beides II, 1917. ff.; das natürliche Herz hört lieber V. als D. I, 1185. f. — S. \gesperrt{Prophezeiungen} 1.

Dudaim. Was für eine Frucht II, 565. f.

Dürre im Jahre 1540. I, 594.

\section{E}

Ebenbild, das göttliche. Unterschied zwischen „Bild“ und „Gleichniß“ I, 72. f. 410. f.; worin es bestand I, 74. ff. 145. 172. 411.; der erste Mensch hatte es wesentlich an sich I, 76.; wohl an seinem Angesicht zu sehen I, 104.; was es gewesen, werden wir erst im ewigen Leben eigentlich erfahren I, 79.; ist verloren I, 75. f. 411.; durch den Glauben, doch nicht völlig, wieder hergestellt I, 78. ff. — S. \gesperrt{Erbgerechtigkeit}, \gesperrt{Erneuerung}.

Eber. Sems Urenkel I, 680.; woher der Name I, 885.; als Vorfahre Christi hervorgehoben I, 679. f. 715.; standhaft gegen Nimrods Kirche I, 715. 885.; lehrte und regierte die Gemeinde II, 89.; rieth Rebekka II, 262. f.; lebte mit Jakob 80 Jahre I, 716.; II, 89.; viele seiner Nachkommen Priester II, 456. f.

Eck, Dr. Frug nichts nach der Wahrheit des Evangeliums II, 671. f.; lästerte Luther II, 1274.; schreckliches Ende II, 1274.

Eden. Heißt Lust I, 106. f.; Name eines Orts I, 106. f., nicht der ganzen Erde I, 119.; weiter Raum, zur steten Wohnung der Menschen bestimmt I, 118.; ob die Flüsse E.’s Eine Quelle hatten I, 119.; daß die Beschreibung E.’s. jetzt auf keinen Ort der Erde paßt, kommt von der Sintfluth I, 120. f. — S. \gesperrt{Paradies} 1.

Eder. Bedeutung des Namens, Lage dieses Thurmes, später in eine christliche Kirche umgewandelt II, 972.

Edom. 1. \gesperrt{Name Esaus}. — S. \gesperrt{Esau} 1.

2. \gesperrt{Name eines Landes} II, 984. 993.; wie lange das Reich Edom gewähret II, 1018.

3. \gesperrt{Edom} nennen die Juden die christliche Kirche und das Römische Reich II, 984.

Edomiter. Wohnsitz I, 880.; Israel unterworfen II, 822. f. 1009.; abgefallen II, 1018.; viele selig II, 828. — S. \gesperrt{Edom} 2, \gesperrt{Esau} 2.

Egypten. Anfang des Reichs I, 712.; wohlgeordnet II, 1782. f. 1808.; hatte viele herrliche Würze II, 2073.; reich und fruchtbar II, 1613. f.; Zustand zu Abrahams Zeit I, 789. 804. 814., zu Josephs Zeit II, 1223. 1293. 1350. f. 1358.; Fruchtbarkeit verloren I, 852.; II, 1794.

Egypter. Religion: Priesterstand in Ansehen II, 1413.; hatten Erkenntniß Gottes I, 812. f., zu Josephs Zeit noch Ueberbleibsel der Lehre der Väter II, 1417.; hatten einen seltsamen Aberglauben II, 1825., beteten Katzen und Mäuse etc. an II, 1829.; durften keinen Ochsen schlachten, der im Pfluge gezogen hatte II, 1615.; verabscheueten andere Völker II, 1613. f. 1776., die Juden II, 1613. f. 1780., die Viehhirten auch ihres eigenen Volkes II, 1614. f. 1785. Berühmt in der Astronomie I, 805.; bedienten sich der Hieroglyphen I, 1476.; verstanden die Kunst, Leichname zu balsamiren II, 2073.; Gebräuche bei Mahlzeiten II, 1615. ff. 1619. ff.; mäßig II, 1809.; Strafe der Ehebrecher I, 814. — S. \gesperrt{Abraham} 21, \gesperrt{Joseph} 15.

Ehebrecher. Bei ihnen hört die Brunst des Mannes zum Weibe und umgekehrt auf II, 1265.; wie man den sicheren, wie den erschrockenen E. predigen soll I, 1188. f.

Ehebruch. Allezeit Sünde I, 818.

Ehegrade. Zu Abrahams Zeit durfte man noch seines Bruders Tochter heirathen I, 1601.; Bruders Stieftochter heirathen ist nicht verboten I, 1351. f.

Eheleute. 1. \gesperrt{Ehre}. Gott hat Wohlgefallen an ihnen II, 568.; haben größeres Lob in der Schrift denn Jungfrauen II, 571.; bei den Heiden angenehmer denn alle Philosophen II, 568. f.; sollen auch bei den Christen angenehm sein II, 569.

2. \gesperrt{Arten}. Verschiedene Arten der E. I, 1749.; was von denen zu halten, die keine Kinder wollen II, 579.

3. \gesperrt{Verhältniß zu einander}. Die Männer haben keine Ursache, sich über die Weiber zu erheben I, 314.; haben Alles gemein I, 167. f. Der \gesperrt{Mann} ist die Sonne, das Weib der Mond I, 84.; Herrschaft ist bei dem Manne I, 248. f.; ist Herr des Weibes I, 268. f.; regiere freundlich, mit Liebe I, 794. 987., sei nicht herrisch I, 1353. Das \gesperrt{Weib} soll nicht herrschen I, 1676.; dem Mann um der Sünde willen unterworfen I, 141. 168. 248.; soll gehorchen I, 793. 1388.; ist eine Gehülfin, was von weiblichen Thieren nicht gilt I, 143.; soll dem Mann folgen, wenn er zieht I, 269. 382. f.; II, 665. f.; Wenige folgen Sarahs Exempel I, 1170. f. — S. \gesperrt{Argwohn}, \gesperrt{Liebe} 5.

4. \gesperrt{Gegenseitiges Verhalten}. Nichts schöner, denn Liebe und Einigkeit unter E. I, 800. f.; sollen freundlich und gelassen sein II, 159. ff.; mögen scherzen und sich küssen II, 158. f.; gegenseitig verzeihen I, 1354.; II, 161. \gesperrt{Vertrauen} muß sein Maß haben I, 801.; kein Argwohn soll sein I, 800. f.; Vertrauen die vornehmste Tugend I, 800. Sollen sich mäßig und \gesperrt{keusch} halten II, 167. f. — S. \gesperrt{Liebe} 5.

5. \gesperrt{Verhalten des Mannes}. Soll sein Weib als Gebäude Gottes erkennen I, 163.; mehr an ihr, als an den Eltern hängen I, 169.; sie lieben, nicht bloß sich ihrer gelüsten lassen I, 1743.; Freundlichkeit eine schöne Tugend am Manne I, 159. ff., Bitterkeit verächtlich und sündlich II, 162. f.; 602. ff.; Jeder lasse sich an seiner Eva genügen I, 294.; wohne besonders in der Schwangerschaft mit Vernunft bei ihr II, 602. ff.

6. \gesperrt{Verhalten des Weibes}. Frommes Weib hat Mitleid mit ihrem Mann, wenn er beschweret wird II, 969.; ist recht, wenn sie durch Fruchtbarkeit begehrt, dem Manne zu gefallen II, 564.

Vgl. \gesperrt{Ehestand}, \gesperrt{Ehemann}, \gesperrt{Eheweib}.

Ehelosigkeit. 1. \gesperrt{Im Allgemeinen}. Vom Ehestande übertroffen I, 1165.; mit Unreinigkeit behaftet II, 530.

2. \gesperrt{Im Pabstthum}. Warum angeordnet I, 293. 1374.; nicht von Gott I, 1164., wider Gottes Wort I, 165. 1033., eine Versuchung Gottes I, 525.; stinkt vor Gott II, 562.; greuliche Sünden und Unzucht sind daraus gefolgt I, 292. f.; II, 539. 1198.; das war Gottes Strafe II, 163.; daher viel uneheliche Kinder zu Rom I, 293., zahlloser Kindermord I, 293., Verachtung des Ehestandes I, 1378.; sollte aufgehoben werden II, 364.

Ehemann. Seine Werke herrlich, aber mühevoll II, 249.; nichts ihm so unleidlich, als wenn ein Anderer nach seines Weibes Ehre steht oder ihr Gewalt thut II, 1278. f.; kann eher keusch leben, als der noch nicht verehelicht war II, 1264. — S. \gesperrt{Eheleute}, \gesperrt{Geschlechter}, \gesperrt{Mann}.

Ehescheidung. Gott hat Mißfallen daran II, 162.; warum von Gott den Juden gestattet I, 169. 1230. f.; papistische Ehescheidungsgründe I, 1729. f.

Ehestand. 1. Die \gesperrt{Lehre vom} E. ist nach dem Evangelium das wichtigste Stück II, 359.; fleißig zu treiben II, 358.; züchtig zu handeln I, 1644.; Jugend darin zu unterweisen I, 976.; dafür Gott zu danken II, 167.

2. \gesperrt{Was der} E. \gesperrt{sei} II, 359., falsche und rechte \gesperrt{Definition} I, 1645. 1653. 1673. 1677.

3. \gesperrt{Stifter und Erhalter} des E.’s. ist Gott I, 1658. f.; II, 360. 538.; er hat ihn eingesetzt I, 163. 1373. f.; nach der Sintfluth bestätigt I, 588.; hat ihn geordnet und regiert ihn I, 972., liebt und erhält ihn I, 990., schützt ihn I, 1323., hat Gefallen daran I, 981. 1730.; II, 477. f.; führt Mann und Weib zusammen I, 1737. f. 1740.; II, 161. f., durch den Dienst der Engel I, 1684. ff.

4. \gesperrt{Die Personen}, die ehelich werden. Wen man erwählen soll I, 1675. ff.; keine, die Einem nicht gefällt II, 487.; Starke sollten mit Starken ehelich werden II, 486.; Alte mögen Junge ehelichen I, 1752.; Eltern sollen ihre Kinder nicht mit Gottlosen verloben I, 1243. f.; die Ehe mit ungläubigen Weibern hat Gefahr I, 468.; die Frommen der ersten Welt verboten ihren Kindern die Ehe mit Kainiten I, 445. f. 470.; die Juden sollten sich nicht mit den Cananäern vermischen I, 445. f.; im A.\,T. keine Sünde, eine Heidin heirathen II, 1416.; die Heiden hatten ein Gesetz, daß sie die nächsten Blutsfreundinnen heirathen, oder ihnen die Aussteuer geben mußten II, 1197.; eine zweite Ehe ist statthaft I, 1748. f.

5. \gesperrt{Unauflöslichkeit} des E.’s. II, 487.

6. \gesperrt{Zweck} des E.’s. Endursache desselben I, 1673.; Mehrung und Fortpflanzung II, 579.; christliche Erziehung der Kinder II, 360.; Meidung der Unzucht II, 486.; Fortpflanzung und Meidung der Unzucht I, 142. 144. f.; II, 360.

7. \gesperrt{Heiligkeit} des E.’s. Ein heiliger, ehrlicher Stand, nicht der Hurerei gleich II, 365.; hindert nicht heilige Werke I, 1162.; nicht fleischlich und unfläthig I, 164. 1684. 1691.; nicht der E., sondern der Mißbrauch des E.’s vor der Sintfluth gestraft I, 470.; Verderben des E.’s kommt von der Sünde I, 169. 205. f.; Gebrechen des E.’s soll man nicht groß machen und verlachen I, 293. f.; ist geheiligt durch Christum I, 1659. f., durch Gottes Wort II, 125.

8. \gesperrt{Werke} des E.’s. Gott sieht im E. auf das Kinderzeugen und die Keuschheit I, 1752.; die Werke des E.’s gehören zu den allergrößten I, 528.; sind heilig I, 1164. f., durch Christum geheiligt I, 1665. f.; Beischlaf keine Unreinigkeit II, 168., doch haftet ihm Sünde an II, 168., um derentwillen man Vergebung bedarf II, 487. f.;
\quellseite{547}
Gott deckt die Sünden und Gebrechen zu I, 1648.; II, 157. 163. f. 167. f.

9. \gesperrt{Ehre und Herrlichkeit} des E.’s. Ist verachtet, doch herrlich I, 1372. f.; deßselben braucht man sich nicht zu schämen II, 365.; nichts Trefflicheres in diesem Leben I, 1644.; ein Ueberbleibsel aus dem Stande der Unschuld I, 162. 167. f.; Gott rühmt ihn im vierten Gebot I, 1738. f.; man soll ihn ehren I, 293. f. 1680.; II, 1198.; wer die Gabe der Keuschheit hat, soll den E. nicht verachten I, 1172. f. 1735.; II, 1073.; ihn zu Ehren bringen, eine nöthige und nützliche Arbeit I, 165. 1681.; die heiligen Väter ehrten ihn I, 1746. ff.; die Heiden schätzten ihn I, 969. 976. 1680. f.; wir sollen über dem E. die Gabe der Keuschheit nicht verachten II, 538. f.

10. \gesperrt{Feinde und Verächter} des E.’s. Satan ein Feind des E.’s I, 292.; II, 359. 487. 1264. f.; Welt und Fleisch verstehen ihn nicht zu schätzen I, 294.; II, 358.; bei manchen Heiden verachtet I, 976. Es verachten ihn Papst I, 1063. 1680., Mönche II, 166. f., Papisten I, 1162. 1172. f.; II, 1175., wälzen sich in Unzucht II, 556.; sie halten ihn für Hurerei I, 1652.; haben dagegen gewüthet I, 292.; versündigen sich schwer I, 291. f. Verächter des E.’s., Ungeheuer I, 163.; wie ihnen zu begegnen I, 1733. Ursachen der \gesperrt{Verachtung}: Erbsünde I, 144. 975. f., Cölibat I, 1063. 1378., man will nicht an Eine Person gebunden sein I, 975.; II, 487. Ursachen der \gesperrt{Ehescheu}: die Beschwerden I, 246. 1679. f.; II, 361., Armuth II, 542. Aus der Verachtung folgt großer Schade I, 975.

11. \gesperrt{Nothwendigkeit und Nutzen}. Der E. ist \gesperrt{nöthig} I, 1176.; das erkannten die Heiden I, 169.; nach dem Fall nöthiger als vorher I, 164. Er ist \gesperrt{nützlich} I, 976. 1680.; wo er fällt, fällt alle Ordnung I, 1227.; ist eine Pflanzschule der drei Stände I, 294.; ohne ihn müßte das menschliche Geschlecht vergehen II, 1073.; zu loben um des Nutzens willen I, 1162.; Nutzen erkannten auch die Heiden I, 169. Ist ein \gesperrt{guter} Stand, darum eine Freundschaft, wenn man Einem seine Tochter gibt I, 1736.; hätte Gott ihn nicht eingesetzt, Alle würden ihn begehren II, 163. f.; wir sind des E.’s nicht werth II, 168. f.; ein Zeichen der Liebe Gottes I, 588. f.

12. \gesperrt{Fruchtbarkeit und Unfruchtbarkeit} des E.’s. Die eheliche Beiwohnung wäre im Anfang nicht ohne Frucht geschehen I, 295.; trotz der Sünde bleibt dieser Segen Gottes I, 294. \gesperrt{Fruchtbarkeit} ein Werk Gottes II, 518. ff.; galt allezeit für einen großen Segen II, 523. 532.; die heiligen Weiber begehrten sie II, 1176.; nicht geachtet zu Luthers Zeit II, 532.; man sollte sie am Menschen nicht weniger als am Vieh begehren II, 1176., das Gegentheil geschieht II, 1197.; es ist unmenschlich, dieselbe verachten II, 532. f.; als man sie zu verachten begann, hat sich die ganze Welt geändert I, 1378. \gesperrt{Unfruchtbarkeit} eine Folge der Sünde I, 163. f.; galt im A. T. für eine Schmach und Fluch I, 724. 1079. f. 1378.; II, 25. 545. 577. ff. 1197., jetzt nicht mehr I, 1079. f. 1378.; II, 578.; was von denen zu halten, die Unfruchtbarkeit begehren II, 579.; berechtigt nicht, ein zweites Weib zu nehmen I, 972. ff.

13. \gesperrt{Wehe} des E.’s. \gesperrt{Gute Ehen} sind selten I, 1353. f. 1678. 1680.; II, 1264.; woher das kommt I, 1353. f.; ist des Teufels Schuld II, 160.; kommt daher, wenn man nur Wohllust sucht I, 1672. f., wenn man den Stand nicht recht anfängt II, 1344., ohne den Willen der Eltern I, 1457.; ohne Christum und den Glauben kann die Ehe nur übel gerathen I, 1672.; ehelich werden kein Scherz I, 1644. f. Der E. ist nicht eitel \gesperrt{Wohllust} II, 137.; man sei auf Anfechtungen gefaßt I, 984. f.; da sind vielerlei Uebel, Gebrechen und der Tod I, 1730. f.; man entdeckt allmählich Untugenden I, 979.; viel Sorge und Mühe II, 1199.; viel Unlust und Anfechtungen I, 983. f. 1392. f.; da gibts Zank I, 986. f. 1386. ff. 1392., zu dem Satan reizt II, 1180.; die erste Liebe erkaltet II, 488.; oft erkennt man die Ursache der Noth nicht I, 821. f. Diese Leiden sind Uebungen des Glaubens, der Hoffnung etc. I, 1392. f. Wie man die Mühseligkeiten überwinden soll II, 360. f. 363. 368. f.; man tröste sich damit, daß Gott den E. will I, 1644. f., der E. ihm gefällt I, 1486., die Eltern eingewilligt haben II, 366. 1107. f.; man denke an das Gute, die Kinder I, 1734.; man bete I, 1679.; II, 364.; Gottes Wort macht geduldig II, 126. f. Das Kreuz überwindet man leichter, wenn man den Stand recht anfängt I, 1652. f. 1678. f.

14. \gesperrt{Rechter Anfang} des E.’s. Rechte Gesinnung bei Antritt II, 1344.; nicht bloß Wohllust suchen I, 1672.; Vorsicht bei der Wahl II, 1415.; man bete I, 1354. 1650. ff. 1677. f. 1682. f. 1707.; man sehe nicht auf Geld, sondern Frömmigkeit I, 1707.; ob man auf Schönheit und Stärke sehen mag II, 484. ff.; Einwilligung der Eltern ist nöthig I, 1455. f. 1682.; II, 364. f. 366. ff. 1161. f. und der Kinder I, 976. 1647. f.; II, 361. 367. 506. — S. \gesperrt{Verlöbnisse} 2.

15. \gesperrt{Wer und wann} man ehelich werden soll. Wer die Gabe der Keuschheit hat, braucht es nicht I, 1735.; II, 1073., sondern wer dieselbe nicht hat II, 165. 361. 368. Man soll nicht gleich zur Ehe greifen, wenn man die Flammen der Unzucht spürt II, 21. f.

Vgl. \gesperrt{Eheleute, -mann, -weib, Hausstand, Jakob} 9, \gesperrt{Leviratsehe, Polygamie}.

Eheweib. Frommes E. kommt vom HErrn II, 22.; Nutzen des E.’s I, 1614.; Gebrechen durch die Mutterschaft aufgewogen I, 1731. ff.; auch ein gottseliges will geliebt, oder doch nicht verachtet sein II, 519. ff.; sie sei häuslich I, 248.; Tod eines gottseligen E.’s, aber nicht eines bösen, ist dem Mann ein großer Schmerz II, 964. f. — S. \gesperrt{Eheleute, Geschlechter, Weiber}.

Ehre Gottes. Dazu dienen alle unsere Werke II, 1191.

Ehrenworte. Was man E. nennt I, 1215.

Ehrerbietung. Geboten II, 824. ff.; Gott will E. gegen die weltliche Obrigkeit II, 829.; Einer erweise sie dem Andern I, 1444. 1634.; auch Geringeren erweist die Liebe E. II, 824.; ob man Fremdlingen solche E. erweisen soll, wie einst Abraham gethan I, 1146.; E. gegen die Eltern wirkt der Heilige Geist II, 1676. — S. \gesperrt{Eltern} 2.

Ehrgeiz. Welt in E. ersoffen, daher Unfriede I, 844.

Eid. 1. E. \gesperrt{Gottes}. Nur einmal steht in der Schrift, daß Gott geschworen habe I, 1557.; Gott schwört bei sich selbst I, 1547.; was das heißt I, 1556.; Gottes E. ist heilig und wichtig I, 1548. f., tröstlich I, 1549. f.; sein E. bei den Verheißungen, ein Zeichen der Liebe I, 1548. f.; Zweck des E.’s: wir sollen glauben I, 1554. f.

2. E. \gesperrt{der Menschen}. Jeder E. enthält eine Anrufung Gottes und einen Fluch I, 1674.; II, 1483.; ein sehr heilig Ding und Gottesdienst I, 1656. f.; die Schrift macht daraus gleichsam ein Sacrament I, 1465.; Gott und Menschen ein angenehmer Dienst I, 1464.; ihn halten ist Gottesdienst I, 1674.; welcher E. verwerflich und welcher gut I, 1464. f.; man mag beim Himmel etc. schwören II, 1483.; bei der Seele schwören ist bei Gott schwören I, 1547. f.; Gebräuche beim E. I, 1656. ff. 1673., bei Josephs E. II, 1834. ff. 1840.; \gesperrt{Christen} dürfen schwören I, 1692. f., wann I, 1464.; II, 1483., sollen es nicht ohne Noth I, 1675.; Christus straft leichtfertige E. II, 1483.; \gesperrt{falscher} E. Verleugnung u. Verachtung Gottes I, 1674. f. — S. \gesperrt{Schwören}.

Eifersucht. Macht die Ehe zur Hölle II, 161.

Eigenlob. Jakob lobt seine Treue etc. II, 585.

Eigenthum. Ob man E. haben darf I, 835. f.

Eingebung. Man folge nicht jeder II, 95. ff.; Kennzeichen, ob sie vom Heiligen Geiste II, 96. f.

Einigkeit. Kann nicht sein, wo man nicht Aussöhnung sucht II, 871. f.; sie zu erhalten gib von deinem Rechte nach I, 842. 847. f., doch nicht von Gottes Gesetz I, 848.

Einsamkeit. Thut wehe I, 841.; man ziehe sich nicht in die E. zurück II, 1642.

Einsiedler. Entstanden in den Verfolgungen I, 1320.; sind kleine Heilige I, 1320.

Elam. Vater der Perser I, 679. 713.; vertrieben I, 713.

Elassar. Wohl Assyrien I, 878.

El Bethel. Ortsname II, 937.

Elemente. Veränderlich, ob auch verweslich I, 35.; Lehre der Philosophen I, 32. f.

Elia. Abkunft II, 2033.; hatte menschliche Affecten I, 1308. f., Erbsünde II, 843.; ward getröstet I, 1341.

Elieser. Bedeutung des Namens I, 937.; guter Theologe I, 1699. f.; fromm I, 1655.; treuer Knecht I, 1722. ff.; vielleicht später frei I, 1655.; hat gebetet I, 1695. ff. 1702. und gedankt I, 1718. f.; ob sein Gebet eine Versuchung Gottes war I, 1695. ff. — S. \gesperrt{Abraham} 22.

Elihu. Ob er Bileam war I, 1601. f.

Eliphas, Esaus Sohn. II, 990. ff.; guter, frommer Mann II, 991.; behielt Beschneidung und göttliche Lehre II, 999.; ob er der Freund Hiobs war II, 990. f.

Elisa. Japheths Enkel, Völkername I, 661.

Ellipse. Beispiel einer E. II, 596.

Eltern. 1. \gesperrt{Eltern}. In ihnen ehrt man Gott I, 1635. und gehorcht ihm I, 346. 762.; haben ein priesterlich Amt I, 302. f.

2. \gesperrt{E. und Kinder}. \gesperrt{Liebe} der E. zu den K. ist natürlich II, 627., eine große, heftige Neigung II, 328.; wie groß sie sei und wie groß der Schmerz, wenn K. die E. betrüben, verstehen die K. nicht II, 328. ff.; sie können an den E. leicht zu Mördern werden II, 328. ff. \gesperrt{Hoffnungen} der E. schlagen oft fehl I, 655.; K. gerathen nicht immer nach den E. II, 1415.; wie sich E. über mißrathene K. trösten sollen II, 979. Die Mutter gibt dem Kindlein mehr Nahrung als der Vater II, 1173. E. sollen die K. in Zucht halten II, 864. f. 870. f., von Amts wegen \gesperrt{strafen} I, 347., aus Liebe I, 800.; II, 807. 1460. f. 1462. f., so ihr Bestes suchen II, 1491.; das sollen die K. merken können II, 1491. f.; E. sollen die K. nach der Strafe wieder versöhnen I, 577. Welche K. vom Vater, welche von der Mutter \gesperrt{vorgezogen} werden II, 91. f.; die E. pflegen fromme K. den gottlosen vorzuziehen II, 1036. E. sollen den K. zu einem gottseligen \gesperrt{Ehestande} helfen I, 1682.; haben keine Gewalt über verehelichte Töchter II, 682. \gesperrt{Kinder} sollen \gesperrt{gehorchen} II, 259.; so gehorchen sie Gott I, 1456.; Gehorsam segnet Gott I, 1456. f.; Ungehorsam schreckliche Sünde, Ursache vielen Unglücks II, 1102.; sollen nicht gehorchen, wenn die E. etwas gegen Gottes Wort gebieten II, 267. 943.; sollen die E. \gesperrt{ehren} I, 1597.; II, 328. ff., beten, damit sie dies thun können II, 330.; solche Ehrerbietung Gott das angenehmste Opfer I, 638.; Sem und Japheth Exempel der Ehrerbietung I, 638.; E. von den K. leicht durch Traurigkeit getödtet II, 328.; betrüben die K. die E., so fluchen sie denselben II, 329.; Verachtung straft Gott I, 639.; II, 330.; man soll die E. nicht verachten, außer wo es der Beruf fordert II, 943.; wie K. die Sünden der E. strafen sollen II, 232. — S. \gesperrt{Ehestand} 14., \gesperrt{Vater, Töchter, Verlöbnisse} 2.

Emim. I, 880.

Empfängniß. Dabei bekommt die Frucht bei Menschen und Thieren leicht eine besondere Gestalt und Flecken II, 601. — S. \gesperrt{Fortpflanzung, Geburt, Schwangerschaft, Zeugung}; E. \gesperrt{Christi} s. \gesperrt{Christus} 2.

Enallage. Vertauschung, Beispiel II, 1804.

Engel. 1. \gesperrt{Allgemeines}. Die Schrift schweigt von ihrer Schöpfung I, 27.; sind Geister II, 394. 735.; nicht wie Menschen gezeugt I, 137.; alle gut geschaffen, es gibt gute und böse I, 28.; konnten wie die Menschen fallen I, 137.; halfen nicht bei der Schöpfung I, 14.; können wohl die Sprache der Menschen an sich nehmen, aber nicht sie verwirren I, 707.

2. \gesperrt{Die guten} E. \gesperrt{Natur und Wesen}: können nicht mehr fallen I, 136. 138.; stark I, 1244. f. 1540. f.; II, 661.; verstehen die wunderbare Regierung Gottes II, 731.; demüthig und freundlich, voll Erkenntniß Gottes und seiner Güte I, 1687. f.; beste Freunde der Gläubigen I, 1689. \gesperrt{Amt und Werke}: Vernunft kennt die Lehre vom Amt der E. nicht II, 726.; Gott bedient sich ihrer I, 1540. f.; Amt im Himmel und auf Erden II, 724. 730. f.; sind Boten und Kriegsknechte II, 735.; sind bei uns, sehen unsern Werken zu I, 1522.; II, 200.; dienen uns I, 1685.; haben Lust daran I, 1690. f.; regieren unsere Werke I, 992.; behülflich bei der Ehestiftung I, 1684. ff. 1691. 1694.; strafen die Gottlosen, schützen die Frommen I, 1245. ff.; II, 661.; wachen, regieren, streiten für Fromme und Gottlose II, 724. ff.; in allen Ständen unsere Hüter I, 1686. f.; ihr Schutz gewiß, wenns auch nicht so scheint II, 733. f.; ihres Schutzes bedürfen wir I, 1246.; II, 729., sollen wir uns trösten I, 1246. f.; II, 200. f. 730., dafür Gott danken, den E. nachfolgen I, 1688. f.; Erkenntniß vom Schutz der E. den Gläubigen nöthig und nützlich II, 732.; ihr Dienst eine Wohlthat I, 1690., von den Wenigsten erkannt I, 1686. f. Man soll die E. \gesperrt{nicht anbeten} I, 1247. — S. \gesperrt{Arzenei, Freunde}.

3. \gesperrt{Die bösen} E. Von ihrem \gesperrt{Fall} schweigt die Schrift, mancherlei Gedichte darüber I, 27.; was die Schrift meldet I, 28.; wohl ein Aufruhr und Kampf zwischen den E. I, 28. 137.; wohl durch Hoffahrt gefallen I, 28. 137. f. 299. 1689. f.; wann gefallen, ungewiß I, 184.; viele I, 138.; haben eine Art Monarchie II, 724.; stark I, 1245.; richten allerlei Unglück und Sünde an II, 724. ff. — S. \gesperrt{Teufel}.

4. \gesperrt{Die guten und die bösen} E. Streiten beständig in allen göttlichen Werken I, 1701. ff.; II, 725. ff., um die Menschen II, 724. ff.; die guten schützen uns wider die Anschläge der bösen I, 1246.; woher dann doch so viel Böses II, 726. ff.; Gott bedient sich des Dienstes beider II, 1051. — S. \gesperrt{Erscheinungen, Teufel}.

Enos. Bedeutung des Namens I, 398. f.; Stand der Kirche zu seiner Zeit I, 398. ff. 401. 417. f.

Entzückung. Was das für ein Zustand sei II, 2070.

Epha. I, 1751.

Ephraim. 1. \gesperrt{Josephs Sohn}. Bedeutung des Namens II, 1430. 1432. 1853. — S. \gesperrt{Jakob} I, 21, \gesperrt{Joseph} 4.

2. \gesperrt{Der Stamm}. Wohnsitz II, 2018.; stolz auf Herkunft II, 1368. f. 1848. f. 2053.; deshalb hart gestraft II, 1848. f. 1887. f.; nach Assyrien geführt II,
\quellseite{548}
2051. Die \gesperrt{Stämme} E. und Manasse hatten viel Volk II, 2041., treffliche Männer II, 2041. 2047.; David hielt viel von ihnen II, 2041.; wie sich an ihnen Jakobs Segen erfüllt II, 2046.

Ephratha. Bedeutung des Namens II, 971.

Ephron. I, 1628. ff.

Epictet. Philosoph, ging betteln I, 830.

Epicur. Meinte, der Mensch sei nur zum Essen und Trinken geschaffen I, 68.

Epicuräer. Meinung über Entstehung der Welt und des Menschen I, 150.; verlachen die Unsterblichkeit II, 15. 1305.; verachten den Himmel II, 115., Wort und Sacrament II, 116.; viehisches Leben II, 1305.; Allegorie auf sie I, 618.; nahmen zu Luthers Zeit zu I, 1764.

Erasmus. 1. \gesperrt{Der Heilige} der Papisten II, 1919.

2. \gesperrt{Desiderius} E. Gelehrt und beredt, versteht in Glaubenssachen nichts I, 998.; Meinung über: Erbsünde I, 132., eine Kindesmörderin II, 1515., Socrates I, 581., Tugenden der Heiden II, 461. f.

Erbe. Darüber entzweien sich Brüder leichter als Fremde II, 1110.

Erbgerechtigkeit. Etwas Natürliches I, 201. ff.; was sie war I, 138. f.; aus der Erbsünde zu erkennen I, 139. — S. \gesperrt{Ebenbild}.

Erbsünde. Ist ein großes Verderben der ganzen menschlichen Natur I, 139. f. 203. f., welche nicht nur zum Bösen geneigt ist I, 577., auch nicht nur übernatürliche Kräfte verloren hat I, 201. ff.; eine erschreckliche Finsterniß und Abwendung von Gott II, 1463.; das ist die E., wenn man wie Gott werden will I, 1086. Nicht etwa nur eine Strafe, sondern Sünde I, 131. ff. 816. f.; verdammlich II, 1523.; die größte und schwerste Sünde I, 77.; man soll sie nicht verringern, sondern groß machen I, 174. ff. 575. f.; steckt tief in unserer Natur II, 1524.; ein großer Schade II, 1525. f.; ihre Größe erkennt man an den Früchten II, 1527., Adams Furcht I, 209. und aus Betrachtung der verlornen Gaben I, 173. Klebt \gesperrt{Allen} an II, 1143., auch den Wiedergebornen I, 576. 816.; II, 148. 1496., bis an den Tod II, 1486.; die Papisten lehren, sie werde durch die Taufe aufgehoben, und bleibe nur noch Schwachheit übrig II, 1464. f., das ist falsch II, 1523. f. 1527. f. Die E. will nicht Sünde sein II, 1479. 1523. 1526., nicht offenbar werden II, 1514. 1517. 1714., will schön sein II, 1524., ist listig sich zu schmücken II, 1531.; ist fein abgemalt in den Brüdern Josephs II, 1523. Ist der Ursprung aller Sünde I, 630. f.; wie sie allmählich hervorbricht I, 583. f.; wie sie sich in der Jugend und in späteren Jahren zeigt I, 575.; aus ihr kommt Unglaube I, 1552. f.; II, 146. f., Stolz und Vermessenheit I, 631., Verzweiflung I, 809., Mißbrauch der irdischen Gaben und göttlichen Segens I, 852. f. 855. f., Verachtung des Ehestandes und Neigung zu Hurerei I, 144. 205. f. 975. f.; II, 486. f.; ohne Erkenntniß der E. kein Verständniß der Schrift II, 1523. Eine \gesperrt{Strafe} derselben ist Regent sein II, 1409. ff. Ihr kann wohl gewehrt, aber sie kann nicht ausgerottet werden I, 585.; wir sollen gegen sie kämpfen II, 1143.; das thun die Gläubigen, die Gottlosen folgen ihr I, 576. Warum Gott die E. zugelassen habe, ist eine vorwitzige Frage II, 1526. — S. \gesperrt{Fleisch} 2, 3, \gesperrt{Sünde, Wille} II, 2.

Erde. 1. \gesperrt{Die Erde}. Das Centrum der Schöpfung I, 42. f.; heißt „trocken“ I, 41. f.; warum nicht vom Meere überfluthet I, 41. ff. 545. ff.; zum Wohnsitz der Menschen zubereitet I, 43. ff.; vor dem Fall überall frisch und fruchtbar I, 251.; verflucht um der Sünde willen I, 94. f. 249. f. 355.; worin der Fluch besteht I, 125. 250.; er ist durch die Sintfluth vermehrt I, 250. ff., nimmt mit der Sünde zu I, 252.; die E. soll bebauet werden Menschen und Vieh zum Nutzen I, 1479.; wird einmal aufhören I, 587. — S. \gesperrt{Himmel} 5., \gesperrt{Schöpfung, Sintfluth}.

2. \gesperrt{Der Boden}. Die röthliche Erde soll die beste sein, darnach die schwarze und gelbe II, 69.

Erhaltung. Schaffen und Erhalten Ein Ding I, 1539.; was Gott erschaffen, kann er auch erhalten und vermehren I, 63. 1540. f.; Gott erhält die Creaturen durch sein Wort I, 61. f. 91. f. 1480. f.; II, 370.; erhält uns so jeden Augenblick I, 545.; E. ist nöthig I, 61. f.

Erinnerung. E. vergangener Leiden thut wohl II, 1779.

Erisapfel. Griechische Fabel I, 706.

Erkennen. „Ein Weib erkennen“ I, 294. f.

Erkenntniß. 1. E., \gesperrt{die Gott hat} in der Ewigkeit und in der Zeit I, 1535.

2. E. \gesperrt{von Gott}. Findet sich auch bei den Heiden I, 1058.; die E. Gottes ist rechte Weisheit II, 1688.; welches die rechte E. Gottes sei II, 1674. f.; wir sollen Gott erkennen, nicht wie er in sich ist, sondern wie er sich offenbart hat I, 1551.; II, 177. f., aus dem Wort, nicht durch Speculation I, 1063. f.; aus der rechten E. folgt der rechte Gottesdienst I, 1058.

3. E. \gesperrt{der Sünde}. Kommt durch das Gesetz II, 1494.; nöthig und wichtig II, 1491. 1523. ff., der Anfang des Heils I, 482. f., aber allein nicht genug zur Seligkeit II, 1491.

4. E. \gesperrt{der Worte und Sachen}. I, 998. f. 1359. Diese zweifache E. macht zu einem guten Lehrer I, 1000.; die E. der Sachen ist die wichtigere I, 998. 1002. 1380. f., sie ist E. und Verstand des N. T.’s I, 1004.

5. \gesperrt{Wachsthum in der} E. Worin es bestehe II, 1334.; ist nöthig II, 77. ff. 740. 1334. — S. \gesperrt{Abraham} 10, \gesperrt{Isaak} 6, \gesperrt{Jakob} 6, \gesperrt{Väter} II, 3.

Erlösung. S. \gesperrt{Christus} 3.

Erneuerung. Wie sie geschieht II, 1385. — S. \gesperrt{Ebenbild}.

Ernst. S. \gesperrt{Churfürsten zu Sachsen}.

Erregung. Bei plötzlicher E. erstarrt das Herz I, 1104. — S. \gesperrt{Bestürzung}.

Erscheinungen. 1. \gesperrt{Allgemeines}. Es gibt E. I, 1120.; prüfe sie nach Gottes Wort I, 1120. 1122.; II, 782. f.; viel trügerische im Papstthum I, 1119. 1121. 1528.; II, 1047.; im N. T. auch E. II, 145. f.; in den Gnadenmitteln haben wir herrlichere als Abraham I, 1119. ff., als das jüdische Volk II, 1842., es gibt keine besseren I, 1118. f. 1527.; andere neue E. mit diesen nicht zu vergleichen I, 1120.; wir sollen keine anderen begehren I, 1120. f.; Augustin, die Märtyrer, Luther begehrten keine I, 1526. f. In welchem Falle die E. in der Schrift als E. von Engeln auszulegen sind II, 775. f.

2. E. \gesperrt{Gottes}. In Traum und Gesicht II, 1145.

3. E. \gesperrt{der Engel}. Oft erschienen I, 991.; II, 723., in weißen Kleidern II, 1038.; redeten und handelten wie Menschen I, 1167. f.; nehmen die Gestalt an sich I, 991., nicht immer dieselbe Gestalt I, 992.; wer der Hagar erschien I, 991., in welcher Gestalt I, 1439.; wer die Drei waren, die Abraham erschienen I, 1134. ff. 1198.; E. bei Isaaks Opferung I, 1522. ff.; E., die Jakob hatte II, 723.; im N. T. nicht so häufig I, 1525.; Kennzeichen guter und böser Engel I, 1529. ff. 1532.

Vgl. \gesperrt{Abraham} 8, \gesperrt{Gespenster, Isaak} 3, \gesperrt{Jakob} 5.

Erschrockene. Nutzen nicht, bis sie hören, was ihre Sünde sei, damit sie Hülfe finden mögen I, 1347.; Gott verwirft sie nicht I, 1437. f. — S. \gesperrt{Gewissen} 2.

Erstgeborne. Heißen des Vaters erste Kraft II, 1923. f.; oft von Gott verworfen I, 298. f.; II, 987. 1862. 1881. ff.

Erstgeburt. Bei Juden und Heiden hochgehalten I, 298.; II, 1924.; Rechte der E. II, 92. f. 102. f. 1853. 1861. f. 1886. 1924., in Frankreich und Sachsen II, 1886.; gilt wohl im weltlichen, nicht in Christi Reich II, 1925.; über die E. oft Haß entstanden II, 1110. 1506.; was E. im geistlichen Sinne sei II, 1506.

Erwählung. Gott erwählt aus dem ganzen menschlichen Geschlecht Etliche, die er selig macht II, 1952. — S. \gesperrt{Beruf} 4, \gesperrt{Versehung}.

Esau. 1. \gesperrt{Esau}. Ist ein \gesperrt{Bild} der falschen Kirche II, 118., aller Juden II, 320. f., solcher, die falsche Buße thun II, 317., der Werkheiligen und Heuchler II, 99. Von seinen drei \gesperrt{Namen} II, 101. f.; von dem Namen Esau II, 75. War \gesperrt{röthlich und rauh}, was das bedeutet II, 76. f. Hatte große \gesperrt{Macht} und Gewalt, da er Jakob entgegenzog II, 765. f. \gesperrt{Sünden} E.’s.: gerieth als Jäger in viele Laster II, 83. 85.; war ein Heuchler II, 87. 90. 257.; hoffährtig II, 81. ff. 85. 240. f. 307. ff., und ward deshalb verworfen II, 81. 242. 382. f.; irdisch gesinnt II, 88. ff. 102. f. 115. f.; unbeständig und tückisch II, 87. f. 97. f. 115. f. 121. f. E.’s \gesperrt{Buße}, da er den Segen mit Thränen suchte, war nicht aufrichtig II, 313. ff. 386. f. \gesperrt{Tugenden und Lob} E.’s: war ein großer, frommer, tapferer Mann, hatte aus Isaaks Predigten gelernt, seine bösen Begierden im Zaum halten II, 834.; ehrerbietig gegen den Vater II, 982. f. \gesperrt{Werke}: hat von den ihm gewordenen Verheißungen gepredigt II, 826. f.; hielt sich zur rechten Kirche II, 983., und ward wohl selig II, 123. 828. 830. — S. \gesperrt{Isaak} 16., \gesperrt{Jakob} 23.

2. E. \gesperrt{und seine Nachkommen}. II, 983. ff. 997. ff. 998. ff. \gesperrt{Religion}: haben Abrahams Lehre, Opferweise und Beschneidung beibehalten II, 62. 984. 992. 1009. \gesperrt{Sünden und Tugenden}: Die meisten der Nachkommen E.’s. waren gottlos und hoffährtig II, 1013. f.; unter ihren Königen waren etliche heilig, die anderen doch weltlicher Weise fromm II, 1014.; unter ihnen waren viel vortreffliche und herrliche Männer II, 984. 1010. \gesperrt{Seligkeit}: sie waren nicht alle verworfen II, 999. und von der Gnade ausgeschlossen II, 984. ff.; wohl viele selig I, 1032.; II, 65. f. 828.

3. E. \gesperrt{und seine Weiber}. II, 235. ff. 372. f. 987. ff. Die beiden \gesperrt{Hethitischen} Weiber waren stolz II, 240.; herrschten im Hause Isaaks II, 344.; wollten Abgötterei in der Kirche Isaaks einführen II, 244. f. 998.; haben Isaak geschmiert und Rebekka betrübt II, 257., Rebekka verachtet II, 331. f.; wurden wohl mit der Zeit etwas gelinder II, 991. 993. — S. \gesperrt{Ada, Ahalibama, Basmath, Judith, Mahalath}.

Escol. I, 882. ff. Leutselig, fromm I, 886. 888.; II, 1069.

Esel. Nützlich, zum Kriege nicht tauglich II, 1985. f.

Essen. Heißt in der Schrift oft „glauben“ II, 813.

Esther. Hat Viele bekehrt II, 1773.

Euclio. Sprüchwörtlich für Geiz II, 612.

Eunuchen. Was darunter zu verstehen II, 1220. f.

Euphrat. Trennt Mesopotamien und Canaan II, 654.

Eva. I, 140. ff. Ihr \gesperrt{Name} ein fröhlicher und lieblicher I, 269.; Bedeutung desselben I, 383. Ein \gesperrt{heiliges} Weib, voll Glaubens und Liebe I, 296. 395. f.; \gesperrt{irrte}, als sie Kain für den verheißenen Samen hielt I, 296. f.; was sie \gesperrt{gethan und gelitten} I, 417.; mußte viel von den Kainiten leiden I, 395.; wie lange sie wohl gelebt I, 417.; warum die Römische Kirche sie nicht canonisirt hat I, 397. — S. \gesperrt{Adam} 2.

Evangelisten. Geben oft nur die Summa der Gespräche an I, 1165.

Evangelium. 1. \gesperrt{Das Evangelium} war vom Falle Adams in der Welt, bleibt ewig I, 625.; rühmt gemeine Werke I, 527. \gesperrt{Wirkung}: nie ohne Frucht I, 625.; dadurch verlieren wir nur das Böse, erlangen Erlösung I, 235.; ohne das E. gerathen wir nur tiefer in Feindschaft wider Gott I, 220.; sondert den Weizen von der Spreu I, 462. \gesperrt{Kam} zu den Arabern I, 997. \gesperrt{Aufnahme}: unsere Herzen sind hart und fühllos dagegen I, 527. 1105. f.; muß Schuld sein an Aufruhr II, 641. und Theurung I, 527. Vom E. lehren die Papisten nichts Rechtes I, 1386. — S. \gesperrt{Gesetz} 2., \gesperrt{Gnade, Verachtung} 2., \gesperrt{Wort} 3.

2. \gesperrt{Das E. und die weltlichen Reiche}. Es hebt das bürgerliche Recht nicht auf I, 1465.

Ewig. Bedeutet oft eine endliche Zeit I, 873. f.

\section{F}

Faber. Verlästerte Luther und seine Lehre II, 1273.

Fallacia compositionis et divisionis I, 132. ff.

Fallen. Helden der Kirche fallen oft I, 723. 789.

Fatum. Hebt den Glauben auf II, 1313.

Faulheit. Die Welt ist voll fauler Drohnen II, 1244.; Gott segnet F. nicht II, 1238.

Fegfeuer. Schrift lehrt nichts davon II, 1543. 1548. f. 2068.; Entstehung der Lüge vom F. II, 1547. f.; wer dahin nach des Papsts Lehre kommt I, 1761.; aus der Lehre vom F. folgt die vom Ablaß I, 1761. f.; wer im F. ist, ist in der Hölle I, 1762.

Feierkleider. Was Feierkleider sind II, 1739.

Feiertage. Woher die Gewohnheit, an F. besser zu essen, sich reinlicher zu kleiden I, 1598.

Feinde. 1. \gesperrt{Menschen}. Die heidnische Philosophie lehrt, man soll sie hassen I, 801.; thue ihnen Gutes I, 978.; laß dir ihr Heil angelegen sein I, 1254. f.; leide das Unrecht geduldig, doch strafe die Sünde II, 1148. f.; Gott kann sie zu Freunden machen II, 1295.; man kann aus F. Freunde machen durch ernstes Gebet II, 828. f., durch Geduld II, 226. ff.

2. \gesperrt{Der Teufel} ist unser ärgster F. I, 1689.

Feldherren. Wenige recht tapfer II, 1375.; ein christlicher F. kann mit geringer Mannschaft viel ausrichten II, 1375. ff. — S. \gesperrt{Kriegsleute}.

Fenster des Himmels. I, 546. f.

Festungen. Helfen ohne Gott nichts II, 932. f.

Feuchtes. Definition des F. II, 1929.

Feueranbetung. Woher entstanden I, 725. f. 1665.; bei den Chaldäern I, 1665. und Persern I, 727.

Feuerhimmel. Was man davon dichtet I, 34.

Feuerzeichen in der Luft. Meinungen darüber I, 606. f.

Filial. Entstehung dieses Ausdrucks II, 2042.

Finsterniß. 1. \gesperrt{Dunkelheit}. Macht Menschen und Thiere betrübt und verzagt II, 771.

2. \gesperrt{Verfinsterung}. Zeichen zukünftiger Dinge I, 52.

Firmament. S. \gesperrt{Himmel} 2, \gesperrt{Veste}.

Fische. Den Vögeln nicht ungleich I, 59.; jetzt noch kraft göttlichen Worts aus dem Wasser gezeugt I, 66.; die fruchtbarsten Thiere II, 1872.

Fleisch. 1. \gesperrt{Speise}. F. essen kann keine Sünde sein I, 591.; Anfangs aß man nur das F. reiner Thiere I, 592.; Manche essen das F. unreiner Thiere I, 592. f.;
\quellseite{549}
nach der Erlaubniß, F. zu essen, ward der Menschen Leben kürzer I, 712.

2. \gesperrt{Erbsünde}. Was das F. sei I, 460. f.; gilt vor Gott gar nichts II, 71. ff. 77.; ist wider Gott I, 460.; ihm gefallen nicht göttliche, sondern teuflische Verheißungen I, 757. f.; verachtet die reine Lehre I, 461.; hält sich im Glück und Unglück nicht recht I, 809. f.; ist ungeduldig II, 1228. 1323. f.; liebt die zeitlichen, verachtet die ewigen Güter I, 424.; ist hoffährtig II, 71. ff.; läßt bald alle Hoffnung fallen II, 741.; hindert am Gebet II, 752. Auch die Heiligen haben noch F. II, 736. 741. 836. 1457. 1465., trotz der Vergebung der Sünden II, 1457. 1465.; daher die Faulheit und Schläfrigkeit der Christen II, 1492.; F. ist F. auch an den Heiligen II, 380.; sie behalten noch F., damit sie nicht müßig seien I, 678.; leben im F., nicht nach dem F. II, 803.; müssen dawider kämpfen II, 22. 146. f. 201., es im Zaum halten II, 1318.; sie streiten dawider II, 736. 741. ff.; muß immer getödtet werden I, 1350.; II, 1659. — S. \gesperrt{Erbsünde, Gerechtigkeit} 4.

3. \gesperrt{Fleisch und Geist}. Kämpfen stets wider einander in den Gerechtfertigten I, 1670.; II, 526. f. 1871., durch das ganze Leben II, 1317. f., besonders wenn wir im Unglück sind I, 809. f.; das F. hindert die Freude im Heiligen Geist I, 1105. f. — S. \gesperrt{Glaube} 11, \gesperrt{Kreuzigung}.

Fleiß. Gott segnet ihn II, 1238. ff.

Fluch. Folgt der Sünde, insonderheit dem Mord I, 355. f. 594. f.; trifft oft später ein, als man nach der Drohung erwarten sollte I, 643.; das Land wird schlechter und geringer I, 594. f. 852.; II, 1794.; durch die Sünde kam der F. über die Erde I, 95. 121., und ist die Bebauung und Bewahrung des Landes beschwerlich I, 124. f. Unter dem F. sollen wir geduldig sein II, 306. — S. \gesperrt{Erde} 1, \gesperrt{Länder}.

Fortpflanzung. Ist etwas Wunderbares I, 152. ff. und Gutes I, 141. f.; geschieht kraft des göttlichen Worts I, 65. f. 91. f. 155.; geschieht auf besondere Weise bei den lebendigen Körpern I, 64. f. — S. \gesperrt{Zeugung}.

Fragen, vorwitzige. Was davon zu halten I, 13.; Warnung davor I, 1085.; dadurch fiel wohl der Teufel I, 1085.; Beispiele: I, 13. 21. 46. 53. 111. 136. f. 1084. f. 1540. — \gesperrt{Vorwitz}.

Franciscaner. Folgen Franciscus nicht nach, sorgen nur für ihre Küche I, 831.; von Johannes XXIII. verdammt I, 831. f.; warum sie den Geiz der Päbste nicht mehr straften I, 832.

Franciscus, St. Verachtete irdische Güter, bettelte, theilte den Armen mit I, 831.; verbot etwas Eigenes zu haben I, 835.; was von ihm zu halten I, 831.

Frankreich. Fruchtbar II, 1614.

Franzosen. Hoffährtig II, 1614.; verachten die Deutschen II, 1786.

Freiheit. Worin der Mensch F. hat I, 859. f.

Fremdling. Bedeutung des Worts I, 1617. f.; wer ein rechter F. ist I, 1136.

Freude. 1. \gesperrt{Allgemeines}. Man kann vor F. sterben I, 1104.; keine wahre F., außer der, die das Wort bringt I, 1370.

2. F. \gesperrt{im Heil. Geist}. Ist unaussprechlich I, 1104.; ein Vorschmack des ewigen Lebens I, 1106.; man erfährt sie nur durch den Glauben I, 1106.; Abraham hatte sie I, 1103. ff. 1370. und Sarah I, 1378. f.; Christen haben mehr Ursache dazu denn Abraham I, 1105. — S. \gesperrt{Leiden} 2, \gesperrt{Trunkenheit}.

Freunde. Ein treuer F. der allerköstlichste Schatz im ganzen Leben I, 840. f.; vergessen einander nicht I, 1286. f.; zanken sich zuweilen I, 1392.; die besten F. der Christen, die heiligen Engel I, 1689. — S. \gesperrt{Feinde}.

Freundlichkeit Gottes. Beispiel I, 1331. f.

Freundschaft. Kann nicht bei Argwohn sein I, 795.; schwindet oft I, 797.; von der wahren weiß die Philosophie nichts I, 798., sie entspringt aus dem Heiligen Geist I, 798.

Frevel. Bedeutung des Worts I, 504. f. 511.; bewahrt den Schein der Heiligkeit I, 504. f. 506. f.; folgt, wo man Gottes Wort verliert I, 504.

Friede. Nöthig, damit man Gottes Wort treiben könne II, 223.; um des F.’s willen soll man auch Opfer bringen II, 2020.; was „Friede“ in der Schrift heißt II, 1604. — S. \gesperrt{Gerechtigkeit} 5.

Friedrich, Churfürst. Vergütete den Bauern den durch die Jagd angerichteten Schaden II, 86.; traf Vorsorge gegen Theurung II, 1360.; ward des Regierens müde II, 1407.; Ausspruch über den Stand der Ackerleute II, 553. 1407.

Fromm, gerecht. Unterschied I, 499.

Fromme. 1. \gesperrt{Fromme}. Sind nicht zweizüngig, sind entweder nichts oder sie sinds gar II, 255.; gibt allezeit etliche II, 773. — S. \gesperrt{Christ}.

2. F. \gesperrt{und Gottlose}. Die F. \gesperrt{sind} die Kirche, die G. des Satans Schule I, 1193. f.; die F. werden für Sünder, die Sünder für F. gehalten II, 670. 675. 708. 711.; die G. sehen nur auf die Sünden der F. und ärgern sich daran I, 634.; den Nutzen der Creaturen erkennen nur die F. I, 1613. f.; F. und G. sind einander \gesperrt{äußerlich gleich bei innerlicher Verschiedenheit}, wenn sie meinen, Gott schlafe I, 697. f. 1195. f., wenn sie hoffen und fürchten I, 698. f., in ihren Werken und Tugenden II, 461. ff. 472. 479. ff., wenn sie sündigen I, 1326., wenn sie in gleiche Strafe gerathen I, 823. 825. 894. ff., 896.; II, 337. f. 1534. 1571. 1576. \gesperrt{Glück und Unglück}: Der Teufel versucht die F. mehr als die G. II, 1260.; die F. haben ein Reich, aber im Glauben, die G. besitzen die Reiche dieser Welt I, 644.; die G. haben leiblichen Segen, die F. scheinen unter dem Fluch I, 648. Die F. sind im Unglück, die G. im Glück I, 665. ff. 677. f. 860., so scheint es vor der Vernunft II, 861.; so war es stets I, 667.; woher es kommt II, 726. ff., das verstehen die Engel, wir werden es einst auch verstehen II, 731. An dem Glück der G. und Unglück der F. stößt sich die Vernunft I, 807. f. 854. f., ebenso die F. I, 360. f. 667. 807. 853. f., das kommt von der Erbsünde I, 809.; daran sollen wir uns nicht stoßen I, 824.; II, 1013.; wie man dies Aergerniß überwinden soll I, 808. f. 855. 860. und sich trösten I, 667. f. \gesperrt{Von den} F. haben die G. viel Gutes I, 383. 892. ff.; II, 895. ff. 1234. f. 1238. \gesperrt{Verhalten im Glück und Unglück}: Die G. mißbrauchen auch die allerbesten Dinge, die F. können auch die bösen Dinge wohl gebrauchen I, 857.; die G. trösten sich mit den Verheißungen, die sie nichts angehen, die F. erschrecken vor Drohungen, die sie nicht treffen II, 372.; die G. sind sicher, während die F. seufzen und erschrecken I, 933. f. 1195. f. 1238. f.; II, 1263.; die G. erschrecken, während sich die F. freuen I, 701. f.; die F. trösten sich der göttlichen Wohlthaten, die G. vergessen sie I, 1343. f.; verschiedenes Verhalten in der Züchtigung II, 1664.; Sodoms Untergang soll die F. trösten, die G. schrecken I, 1276. \gesperrt{Gegenseitiges Verhalten}: die F. müssen unter G. leben, die sich überheben I, 937. f., undankbar sind II, 1235.; die F. sollen zu den Anschlägen der G. nicht schweigen I, 671.; entschuldigen sich die F. gegen die Anklagen der G., so klagen sie die G. damit aufs allerschwerste an II, 681. \gesperrt{Erfolg und Ausgang}: Die F. sind seliger in ihrem Unglück als die G. in ihrem Glück II, 1085.; den F. dient Alles zum Besten, den G. zum Schaden I, 1326. ff.; die G. thun den F. Gutes, auch wenn sie ihnen Böses thun II, 1150.; das Unglück der F. wandelt sich in Glück, der G. Glück in Unglück I, 648. 855. 860.; II, 1011. ff. 1018. f.

Fruchtbarkeit. S. \gesperrt{Ehestand} 12.

Führung. Es ist besser, Gott führe uns, als daß wir uns selbst führen II, 1635. — S. \gesperrt{Regierung}.

Furcht. 1. \gesperrt{Furcht}. Eine Folge des Falls I, 76. f.; II, 703.; es gibt eine rechtmäßige F. II, 174.

2. \gesperrt{Gott} heißt unsere F., warum II, 721. f.

3. F. \gesperrt{Gottes}. Gott will gefürchtet sein I, 1204.; ist der höchste Gottesdienst I, 1538.; vertreibt alle andere Furcht II, 721.; wo sie fehlt, folgen die greulichsten Sünden I, 630. f. — S. \gesperrt{Fürchten}.

4. F. \gesperrt{Isaaks}. Bedeutet den Gott Isaaks II, 701. ff.; ist Christus II, 703. f.

5. F. \gesperrt{und Hoffnung}. In Christen ist Beides I, 699.; man hat gemeiniglich mehr F. als H. II, 1338.; die Gottlosen spotten Gottes in F. und H. I, 698.

Fürchten. Gott fürchten heißt Gott dienen I, 1493. 1537.; man soll Gott nur in seinem Worte fürchten I, 1537. f. — S. \gesperrt{Furcht} 3.

Fürsorge Gottes. Gott sorgt I, 1511. lange voraus II, 1068., ehe wir sorgen können I, 48., wie eine Mutter II, 1586.; sorgt für unser Bestes auch wider unsern Willen II, 1068.; seine F. geht über Wachende und Schlafende II, 1332. f., über die Christen II, 660. f. 716., über jede einzelne Person seiner Gläubigen II, 1768. f., die Gott wider Welt und Teufel schützt II, 657. 659.; derselben können sich Christen auch im Irdischen getrösten II, 617.; erstreckt sich auf Großes und Kleines I, 1693.; II, 469. ff. 473. 518. ff. 617. 1074. 1769., auf unsere Leiden und Thränen II, 1333. f.; Gott sorgt mehr um uns, als um sich I, 43.; gibt sich kund in der Schöpfung I, 47. f. 87., in der Erhaltung der Menschen und Thiere in der Arche I, 524.; nur der Glaube kann sich derselben trösten II, 617. f.

Fußwaschen. Ein alter Gebrauch, namentlich auch in der Kirche II, 1600. f.; kommt von den Patriarchen I, 1146.; was uns Christus damit lehrt I, 1215.; päbstliches F. ist bloßes Schauspiel I, 1215.; II, 1601. und Heuchelei II, 1601. 1603.

\section{G}

Gaben Gottes. 1. \gesperrt{Art und Natur}. Alles, was wir haben, sind G. Gottes I, 1654. f., auch die leiblichen Güter I, 1710. 1718.; II, 296.; der leiblichen G. bedürfen wir um der geistlichen willen II, 294. ff.; die höchste G. ist, daß man ein Glied der rechten Kirche ist I, 448.; zu den G. des Heiligen Geistes gehört auch das Blut Christi II, 1992. G. erben nicht fort II, 1368. ff.; Gott theilt sie aus, wie er will II, 1371., aus Gnaden, nicht aus Verdienst I, 1418. ff., die allerbesten unter die Gottlosen I, 996.; sie gereuen ihn nicht II, 309. f.

2. \gesperrt{Zweck} der G. Wozu sie Gott gibt I, 299. f.; daß wir Andern dienen I, 722.; je mehr Gaben, desto mehr dienen I, 985.

3. \gesperrt{Erkenntniß} der G., \gesperrt{Gebrauch und Mißbrauch}. Man soll die G. \gesperrt{erkennen} I, 985.; man verachtet sie, weil sie alltäglich sind II, 2049.; wie man sie an Menschen loben soll I, 759. \gesperrt{Gebrauch}: daß wir Gottes Willen dienen II, 1390. f., sie zu Gottes Ehre und des Nächsten Nutzen anwenden II, 263. f., Gott dafür danken I, 1718., uns derselben nicht überheben I, 448. 911.; II, 819., auch nicht der G., daß man ein Glied der rechten Kirche ist I, 448. f.; daß wir der leiblichen G. genießen I, 1710. 1714. \gesperrt{Mißbrauch}: Die G. gerathen den Gottlosen zum Schaden I, 1330.; sie mißbrauchen die G. I, 1328. f., schmähen und verspotten Gott I, 299. f.; man überhebt sich im Herzen I, 985.; je mehr Güter, desto größer die Hoffahrt I, 438. 923.; so überhoben sich die Menschen vor der Sintfluth I, 441.; solche Hoffahrt ist teuflisch I, 985.; Gott straft den Mißbrauch I, 300. 437. f. — Gott kanns nicht unterlassen, uns mit G. zu zieren, wir nicht, uns derselben zu überheben I, 440. f. 477.

Gabriel. Ableitung des Namens I, 1013.; Bedeutung desselben I, 1198.; welcher Engel so geheißen haben soll I, 1198.

Gad. 1. \gesperrt{Jakobs Sohn}. Bedeutung des Namens, Anlaß dazu II, 562. f. Jakobs Segen II, 2030. ff. — S. \gesperrt{Jakob} I, 19.; II.; \gesperrt{Joseph} 12.

2. \gesperrt{Der Stamm}. II, 2030. ff.

Galater. Waren Deutsche I, 657.

Galgenreue. Was G. ist II, 313. ff. 336. 663. 1467. — S. \gesperrt{Gewissen} 3.

Galiläa. Fruchtbares Land II, 2034.

Gallen. Verschnittene Priester der Cybele II, 1222.

Garizim. Berg, Lage I, 779.; darauf ein Tempel I, 779.

Gast. Sein wird man bald überdrüssig II, 482.

Gastfreiheit. \gesperrt{Ruhm und Lob}: übertrifft alle Werke der Carthäuser I, 1162. f.; wo sie geübt wird, übertrifft ein Hüttlein einen Palast I, 1158.; es lobt sie Christus I, 1213., der Hebräerbrief I, 1133. \gesperrt{Pflicht}: Christen sollen G. üben I, 1133. f. 1136., wenn auch Landstreicher sie oft mißbrauchen I, 1139. \gesperrt{An wem} zu üben I, 1720. f., nicht bloß an Eltern, auch an Brüdern I, 1158., vornehmlich, doch nicht ausschließlich an des Glaubens Genossen I, 1140. ff.; welches die rechten Fremdlinge, an denen sie zu üben I, 1135. f. \gesperrt{Wie} zu üben: mit Demuth und Ehrerbietung I, 1142. f. 1148.; gern und mit dem Besten I, 1156. f.; mit Vorsicht II, 1601. f.; Viele stellen sich nur gastfrei I, 1215. f. \gesperrt{Warum} zu üben: man nimmt dabei Gott auf I, 1143. f. 1721., Gottes Sohn I, 1134. 1146. 1159.; wer nicht gastfrei ist, verstößt den Sohn Gottes I, 1134. G. ist \gesperrt{selten}: es mangelt daran I, 1159.; nicht so häufig wie ehedem I, 888.; wir sind träge dazu I, 1145., woher das kommt I, 1147. 1154. f. 1158. f. \gesperrt{Wo} sie sich findet: an allen Orten, da die Kirche ist I, 1133., da ist viel Gelegenheit I, 1135.; nur in der Kirche I, 1138. 1146. f., nicht im Papstthum I, 1138., nicht in der Welt I, 1146. f. Wie man sich \gesperrt{trösten} soll, wenn G. übel belohnt wird I, 1721.

Gastmähler. Man veranstaltet Gastmähler in Glück und Fröhlichkeit I, 902. 1381. ff. 1154. ff.; sind recht, um die Herzen fröhlich zu machen I, 1383. f.; angestellt, den Ehestand zu fördern I, 1682.; es soll nicht stillschweigend dabei zugehen I, 1159. f.; man hüte sich vor Völlerei I, 1715. — S. \gesperrt{Hochzeitsmahl}.

Gaza. Lage I, 678.

Gebären. „Auf den Schooß gebären“ II, 544.

Geberden. Zorn und Liebe offenbaren sich darin I, 319.

Gebet. 1. \gesperrt{Gottes Wille}. Wir sollen beten I, 1110., Gott will gebeten sein I, 1267.; II, 903. 918.; Christus hat uns eine Weise zu beten vorgeschrieben I, 1267.; wir sollen, weil dies der geordnete Wille Gottes ist I, 1268.; warum Gott will, daß wir beten II, 760. f.; G. ist Gott
\quellseite{550}
ein angenehmer Dienst und Opfer I, 1332. 1334. 1699.; II, 755.; Unterlassung des G.’s erzürnt Gott I, 1272. f.

2. \gesperrt{Was es sei}. Was beten heißt I, 1112.; G. ein Stück des prophetischen Amts I, 1366.; vom Wesen des G.’s wissen die Papisten nichts I, 1366.

3. \gesperrt{Theile} des G.’s. I, 1268. f.; II, 756.

4. \gesperrt{Eigenschaften} des rechten G.’s. Wie das rechte G. beschaffen sei II, 28. f.; wann es vollkommen und stark ist II, 575. ff. Eigenschaften des rechten G.’s II, 760.; man muß sich an die Verheißung halten II, 57. 756.; im Glauben beten I, 1110. 1364.; II, 31.; beherzt beten I, 1203.; glauben, man sei erhört I, 1704.; im Namen JEsu I, 1705. f., ohne Murmeln und Zweifel I, 1271. G. soll sein demüthig II, 759. f., andächtig II, 31. f., anhaltend I, 1204.; II, 28. f. 60. 818. 1317.; wann wir anhalten sollen I, 1702.; II, 47. f., wie lange II, 794., solch Anhalten gefällt Gott wohl I, 1204.; II, 58. 796. 818.; G. soll brünstig sein II, 1643. f., ist nicht immer gleich brünstig II, 1644. ff. 1653., wie man Inbrunst erwecken kann II, 1645. f. Das rechte G. fragt nicht nach viel Worten, macht viel Seufzens II, 756.

5. \gesperrt{Nothwendigkeit} des G.’s. Ist nöthig in allen Ständen II, 277. ff., sonderlich im Regiment II, 276. f.; uns nöthig, Gott ein Dienst I, 1332. 1334. 1699.; nöthig, weil der Satan die Menschen von Sünde zu Sünde treibt I, 1226.

6. \gesperrt{Kraft und Würde} des G.’s. G. hat eine große Kraft I, 1705. f.; II, 64. 577. 817., besonders das gemeinschaftliche I, 1705.; richtet große Dinge aus I, 1547.; II, 29. ff.; dadurch überwinden wir Gott II, 952., unsere Feinde I, 1581. f.; II, 828. f. 1645. 1647. 2016. f., lindern wir die gemeine Noth II, 818.; das schwache G. soll man nicht verachten II, 58.; von der Kraft des G.’s wissen die Papisten nichts I, 1366. Rechtes G. unbezahlbar I, 1363. f.; wir verstehen die Größe der Dinge nicht, um die wir bitten I, 1109. 1112. f.

7. \gesperrt{Erhörung} des G.’s. Daran nicht zu zweifeln I, 1113.; II, 57. f. G. kann nicht vergeblich sein I, 1361.; Erhörung verheißen I, 1267., ist \gesperrt{gewiß} II, 57. f. 794. 918. f.; Gott erhört I, 1269. 1288. f. 1319. 1362. f.; II, 58. 61. 64. 1644. 1652. f., sobald wir beten I, 1703. f.; II, 5. 1644. f.; Gott erhört auch der Frommen Verlangen und Wünschen I, 1697.; Gottes rechter Titel ist: Erhörer des Gebets I, 1109. \gesperrt{Art und Weise} der Erhörung: Gott gibt mehr, denn wir bitten I, 1107. ff. 1111. ff.; II, 818., über Bitten und Verstehen I, 1705. 1711.; II, 61. f. 574. 577. 1328. f. 1381. f., wir bitten immer zu wenig II, 818. f.; Gott erhört auf andere Weise, als wir hoffen I, 1111., gibt etwas Besseres I, 1272. f. \gesperrt{Verzug} der Erhörung: Gott scheint manchmal zu schlafen I, 700.; durch Verzug soll unser G. kräftiger und brünstiger werden II, 575. ff.; rechtes Verhalten beim Verzug II, 576. f.; wie zu trösten, wenn man sieht, daß das G. für die Welt vergeblich ist I, 1395. — S. \gesperrt{Hülfe}.

8. \gesperrt{Um was} man beten soll. Auch um leibliche Dinge II, 29. f.

9. \gesperrt{Zeit, Ort und Weise} des G.’s. Wir sollen beten \gesperrt{ohne Unterlaß}, und wie das geschieht I, 1365. f., bei allem, das wir anfangen I, 1700.; II, 1299., in Noth II, 750. 754. 920. Ein einsamer \gesperrt{Ort} eignet sich zum Gebet II, 771.; warum man den aufsucht II, 771. \gesperrt{Weise}: man bitte in leiblichen Dingen mit Bedingung I, 1702.; wie man dabei Gott Zeit, Stätte und Person vorschreiben darf I, 1695. ff. 1702.; G. geschieht auch ohne äußerliche Geberde II, 1299.

10. \gesperrt{Hindernisse und Ermunterungen}. Das Gebet \gesperrt{hindert}: Satan I, 1704. f., unser Fleisch II, 752., unsere Geschäfte und mancherlei Sorgen II, 1653.; haben dabei immer mit Unglauben zu kämpfen I, 1109.; das G. hindert ein böses Gewissen II, 1141., Haß II, 1141. f., wenn man in Sünden liegt II, 1194.; vom G. soll uns nicht abschrecken die Größe der Noth und erbetenen Wohlthat I, 1108. ff., auch nicht unsere Unwürdigkeit und Gottes Majestät I, 1110. 1271. f. \gesperrt{Ermuntern} sollen uns die Exempel anderer Beter I, 1704. f.

11. \gesperrt{Wer} allein beten kann. G. ist eine große Kunst I, 1109., schwere Arbeit I, 1111. f. 1272.; II, 30. f.; beten kann nicht: wer das Wort nicht hat I, 456.; II, 1644., ein Gottloser I, 494., sondern nur ein Prophet I, 1366., ein Gläubiger II, 2014. f.; das G. der Papisten ist ein Spott II, 58., Gott ein Greuel I, 1364.; das G. der Mönche taugt nichts II, 756., ist eine Lästerung Christi I, 1364. f. — S. \gesperrt{Schöpfer}.

12. G. \gesperrt{und Arbeit}. Bete und arbeite I, 1270.

13. G. \gesperrt{und Predigt}. Gehören zusammen I, 838.

14. G. \gesperrt{und Vorherbestimmung} Gottes. Bete, obwohl Alles geschieht, wie Gott beschlossen hat II, 47. f.

15. G. u. \gesperrt{Wort} Gottes. G. folgt dem Wort II, 137.

Vgl. \gesperrt{Beter, Glaube} 7, \gesperrt{Jakob} 10, \gesperrt{Mönche} 2, \gesperrt{Seufzen, Vater Unser}

Gebote Gottes. 1. \gesperrt{Allgemeines}. Unser Haus, Leib, Herz ist voll G. I, 1501.; Gott gebietet oft geringe, lächerliche und ärgerliche Dinge I, 527., da sehe man nicht auf das Gebotene, sondern den Gebietenden I, 527., der macht alles Gebotene groß I, 528. ff. 1497. f.; nicht jeder Befehl Gottes gilt allen Menschen I, 568. \gesperrt{Verhalten}: man soll nichts davon thun noch dazu thun I, 526. f. 847. 1495. f.; ohne Gottes G. nichts vornehmen I, 567., mit demselben mag mans frei wagen I, 1499.; vom Gehorsam gegen dieselben I, 1261. ff.; wir sollen Gehorsam leisten ohne Verzug I, 1413. f. 1417. 1497., ohne zu fragen warum? und wozu? I, 1125. 1127. f. 1260. f. 1496. f., welcher Vorwitz ins Verderben stürzt I, 1125., ohne Bedenken, und wo das Vermögen fehlt, doch zum Gehorsam bereit sein I, 1261. f. 1497. f.; gehorchen, wie närrisch und ärgerlich das Gebotene auch sein mag I, 1129., wenn wir auch verspottet werden I, 1131.

2. \gesperrt{Die Zehen Gebote}. Die zweite Tafel muß der ersten weichen I, 1421.; II, 266. 640. ff. 645. f.; die erste Tafel soll man zuerst treiben, dann werden die Werke der zweiten nicht fehlen I, 399. ff.; den Sünden wider die erste Tafel folgen die Sünden wider die zweite I, 441. f. 443. ff. 469.; die Sünden wider die erste Tafel bleiben gemeiniglich verborgen unter dem Schein der Heiligkeit, nicht so die wider die zweite Tafel I, 503. f. \gesperrt{Inhalt}: die darin geforderten Werke sind nicht leicht I, 1180.; die G. verdammen auch die böse Lust im Herzen I, 1348.; Inhalt der ersten Tafel I, 400. f.; vom 1. Gebot I, 1057. f. 1452. 1517. f.; II, 8.; vom 2. Gebot I, 1402. ff.; vom 4. Gebot I, 1735.; II, 329. f.; vom 5. Gebot II, 1115. 1127. f.; vom 6. Gebot I, 1323.

Vgl. \gesperrt{Gesetz} 1. 2.

Geburt. 1. \gesperrt{Die leibliche} G. des Menschen ist ein großes Werk I, 1371.; II, 16., trotz der anhaftenden Sünde I, 1374. f.; jetzt noch so wundersam wie einst Isaaks G. I, 1370.; der Heilige Geist hat Lust daran II, 1214., die Heiden hatten sonderliche Hochachtung davor II, 1214., wir haben keine Freude daran II, 1214. f., verachten sie I, 1371. ff.; Beispiele gefährlicher Geburten
II, 68. ff. 1216. f.; auf die leibliche G. und Herkunft stolz sein ist thöricht I, 312. f.; II, 1368. f. 1371. — S. \gesperrt{Zeugung}.

2. \gesperrt{Fleischliche und geistliche} G. Es ist zweierlei G.: nach dem Fleisch und nach der Verheißung II, 33. ff. 65. Der fleischlichen soll man sich nicht rühmen I, 1031. f.; der natürliche Mensch verläßt sich auf die fleischliche G., um selig zu werden I, 1430. f.; darauf verläßt sich die falsche Kirche II, 35., Türken, Juden und Papisten II, 36. ff.; gilt nichts vor Gott I, 328. 1132. 1430. f.; II, 70. ff. 1177. ff., nützt nichts zur Seligkeit I, 1030. f. 1398. f. 1427. f. 1430. f.; II, 23. ff. 66. f., ist ohne die geistliche G. schädlich zur Seligkeit II, 43. f.; auch die fleischliche Verwandtschaft mit Christo hilft nichts zur Seligkeit II, 1201.; Gott hat aus allen Völkern Etliche angenommen II, 985. f.; das Pochen auf die fleischliche G. hat die Juden um Leib und Seele gebracht I, 1032. — S. \gesperrt{Mensch} 2.

Geburtstag. Feier desselben I, 1382.

Gedanken. G. nennt die Schrift oft Worte I, 1696.; \gesperrt{eigenen} G. sollen wir nicht folgen II, 211. f.; \gesperrt{böse} G. können wir nicht hindern, aber man soll sie nicht ausführen II, 785. f.; Zustand eines in Gedanken Versunkenen II, 290. f. 500. f. — S. \gesperrt{Bestürzung, Erregung}.

Gedud. Bild in Mansfeld II, 2031.

Geduld. 1. G. \gesperrt{Gottes}. Zeigt sich darin, daß er uns täglich Sünden vergibt I, 1553., daß er den Gottlosen noch Zeit zur Buße gibt I, 965. f., die Sünder trägt II, 729. — S. \gesperrt{Langmuth}.

2. G. \gesperrt{der Kinder Gottes}. Ist eine der Vernunft verborgene Tugend II, 27.; was die Philosophie von G. lehrt II, 1666.; durch G. sündigt man nicht II, 1152.; trotz aller G. soll man die Sünde strafen I, 986.; II, 231. ff.; G. muß aufhören, wenn man Gott verlieren würde II, 243. f.; läßt nicht zu Schanden werden I, 827., trägt reiche Früchte II, 1421. 1430. f., hat eine allmächtige Kraft, macht aus Feinden Freunde II, 226. ff., Gott belohnt sie herrlich II, 1381. f.; soll nicht in Ungeduld umschlagen II, 1314.; an G. soll sich gewöhnen, wer Gewalt zu strafen hat II, 1152.; G. sollen wir an der Väter Exempeln lernen II, 380.

3. G. u. \gesperrt{Bekenntniß}. Soll bei einander sein I, 986.

4. G. u. \gesperrt{Hoffnung}. Eigentlich dasselbe II, 1293.

Gefahr. Man soll sich nicht G. auf den Hals laden, wenn man auch Verheißung hat II, 637. f.; darin soll man sich schützen, so gut man kann I, 764., nicht verzagen, wenn auch keine Hoffnung mehr ist I, 615.

Gefängniß. Die Könige hatten G., über die sie ihre Edelleute setzten II, 1296. Die Juden wußten, wie lange das \gesperrt{babylonische} G. währen sollte I, 961. f.; in ihrem Gefängniß klagten sie vergeblich I, 989. f.

Gefühllosigkeit. Ist Raserei, nicht Stärke II, 944. f., unnatürlich und unmenschlich II, 1608., unchristlich II, 945. 948. 1500. f. 1573., von Gott gehaßt und verdammt II, 1572. 1625., teuflisch II, 1679.; ward von den Mönchen in Egypten erdacht, vom Pabst bestätigt II, 1679.; Exempel ihrer G. II, 1608. f.; wie man der G. der Mönche begegnen soll II, 1608. — S. \gesperrt{Lieblosigkeit, Mönche} 2.

Gegensätze. Es gibt mancherlei I, 267. f.

Geheimnisse der Schrift. Keine Schande, wenn wir etliche derselben nicht wissen I, 536. — S. \gesperrt{Widersprüche}.

Gehenna. Ursprung des Worts II, 1550. ff.; so nennen die Scholastiker die Hölle II, 1550.; ist der Ort der Gottlosen II, 1555., verschieden vom Scheol der Gottseligen II, 1555. — S. \gesperrt{Hölle}.

Gehorsam. \gesperrt{Art und Natur}: Wort und G. gehören zusammen I, 767.; Wort und Beruf macht allein rechten G. I, 764.; der rechte bleibt beim Wort, thut nichts davon, nichts dazu I, 526. f. 1493.; thut, was Gott befohlen hat I, 761. f., und zwar ihm befohlen hat I, 762.; fragt nicht lange warum? I, 1125. f.; langsamer G. ist kein G. I, 1503.; G. den Menschen geleistet, ist Gott geleistet I, 1725. f.; rechter G. ist der G. gegen die Eltern I, 762. und Obrigkeit I, 763.; der äußerliche G. muß dem innerlichen folgen, nicht umgekehrt I, 760.; G. ohne Gottes Gebot ist Teufelsgehorsam I, 764. f. 767. G. preist Gottes Weisheit und Güte I, 529.; ist tröstlich im Unglück I, 763. f.; von Gott gesegnet I, 1726.; Gott ein angenehmes Opfer I, 1523. G. beweisen Christen, indem sie sich willig in Leiden schicken I, 375. ff.; am rechten G. mangelt es I, 1726.; der Teufel ist ihm feind I, 1126. Im \gesperrt{Pabstthum} viel Gerede von G., doch wußte man nicht, was G. ist I, 761.; der vom Pabst geforderte selbsterwählte ist schändlich I, 1523., führt zur Hölle I, 764. f.; wie dem Pabst zu antworten, wenn er G. fordert I, 765. f.; der Mönche Unterscheidung zwischen G., der etwas, und G., der wenig von sich selbst habe I, 1124. — S. \gesperrt{Eltern} 2, \gesperrt{Mönche} 2, \gesperrt{Pabst} 1.

Geist. 1. \gesperrt{Der Heilige} G. Kommt nicht ohne das \gesperrt{Wort} II, 263.; sein \gesperrt{Amt} ist: die Welt strafen I, 458. 482., lebendig machen I, 12.; ist den Sündern nahe, um sie zu heiligen II, 1202.; kann allein unsere Unart und Schaden heilen I, 483.; lehrt allein die Schrift recht verstehen I, 452.; wirkt wahre, beständige Freundschaft I, 798.; macht sanftmüthig und demüthig I, 1420.; treibt zu ehrlichen, gottseligen Werken II, 483.; bewegt die Herzen durch äußerliche Dinge: Wort und Ceremonien II, 293.; wird betrübt, wenn sich die Frommen über die Welt betrüben I, 454.; des Heiligen Geistes Gabe kann man an äußerlichen Geberden merken II, 1286. f.

2. G. \gesperrt{Christi}. Nöthig die Schrift zu verstehen I, 456.

3. G. \gesperrt{der Gnade}. Was er sei I, 456.; daran hängt der Geist des Gebets I, 456.

4. G. \gesperrt{des HErrn}. Der Heilige G., nicht der Wind I, 11.

5. \gesperrt{Gewisser} G. Was das sei II, 1353. ff.; nöthig II, 1931.

Vgl. \gesperrt{Fleisch} 3.

Geister. Wie man sie prüfen soll I, 1531.

Geistliche. Wen die Papisten so nennen und wer wirklich ein G. ist I, 1620.

Geiz. Der scheußlichste aller Greuel II, 612., eine Welt aller Schande und Laster II, 614., ein schreckliches Ungeheuer, das keine Ruhe läßt II, 654., Thorheit I, 1627., argwöhnisch II, 686.; was er am Menschen ausrichtet, zeigt Laban II, 508., macht ihn zum Unmenschen II, 595. 597. 627. f., ist Einzelnen und der ganzen Gemeinde schädlich II, 614.; wie Gott ihn straft II, 586. ff., Neid ist die gerechte Strafe II, 612.; haftet Allen an I, 1627.; davor muß man warnen, wenns auch vergeblich scheint II, 588. 590. f.

Geizhälse. Sind wie silberne und goldene Götzen, die keine Empfindung haben II, 597. f., eine schädliche Erdlast II, 598. 610.; führen den Namen Gottes, aber unnützlich II, 582. f., gönnen sich selbst nichts II, 612., thun nichts Gutes, außer wenn sie sterben II, 598., gehen nur mit List und Betrug um II, 588., gönnen dem Nächsten
\quellseite{551}
nicht, was ihm mit Recht gebührt II, 610. ff.; werden dem Gerichte Gottes nicht entfliehen II, 592., ihr Reichthum ihr Verderben II, 592., unglücklich bei all ihrem Gut II, 612., behalten nicht, was sie zusammenscharren II, 205., ihre Kinder darben und betteln II, 590. f.; von allen Völkern verdammt II, 610.; man soll sie nicht zur Communion lassen II, 610.

Geld und Gut macht stolz I, 440. f., macht, nach d. Sprüchwort, edel und schön II, 1787.; G. vergraben eine alte Gewohnheit, so ist viel G. verloren II, 1597.

Gelegenheit. Ist wahrzunehmen I, 1262. f.; II, 98. f.

Gelübde. G. des Glaubens und der Liebe II, 452. ff.; G. heißen in der Schrift Danksagungen II, 454.; stehen frei II, 453. 455. f.; wann G. ein natürliches und sittliches Gebot werden II, 453.; G. soll man halten II, 450. f. 454. f.; „Gott ein G. thun“ II, 450. ff. — S. \gesperrt{Jakob 9, Klostergelübde}.

Gemeinschaft der Eigenschaften in Christo. Was das sei II, 398.; an diesem Artikel ärgert sich Welt, Vernunft und Teufel II, 398. f.; diesem Geheimniß ist der Teufel feind II, 403. f.

Generatio aequivoca. I, 66.

Genugthuung für die Sünde. Das Wort „Genugthuung“ sollte man aus der Kirche thun II, 1455.; G. ist uns unmöglich II, 1463.; Lehre der Scholastiker II, 1456.

Georg. 1. \gesperrt{Schutzheiliger} der Bauern II, 1919.

2. \gesperrt{Herzog} G. Befahl, man solle in seinem Lande die papistische Religion halten II, 652.; sah voraus, daß seine beiden Söhne vor ihm sterben würden II, 1311.; Manche meinen, er habe sich noch vor seinem Ende bekehrt II, 116.

3. \gesperrt{Magister} G. Meuchlings ermordet II, 1108.

Ger. Judas Sohn II, 1177. f.; mit dem 12. Jahre Ehemann II, 1166. 1177.; seine Bosheit war groß II, 1180. — S. \gesperrt{Juda} I, 2.

Gerar. Lage I, 1306.; II, 199.; nahm Gottes Wort an I, 1340. f.; Isaaks Geburtsort I, 1358.

Gerecht. „Gerecht vor Gott“ I, 537., wird man nicht durch Werke I, 620., ist, wer an Christum glaubt II, 1544.; die Gerechten vor Gott haben auch noch Sünde, die sie aber nicht verdammt I, 1670. f. — S. \gesperrt{Fromm, gerecht}.

Gerechtigkeit. 1. \gesperrt{Im Allgemeinen.} Bei jeder Art G. ist Zuversicht und Vertrauen I, 1575. f.

2. G. \gesperrt{Gottes.} Ist nicht, daß er die Leute verdammt, sondern die Elenden schützt, die Schuldigen straft I, 1323. f.; tröstlich, wenn das Gewissen uns schreckt I, 1324.; ist die G., dadurch wir gerecht werden II, 1487.

3. G. \gesperrt{vor Gott.} Ist zweierlei: die vollkommene durch den Glauben, die unvollkommene durch das Leben I, 1011.; Gott fordert über die fleischliche Geburt eine persönliche G. I, 1451.; ist in Christo zu finden I, 760. 1022. 1025., durch den Glauben erlangt I, 1451., an den Werken erkannt I, 1537. — S. \gesperrt{Abraham} 13, \gesperrt{Gerecht, Heiligkeit, Rechtfertigung}.

4. G. \gesperrt{des Fleisches.} Gilt vor Gott nicht I, 1572. f.; II, 70. ff.

5. G. und \gesperrt{Friede} küssen einander II, 1810. f.

Gerichte. 1. G. \gesperrt{Gottes.} G. a priori und a posteriori, was das sei II, 70. f.; Gott verzieht damit, sieht endlich drein I, 506. f. 540.; sind unbegreiflich, der Vernunft ärgerlich I, 1128. f. 1285.; II, 71. f.; G. und Anklage Gottes kann kein Gottloser oder Teufel ertragen I, 362.

2. \gesperrt{Bürgerliche} G. Wozu geordnet I, 599.

Gerson. Aergerte sich am ehelichen Leben der Väter II, 556.; seine Bücher lese man, doch mit Bedacht I, 1253.; Meinung über: das beschauliche Leben I, 1252. f., unbillige Strenge der Mönche I, 845.

Geruch. „Lieblicher Geruch“ I, 572. f.

Gerücht. S. \gesperrt{Name} 4.

Gerundensis. Fehlt in der Auslegung, hat wohl gute Wort-, keine Sachkenntniß I, 322. f.

Geschenke. Im A.\,T. verboten, im Regiment schädlich I, 1473.; ob man sie machen oder annehmen darf, hängt von der Absicht ab I, 1473. f.; G. versöhnen II, 766. 1581. ff.; G. gemacht haben Elieser I, 1694. 1711. ff. 1738., Jakob II, 766. ff. 832. ff. — S. \gesperrt{Brautgeschenke, Hochzeitsgeschenke}.

Geschichten. \gesperrt{Abdruck} von fünf Fingern auf einem Altar II, 790. \gesperrt{Agatha}, die Märtyrerin II, 1994. \gesperrt{Amiris} flieht aus Sibaris I, 1258. f. \gesperrt{Anastasia} vor dem Richter II, 1982. f. \gesperrt{Anfechtung:} wie eine Nonne geistliche A. überwand II, 182.; eine angefochtene meint, daß sie nicht glaube II, 179. \gesperrt{Antonius} zu einem Gerber gewiesen I, 1181. \gesperrt{Aussteuer} von 30 Gulden II, 206. \gesperrt{Bauer} spitzt die Reifen seiner Räder zu II, 275. f. \gesperrt{Bischofsamt} will ein Mönch nicht annehmen I, 763. \gesperrt{Blutschande} begeht eine Mutter mit ihrem Sohne II, 1003. \gesperrt{Dompfaff} zu Erfurt hat 20 Lehen II, 1732. \gesperrt{Ehebrecherin} bekennt auf dem Sterbebette II, 1006. \gesperrt{Ehelosigkeit} der Priester durch Gregor aufgehoben I, 292. f. \gesperrt{Ehestand:} Censorius empfiehlt ihn I, 1176.; ein Mann nennt sein fruchtbares Weib eine Sau II, 532. f. \gesperrt{Ehrsucht:} ein Bischof betet, er möge vor E. bewahrt bleiben I, 923. \gesperrt{Einsiedler} will seinen weinenden Sohn ersäufen I, 1523. \gesperrt{Fegfeuer:} Entstehung der Lüge vom F. und Opfermessen II, 1547. \gesperrt{Gedanken:} Rath böse G. betreffend II, 785. f. \gesperrt{Gefühllosigkeit} der Mönche II, 1608. f. \gesperrt{Geizhals}, der sich selbst nichts gönnt II, 602. \gesperrt{Gnade Gottes:} ob man von irdischem Glück darauf schließen kann I, 1430. f. \gesperrt{Hannibals} Traum II, 1049. f. \gesperrt{Hauptmann} zu Antiochien mißbraucht seine Gewalt I, 820. \gesperrt{Henning}, Dr., als Disputator I, 1178. \gesperrt{Hippocrates} rettet eine des Ehebruchs Angeklagte von Verurtheilung II, 601. \gesperrt{Hirten:} Untreue und Bosheit derselben II, 700. \gesperrt{Iphigeniens} Historie von einem Maler dargestellt I, 342. \gesperrt{Kindesmörderin:} Eine Lacedämonische K. II, 1515.; eine K. verräth sich durch ihr Betragen I, 335. \gesperrt{Klosterneuburg:} todte Kindlein zu K. I, 292. \gesperrt{Leben:} eine Nonne, die ein \gesperrt{beschauliches} Leben führte, vom Teufel verspottet I, 1253. f. \gesperrt{Lucia}, St. II, 505. \gesperrt{Luther} und der Rottengeist, der behauptete, vor dem Fall habe es kein Gesetz gegeben I, 130. f. \gesperrt{Martin}, St., vom Teufel gestraft I, 336.; II, 1515. \gesperrt{Märtyrer} beißt sich die Zunge ab II, 1259. f. \gesperrt{Mechtilde} überwindet in Anfechtung II, 783. \gesperrt{Mißgeburt:} eine Frau gebiert eine Ratte II, 602. \gesperrt{Mönch:} ein M., der nicht arbeiten wollte II, 1244.; ein M. setzt sich Läuse in die Kappe II, 1878. f.; ein M., der im Pabstthum nie von Verheißung gehört II, 1889. \gesperrt{Monica:} Traum von Augustins Bekehrung II, 252. f. 1308. f.; wird erhört I, 1111.; stiftet Frieden unter den Weibern II, 350. \gesperrt{Muttermale:} ein Mann, deß Gesicht wie das eines Todten war II, 602.; Königin gebiert ein Mohrenkind II, 601. \gesperrt{Pabst:} Leos Meinung über Unsterblichkeit I, 1243.; Sixtus gebietet dem Meere zu schweigen I, 1545. \gesperrt{Papisten} suchen eine Königin am Leben zu erhalten II, 1625.; \gesperrt{Riese}, der ein guter Kämpfer war II, 791. f. \gesperrt{Schiffers} Gelübde, dem St. Nicolaus gelobt I, 895. \gesperrt{Staupitz} setzt einen klagenden Prior ab II, 542. f. \gesperrt{Strauchdieb} II, 1263. \gesperrt{Taufe:} ein Arzt tröstet sich derselben und begegnet muthig dem Teufel II, 513. f. \gesperrt{Türken} spotten über in der Schlacht gefallene Christen II, 51. \gesperrt{Undankbarer} bringt seinen Wohlthäter an den Galgen I, 977. \gesperrt{Unsterblichkeit:} ein Arzt wird von derselben überzeugt I, 930. f. \gesperrt{Vater Unser:} Werth Eines V.\,U. I, 1363. f.; wie Einer beim Beten desselben seine Gedanken umher irren läßt II, 32. \gesperrt{Versuchen, Gott:} zwei Brüder, die in fremden Landen reisten II, 1563. \gesperrt{Vincentius}, der Märtyrer II, 1989. \gesperrt{Waldbruder}, der die Sonne nicht ansehen wollte II, 1641. \gesperrt{Wiedertäufer} deuten einen Bogen in den Wolken, dabei eine blutige Hand war II, 423. \gesperrt{Zorn:} ein Mönch in der Wüste vom Zorn übermannt II, 1502. \gesperrt{Züchtigung:} ein Bauer will nicht von Gott gezüchtigt sein II, 809.

Geschlecht, männliches und weibliches. Knaben und Männer haben eine andere Gemüthsart als Mädchen und Frauen II, 91. f. 159. 497.; die Weiber haben besondere Gaben vor den Männern I, 147., sind von Natur barmherziger I, 1362.; haben besondere Gaben zu trösten, übertreffen auch sonst zuweilen die Männer II, 946.; finden schnell guten Rath, taugen aber doch nicht zu öffentlichen Aemtern II, 686. f.; man sieht an ihnen größeren Nutzen II, 2049.; haben nicht so viel und große Gebrechen wie die Männer I, 1730. f. Jedes hat seinen Samen I, 1169.; das männliche G. ist geschaffen zu zeugen, das weibliche zu gebären II, 569. f., Andere zu ernähren und erziehen II, 1585.; die natürliche Neigung der G. zu einander ist Gottes Werk II, 477. — S. \gesperrt{Mädchen, Mann, Weiber}.

Geschlechtsregister. I, 404. ff. 708. ff. Warum in der Schrift verworren II, 845. f.; geben zu unzähligen Fragen Veranlassung I, 719.; man soll darin nicht forschen II, 845. f.

Geschlechtstheile. Durch Christum geheiligt I, 1666.

Geschlechtstrieb. Durch den Fall verderbt I, 76.

Gesetz. 1. \gesperrt{Gesetz. Art und Natur:} G. gab es schon im Stande der Unschuld I, 130. ff.; man unterscheide zwischen dem G. vor und nach dem Fall I, 133., zwischen dem G. Mose und der Väter II, 144.; was eigentlich zum G. Mose gehört II, 144.; das geschriebene G. geht nur die Juden an, den Heiden ist es ins Herz geschrieben I, 1018., so z.\,B. das Gebot der Nächstenliebe I, 339.; das G. geht mit den 4 Stücken um: Laster und Tugend, Strafe und Belohnung I, 1077.; zum G. gehört Strafe und Schwert I, 988. \gesperrt{Zweck, Nutzen, Wirkung:} wozu das G. vor dem Falle gedient I, 133. f., wozu nach dem Falle I, 133.; es offenbart die Sünde, richtet Zorn an II, 1491., macht die Sünde lebendig I, 201.; II, 1494. und mächtig II, 2083., erschreckt I, 992., schlägt darnieder I, 1444., ist ein Amt des Todes I, 621. 1429., bläst auf I, 988.; II, 1883., macht keine Christen I, 1077., kann das Gewissen nicht trösten I, 619., macht nicht gerecht I, 769. 949.; II, 1018. \gesperrt{Verbindlichkeit:} G. geht Alle an I, 1017.; Gott ist nicht ans G. gebunden II, 1862. 1876. 1924., kann vom G. dispensiren I, 1231., auch das Sittengesetz ändern II, 515. \gesperrt{Erfüllung:} selbst den Gläubigen unmöglich I, 948. \gesperrt{Ende:} Christus ist das Ende des G.’s I, 1010. 1024. f., des G.’s Moses II, 144. f. — S. \gesperrt{Gebote, Gnade} 6, \gesperrt{Naturgesetze}.

2. \gesperrt{Gesetz und Evangelium. Scheidung und Vermischung:} man scheide G. und E. recht II, 392.; darin soll ein Prediger geschickt sein I, 1430.; rechte Scheidung ist Weisheit über alle Weisheit I, 1068.; man hat beides früh vermischt I, 1009., so auch im Koran I, 1009. \gesperrt{Amt und Wirkungen:} von beider Amt I, 1427. ff.; nur das E. kann der sündigen Natur helfen I, 223. f.; G. ohne E. richtet nur Zorn an I, 218. 221. f.; wozu das G. in der Kirche dienlich I, 1189. f.; G. zu lehren, damit wir unser Elend erkennen und Gnade begehren I, 949.; Gott predigt erst G., dann E. I, 1330., so sollen auch wir thun I, 1426. ff.; beides ist nöthig, G. und E. II, 1875. f. 1890. 1922.; erst muß das G. demüthigen und schrecken, dann tröstet das E. und richtet auf I, 992. 1444. ff.; II, 1705. 1710. 1876. 1981. f.; G. tödtet, E. macht lebendig II, 1494.; wenn G. und E. zusammenkommen, dann werden Christen I, 1077. \gesperrt{Wem zu predigen:} für wen beides ist I, 1186. ff. 1205. 1208. f. 1210. f.; G. für die Sicheren I, 1205.; E. den Zerschlagenen I, 1427., den Betrübten, Angefochtenen II, 616., den Gläubigen I, 1239. Der natürliche Mensch läßt sich das E. gefallen, wenn er nur nicht seine Sünden bekennen soll I, 1569. \gesperrt{Beides} muß in der Kirche gelehrt werden I, 1186. ff. 1204. f. 1208. f. 1210. f.; auch das G. I, 1190. 1301.; das G. in der Kirche nöthig um der Bösen willen und weil auch die Heiligen noch Fleisch haben I, 1190. f. 1244. — S. \gesperrt{Antinomer, Gesetzprediger, Gnade} 4.

3. \gesperrt{Jüdisches} G. Gilt nicht mehr I, 1221.

4. \gesperrt{Weltliches} G. G. ist eines weisen Mannes Rath I, 846.; G., das wider die Liebe streitet, ist kein G. I, 847.; alle G. und Rechte sollen zum Frieden gebraucht und dahin gerichtet werden I, 843. ff. 846.; bedürfen eines weisen Mannes, sie zu handhaben I, 846. — S. \gesperrt{Liebe} 6, \gesperrt{Recht}.

Vgl. \gesperrt{Recht}.

Gesetzprediger. Abgebildet durch Noahs Raben I, 619. ff.; was sie predigen I, 619. f.; von Cato, Cicero in Sittenpredigten übertroffen I, 620. f.

Gesichte. Was ein G. sei I, 930.; Satan kann sie auch hervorbringen I, 931.; II, 1049.; göttliche G. führen die Leute aus sich, daß sie meinen, sie träumen I, 939. — S. \gesperrt{Offenbarungen}.

Gesinde. \gesperrt{Stand:} dieser Stand hat dieselbe Ehre wie der Herren Stand I, 1728.; Gott hat Knechte und Herren gleich lieb I, 996. \gesperrt{Werke:} man diene um Gottes willen getreulich II, 483.; die geringsten Werke, im Glauben gethan, gefallen Gott II, 1742. f.; G. will Herr im Hause sein II, 554.; sündigt oft, davon der Herr nichts weiß I, 1467. \gesperrt{Getreues} G. ist ein Segen I, 1655.; II, 1231.; \gesperrt{gläubiges} dient nicht Menschen, sondern Gott I, 1725.; getreues ist selten II, 483.; durch den Glauben hatte Isaak frommes G. II, 129.; \gesperrt{untreues} G. eine Strafe I, 1655., macht den Hausstand beschwerlich II, 700.; \gesperrt{Klage} über das G. ist allgemein I, 1655. 1722. f. Wie Knechte, die einen geizigen Herrn haben, handeln mögen, um zu ihrem gebührlichen Lohn zu kommen II, 606. — S. \gesperrt{Abraham} 18, \gesperrt{Aufsicht, Beruf} 1, \gesperrt{Jakob} 22.

Gespenster. Dadurch hat Satan Viele betrogen I, 1525. — S. \gesperrt{Erscheinungen}.

Gespräche. Man soll G. bei Gastmählern führen I, 1159. f.; christliche G. erquicken das Herz I, 1159. f.

Gestirndienst. Wie entstanden I, 18. 663. 1663. f. — S. \gesperrt{Himmelskörper}.

Gethsemane. Bedeutung des Namens II, 1554.; da hat Christus die Schmerzen der Hölle erlitten II, 1554.; Christus nennt G. Gehenna II, 1554.

Gewaltige in der Welt. Warum die Tyrannen vor der
\quellseite{552}
Sintfluth so heißen I, 477. f.; hatten ihre Gewalt vom Teufel I, 478. — S. \gesperrt{Tyrannen}.

Gewissen. 1. \gesperrt{Gutes} G. Ein fein Ding II, 1106.; macht getrost und fröhlich II, 1595. f. 1684. 1721.; ein Fels in Trübsal II, 1147. f.; wie glückselig, wer ein gut Gewissen hat II, 1107. f.; leicht verletzt II, 1723.

2. \gesperrt{Erschrockenes} G. Der elende Zustand eines erschrockenen G.’s I, 210.; II, 1960.; ist schwer zu trösten II, 2077. ff.; wie zu trösten I, 1351.; II, 1900.; der Glaube beruhigt es II, 2080. f.; wenn es uns straft, nehme man zu Gott Zuflucht I, 339.; der Schmerz des G.’s kehrt trotz des Trostes der Vergebung immer wieder II, 1107. — S. \gesperrt{Anfechtung} 9, \gesperrt{Erschrockene, Traurigkeit}.

3. \gesperrt{Böses} G. Ist ein schrecklich Ding I, 351., ein groß Unglück II, 1106., die größte Strafe I, 719., eine böse Bestie und Erinnye II, 1514. ff. 1521., eine Wunde, die nicht zu heilen ist II, 2076., ein Beweis, daß man gesündigt hat I, 215. \gesperrt{Was es thut:} es fürchtet sich, wo es sich nicht fürchten sollte, und umgekehrt II, 1590. ff. 1684., flieht vor Gott II, 1599., legt auch das Beste aufs übelste aus II, 1594. f. 1597. f., plagt uns wegen erdichteter Sünden II, 1520. f., bereitet große Qual II, 1514. ff., mehrt die Leiden II, 1148., macht furchtsam I, 207. ff.; II, 712. 1215. ff., macht zu Feinden Gottes II, 1721., hindert das Gebet II, 1141., macht die Hölle zur Hölle II, 1595., zündet das höllische Feuer an II, 1720. \gesperrt{Es folgt} auf allerlei Sünden II, 1106.; ist Strafe der Sünde I, 207. f. 209.; II, 1106.; Todtschläger und die zum Tode sündigen, werden von Gewissensbissen gepeinigt I, 351. f. \gesperrt{Hülfe:} das Blut des Sohnes Gottes ist die rechte Arznei II, 2079.; Gottes Wort beruhigt es II, 1593. ff.; Bekenntniß der Sünde ist das Mittel dagegen II, 1522.; man nehme Christi Blut im Glauben an II, 2080. f. — S. \gesperrt{Anfechtung} 9, \gesperrt{Galgenreue, Reue} 1.

Gewißheit. G. in der Lehre nöthig II, 1352. ff.; G. über Vergebung der Sünde und Seligkeit soll ein Christ haben II, 1353. f.

Gideon. Aus dem Stamm Manasse II, 2041.; G.’s Glaube nicht wie Abrahams I, 891.; G.’s Gebet keine Versuchung Gottes I, 1698.

Giganten. Ursprung der Fabel von ihnen I, 687.

Gihon. Der Nil I, 119. 123.

Gilead. Es gab zweierlei G. II, 2033.; berühmter Ort II, 715.; hieß anders zu Jakobs Zeit II, 663.; metaphorisch: Jerusalem II, 715.

Glaube. I. \gesperrt{Das Credo.} Von den Worten: Ich glaube eine heilige christliche Kirche I, 1408.

II. \gesperrt{Glaube} (fides). 1. \gesperrt{Wichtigkeit} des G.’s. Er ist das Hauptstück der christlichen Lehre I, 1315.; wird 1 Mos. 15, 6. zum erstenmal in der Schrift lobend erwähnt I, 940.

2. \gesperrt{Art und Natur} des Glaubens. Der G. ist eine Gabe Gottes I, 1554.; II, 2089.; nicht eine lächerliche, kalte Eigenschaft und Geschicklichkeit II, 381., wie die Papisten meinen II, 418., sondern etwas Lebendiges II, 381., ein schweres Ding II, 618.; I, 556.; II, 146. f.; eine wunderbare Kunst II, 1963., von allen Künsten verschieden II, 284., nicht jedermanns Ding I, 757. f.; was der G. sei I, 941. f. 943. f.; ist mehr ein Leiden, denn eine Wirkung I, 756.; ist eine Veränderung der ganzen Natur I, 756. f.; „glauben“ heißt, sich in Christo segnen I, 1567.; ist Gott ein angenehmes Opfer I, 1319., der einige und angenehmste Gottesdienst I, 1557. \gesperrt{Unser} G. ist ein anderer, denn der, den \gesperrt{Adam} vor dem Fall hatte I, 188. f. Was \gesperrt{historischer} Glaube sei II, 287.; derselbe setzt sein Vertrauen nicht aufs Wort, besteht nicht in Anfechtungen II, 1905.; historischen G. hat auch der Teufel II, 1905. Der G. ist nicht immer gleich stark II, 288. f., bald \gesperrt{stark}, bald \gesperrt{schwach} I, 530. f. 789. f. 927. 1553. II, 151. ff. 897.; wird oft schwach II, 1556., wenn er die Anschläge der Feinde wider die Kirche sieht und Gott verzieht I, 704. f.; die Schwachheit des G.’s ist bei den erfahrenen Christen anderer Art als bei den jungen Christen II, 927. f.; der schwache G. macht das Herz verzagt II, 927., ist dann am stärksten, wenn er am schwächsten ist II, 750.; Gott duldet mit Langmuth den schwachen G. II, 836., verachtet ihn nicht II, 412., verwirft ihn nicht II, 288.; daß der G. oft schwach wird, läßt Gott zu, damit wir nicht stolz werden II, 928., kommt vom Fleisch II, 411. f.; unseres schwachen G.’s sollen wir uns schämen II, 2057. f.; haben wir keinen starken G., so sollen wir im schwachen aushalten und kämpfen II, 305. f. — S. \gesperrt{Kleinglaube, Schwachgläubige}.

3. \gesperrt{Gegenstand} des G.’s. Er hat es nicht mit gegenwärtigen und sichtbaren, sondern mit zukünftigen und unsichtbaren Dingen zu thun II, 4. 353. 757. 1292.; hängt sich stracks an das, das Nichts ist, und wartet, bis daraus Alles werde II, 284.; nimmt Dinge an, die stracks wider einander sind II, 1756. f.; glaubt die Verheißungen I, 1556. f.; ergreift den Segen Abrahams zuversichtlich I, 1572. f. 1583.

4. \gesperrt{Der} G. u. \gesperrt{das Wort Gottes.} Durch das Wort Gottes, nicht durch Speculation wirkt der Heilige Geist den G. II, 290.; dem Wort sollen wir glauben I, 1552.; Wort und G. gehören zusammen I, 944. 1495.; II, 1904. f. 1914. 1962.; lassen sich nicht scheiden I, 567.; der G. ist nöthig, wo das Wort ist I, 188.; alle Drohungen und Verheißungen fordern ihn I, 940.; G. und Furcht sollen bei einander sein II, 1917. f. 1920.; der rechte G. vertraut aufs Wort II, 1905.; ist nöthig, damit man sich der Fürsorge Gottes trösten kann II, 617. f.; der G. hat ein gutes Gedächtniß für die Verheißung II, 1796., hält sich fest daran II, 284., in den größten Nöthen I, 1318. f., im Kampf wider das Fleisch II, 287. f., ob er auch das Gegentheil empfindet II, 381. f., wenn Gottes Rath unseren Gedanken zuwider läuft II, 1290. f., wenn auch nichts zu hoffen ist II, 1962. f.; der G. thut nichts zum oder vom Wort I, 526.; der G., der auf dem Worte steht, kann nicht fehlen I, 792.; von dem G. und den Verheißungen muß man oft predigen II, 382. — S. \gesperrt{Verheißung} 3.

5. \gesperrt{Uebung und Wachsthum} des G.’s. Der G. übt sich immer im Wort und Gebet II, 1339.; wächst nach und nach II, 1334. f.; ist stärker in Noth denn in Glück II, 1729.; bedarf der Uebung durch Werke I, 1123., durch Kreuz und Widerwärtigkeiten I, 1422., deren Zweck solche Uebung ist II, 755. f. 1011. f.; wächst auch dadurch, daß wir in Sünden fallen I, 1327.; wird gestärkt, wenn man die Exempel der göttlichen Wohlthaten zusammen lieset II, 756. Wir werden nie Meister im G. I, 1554. f.; um Mehrung des G.’s sollen wir bitten II, 1894.

6. \gesperrt{Kraft und Wirkungen} des G.’s: er thut Dinge, die unglaublich und unmöglich scheinen II, 289.; ihm muß Alles gelingen I, 891.; II, 271. ff. 286.; erlangt, was er hofft I, 809.; II, 1870.; macht dem Gläubigen alle Dinge möglich II, 283. f.; verzagt nicht, wenn ein Anschlag nicht sogleich glücken will II, 274. f.; heiligt alle Dinge II, 468.; kann bei Gott nichts verlieren,
aller Verlust ist Gewinn II, 1697.; ernährt auch in der Theurung II, 303. Durch den G. wird das \gesperrt{Ebenbild Gottes} etlichermaßen wieder hergestellt I, 411. und alles erstattet, was wir verloren haben I, 1690. f. Der G. stiftet \gesperrt{Frieden}, wo sonst Feindschaft war I, 691. f.; durch den G. erlangt man \gesperrt{Gerechtigkeit vor Gott}, er macht gerecht und selig I, 941. ff. 952. f. 1018., er allein I, 1577. f. 1584., nicht sofern er durch die Liebe erst sein Wesen erhält (fides formata) 945. ff., nicht als Werk, sondern weil er Gottes Gnade in Christo ergreift I, 316. f., nicht als unser, sondern als Gottes Werk I, 945.; er macht gerecht, das erkannten David und die Propheten I, 1590. f.; er hat Einfluß auf die leiblichen Kräfte I, 414. f.; erlangt Gottes \gesperrt{Gnade} I, 324. f.; II, 1893. f.; ist das rechte wahre \gesperrt{Leben} in Gott II, 1901.; macht Einen zu einem himmlischen Mann, Ueberwinder des Todes, Erben des Lebens II, 1871.; macht uns, wenn wir gestorben, lebendig II, 1901.; beweist sich in einem herzlichen \gesperrt{Verlangen} nach Christo und nach der Erlösung von Sünde und Tod I, 430. f. — S. \gesperrt{Ebenbild}.

7. G. \gesperrt{und Gebet.} Der Glaube geht voraus, Gebet folgt II, 571.; man soll täglich um den Glauben bitten II, 147.; beten, wenn auch der G. schwach ist II, 1586.; wo Gebet ist, da ist auch G. II, 1578.

8. G. \gesperrt{und Bekenntniß.} Gehören zusammen II, 1297. f.; Bekenntniß folgt, wo G. ist II, 1203. 1228. f.

9. G. \gesperrt{und Werke.} Gehören zusammen, aber man soll sie nicht unter einander mengen I, 1577., sondern weit von einander scheiden I, 945. ff. 1457.; sind unzertrennlich I, 1582. f.; der G. soll vorausgehen I, 1706., die Werke als Danksagung folgen II, 446. f.; der G. ist die Hauptsache, dann sind die Werke gut II, 471. f.; der G. kann nicht ohne Werke sein I, 1510.; die Werke sollen folgen II, 467.; aus G. folgt der Gehorsam des Weibes und Gesindes I, 773.; Alles soll man im G. thun I, 1684.; der G. ist in allen Ständen nöthig II, 542. f.; Gott gefällt nichts, was nicht aus dem G. geht I, 323. 1706.; II, 2016.; alle Werke, die aus dem G. gehen, gefallen Gott wohl I, 316. 858. ff. 866. 954. 1367. 1667., der G. macht sie köstlich II, 463. 471. 1750., auch die geringen Werke II, 479. ff. und solche, darin dem Menschen Freiheit gelassen I, 860.; Gott sieht auf den G. I, 316. 328.; die Werke offenbaren und üben den G. I, 1123. f., sind Früchte I, 1537. und Zeugnisse des G.’s I, 317.

10. \gesperrt{Tugenden,} die aus dem G. hervorgehen. Der G. ist nicht ohne Liebe, Hoffnung u. andere T. I, 946. 948.; wirkt herrl. Tugenden. I, 948. f. G. \gesperrt{und Liebe} muß man unterscheiden I, 947. f.; der G. ist nicht ohne die Liebe I, 946. 948., ist die Mutter der Liebe I, 949.; fides informis II, 381.; der G. ist das beschauliche, die Liebe das wirkliche Leben II, 560. f.; der G., nicht die Liebe, macht gerecht I, 948.

11. \gesperrt{Kampf und Sieg} des G.’s. Wir sollen uns auf Kampf und Uebung des G.’s rüsten I, 783.; der G. macht muthig I, 1554., verrichtet große Wunderwerke I, 1498. f.; II, 129., sieht auf Gottes Gebot und überwindet so alle Feinde II, 286. Der G. wird durch \gesperrt{Anfechtungen} geplagt II, 381., versucht I, 786., durch Leiden geprüft II, 772., durch die Noth bewährt I, 784.; thut im Frieden gute Werke, in Anfechtungen kämpft er II, 658. f. 762. und überwindet II, 897. f.; kann böse Gedanken nicht hindern II, 785. Wir sollen den G. in Anfechtungen stärken I, 771.; wenn nicht auch Trost käme, würde er sinken I, 828.; der G. ist ein trefflich Ding, macht alle Leiden zum Kinderspiel I, 1067. Der G. ist ein allmächtiger Kampf wider das \gesperrt{Fleisch} II, 287. f.; wird vielfach durch die Erbsünde gehindert I, 1567. f., durch Unzucht verletzt II, 1272. Der G. siegt über \gesperrt{Gesetz, Sünde und Tod}, wie schon im Paradies verheißen I, 232. f.; überwindet \gesperrt{Gott}, wenn er anhält II, 793. 795. f. 811. f.; Gott widersteht ihm nicht gern II, 790.; durch den G. wird man der \gesperrt{Sünde} Herr I, 329.; er überwindet den \gesperrt{Tod} I, 1513. ff.; II, 286. f. 1905.; beweist seine Kraft besonders im Tode II, 381.; je stärker der G., desto schwächer der Tod, und umgekehrt I, 425. Die \gesperrt{Vernunft} muß getödtet werden, wo der G. Raum haben soll II, 56.; der G. ist der Vernunft ein unbekannt und verborgen Ding II, 706.; erscheint ihr gering und wie ein eitler Wahn I, 556.; II, 289.; wir sollen uns im G. durch die Vernunft nicht irre machen lassen II, 285. f., wider ihren Zweifel streiten I, 756., sie gefangen nehmen unter den Gehorsam des G.’s. I, 193.; der G. kann Dinge, die stracks wider einander sind, vereinigen, wie, Leben und Tod I, 1515. 1552. f. 1556. 1661. f., und glaubt auch, was schwer zu glauben ist I, 1513. f. II, 618. Der G. steht allein wider die ganze \gesperrt{Welt} auch unter schweren Aergernissen I, 466. ff. und wider ihr Urtheil I, 537. f. und überwindet II, 382., überwindet \gesperrt{Welt und Teufel} II, 412. 906. — S. \gesperrt{Anfechtungen} 6, \gesperrt{Beständigkeit}.

12. \gesperrt{Irrlehrer.} Die Papisten wissen von dem rechten Glauben nichts I, 1364. 1556. f.; II, 1569. 1904.; verachten ihn als ein gering Ding I, 556. 1310., ebenso die Welt I, 1306.

13. \gesperrt{Der} G. \gesperrt{Gottes.} Ist die Gabe des G.’s II, 1487.

Vgl. \gesperrt{Abraham} 10, \gesperrt{Beständigkeit, David, Isaak} 6, \gesperrt{Jakob} 6, \gesperrt{Joseph} 7, \gesperrt{Lot} 1, \gesperrt{Noah} 6, \gesperrt{Rebekka, Religion, Sarah, Väter} I; II, 3.

Glauben sollen wir dem Gott, der sich im Wort offenbart hat I, 1552.

Glaubensgenossen. Ihnen vor Andern wohlzuthun I, 1140. ff.

Gläubig werden heißt das Wort annehmen II, 1830.

Gläubige. 1. \gesperrt{Im Allgemeinen.} Sind Israeliter II, 811. f.; ihnen wird, was Zorn ist, Gnade und aus dem Tod das Leben I, 615. f. 1519. f.; Gott liebt sie I, 1342. — S. \gesperrt{Christ}.

2. G. \gesperrt{des} A.\,T.’s. Hatten dieselbe Lehre wie wir, predigten von dem zukünftigen Christus II, 1751.; glaubten die Lehre von der Dreieinigkeit I, 708.; wie weit sie die Verheißung vom Weibessamen verstanden I, 236. ff.; leben durch die Hoffnung auf Erlösung durch den Weibessamen, eben wie wir I, 241. — S. \gesperrt{Väter} I; II, 3; \gesperrt{Vorbilder}.

3. G. u. \gesperrt{Ungläubige.} S. \gesperrt{Fromme} 2.

Gleichkauf. Was G. ist II, 117.

Glossa ordinaria. Von der Himmelsleiter II, 394 f.; Auslegung von 1 Mos. 49, 11. 12. II, 1992.

Glück. 1. \gesperrt{Glück.} Welt schließt vom G. auf Gottes \gesperrt{Gnade} I, 1218. f. 1430. f.; II, 343.; ist ein falscher Schluß I, 1217. 1400. ff.; II, 1615.; G. beweist nicht, daß ein Volk Gottes Volk sei I, 1400. f. Ohne den Heil. Geist \gesperrt{wirkt} G. Hoffahrt I, 1218.; II, 1386. ff. 1390. 1729.; macht Narren II, 2038. \gesperrt{Verhalten:} im G. denkt man nicht an das mögliche Unglück II, 1539.; G. zu ertragen, dazu gehört ein Mann II, 1386. 1729.; im G. sollen wir unser Verdienst nicht rühmen II, 1631., uns nicht überheben, sondern demüthigen I, 1217.; die Gottlosen
\quellseite{553}
werden trotzig und vermessen im G. I, 687., halten sich selbst für die Ursache alles Segens I, 894. — S. \gesperrt{Gnade}.

2. G. u. \gesperrt{Unglück.} Sind bei einander I, 1354.; G. schwerer zu tragen denn U. II, 1386. 1395. f. 1849.; wir sind in G. und U. arme Leute II, 1539. f.; durch beides versucht Gott die Christen II, 1622. f.; Gott ist in G. und U. bei uns II, 1571. f., soll in beiden unsere Zuflucht sein II, 1631. f.; sollen uns in beidem an Gottes Verheißung halten II, 194.; es ist besser, im U. sein mit einem gnädigen Gott, als im G., wenn Gott zürnt II, 526.; im G. demüthig, im U. unverzagt sein, erfordert ein gläubiges Herz II, 1729., so sollen sich Christen verhalten II, 154. — S. \gesperrt{Fromme} 2.

Gnade. 1. \gesperrt{Gnade Gottes.} Ist allein in Christo zu finden I, 1025.; durch die \gesperrt{zufällige} G. wurden viele Heiden im A.\,T. selig I, 368. f.; welche große Wohlthat die im Evangelium geoffenbarte und geschenkte G. ist II, 410. f.; es ist G., daß Gott im Predigtamt zu uns redet I, 869.; sie leuchtet hervor aus dem Fluch über Satan im Paradies I, 231. ff.; bleibt, auch wenn Gott zu zürnen scheint II, 1490.; bringt abgöttische Menschen wieder zurecht I, 723.; unser Herz versteht die überschwängliche G. Gottes nicht I, 1109.; wie wir zur Gnade und Seligkeit kommen I, 1433. f.; die G. zu ergreifen, steht nicht in des Menschen Gewalt II, 2084. ff.; von Natur möchte der Mensch sie verdienen II, 108. ff.; G. muß man im Glauben annehmen, sonst sind wir verloren II, 411.; an der G. sollen wir nicht zweifeln, wie doch die Papisten lehren I, 1550.; aus den göttlichen Verheißungen soll man schließen, daß Gott gnädig sei I, 1402., nicht aus irdischem Glück I, 1217. ff. 1400. ff. 1430. f.; II, 343. 1615.; man kann aus der G. fallen I, 520.; G. mißbrauchen die Menschen zum Geiz II, 109. f.; Exempel der G. dienen den Furchtsamen zum Trost I, 732. f. — S. \gesperrt{Glück} 1; \gesperrt{Natur} 3, \gesperrt{Zeichen} 2.

2. G. \gesperrt{finden.} Diese Redeweise schließt alles Verdienst aus I, 497. f.

3. G. u. \gesperrt{Verdienst.} Können nicht neben einander bestehen II, 448.; bei Gott gilt kein V., sondern nur G. I, 1420.; Lehre von der G. ohne V. durch Abrahams Berufung bestätigt I, 730. f.; ob der Mensch durch seine natürlichen Kräfte Gnade verdienen könne, si non de congruo, tamen de condigno I, 580.

4. G. u. \gesperrt{Zorn.} G. den Betrübten, Z. den Sicheren I, 1191. f. — S. \gesperrt{Gesetz} 2.

5. \gesperrt{Auf} G. \gesperrt{sündigen} sollen wir nicht I, 376.; II, 2083. ff.; Gott straft es II, 1703. ff.

6. G. \gesperrt{und Gesetz.} G. ist Meisterin über das Gesetz II, 1862.

Gnadenmittel. Ihre Herrlichkeit zu erkennen, dazu gehören geistliche Augen II, 434. f.; wo sie sind, da finden wir Gott I, 1048.; II, 2000.; darin ist der Himmel über uns ausgebreitet II, 1098. f.; Gott könnte ohne sie selig machen II, 778., ohne sie wirkt der Heil. Geist nicht I, 1251.; Gott gibt sie uns zur Seligkeit II, 1049. 2000.; sind nicht um des äußerlichen Werkes willen gegeben, sondern daß wir der Wohlthaten Christi gedenken II, 918.; dadurch will Christus den Zweifel an der Seligkeit aus den Herzen reißen I, 1066.; werden auch von Gottlosen kräftig verwaltet I, 1405. 1408., z.\,B. unter den Papisten I, 1405.; weil alltäglich, werden sie verachtet I, 1118. f.; die menschliche Natur will sich daran nicht genügen lassen II, 561.; wir sollen uns zu denselben halten II, 778. f., sie ehrerbietig brauchen II, 1842. — S. \gesperrt{Erscheinungen} 1, \gesperrt{Predigtamt, Sacramente, Schlüssel, Wort} 3, 4, \gesperrt{Zeichen} 2.

Gog und Magog. Gog heißt Dach I, 658.; welche Völker G. und M. seien I, 658.

Gojim. Kein schändlicher Name I, 744.

Golgatha. S. \gesperrt{Schädelstätte}.

Gomer. Bedeutung des Namens, Anlaß dazu I, 664.; Wohnsitz I, 657. f.; welcher Völker Vater I, 657. f.

Gosen. Lage II, 1715. 1776. f. 1794.; fruchtbar II, 1776. f. 1794.

Gott. 1. \gesperrt{Gott.} Alle halten G. für eine \gesperrt{göttliche Gewalt,} die da wohlthätig ist II, 758. Wenn man von dem \gesperrt{Wesen} der Gottheit handelt, unterscheide man zwischen dem verborgenen und geoffenbarten G. II, 176.; von dem verborgenen und geoffenbarten G. II, 184. f.; G. ist ein Geist I, 487., einig und dreieinig II, 213. ff., welche Lehre man freilich nicht begreifen kann II, 214.; das göttliche Wesen an sich ist unbegreiflich I, 21.; G. ist in seinem Wesen ganz unerkenntlich und unbegreiflich I, 487. f. vor wie nach der Schöpfung I, 13.; der Mensch kann G. durch seine Vernunft nicht finden I, 1046. f.; über ihn sollen wir nicht forschen II, 176. ff.; wer von G. fragen will, der sehe auf den Gottesdienst I, 1059.; außer dem Wort hat man den rechten G. nicht I, 1718.; man halte sich ans Wort, speculire nicht darüber hinaus I, 17.; G. ist, wo die Gnadenmittel sind I, 780. 1048. 1116.; nur von dem im Wort geoffenbarten G. ist Leben zu hoffen I, 1064. Ist ein G. nicht nur der Juden, sondern auch der Heiden I, 1023. 1029.; II, 909. 1408. f.; was das heißt: „jemandes G. sein“ I, 1057. f. 1064. 1067. 1518.; II, 758.; war in besonderem Sinne der Juden Ernährer und Beschützer I, 1016.; ist unser Vater I, 155.; sieht die Person nicht an I, 328. Wir haben einen G., der mit uns \gesperrt{redet} und umgeht I, 1118.; er handelt mit uns nicht nach seiner Majestät, sondern er läßt sich zu uns herab I, 1442. 1448. f. 1549.; das ist nöthig, doch verachten wir G. grade deshalb I, 1442. f. 1448. f. Juden, Türken und Papisten haben einen Aristotelischen G. I, 1551. — \gesperrt{Eigenschaften} G.’s siehe unter: \gesperrt{Allmacht, Barmherzigkeit} etc. — S. ferner: \gesperrt{Angesicht, Drohungen, Ehre, Eid} 1, \gesperrt{Erkenntniß} 1, 2, \gesperrt{Erscheinungen} 2, \gesperrt{Gaben, Grammatik} 2, \gesperrt{Hülfe, Leben, Menschenworte, Name} 1, \gesperrt{Reue, Ruhe, Schöpfer, Spott, Strafe, Versehung, Vorsehung, Werke} I, \gesperrt{Wille} I, \gesperrt{Wohlthaten} 2, \gesperrt{Wort, Zorn} 1.

2. \gesperrt{Der starke} G. \gesperrt{Isaaks.} Ist Christus II, 853.

3. \gesperrt{Baum, Garten, Stadt} G.’s. Was das heiße I, 849. ff.; II, 546. f.

Gottesdienst. 1. \gesperrt{Rechter} G. Was es heißt, Gott dienen I, 398.; besteht nicht in Ceremonien, sondern in Anrufung und Dank I, 400. f. 778. 1058. f.; II, 921., in Ehre und Furcht I, 1493. 1537. f.; II, 721., in Glauben, Hoffnung, Liebe II, 4. 8., im Warten des Berufs I, 1058. f.; der höchste G. ist, daß man Gott für wahrhaftig hält I, 1182. 1537.; II, 64. 129. ff. 183., auf seine Hülfe harrt II, 4. ff. 11. f.; der rechte, höchste G. ist lehren und beten, darnach folgt die Tödtung des Fleisches II, 936. f.

2. \gesperrt{Falscher, selbsterwählter} G. \gesperrt{Falscher} G. besteht zu allen Zeiten I, 727. f.; der G. sei nicht Sadducäisch oder Epicurisch I, 1059.; der G. der Heuchler gefällt Gott nicht II, 721.; der G. der Gottlosen ist gezwungen und vergeblich I, 1537. f. \gesperrt{Selbsterwählter:} ohne Gottes Befehl richte keinen G. an I, 567. ff.; II, 906. f.; verwirf allen selbsterwählten G. I, 1538. 1592. ff. 1597. f.; er gefällt Gott nicht I, 1437.; II, 778. f. 2014. ff.; ist Sünde, Abgötterei I, 567. 778.; ob Noahs Opfer ohne Gottes Gebot war I, 567. ff.; selbsterwählter G. entsteht, wenn man Gottes Wort verläßt II, 779., gefällt den Leuten II, 912., findet sich bei Türken, Juden, Mönchen, Papisten I, 567.; II, 777. 907. 1109.; der Eifer für selbsterwählten G. ist immer größer als der Eifer für den rechten II, 2016.

3. \gesperrt{Ort, Zeit, Ceremonien} des G.’s. \gesperrt{Ort:} man kann Gott nicht überall dienen, wenn er einen gewissen Ort dazu bestimmt hat II, 777.; G. im A.\,T. an Jerusalem gebunden II, 776. ff.; G. im Freien wäre um vieler Ursachen willen gut I, 848. Die \gesperrt{Morgenzeit} war in der Schrift und Gewohnheit zu Gebet und Predigt verordnet I, 99. \gesperrt{Ceremonien} im Stande der Unschuld einfach, G. ohne viel Mühe und Kosten I, 130. — S. \gesperrt{Ceremonien}.

Gottesfurcht. S. \gesperrt{Furcht} 3.

Gottesfürchtige. S. \gesperrt{Christ}.

Gottlose. Gott gibt ihnen das Allerbeste im Leiblichen I, 853. f.; Gottes Gaben gereichen ihnen zum Schaden I, 1330.; Gott lacht ihrer Anschläge I, 703., macht sie zu nichte II, 1949. f.; an ihnen mißfällt Gott Alles II, 848.; reißen den Namen der Kirche an sich I, 933.; ziehen die Verheißungen auf sich, verachten die Drohungen I, 497. 936. f. 1274. f.; bekennen Gott nur mit den Worten II, 675.; fürchten Gott nicht, wo er zu fürchten ist I, 1538.; sind mit bösem Gewissen gestraft I, 373. f.; II, 712. 719.; können nicht beten I, 494.; prophezeien sich ihr Unglück I, 695. 699.; es trifft sie, was sie fürchten I, 699.; das Unglück, das sie Andern drohen, trifft sie selbst II, 695.; sündigen, um nur ihren Muthwillen auszuüben I, 1194., ob sie wohl wissen, die Strafe ist nicht weit I, 697.; in der Noth sind sie verzagt, darnach aber wieder die Alten I, 895.; fallen immer wieder in Sünde I, 1334.; ein einziger kann wohl eine ganze Stadt irre machen I, 1100.; der Tod kommt ihnen immer unerwartet I, 544.; ihr Glück nimmt kein gutes Ende II, 373. — S. \gesperrt{Fromme} 2, \gesperrt{Leben} 10, \gesperrt{Ungläubige}.

Gottlosigkeit. G. und Aberglaube bleibt eine Zeit wie die andere I, 727.; G. und Thorheit sind bei einander I, 340.; G. folgt aus Verachtung des Worts II, 98.; G. um so größer, je reiner und reichlicher das Wort gepredigt I, 1200. f.; mehrt sich von Tag zu Tage I, 506.

Gottselige. Sind Wunderwerke und Gaben Gottes II, 773.; warten nicht nur der geistlichen Uebungen II, 1076.; müssen Verfolgung leiden I, 1393., Anfechtung und Kreuz erwarten II, 126. f. 139., das ist nöthig und heilsam II, 807. f.; müssen im Kampf doch endlich siegen II, 800.; Gott hilft ihnen II, 899., zählt ihre Thränen II, 772.; sollen vorsichtig wandeln um der Feinde willen II, 1028. — S. \gesperrt{Christ, Leben} 6.

2. G. u. \gesperrt{Heuchler.} H. haben einen größeren Schein, denn die G. II, 308.; reden eben wie die G., meinen es aber nicht aufrichtig II, 709.; suchen die G. leiblich und geistlich zu betrügen II, 712. f.; morden die G. unter einem guten Schein I, 331.; es ist G. beschwerlich, mit den Heuchlern umzugehen II, 708. ff. 713.; G. sollen die Gemeinschaft der H. fliehen II, 713.

Gottseligkeit. Man übt G., wenn man gläubig in seinem Beruf einhergeht I, 867.

Götzenbilder. Woraus verfertigt II, 923. f.; gab solche unter den Juden zur Zeit Moses II, 651., Davids II, 650. f., unter den Deutschen zu Luthers Zeit II, 650. f.; Jakob vergräbt sie II, 924. f., Mose gebietet sie zu verbrennen II, 924. f.; ob das jetzt noch nöthig II, 925. f.; die Kirche hat sie eingeschmolzen II, 925. f.

Götzendiener. Wer ein G. ist I, 183. f.

Götzendienst. Was dabei die Hauptsache II, 913.; man beabsichtigte nicht Holz und Steine anzubeten II, 674. f. 922.; an den lustigsten Orten getrieben II, 939.; getrieben von den Juden II, 960. im Thal Thopheth II, 1550. ff., auch im Pabstthum II, 913. f. — S. \gesperrt{Abgötterei}.

Grab Christi. Ob da, wo Jakob die Himmelsleiter sah II, 428.

Grabmale. Damit haben die Väter im A.\,T. die Gräber geziert II, 971.; erinnern an das ewige Leben II, 971.; die Heiden haben die Gräber herrlich zugerichtet II, 1836.

Grade in den Hochschulen. Zweck II, 1404.

Grammatik. 1. \gesperrt{Grammatik.} Allein genügt sie nicht zur Schriftauslegung I, 364. 451. f. 998. ff. 1002.; ungeschickte Auslegung kommt von der G. I, 451. f.

2. G. \gesperrt{Gottes.} Darin alle Creaturen Worte I, 60.

Grammatiker. Aufgabe eines guten G.’s. II, 849.

Gregorius. Erkannte Christum nicht sehr wohl II, 1977.; verstand nicht, was wirkliches und beschauliches Leben ist II, 956.; glaubte alle Erscheinungen I, 1527. f.; durch seine angeblichen Erscheinungen richtete er viel Unheil in der Kirche an I, 1121. f.; Stifter der Lüge vom Fegfeuer und Messen II, 1547.; sein Decret von der Priester Ehelosigkeit I, 293. \gesperrt{Meinungen und Aussprüche} über: die Erkenntniß der Schuld I, 1347. f.; die Gottlosen II, 1150.; die Himmelsleiter II, 395.; die heilige Schrift I, 2.; ubicunque serpens caput etc. II, 884.

Griechen. Die allerärgsten Leute I, 854.; ihre Könige sehr zu Unzucht geneigt I, 804.

Gritel. S. \gesperrt{Agricola}.

Grimmwalt. II, 1287.

Großgriechenland. I, 658. f.

Gründe. In G. ist das Land dürr II, 199.

Gut. 1. \gesperrt{Gut.} Man unterscheide zwischen dem, was weltlich, und dem, was theologisch gut ist I, 483.

2. \gesperrt{Gut und böse.} Gottes Wort macht, daß etwas gut oder böse sei I, 117. f.

Gut. \gesperrt{Ehrlich} erworbenes G. segnet Gott II, 206.; \gesperrt{unrecht} G. gedeiht nicht, bringt Strafe II, 205. f. 586. ff. 590. f. — S. \gesperrt{Güter, Schweiß}.

Güte Gottes. Ist unermeßlich, unendlich, unbegreiflich II, 729., unerschöpflich II, 1698. 1700.; übertrifft der bösen Engel Bosheit II, 731. f.; wo Gott ist, da ist Alles II, 569.; Gott ist so gütig, daß er nichts Böses würde geschehen lassen, wenn er es nicht wieder gut machen könnte II, 685.; Gott freut sich, uns Gutes zu thun I, 573.; Gottes G. muß man erfahren, kann man durch Speculation nicht erkennen II, 1383.; hat sich mannigfach offenbart I, 47. f. 50. 87. 1087.

Güter. 1. \gesperrt{Zeitliche, leibliche, irdische.} Sind gute Gaben und Creaturen, werden aber mißbraucht I, 830. 832. 855. f.; II, 296., deshalb aber nicht zu verwerfen I, 857.; Gott theilt sie aus, wie er will II, 1071., gibt sie uns zum Trost I, 828. f.; man soll sie nicht verachten II, 297., sondern gebrauchen I, 833. f., dankbar dafür sein II, 296. f.; zum Gebrauch bringe man ein frommes Herz I, 833.; der natürliche Mensch mißbraucht sie, der erneuerte braucht sie recht I, 856. f.; wie die Gläubigen sie ansehen und gebrauchen I, 739.; II, 121.; man hänge das Herz nicht daran I, 918. — S. \gesperrt{Gaben, Reichthum}.

2. \gesperrt{Geistliche} G. Was darunter zu verstehen I,
\quellseite{554}
1619. f.; die höchsten sind: Erkenntniß des ewigen Lebens, Glaube, Hoffnung der Auferstehung I, 1421. f.; was nicht dazu gehört II, 103. f.; man kann sie nicht kaufen oder verkaufen II, 107. ff.; wie zu gebrauchen II, 111. f.; was die Papisten „g. G.“ nennen I, 1619. f. — S. \gesperrt{Gaben, Simonie, Zeit} 2.

3. \gesperrt{Himmlische} G. Sind uns durch den Glauben gewisser als das irdische Leben II, 1869.

Gutes und Böses. Viel mehr G. denn B. in der Welt II, 728. ff. — S. Gut.

\section{H}

Haare. Ursachen grauer H. II, 1799. f.

Habsucht und Ehrgeiz findet man bei Edelleuten, Herren und Fürsten I, 915.

Hadad. II, 1016. f.; fromm ib.

Hadar. II, 1017. f.

Hagar. I, 990. ff. 997. ff. 1435. ff. War gläubig und gerecht I, 1001. f. 1006.; eine treffliche Matrone II, 9.; hatte leibliche und geistliche Anfechtungen I, 1424. f. 1435. ff. 1448. f.; ward eine Mutter in der Kirche und was sie gelehrt I, 1439.; ein Bild der fleischlichen Menschen I, 990. und der falschen Synagoge I, 1000. f. — S. \gesperrt{Abraham} 23, \gesperrt{Ketura}.

Halberstadt. Woher der Name I, 1610.

Halle. Salzbrunnen bei H. I, 546.

Ham. 1. \gesperrt{Noahs Sohn.} I, 627. ff. 640. ff. Bedeutung des Namens und Anlaß dazu I, 654. f.; Noahs jüngster Sohn I, 435. 502. 664. f.; vor der Sintfluth gläubig, nachher verworfen I, 519. f.; störte zuerst den seligen Zustand nach der Sintfluth I, 684.; war allmählich gefallen I, 630. f. 640.; die Erbsünde machte ihn vermessen und stolz I, 630. ff.; war ein Verächter Gottes, daher auch seines Vaters I, 638.; von der Größe seiner Versündigung an Noah I, 639. 687.; warum sein Name in Noahs Weissagung nicht vorkommt I, 640. 650. f.; verachtete den Fluch I, 641. f.; von der Erfüllung des Fluchs I, 642. f. 647. f. 665. f.; hatte die besten Länder inne I, 666.; Aehnlichkeit mit Kain I, 648. 655.; sein schreckliches Exempel soll man der Jugend fleißig vorhalten I, 639. — S. \gesperrt{Noah} 14.

2. H.’s \gesperrt{Geschlecht.} I, 664. ff. Wohnsitz I, 657.; verachtete Noahs Fluch I, 657.; pflanzte die Abgötterei im Morgenlande I, 656.; nicht Alle verdammt I, 644. 1033.

3. \gesperrt{Das Land Ham.} Nicht weiter bekannt I, 880.

Handauflegung. Ein alter Gebrauch II, 1860.

Handelsstädte. Stets die berühmtesten Städte I, 665. f.

Hanna. Verzweifelte daran, daß sie würde Kinder bekommen II, 574.

Hannas. Vom Stamme Simeon II, 117.

Hannibal. Sein Befehl nach der Schlacht bei Cannä II, 749.; sein Traum und dessen Auslegung II, 1049. f.; Gott setzte seinem Zuge gegen Rom ein Ziel II, 815. f. 931.; vincere scis, H., etc. II, 931.; kein Held wie Abraham I, 889.

Hanoch. 1. Bedeutung des Namens I, 383. f.; Lage der \gesperrt{Stadt} I, 668.

2. \gesperrt{Hanoch und Henoch.} Hanoch, Kains Sohn, ein Anfang des leiblichen, Henoch des geistlichen Segens I, 384.

Haran. 1. \gesperrt{Abrahams Bruder.} I, 722.; jüdische Meinung von seinem Märtyrertode I, 722.

2. \gesperrt{Das Land.} I, 772.

3. \gesperrt{Die Stadt.} Soll auch Charran geheißen haben II, 387.; der Römer Niederlage daselbst II, 387. 475.; H. war eine Kirche und treffliche Schule II, 1015.

Harren. Die Lehre, man solle auf Gott h., verlachen die Weltmenschen II, 5.; auf Gott h. ist schwer II, 4., kein Werk des Fleisches oder der Heuchler II, 4.; der rechte Gottesdienst II, 4. f., läßt nicht zu Schanden werden II, 6. f.; Isaak hat geharrt II, 3. ff. — S. \gesperrt{Hülfe}.

Harz. Verschiedene Arten II, 1124.

Haß. Menschlicher H. läßt nach, teuflischer gegen die Heiligen Gottes nimmt täglich zu I, 343.; hindert das Gebet II, 1141.; macht, daß man einen Menschen flieht II, 1070.

Hauptstadt. Nach ihr richten sich gewöhnlich die anderen Städte I, 1223. f. 1227.

Hausarbeit. Dazu Rebekka angehalten I, 1707. f.

Hausmutter. Ihr Amt II, 625. f., höchstes Lob I, 1171.; wird geplagt und gedemüthigt I, 1165.; Sarah eine heilige H. I, 1165.; Lea und Rahel gottselige H. II, 540. — S. \gesperrt{Hausstand}.

Hausstand. Von Gott gestiftet II, 1073., geliebt I, 990., regiert und geschützt I, 992. f.; II, 523. 562. 617.; mit Verheißung geziert II, 1732.; nöthig II, 458. f. 1073. f.; mühevoll I, 979. f.; II, 700. \gesperrt{Werke} des H.’s gefallen Gott wohl, hindern den Gottesdienst nicht II, 1072. ff.; die häuslichen Werke sind von selbsterwählten zu unterscheiden I, 1178., sind geboten I, 1479., wenn im Glauben gethan, gottgefällig I, 1162. ff. 1724.; II, 569., köstlich I, 1172. f., von Christo und den Aposteln hochgehalten I, 1163.; sie übertreffen aller Mönche Werke I, 858. f. 1161. f. 1479.; haben keinen Schein vor der Welt I, 1180. Die rechte Lehre von den Werken des H.’s soll man der Jugend einprägen II, 1074.; H. geht nicht ohne Sünde ab I, 821.; die sich nicht wollen züchtigen lassen, soll man aus dem Hause stoßen I, 1034. f.; dem H. darf man sich nicht entziehen II, 1072.; er fordert gläubige Menschen II, 542.; Christen müssen auch das Hausregiment verwalten II, 458. ff.; die Papisten verachten den H. und seine Werke II, 1072. ff.; Satan, ein Feind des H.’s, richtet darin Unglück an II, 327. f. — S. \gesperrt{Ehestand, Hausmutter, Hausvater, Stände, Väter} II, 7.

Hausvater. Sein Amt II, 88. f. 1240.; verantwortlich für die Thaten seiner Untergebenen I, 1034.; wie er sein Amt führen soll II, 279. — S. \gesperrt{Hausstand}.

Hazelon Thamar I, 881.

Hebammen. S. \gesperrt{Wehmütter}.

Hebräer. 1. \gesperrt{Die Menschen.} Woher der Name I, 679.; Anlaß dazu I, 715. 885.

2. \gesperrt{Der Brief an die} H. Verfasser II, 1886.; sieht Cap. 13, 2. auf 1 Mos. 18, 2–5. I, 1133.

Hebron. Lage II, 199. 853.; königliche Stadt I, 782.; viertheilig I, 672.; frühere Namen H.’s I, 1609.

Hebroniter. Fielen nach Abrahams Zeit ab I, 1631. 1635. f. — S. \gesperrt{Abraham} 24.

Hehler und Stehler gleich II, 885.

Heiden. 1. \gesperrt{Im Allgemeinen.} Wie weit sie Gott erkennen I, 1056.; haben keine Hoffnung, weil sie nicht glauben I, 1514.; Gott hat sie nicht verworfen II, 1172. ff., sollen selig werden I, 1114.; auch sie geht die Verheißung von Abrahams Samen an I, 1560.; hielten Sünde für Gottesdienst I, 455.; „Heiden“ sind nach dem Sprachgebrauch der Schrift die Völker am Mittelmeer I, 878. — S. \gesperrt{Juden} 3.

2. \gesperrt{Im} A.\,T. Sollten wohl nicht beschnitten, doch aber des Segens Abrahams theilhaftig werden I, 749. 1028. f. 1032. f. 1036. 1044. 1046.; II, 228. f. 1418. f.; wozu der Juden Beschneidung ihnen diente I, 1028. f. 1031. 1043. f. 1055. 1068. 1093.; konnten Gottes Volk nicht durch die Beschneidung werden, sondern durch den Glauben I, 1052.; viele H. sind des geistlichen Segens theilhaftig geworden I, 767. f. 875. 884. 1014. 1023. f. 1026. 1345. und selig I, 356. f. 1031. ff. 1044. 1055.; II, 1952., aus zufälliger Gnade I, 368. f. 644. f. 1033.; II, 984. f., aber nicht als Heiden, sondern durch das Wort II, 1705. f. 1827. ff., durch den Glauben I, 1078.; ihre Berufung ist von Noah geweissagt I, 646. f. 653., wirklich geschehen I, 654. 706. 876. f. 909. 920. 1302. 1307. 1332. f. 1631.; II, 908. ff.; haben von den Juden das Wort empfangen II, 1758. 1831.

Heil. Kommt von den Juden I, 646.; H. bedeutet oft Sieg II, 2027.

Heilige. 1. \gesperrt{Die wahren} H. Wer heilig ist II, 459. 705.; Christen sollen sich nicht schämen, sich H. zu nennen II, 391.; haben noch Sünde, Gebrechen und schwere Seuchen I, 1669. f.; II, 664.; auch die besten fallen II, 843.; Gott bewahrt sie vor Sünden I, 1331. f.; wie man von ihnen predigen soll I, 734.; die H. waren Menschen wie wir I, 1309. — S. \gesperrt{Christ, Fleisch} 2.

2. \gesperrt{Die} H. \gesperrt{des Pabstes.} Der Pabst macht sie zu reinen Engeln I, 723.; sind den wahren H. in der Schrift sehr unähnlich I, 1309. f.; diesen, nicht den H. des Pabstes, sollen wir nachfolgen I, 1310. 1315. f. — S. \gesperrt{Mönche} 2.

Heiligen. Heißt: zum Gottesdienst aussondern I, 96. f.

Heiligenschein. Der rechte H., das Wort Gottes II, 1760. f.

Heiligkeit. Worin rechte H. besteht II, 456. ff. 1566.; ist zweierlei: die vollkommene durch das Wort, die unvollkommene oder Werkheiligkeit durch Liebe und Werke II, 390. ff.; beiderlei H. zu unterscheiden II, 392. 413. ff.; beiderlei soll zusammen sein, die Werkheiligkeit folgt der Heiligkeit des Worts II, 415. f.; Werkheiligkeit ist nöthig II, 414. f., daran hat die Welt Gefallen II, 413. f., man darf sie der H. des Worts nicht vorziehen II, 414. f.; was für einer H. sich die Papisten rühmen I, 528.; sie halten nur Ehelosigkeit für H. II, 238. — S. \gesperrt{Mönche} 2, \gesperrt{Werkheiligkeit}.

Heiligung. Muß zunehmen II, 2003. ff.; wird in diesem Leben nicht vollkommen I, 241. 411. — S. \gesperrt{Heilige} 1, \gesperrt{Natur} 3.

Heimlichkeiten. Soll man nicht offenbaren II, 1002.

Heinrich, Herzog zu Sachsen. Unterstützte Kirche und Schule II, 450.

Heinz Mordbrenner. Wollte für heilig und gerecht gelten II, 708.

Helena. II, 882.

Hemath. Ist Antiochien I, 677.

Hemor. II, 864. ff. Heißt Esel II, 883.; stolz II, 870. f. 881. f.; Verächter Gottes und der Menschen II, 872.; macht sich der Sünde Sichems theilhaftig II, 872.; heuchelt II, 872.

Henning, Dr. Konnte gut disputiren I, 1178.

Henoch. I, 417. ff. Warum seine \gesperrt{Geschichte} in die Schrift aufgenommen worden I, 427.; sie lehrt Unsterblichkeit, Sieg über Tod und Sünde I, 403. 405. 419. 427.; welche Aufnahme sie bei seinen Zeitgenossen fand, den frommen I, 419. ff. 424. f. 428., den gottlosen I, 427. H. \gesperrt{war} auch ein Sünder wie wir I, 405. 422. 426.; wie er vor Gott gewandelt I, 421. f.; ein Ehemann und Hausvater, kein Mönch I, 419. 428.; führte sein göttliches Leben im öffentlichen Predigtamt I, 418.; war groß seines Berufs und Predigtamts halben I, 419.; vor Anderen muthig gegen Satan und die Kainiten I, 418. f.; damals Adam todt, Seth Hoherpriester I, 422.; lebt nun ein göttliches, ewiges Leben I, 421. f. 426. — S. \gesperrt{Hanoch} 2.

Hercules. II, 516. — S. \gesperrt{Theseus}.

Herkunft. S. \gesperrt{Geburt} 1.

Herodes. Herrschte widerrechtlich und tyrannisch über die Juden II, 324. f. — S. \gesperrt{Bethlehem}.

Herodot. Was er von den Egyptern sagt II, 1613.

Herrlichkeit. S. \gesperrt{Leiden} 3, \gesperrt{Niedrigkeit}.

Herrschaft. über die Thiere. Wir haben nur noch ein Ueberbleibsel I, 82. 162.; nach der Sintfluth von Gott vermehrt I, 589. f.; wir haben sie, obwohl wir uns vor Thieren fürchten, diese wider uns wüthen I, 593. f.; wenn die Sünden zunehmen, nimmt die H. ab I, 593. f.; daß Satan sie nicht hat, eine große Wohlthat I, 162. — S. \gesperrt{Thiere} 1.

Herrschsucht. Kommt von der Erbsünde II, 1410.

Herz. Gott sieht ins H. I, 259., will vor Allem das H. haben II, 258., hat des Menschen H. in seiner Gewalt II, 815.; er allein kann es ändern II, 815.; Herz und Willen des Menschen ändern ist der herrlichste Sieg II, 816. f. Ein \gesperrt{betrübtes} H. wiederholt die Klage oft I, 541. ff.; „Einem das Herz stehlen“ II, 653. ff.

Hesiod. Ausspruch: vom Ehelichwerden II, 488., von dreierlei Leuten II, 1365.

Hethiter. Woher der Name II, 993.; gewaltiges Volk, das ärgste unter den Cananitern II, 993.

Heuchelei. Ist argwöhnisch II, 686.; kann erschreckliche Dinge ausrichten II, 1142., kann Gott nicht betrügen II, 259.; Gottselige sollen der H. feind sein II, 258.

Heuchler. II, 481. ff. Haben keinen Gott II, 679.; sind tückisch und listig II, 132., Wölfe in Schafskleidern, betrügen Andere II, 258.; ihr Haß nicht menschlich, sondern teuflisch I, 343.; machen viel Worte, die einen großen Schein haben II, 307. 710. f. 715. 720.; verstellen sich II, 1612., als wollten sie Gott und dem Nächsten dienen I, 333.; gehen wohl zur Kirche, verachten aber innerlich Gottes Wort II, 90. 721.; rühmen sich Gottes und der Gottseligkeit II, 716.; stellen sich freundlich, bis Gelegenheit zu schaden kommt I, 331. f.; beschuldigen Andere unrechtmäßig II, 680.; schmücken sich, verachten Andere II, 717.; sind halsstarrig und hoffärtig selbst bei offenbaren Sünden, wollen die Allerheiligsten sein I, 337. f. 339.; bemänteln oder leugnen ihre Sünden I, 336.; halten es für die größte Schande, eine Sünde zu bekennen II, 871.; klagen eher Gott an, als daß sie sich anklagen lassen I, 340.; man kann H. nicht zu rechter Buße bringen II, 699.; ein H. ist ein solch Ungeheuer, das schlechterdings nicht sündigen kann, ist auch nicht zu bekehren II, 708.; von der H. Buße II, 660. ff. Christus straft H. am härtesten II, 1128.; ihr Dienst gefällt Gott nicht II, 721.; ihre Strafe ist ein böses Gewissen II, 719.; Gott läßt sie nicht lange verborgen I, 340.; ein Bild der H. ist Kain I, 331. f., Esau II, 99. — S. \gesperrt{Gottselige} 2.

Heuschrecken. Im Jahre 1542 I, 594.

Hevila. Welche Länder H. umfaßte I, 119. 123.; jetzt noch ein reiches Land I, 123.

Hiddekel. Der Tigris I, 119. 123.; ein geschwindes Wasser I, 123.

Hierarchien. I, 287. f.

Hieroglyphen. S. \gesperrt{Bildersprache}.

Hieronymus. 1. \gesperrt{Der Kirchenvater.} Klagt über unzüchtige Träume II, 530.; betet in einem Hymnus: Ne polluantur etc. II, 1928.; ließ sich durch ein einziges Wort zu einer heftigen Schmähschrift bewegen I, 1356.;
\quellseite{555}
ist wohl selig geworden II, 1678. Als \gesperrt{Uebersetzer u. Ausleger der Schrift}: übersetzt zu frei II, 1261., fehlerhaft II, 1307., verdunkelt den Text II, 1336. 1425.; Uebersetzung von 1 Mos. 49, 4. II, 1929., V. 6. II, 1942.; rühmt die Allegorien I, 284., allegorisirt I, 284.; II, 557., seine Allegorien sind fehlerhaft I, 612. 626.; II, 774. 811. \gesperrt{Meinungen u. Aussprüche} über: 1 Mos. 38, 14. II, 1189.; 1 Mos. 49, 11. 12. II, 1985.; die Anachoreten I, 997.; die Blutschande Judas mit Thamar II, 1168.; die Buße als das andere Brett II, 1465.; „Chesib“ und des Patriarchen Juda Weib II, 1163.; die Ehe Abrahams mit Ketura I, 1744.; das eheliche Leben der Väter II, 556.; die zweite Ehe I, 1748.; „Gehenna“ II, 1550.; die Geizigen II, 612.; Gosen II, 1777.; „Israel“ II, 795.; Jakob II, 1745.; Josephs Ehrentitel II, 1402., schöne Gestalt II, 2041. f.; Judas Ischarioth II, 104. 1717.; kibrath II, 961. f.; einen Märtyrer, der sich die Zunge abbiß II, 1259. f.; das Mönchsleben I, 1596. f.; II, 642. 943. 1678.; die Pferdezucht in Spanien II, 602.; Potiphar II, 1222. 1414.; reeh und roeh II, 1186. 1202.; sar hattabachim II, 1222.; die Siebenzig II, 593. 1847. f.; den Stamm Simeon II, 1117. 1717.; Thabor II, 2018.; den Thurm Eder II, 972.

2. \gesperrt{Hieronymus von Prag}. Ein rechtschaffenes, heiliges Glied der Kirche I, 459.; hatte unverzagten Muth und Beständigkeit I, 538.; war trunken vom Heil. Geist II, 2002. f.

Hilarion. Sein strenges Leben I, 1355.; was davon zu halten I, 865. f.

Hilarius. Meinung: vom Unterschied der Personen in Gott I, 61., von den Schöpfungstagen I, 6. 45.

Himmel. 1. \gesperrt{Bedeutungen} des Worts I, 35. ff. 40. 52.

2. \gesperrt{Das Firmament} I, 29. ff. Woher der H. im Hebräischen seinen Namen hat I, 40. 52.; ist ein Wunder, daß er nicht zerschmilzt oder bricht I, 29. f.; dünner und feiner als die Luft I, 30.; warum er blau erscheint I, 30. — S. \gesperrt{Feuerhimmel, Krystallhimmel, Sphären}.

3. \gesperrt{Aller Himmel Himmel}. Das ganze Firmament I, 37. f.

4. \gesperrt{Himmel der Seligen}. Davon können wir nichts Gewisses sagen I, 37.

5. H. \gesperrt{und Erde}. Haben jetzt ein Werkelkleid, werden einst ein Pfingstkleid tragen II, 1739. f.

6. H. \gesperrt{und Hölle}. Lieber mit Gott in der Hölle, als ohne Gott im H. II, 569.

Himmelfahrt Christi. S. \gesperrt{Christus} 2.

Himmelreich. Darin wird unter die Ungleichen Alles gleich getheilt I, 1074.

Himmelsheer. Die Sterne und Planeten I, 90. — S. \gesperrt{Gestirndienst, Zeichen} 1.

Himmelskörper. Lauf derselben durch Gottes Wort geordnet und erhalten I, 36.; wurden für vernünftige Creaturen gehalten I, 57., für ewig und gleichsam Götter I, 50.

Himmelsleiter. Mancherlei Auslegungen II, 394. ff.; wo Jakob sie sah II, 428.

Hiob. Der \gesperrt{Name} H. findet sich beim deutschen Adel II, 1015. Wer H. gewesen I, 1601; II, 383. 1014. ff.; gewaltig, weise, verständig im Wort des HErrn I, 1601.; II, 383. 1015.; Gott prüfte seinen Glauben und Beständigkeit I, 822. H.’s \gesperrt{Weib} hat nicht gleich im Anfang gemurrt II, 1326.; H.’s \gesperrt{Freunde} weise, gelehrte Männer II, 1016.

Hippocrates. Rettet eine des Ehebruchs Beschuldigte von Verurtheilung II, 601.

Hiram. Erkannte den rechten Gott I, 1052.

Hirten. Können viel zur Fruchtbarkeit der Heerde beitragen II, 697.

Hiskia. Sein Gebet keine Versuchung Gottes I, 1698.

Hoba. I, 891.

Hochmuth. S. \gesperrt{Hoffahrt}.

Hochzeit. Vom rechten Maß der Fröhlichkeit und des Gepränges I, 1717. f.

Hochzeitsgeschenke. Nicht verwerflich I, 1694. — S. \gesperrt{Geschenke}.

Hochzeitsmähler. Nicht zu verwerfen I, 1692. f.; Zweck derselben I, 1693. 1714. — S. \gesperrt{Gastmähler}.

Höfe. An den H. der Fürsten herrschen die sieben Todsünden mehr, denn sonstwo II, 1726.; selig, wer sich hüten kann, daß er sich nicht an die H. begebe II, 1789.

Hoffahrt. \gesperrt{Art und Natur}: ist teuflisch I, 985., eine Sünde wider das 1. Gebot I, 1452., thöricht II, 1388. f.; beschönigt sich mit Gottes Wort I, 994.; führt zu fleischlichen Lastern I, 441.; macht die Menschen toll und unsinnig II, 1388. f. 1392. f. \gesperrt{An wem} sie sich findet: alle Menschen sind hoffährtig I, 440. 980.; II, 1370. f. 1390.; würden nicht zu Engelsdiensten taugen I, 1687. f.; H. ist uns eingepflanzt I, 312. f. 985.; II, 1389. 1631.; wir können nicht anders als stolzieren I, 924.; H. findet sich in allen Ständen I, 981. f.; worauf die Menschen stolz sind II, 1388.; wir mißbrauchen die Ehre stets zu H. I, 314.; H. um so größer, je größer die Gaben I, 438. 477. 923. f.; H. folgt, wo die Lehre von der Gnade verlischt I, 1453. f.; H. war die Sünde der ersten Welt I, 439. \gesperrt{Strafe}: H. macht unselige Menschen II, 1389. f. 1392. f.; Gott ist ihr feind I, 299. ff. 477. 994. 1430. 1444. 1453.; II, 1372. 1389. 1393., straft sie I, 313. 437. ff. 1689. f.; II, 81. ff. 263. f. 1370. f.; durch H. verliert man die göttlichen Verheißungen II, 243., werden die meisten Erstgebornen in der Schrift unglückselig II, 987.; Gott stürzt die Hoffährtigen II, 242. \gesperrt{Bekämpfung}: Gott will die H. dämpfen, das wollen wir nicht leiden I, 312.; wie man der H. wehren kann II, 1390.; wir sollen uns die Exempel gestraften Hochmuths merken I, 439. \gesperrt{Erkenntniß}: die Heiden haben nicht erkannt, warum Gott die Hohen erniedrigt I, 313. f.

Hoffnung. Ist tröstlich II, 1217.; H. auf Gottes Verheißungen scheint vergeblich II, 50., ist es nicht I, 695., rechte H. macht nicht zu Schanden I, 810. 827.; Gott ist denen gnädig, die auf ihn hoffen I, 806.; H. sollen wir nicht verlieren, wenn wir auch Gottes Zorn fühlen I, 1004., wenn in Gefahr und Noth auch keine Hoffnung mehr da ist I, 615. 810. f. — S. \gesperrt{Eltern} 2, \gesperrt{Furcht} 5, \gesperrt{Geduld} 4, \gesperrt{Harren}.

Höflichkeit. S. \gesperrt{Sittigkeit}.

Hofschranzen. Finden sich an allen Höfen I, 1474.; suchen nur den eigenen Nutzen I, 1460.; hindern den Willen der Fürsten oft I, 1467. f.

Hofteufel. Ist nicht todt II, 1726.

Hoherpriester. Von seinem Kleide II, 1039., seinem Stirnblatte I, 1713.

Hölle. Was und wo sie sei II, 2067. f.; H. ist der Zorn Gottes II, 1720., des Anschauens Gottes beraubt sein I, 1762.; Strafe um so schwerer, weil man sie als verdient erkennen wird I, 255.; die größte Pein wird sein, daß die Gottlosen nicht fliehen können, ob sie gleich wollen I, 211. — S. \gesperrt{Gehenna, Gewissen} 3, \gesperrt{Himmel} 6.

Höllenfahrt Christi. S. \gesperrt{Christus} 2.

Horiten. Woher der Name II, 1001.; Edomiter im steinigten Arabien I, 880.; von Esaus Geschlecht aus Seir vertrieben II, 999.

Hosea. Weissagte wider Israel I, 779.; arbeitete 80 Jahre, bekehrte Viele II, 2047. — S. \gesperrt{Propheten} 2.

Hugo. Aergert sich am ehelichen Leben der Väter II, 556.

Hülfe Gottes. Gott kann und will helfen II, 1386.; seine H. fängt dann an, wenn Alles verloren scheint II, 1285.; Gott ist ein Helfer I, 777.; II, 893. f. 901. ff., hilft den Gottseligen I, 806. f. 818. 826. 899.; er verzieht wohl, hilft aber gewiß zu rechter Zeit I, 777. 807. 883. f. 891. f.; II, 5. f. 1677.; H. G.’s am nächsten, wenn die Noth am größten II, 1640.; darauf soll man geduldig warten I, 561. 1692.; wie man sich trösten soll, wenn Gott verzieht II, 5. 1324. f. — S. \gesperrt{Gebet} 7, \gesperrt{Harren}.

Hurenhäuser. Nicht zu gestatten I, 1231. f.

Hurer. Haben auch Noth, wie Ehemänner II, 362.

Hurerei. Die Erbsünde macht dazu geneigt I, 975.; mißfällt Gott I, 294.; bringt leiblichen und geistlichen Schaden I, 975.; H. Lediger ist Sünde I, 204. 1225.

Huß. Rechtschaffenes Glied der Kirche I, 459.; hatte unverzagten Muth und Beständigkeit I, 538.; war trunken vom Heil. Geist II, 2002. f.; strafte am Papstthum die Simonie II, 105.; sein Tod hat dem Papst nichts genützt II, 1950.

Hyperbel. Redefigur II, 69. f. 1978.

Hypothesen. Sind nützlich I, 33. f.

Hysteron proteron. Redefigur II, 548. f.; macht oft die Erzählung dunkel II, 572.; sehr gebräuchlich II, 549. ff., auch in der Schrift II, 1066. f. 1158.; Beispiele: II, 982. 1067. 1158. 1220. 1511. f.; auf dieselbe muß man bei Auslegung der Schrift achten II, 549. 1510. ff. — S. \gesperrt{Anticipatio, Recapitulatio}.

\section{I}

Ich kann nicht. Hüte dich vor dem „I. k. n.“ II, 751.

Idumäa. Frühere Bewohner II, 1008. f.

Indien. S. \gesperrt{Hevila}.

Inseln. Ob sich I. aus dem Meer erheben I, 662

Iphigenia. Gemälde ihrer Historie I, 342.

Ir. Theil Ninives I, 674.

Irad. Bedeutung des Namens I, 384. f.

Irrlehren. Eine bringt die andere II, 1277. 1829. — S. \gesperrt{Lehre} 3.

Irrlehrer. Durch sie versucht uns Gott I, 1533. f.; betrübend, wenn sie die Lehreinigkeit stören II, 245.; ihnen bringt man jückende Ohren zu I, 726. f. — S. \gesperrt{Ketzer, Propheten} 3, \gesperrt{Rottengeister}.

Irrthum. Folgt aus vorwitzigen Disputationen I, 16. f. 192. f.; Gottes Gnade bringt davon zurück I, 723. — S. \gesperrt{Irrlehren}.

Isaak. 1. \gesperrt{Geburt} I.’s. I, 1368. ff. I.’s G. war wunderbar, nicht Folge fleischlicher Wohllust I, 1376., geschah kraft der Verheißung I, 1101., er ward erst im Glauben empfangen, ehe er leiblich gezeugt ward II, 2028. f.; ist ein Bild der Auferstehung I, 1096.

2. \gesperrt{Name} I.’s. Woher der Name I, 1104.

3. \gesperrt{Erscheinungen und Offenbarungen}. II, 138. ff. 210. ff. Gott redet nicht so oft mit ihm wie mit Abraham II, 1., nur ein- oder zweimal II, 140.

4. \gesperrt{Verheißungen} I. geschehen. II, 138. ff. 210. ff.

5. \gesperrt{Bund} Gottes mit I. Ein zwiefacher: der Beschneidung und der Verheißung I, 1115.

6. \gesperrt{Glaube} I.’s. II, 3. ff. In der \gesperrt{Erkenntniß} mußte er noch wachsen II, 77. ff. 93. ff.; hat die Lehre von der Dreieinigkeit und Auferstehung der Todten verstanden II, 215. In I. war ein großes Licht des \gesperrt{Glaubens} I, 519., großes Verlangen nach dem verheißenen Samen II, 704.; wie sich sein Glaube gezeigt I, 1508. ff.; II, 128. ff. 131. ff. 136. f. 703.; um seines trefflichen Gottesdienstes willen heißt Gott die Furcht I.’s II, 703.

7. \gesperrt{Tugenden} I.’s. I. hatte sich an \gesperrt{Geduld} gewöhnt II, 130.; war sehr geduldig II, 237. f., nicht rachsüchtig II, 206. f.; wie sich seine Geduld gezeigt II, 128. ff. 136. 195. 204. f. 209. 226. f.; gehorsam I, 1512.; keusch II, 23. f.; liebte seine Mutter I, 1647. 1742.

8. \gesperrt{Werke} I.’s. Häusliche Werke II, 203. ff. 234. f.; I. predigte II, 185. 215. 221. ff. 999., strafte die Sünde II, 231. ff., tröstete Jakob über sein Unglück II, 980.

9. \gesperrt{Reisen} I.’s. II, 122. ff. 198. ff.

10. \gesperrt{Sünden und Schwachheiten} I.’s. I. verleugnet sein Weib II, 150. f. 155. f. 173. f.; war ein Melancholiker II, 291.

11. \gesperrt{Anfechtungen u. Leiden} I.’s. An I. sehen wir, wie wunderbar Gott seine Heiligen regiert II, 1. I.’s \gesperrt{Anfechtungen} übertreffen alle Kasteiungen der Mönche II, 135. ff.; von seinen mancherlei A. II, 10. f. 128. f. 190. f. 194. ff. 199. f. 219. f. 344. f.; hat bis in sein vierzigstes Jahr wider die Flammen der Unzucht gestritten II, 20. f. I. mußte viele und schwere \gesperrt{Widerwärtigkeiten} erfahren II, 236. f. 982. 1023. f. 1158. f., auch viel Mühseligkeit in seiner Haushaltung II, 191. und Hohn um Rebekkas Unfruchtbarkeit willen II, 27. f. In den A. und L. hielt er sich an Gottes Wort II, 198. ff., ward er getröstet II, 7. 191. f., siegte er durch den Glauben II, 189. ff. 210. — S. unten 15.

12. I. \gesperrt{und Abimelech}. II, 150. ff. 169. ff. 203. ff. 225. ff. I. hat A. in der rechten Erkenntniß unterrichtet II, 185. 227. f.; A. demüthigte und bekehrte sich II, 232. f., ward ein Schirmherr der Kirche I.’s II, 230.

13. I. \gesperrt{und Ismael}. I. ist aus der Verheißung geboren, Ismael nicht I, 1029. f.; I. hatte geistliche, Ismael leibliche Verheißungen I, 1395. f.; Ismael hatte größeren leiblichen Segen, als I. II, 10. f.; sind ein Bild der rechten und falschen Kirche II, 67.; der Saracenen Fabeln von ihnen I, 1504.

14. I. \gesperrt{und die Philister}. II, 189. ff. 203. ff. I. bringt die Philister zur rechten Erkenntniß II, 227. f.

15. I. \gesperrt{und Rebekka}. I, 1741. ff.; II, 19. ff. 32. ff. 150. ff. Beide mußten Anfechtungen und Leiden ertragen II, 25. ff. 45. 235. ff. 343. ff. 373. f. 993.; darin nahmen sie ihre Zuflucht zum anhaltenden Gebet II, 27. ff. 33. und frugen bei Sem den HErrn um Rath I, 712.; II, 55. f. 59. Sie lebten keusch II, 80.; hatten freundlichen Streit über die Verheißung, daß der Größere dem Kleineren dienen sollte II, 250. ff.; R. regierte das Haus, da I. blind geworden II, 331. f. 352. — S. \gesperrt{Esau} 3.

16. I. \gesperrt{und Rebekka, Jakob und Esau}. II, 67. ff. 246. ff. 261. ff. 269. ff. 280. ff. 292. ff. 307. ff. 325. ff. 335. ff. 358. ff. 369. ff. Esau wußte sich beim Vater einzuschmeicheln, Jakob hielt sich zur Mutter II, 256. ff.; I. zog Esau vor II, 92. ff., R. den Jakob II, 91. f., wobei sie auf Eingeben des Heil. Geistes handelte II, 95. f.; Esau war Licentiat zum Papstthum, bis Jakob den Segen erlangte II, 281. und führte die Herrschaft II, 330. ff.; ob Betrug, da R. Jakob den Segen verschaffte, Sünde war II, 261. ff. 310. ff., war ein
\quellseite{556}
gottseliger Betrug II, 244.; warum I. denselben nicht bemerkte II, 501. Der Segen Jakobs faßt alle drei Stände in sich II, 317., der Segen Esaus ist eigentlich kein Segen II, 322. ff.; warum sie Jakob wegschickten II, 444. f. und gar zu abgöttischen Leuten II, 382. ff.; in Jakobs Abwesenheit führte Esau die Herrschaft II, 344., verwaltete er das Priesterthum II, 999.; Esau ward von seinen Eltern nicht gar verworfen II, 1001.

17. \gesperrt{Tod und Begräbniß} I.’s. II, 979.

Vgl. \gesperrt{Abraham} 25.

Isaschar. 1. \gesperrt{Jakobs Sohn}. Wann geboren II, 551. 570.; wie der Name zu lesen II, 571.; Jakobs Segen II, 2018. ff. — S. \gesperrt{Jakob} I, 19; II; \gesperrt{Joseph} 12.

2. \gesperrt{Der Stamm}. Wohnsitz II, 2018. ff.; ein guter fauler Pater II, 2019.; nach Assyrien geführt II, 2020.

Ismael. 1. \gesperrt{Ismael}. I, 993. ff. 1451. ff.; II, 9. ff. Des Namens Ableitung I, 1013., Bedeutung I, 996.; war hoffährtig und vermessen I, 1387. ff. 1420. 1428. ff.; II, 242., ein antinomischer Epicurer I, 1429., ward ein Kind der Verheißung I, 1437. f. 1452. 1455., ein geschickter und gelehrter Prediger I, 1452. und großer Theologe II, 12.; behielt Abrahams Gottesdienst bei II, 12. 984.; war ein Kriegsmann, tapfer I, 1454. und fromm II, 13.; ward wohl selig I, 1002. 1030. 1051. 1418.; II, 830. — S. \gesperrt{Abraham} 26, \gesperrt{Isaak} 13.

2. \gesperrt{Ismael und seine Nachkommen}. I, 997. ff. II, 9. ff. I.’s Nachkommen waren nicht von dem Segen, Abraham verheißen, ausgeschlossen I, 1023. f. 1113. f.; viele hatten die rechte Erkenntniß I, 1114. 1418. 1420.; I.’s N. behielten die Beschneidung I, 1114., die Lehre vom Messias II, 1009.; viele wurden selig I, 1002.; II, 12.; sind unter die Türken vermengt I, 1114.; ließen die rechte Lehre fahren II, 13. — S. \gesperrt{Araber, Saracenen}.

Ismaeliter. Trieben Handel mit Gewürz II, 1124. f.; haßten Jakobs Geschlecht II, 1123. — S. \gesperrt{Ismael} 2.

Israel. 1. \gesperrt{Der Name}. Bedeutung desselben II, 795. 952. f.; wird in der Schrift mit besonderem Nachdruck gebraucht II, 976.; setzt einen Unterschied zwischen Fleisch und Geist II, 951. — S. \gesperrt{Jakob} 3.

2. \gesperrt{Das Königreich} I. Hieß auch das K. Joseph II, 1071.

Israeliten. Israeliten nach dem Fleisch und nach dem Geist II, 1938.

Italien. Früher sprach man da Griechisch I, 658. f.

Italiener. Verachten andere Völker II, 1776. 1786.; grüßen sehr höflich II, 737.; spotten über Christen II, 1995.; eheliches Leben II, 156.

\section{J}

Jabal. Bedeutung des Namens I, 386. f.

Jacobus. Macht einen närrischen Schluß in seiner Epistel I, 1536.

Jagd. Ist nicht sündlich II, 86. f.; geht meist nicht ohne Sünde ab II, 85. f.; Zweck derselben II, 86. f.

Jäger. Sind nicht aufrichtig II, 87.; versündigen sich oft II, 83. 85.

Jahr. Anfang desselben vor der Sintfluth I, 543. f.; jüdische Eintheilung I, 543. f. 586.

Jakob. I. J., \gesperrt{der Patriarch}. 1. \gesperrt{Ehre und Bedeutung}. J. war kein ungelehrter Bauer, sondern ein Prophet und Bischof II, 1798. 1959. f., ein Patriarch, hatte Gottes Wort, Gott selbst und seine Engel stets bei sich II, 1731., ein Patriarch und Licht der Welt II, 374. f.; wird durch den Namen „Israel“ geehrt II, 953.; stand bei den Egyptern in Ansehen II, 2072.

2. \gesperrt{Das Leben} J.’s. J. ist ein wunderbares Exempel der göttlichen Regierung II, 1068. 1091. f.; Chronologisches, sein Leben betreffend II, 1024. 1062. f. 1065. f.

3. J.’s \gesperrt{Namen}. Anlaß zu den Namen Jakob und Israel II, 795.; Aehnlichkeit beider N. II, 952.; heißt J. von der Fußsohle II, 78.; bekommt den N. Israel II, 793. ff. 950. ff.; ward billig Israel genannt II, 817.; warum J. später nicht immer Israel genannt worden II, 956. f.; „Jakob“ soll das wirkliche, „Israel“ das beschauliche Leben bezeichnen II, 953. — S. \gesperrt{Israel} 1.

4. \gesperrt{Kenntnisse und Gaben} J.’s. Hatte die \gesperrt{Kunst}, bunte Lämmer zu erzielen, gelernt II, 601. 607.; \gesperrt{war} klug II, 592. f. 604. f.; anfänglich mit seinen vier Weibern sehr arm II, 541.; woher er das Geld hatte, Laban zu bezahlen, was gestohlen oder zerrissen wurde II, 698.; ward durch die Kunst, bunte Lämmer zu erzielen, sehr reich II, 606.; seinen Reichthum hatte er Laban nicht geraubt, von Gott bekommen II, 622. ff. 636.

5. \gesperrt{Erscheinungen und Offenbarungen}. II, 390. ff. 405. ff. 419. ff. 722. f. 900. ff. 948. ff. 1750. ff. J. hatte verschiedene Erscheinungen II, 1850.; Gott redete nicht so oft mit ihm wie mit Abraham II, 1. — S. \gesperrt{Himmelsleiter}.

6. J.’s \gesperrt{Glaube}. J. kannte die Lehre von der heiligen Dreieinigkeit II, 941., von der Menschwerdung, Kreuzigung und Auferstehung des Sohnes Gottes II, 1867., von den guten Engeln II, 725. f. Er \gesperrt{glaubte} die Dreieinigkeit II, 1863. ff., die Auferstehung der Todten und ein ewiges Leben II, 1800. f. 1901. 1903.; war voll Heiliges Geistes, durch dessen Erleuchtung er viel aus den Verheißungen nahm II, 1875.; wuchs in der Erkenntniß II, 1852. f.; verstand anfänglich die Verheißung von der Mehrung seines Samens nicht recht II, 962. f. 977. f.; meinte erst, der verheißene Same werde von Joseph kommen II, 970. 1061. f. Hatte den \gesperrt{wahren Glauben} II, 1905., glaubte an Christum II, 1841.; glaubte, darum waren ihm alle Dinge möglich II, 814.; hatte einen ritterlichen, starken G. II, 379.; sein G. war bald stark, bald schwach II, 1870. f.; war stark II, 750. 840., ward schwach II, 744. 896. ff. 1444. ff. 1556. ff. 1568. ff.; wie sich sein G. bei verschiedenen Gelegenheiten gezeigt II, 541. ff. 619. f. 635. f. 689. 1899. ff.; da er nach Egypten kam, gedachte er der Verheißung 1 Mos. 15, 16. II, 1797.

7. \gesperrt{Tugenden} J.’s. J. war demüthig II, 81. gegen Gott II, 759. 829., Laban II, 580. ff. 590., Esau II, 768. 819. 830. ff. 839.; ehrerbietig gegen die Obrigkeit II, 1797.; fleißig II, 1239.; freundlich gegen das weibliche Geschlecht II, 935.; fromm, aufrichtig, beständig, nicht albern II, 87. f.; geduldig II, 237. f. 494. ff. 498. ff. 580. f. 680. 701. 899. 1090.; gehorsam II, 269. ff. 366.; keusch II, 488. ff. 491. 498. f. 529. ff. 551. f. 555.; liebte die Seinen II, 944.; mitleidig und barmherzig II, 510. 516. 1669.; treu II, 589. 696. ff.

8. \gesperrt{Werke} J.’s. J. war kein Kriegsmann, suchte zuerst das Reich Gottes II, 88. f.; alle seine Werke gefielen Gott wohl II, 616. f. 1098.; mußte nach Isaaks Tode die Last des Regiments allein tragen II, 980.; hat sich des Hausstandes angenommen, um die Seinen und seine Habe gesorgt II, 1071. f. 1076.; bediente sich bei seinen Unternehmungen der dargebotenen und vorhandenen Mittel II, 637. f. 746. ff. 763. ff. \gesperrt{Einzelne Werke}: daß er den Stein von dem Brunnen wälzt, war ein W. des Glaubens und der Liebe II, 476. ff.; schwor bei dem wahren Gott II, 721.; kaufte ein Stück Acker zum Eigenthum II, 852. f.; behielt das Land und anderen Raub der Sichemiten II, 874. f. 877. 890. 1070. f.; pachtete Aecker II, 1069.; besuchte oft Isaak und Esau II, 853. Ein \gesperrt{Exempel} ist uns Jakob nicht in seiner Geduld, da er durch Lea betrogen ward II, 506., in seiner Polygamie II, 510. ff. 533., da er durch Kunst seinen Lohn mehrte II, 608. f.

9. J.’s \gesperrt{Gelübde}. II, 438. ff. Er gelobte eine Kirche und Schule zu bauen und erhalten II, 449. ff.; handelte nach Gottes Wort II, 440.; dies war ein Gelübde der Danksagung II, 449.

10. \gesperrt{Gebet} J.’s. J. hat in den verschiedenen Nöthen gebetet II, 750. 771. 898. 1753. f.; sein Gebet, da ihm Esau entgegenzieht, hat alle Eigenschaften eines guten G.’s II, 752. ff.; dadurch machte er den Feind Esau zum Freunde II, 828. f.; sein Exempel soll uns zum G. reizen II, 919.

11. \gesperrt{Die Kirche} zu J.’s \gesperrt{Zeit}. II, 922. ff. 934. ff. J. war ein Hoherpriester, voll Heil. Geistes II, 1751. f.; beschäftigte sich mit Gebet und Gottes Wort II, 958.; hatte in seiner Gemeinde auch Fremde II, 908. ff.; baute Altäre II, 853. 939.; baute die erste Kirche, wo er die Himmelsleiter gesehen I, 779.; lehrte II, 854. f. 907. ff. 936. ff. 1833.; wie und was er gepredigt II, 875. f.; lehrte die Dreieinigkeit II, 1863. ff.; predigte von Christo II, 563. 1751. f.; strafte II, 891. ff.; tröstete II, 753. f. 762. f. 934. f.; richtete eine Reformation seiner Kirche an und begann mit der Lehre II, 910. ff.; bekehrte im Hause Labans wohl Etliche II, 649.; seine Söhne waren eine böse Frucht II, 1087. f.; gab wohl den Kindern Ebers den Zehnten II, 961.

12. \gesperrt{Ehestand} J.’s. J. enthielt sich des Ehestandes beinahe 70 Jahre II, 90. 248. 484. f.; an seiner Polygamie stoßen sich fleischliche Menschen II, 528. f., die Papisten II, 539. ff.; sein E. war sehr keusch II, 572. f.; hielt sich zu den vier Weibern um der Nachkommenschaft willen II, 491. 531. f.; hielt sich nach Rubens Blutschande zu keinem Weibe mehr II, 977. — S. unten 18.

13. J.’s \gesperrt{Reisen}. II, 375. ff. 456. ff. 475. ff. 608. ff. 636. ff. 722. ff. 763. ff. 844. ff. 848. ff. 922. ff. 1750. ff. 1761. ff. 1782. ff. J.’s Reise nach Haran war unansehnlich und verächtlich II, 465., ein schweres Werk, so aus dem Wort und Glauben gekommen ist II, 473. f.; wie lange er in Suchoth geblieben II, 844. f. 856.; zog als Fremdling hin und her in Canaan, verließ aber Isaak nicht II, 1024. f.; wann er nach Egypten zog I, 959.

14. \gesperrt{Sünden und Schwachheiten} J.’s. Hatte auch mit dem Fleisch und Schwachheiten zu kämpfen II, 805. f. 1503.; von seinem Fleisch kam es, daß er von Esau Grausamkeit befürchtete II, 749., das Gebet zu lange verschob II, 752.; handelte unbedacht und hätte fast Rahel auf die Fleischbank geliefert II, 683. Er sündigte nicht, da er die Erstgeburt kaufte II, 102. f., durch Betrug den Segen erlangte II, 264. ff., zum abgöttischen Laban zog II, 385., Rahel um ihrer Schönheit willen liebte II, 484. ff., zwei Schwestern zur Ehe nahm II, 510. ff., durch Kunst seinen Lohn mehrte II, 606. ff., heimlich von Laban abreiste II, 636. f.; befleckte sich nicht bei Laban mit Abgötterei II, 735. f.; heuchelte nicht in seinem Gespräch mit Esau, ward kein Götzendiener dabei II, 836. ff.; log nicht, da er Esau versprach, gen Seir zu kommen II, 843. f.; hatte kein Theil am Mord der Sichemiten II, 888.

15. \gesperrt{Anfechtungen und Leiden} J.’s. J. mußte Unglück leiden eben wie wir II, 1077. f.; hatte viel und schweres Kreuz und Anfechtungen II, 301. ff. 1062. f. 1094. f. 1135., hatte von Jugend auf Uebung im Kreuz II, 1772. 1800. f.; sein Leben war erbärmlich nach dem Fleisch II, 1095. f.; einem Heiden konnte es nicht übler ergehen denn J. II, 1094. f.; hat mehr erlitten denn ein Märtyrer II, 701.; Anfechtung und Unglück kam über ihn wie ein Sturm II, 900. 918. 966. f. 1094. 1120. 1165.; in seinem Hause hatte der Teufel große Gewalt II, 1094. 1120.; von seinen Leiden hatte er lange graue Haare und ein verwittertes runzlichtes Gesicht, da er zu Pharao kam II, 1799. Er ward \gesperrt{angefochten} II, 380. 386. ff. 389. 392. f. 657. f. 740. ff. 766. 855. ff. 861. f. 867. 896. ff. 942. 946. 962. ff. 973. f. 976. 1752. J.’s \gesperrt{Kampf}: II, 774. ff. 792. ff.; war ein Kampf mit Gottes Sohn II, 780. ff. 799. f. 1866.; währte etwa zwei Stunden II, 771. 788.; war eine hohe Anfechtung II, 774., ein Kampf des Glaubens mit dem Mißglauben II, 787., der höchsten Werke eins II, 1076., ein Kampf mehr des Glaubens als der Arme II, 787. f.; woher es kam, daß J. stärker war als Gott II, 789. f.; mit welchem Segen ihn der Mann gesegnet II, 799.; J. hat wohl nicht immer gehinkt II, 850.; ist in diesem Kampf ein Bild aller Gläubigen II, 811. f. \gesperrt{Josephs Verkauf} war ein schweres Kreuz II, 1021. f. 1094. 1123. 1127., eine schwere Anfechtung II, 1020. f. 1152., stürzte J. in den Abgrund der Hölle II, 1063. f.; dies Leiden gewann einen lieblichen Ausgang II, 1022. J. ward von Esau und dessen Weibern \gesperrt{verachtet} II, 240. f.; mußte bei Laban viel \gesperrt{leiden} II, 608. ff. 615. f.; Judas Blutschande schmerzte ihn II, 1022.; ward durch die Theurung trotz der Verheißung bedrückt II, 1560.; Gott \gesperrt{versuchte} ihn durch seine eigenen Söhne II, 1536.; J. \gesperrt{betrübte} sich bei Dinas Schwächung und Josephs Verkauf mehr über das Aergerniß als das Unglück II, 1089. Gott stärkte und tröstete ihn in den A. II, 396. 424. f. 624. f. 902. ff. 948. ff. 1063.; darin kämpfte sein Glaube II, 380. wider das Fleisch und Unglauben II, 762. 1584. f. 1587. ff.; er hielt sich an das Wort II, 658. 752. ff. 1081. f. 1095. f.; wie er sich getröstet II, 1022. 1092.

16. \gesperrt{Tod, Begräbniß, Auferstehung} J.’s. II, 1832. ff. 1899. ff. 2054. ff. 2072. ff. Starb in großem Glauben II, 2063. f.; warum er in Canaan begraben sein wollte II, 2064. f.; herrlicher begraben als Christus II, 2072.; ist wohl mit Christo auferstanden II, 1836.

17. \gesperrt{Zeitgenossen} J.’s. II, 80. f. 89.

18. \gesperrt{Jakob und seine Weiber}. II, 620. ff. 636. ff. J. hatte die Mägde zu Weibern, nur weil und so lange Lea und Rahel dies wollten II, 551. f. 555. J. \gesperrt{und Lea}: II, 493. ff. 518. ff. 544. ff. 564. ff. Lea ward wohl von Laban gezwungen, J. zu betrügen, war aber selbst auch nicht abgeneigt II, 498. 502. f.; was von ihrem Umgang in der Brautnacht zu halten II, 503. ff.; J. hätte sich nach des Papsts Lehre von Lea dürfen scheiden I, 1729. f.; behielt sie aus Barmherzigkeit, Liebe und Geduld II, 510. 516.; sie gab J. ihre Magd, weil sie nach dem verheißenen Samen Verlangen hatte II, 555. f. J. \gesperrt{und Rahel}: II, 479. ff. 507. ff. 518. ff. 528. ff. 544. ff. 564. ff. 639. ff. 961. ff. Wie sich Rahel verhalten, als Laban dem Jakob Lea beilegte II, 495. 497. f. 502.; J. und Rahel liebten sich herzlich II, 493. 965.; zankten sich II, 538.; R. war eine treue Gesellin J.’s II, 963. 969.; J. wußte nicht, daß R. Labans Götzen gestohlen hatte II, 648. f., brachte sie daher in große Gefahr II, 686. f.; daß sie J. ihre Magd gibt, ist ein Werk des Geistes II, 545. J. \gesperrt{und Bilha}: II, 544. ff. 972. ff. J. \gesperrt{und Silpa}: II, 554. ff. 564.
\quellseite{557}
S. heißt J.’s Kebsweib, war aber sein rechtes Eheweib II, 973.

19. J. \gesperrt{und seine Söhne}. II, 1143. ff. 1153. ff. 1438. ff. 1523. ff. 1535. ff. 1558. ff. 1575. ff. 1738. ff. 1927. ff. 1935. ff. 1959. ff. 2011. ff. 2039. ff. 2054. ff. Jedes hatte seine besondere Haushaltung II, 1742.; J. hätte lieber die Anderen als die Söhne Rahels entbehrt II, 1035. f.; sie waren unnatürlich gefühllos gegen J. und Isaak II, 1054. 1144. f.; J. hatte Verdacht auf seine Söhne II, 1569.; wie sich J.’s letzter Segen im Einzelnen erfüllt hat, weiß man nicht II, 2036. f.

20. \gesperrt{Jakob und Joseph}. II, 1032. ff. 1037. ff. 1770. ff. 1796. ff. 1832. ff. 1846. ff. 1863. ff. 1873. ff. 1881. ff. 1899. ff. 2039. Kein Wunder, daß J. den Joseph lieber hatte als die anderen Söhne II, 1224., das war recht II, 1032. f.; warum er ihn mehr liebte II, 1026. ff. 1034. ff.; nach J.’s Willen gebührte Joseph die Erstgeburt II, 1057. 1060.; J. hoffte von ihm den verheißenen Samen II, 1060.; welcher Art der „bunte“ Rock war I, 1037. f. 1040.; dadurch wollte J. den Joseph als obersten Regenten und Erben auszeichnen II, 1040.; J. läßt sich Josephs ersten Traum gefallen II, 1057.; seine Gedanken über den zweiten Traum II, 1060. ff.; wie sich der zweite Traum erfüllt II, 1066.; warum Joseph keine Gesandten von Egypten aus zu seinem Vater geschickt II, 1427. ff. 1724.; sie haben wohl mit einander über den Stamm Levi berathen II, 1846. f.; J. ehrt Joseph damit, daß er ihm Sichem schenkt II, 1907. — S. oben 15.

21. J. \gesperrt{u. Ephraim u. Manasse}. II, 1846. ff. 1863. ff. 1885. ff. Warum J. sie zu Kindern annimmt II, 1853. ff.; thut es, nicht aus ihrem Verdienst II, 1863.; zieht Ephraim vor, nicht aus Ephraims Verdienst oder Manasses Schuld, sondern aus Eingebung des Heil. Geistes II, 1886. 2039.

22. J.’s \gesperrt{Gesinde}. Hatte viel Gesinde aus der Verwandtschaft Nahors und Tharahs II, 695.; Manche zogen mit J. in seine Heimath II, 695. f.; war fromm durch J.’s Lehre II, 773.

23. J. \gesperrt{u. Esau}. II, 67. ff. 97. ff. 325. ff. 722. ff. 738. ff. 763. ff. 814. ff. 825. ff. 840. ff. J. kommt bei der Geburt in Gefahr II, 70.; durch ihre Geburt will Gott alle Gerechtigkeit des Fleisches abschneiden II, 70. ff.; Gott zieht J. dem Esau vor II, 41.; J. bekommt die Erstgeburt II, 91.; warum J. die Erstgeburt kauft, da sie doch ihm gehört II, 103. 111. ff.; das Linsengericht war der Kaufpreis II, 117. f.; von dem Unterschied ihrer Namen II, 78. f.; wurden von Abraham unterrichtet II, 80. f.; wie sie sich verhalten, da Isaak nach Gerar zog II, 137. f.; Unterschied Beider nach Beruf und Gesinnung II, 81. ff.; E. hatte größeren leiblichen Segen als J. II, 248. 251. f. 1011. 1013. f., doch ist J.’s Herrlichkeit größer als E.’s II, 1012.; E. verachtete den J. II, 82. f.; J.’s Demüthigung vor E. streitet nicht mit der Verheißung II, 769. 821. ff.; J. hat wohl E. besucht II, 844. f. 853.; waren bei Isaaks Tode ausgesöhnt II, 982.; sind ein Bild der rechten und falschen Kirche II, 67.

24. J. \gesperrt{und Laban}. II, 479. ff. 493. ff. 507. ff. 580. ff. 591. ff. 608. ff. 620. ff. 636. ff. 654. ff. 664. ff. 686. ff. 707. ff. 713. ff. Laban beneidete J. II, 615.; begehrte alles Gut J.’s an sich zu bringen II, 635.; schmähte ihn wohl bisweilen als einen Dieb II, 596.; hatte ihm wohl das Haus verboten II, 621.; verdächtigte ihn, als sei er nicht treu und gottesfürchtig II, 711. ff., als sei er ein Ehebrecher II, 717.; J. diente L. treulich II, 483.; sein „Betrug“ an L. mit den Stäben war nöthig und recht II, 582.; J. strafte L. wegen des Geizes II, 585. f., wegen der Härte gegen seine Töchter II, 589. f.; ein Engel tröstete J. über L.’s Ungerechtigkeit und Geiz II, 624. f.; was J. zu erwarten hatte, hätte er L. seine Absicht wegzuziehen mitgetheilt II, 681. ff. L. ist ein Bild eines Heuchlers, J. eines frommen Mannes II, 723.

Vgl. \gesperrt{Isaak} 16.

II. J.’s \gesperrt{Söhne}. 1. \gesperrt{Die zwölf Söhne}. II, 977. ff. Reihenfolge und Zeit ihrer Geburt II, 547. ff. 1035.; nahmen Cananitinnen zu Weibern II, 909. 934. f.; pflegten anfänglich ihre Töchter nur mit Beschnittenen zu verehelichen II, 878. f.; waren Priester unter den Heiden II, 1209.; stolz auf die Verheißung, daß sie Herren des Landes werden sollten II, 876. 879.; die \gesperrt{Söhne Leas} waren stolz, verachteten die Söhne der Mägde II, 1026., wollten von Joseph geehrt sein II, 1031. ff.; Satan hat sie in scheußliche Sünden verwickelt II, 978.; es ist zu verwundern, daß sie so gottlos waren II, 1026.; sind trotz ihrer Sünden selig geworden II, 43.

2. J.’s \gesperrt{Söhne u. die Sichemiten}. II, 864. ff. 879. ff. Sie thaten nicht recht, daß sie betrüglich handelten II, 873. ff. 877. ff.; ihr Angriff auf die S. war eine kühne That II, 878., Andere halfen ihnen II, 879., war unrecht II, 889. f.; wurden deswegen gestraft II, 888.

III. \gesperrt{Jakob}, St., in Spanien. Ein Abgott I, 1531. f.; zu ihm wallfahrtete man I, 864.

Japheth. 1. \gesperrt{Noahs Sohn}. I, 627. ff. 640. ff. Bedeutung des Namens I, 651. ff.; Noahs Erstgeborner I, 435.; wann geboren I, 435. 502.; heißt bei den Heiden Japetus I, 659.; Erfüllung des Segens Noahs über J. I, 653. ff. 665. — S. \gesperrt{Noah} 14.

2. J.’s \gesperrt{Geschlecht}. I, 656. ff. Wer dazu gehört I, 654. 659.; Wohnsitz I, 656. 663. f. 666.; die Namen der Söhne J.’s nicht bedeutungslos I, 664.; unter ihnen die heiligsten Väter I, 662., die uns ein Vorbild sind I, 664.; geriethen allmählich in Abgötterei I, 663.

Javan. 1. \gesperrt{Japheths Sohn}. I, 658. f.

2. J.’s \gesperrt{Nachkommen}. Seine Söhne waren die allergewaltigsten Könige I, 662.; von Japheth in der Lehre unterwiesen I, 662.; seine Nachkommen waren erfahrene Schiffsleute I, 662.; behielten wohl den rechten Gottesdienst I, 662.; welche Völker „Javan“ bezeichnet I, 658. f.

Jehovah. Name Gottes, nicht einer Creatur I, 274.; warum „Jehovah“ Tetragrammaton heißt II, 937.

Jehu. König. Ob aus dem Stamm Gad II, 2033.

Jekel. Toller, unerfahrener Geist I, 1212.; wollte für heilig und gerecht gelten II, 708.

Jephthah. Aus dem Stamm Manasse II, 2041.

Jeremia. Mußte Feindschaft und Gefahr ertragen, weil er strafte I, 457. — S. \gesperrt{Propheten} 2.

Jericho. Lage I, 780.

Jerobeam. Vom Stamm Benjamin II, 1432.; weise, klug II, 425. f.; sein Götzendienst zu Bethel II, 425. f. 959.; wie er denselben vertheidigte II, 907.

Jerusalem. Lage II, 1186.; verschiedene Namen I, 1492.; II, 427. f. 851. 903. 936.; heißt metaphorisch Gilead II, 715.; daselbst regierte Sem I, 1492. — S. \gesperrt{Bethel}.

Jesaia. Las und legte Mose fleißig aus I, 1754.; wiederholt sich oft II, 1356. 1948. — S. \gesperrt{Propheten} 2.

Jethro. Durch wen zur Erkenntniß gebracht II, 1831.

Jisca. Ob J. und Sarah dieselbe Person I, 723.

Jobab. II, 1014. ff. Ob er Hiob war II, 1014. ff.

Johannes. 1. Bedeutung des Namens II, 1553.

2. J. \gesperrt{der Täufer}. That seine besonderen Werke auf Gottes Befehl I, 868.

3. J., \gesperrt{Churfürst}. Bekannte zu Augsburg I, 1336. f.

4. J. \gesperrt{Friedrich}, Churfürst. Fleißig in seinem Regieramt II, 1790. f.; nahm sich der Kirche und Schule an I, 1632.; II, 453. 633.

5. J. XXIII., Pabst. S. \gesperrt{Franciscaner}.

Jona. Hat Viele bei reiner Lehre erhalten II, 2047.

Joseph. 1. \gesperrt{Im Allgemeinen}. II, 1312. ff.

2. \gesperrt{Ehre u. Bedeutung}. Keines Andern Leben wird so genau in der Schrift beschrieben II, 1247.; seine Geschichte ist die schönste im ersten Buche Mose II, 1022.; ein vortrefflicher Mann, trefflicher als Abraham und die anderen Patriarchen II, 1365. 1398.; kann mit den großen Männern des A. und N. T.’s verglichen werden II, 1946.; größer denn Achilles, Scipio, Alexander II, 1396.; ein Lehrer der Welt II, 1365.; eine Figur Christi II, 1119. f. 1127. 1136. f. 1156. 1321. 1328. 1836.; der Morgenstern im Hause Jakobs II, 1028., wo er hochgeehrt war II, 555.; ein Exempel für alle Junggesellen II, 1224. f.

3. \gesperrt{Geburt und Erziehung}. Wann geboren II, 547. ff. 572.; lernte von Jugend auf Gottes Wort II, 1316. f.

4. \gesperrt{Namen}. Die Namen J. und Ephraim gebraucht die Schrift ohne Unterschied II, 1847. 1863.

5. \gesperrt{Gaben}. J. war reichbegabt II, 1036., ein politischer Mensch und vortrefflicher Theologe II, 1688. 1793. 1823.; hatte hohen Verstand, großen Geist, Erfahrung II, 1639. f.; war weise und verständig II, 1341. 1780. f. 1793.; ein erfahrener, muthiger Feldherr II, 1375. ff.; hatte den Heil. Geist II, 1295. f. 1339., der in ihm geschäftig war II, 1297. f.; war voll Weisheit und geistlichen Verstandes II, 1433. f.; verstand den Kampf des Gewissens II, 1685.; hatte einen prophetischen Geist II, 1302. 1348., die Gabe der Weissagung II, 1640. und Träume auszulegen II, 1298.; war Priester und Prophet II, 1297.; von seiner Schönheit II, 2041. f.

6. \gesperrt{Träume} J.’s. II, 1045. ff. Wie der zweite Traum vollkommen erfüllt werden konnte, da Rahel gestorben war II, 1062. f. 1065. ff.; sind erfüllt II, 1066. 1718. 1834.

7. \gesperrt{Glaube} J.’s. J. war ein guter Zunder, fing leicht das Feuer des Worts II, 1096.; woraus seine Bibel bestand II, 1823.; erkannte, daß er von seinem Amte müsse Rechenschaft geben II, 1394.; hatte den Glauben und die Hoffnung der Auferstehung II, 1676.; sein Glaube war brünstig II, 1224., fest II, 1316. f.; er bekannte und übte denselben II, 1228. 1297.; damit überwand er Alles II, 1228. — Siehe unten 11.

8. \gesperrt{Tugenden} J.’s. II, 1025. ff. J. ist ein Wunderwerk II, 1246. ff.; hatte allerlei Tugenden, war ein solcher Mann, den unser HErr Gott gemacht II, 1262.; die Zucht war wohl gerathen II, 1224.; war tugendhaft II, 1036.; hatte alle Tugenden an sich II, 1244. f. 1251., so viel als nur bei dieser unserer Natur möglich ist II, 1225.; hatte größere Tugenden als Achilles, Scipio und Alexander II, 1715.; war ein fester, starker, lebendiger Heiliger II, 1286.; seine Tugenden offenbaren sich in allen Lagen II, 1673. 1677.; ist Bild eines frommen Fürsten II, 1783. f., uns ein Exempel II, 1716. J. war beständig II, 1248. im Glauben und Hoffen II, 1287. 1312., im Glauben und in der Lehre II, 1225. ff. 1305. 1433. 1674. f.; demüthig II, 1026. ff. 1433. 1672. f. 1676. ff. 1715. ff. und ehrerbietig gegen Pharao II, 1339. ff. 1343. 1386. ff. 1403. f. 1413. 1728. 1782. 1784.; ehrbar, wachte über den guten Ruf des Hauses Jakobs II, 1028. ff. 1072.; fleißig II, 1228. 1239. ff. 1437. 1449.; geduldig II, 1027. 1224. f. 1229. f. 1233. 1291. und übte den Psalter II, 1226.; gehorsam II, 1076. 1413.; nicht geizig II, 1243.; gottselig II, 1672. ff.; mäßig in Glück und Unglück II, 1729.; keusch II, 1224. f. 1255. 1266.; liebte Vater und Brüder, war nicht gefühllos II, 1672. f. 1676. f. 1682. 1715. ff. 1743. f.; männlich und muthig II, 1405. f.; nicht rachsüchtig II, 1682. f.; sanftmüthig II, 1256.; sittsam II, 1027. 1715.; treu II, 1228. 1230. ff.; weder verzagt noch vermessen II, 1674.

9. \gesperrt{Werke} J.’s. Hat nichts ohne Gebet gethan II, 1339. f.; hat größere Werke gethan denn Christus II, 1331. — S. unten 12. 13. 14. 15.

10. \gesperrt{Sünden und Buße} J.’s. J. hatte auch noch ungeduldiges Fleisch an sich II, 1317. f. 1323. f., mit dem er kämpfen mußte II, 1152. 1503. Ob er gesündigt dadurch, daß er durch den Schenken frei zu werden suchte II, 1300. ff., eine Heidin heirathete II, 1416. f., seines Vaters Haus vergaß II, 1427. ff., seine Brüder für Kundschafter erklärte II, 1472. f., bei dem Leben Pharaos schwor II, 1483. Wie er über das Murren seines Fleisches Buße gethan II, 1332.

11. \gesperrt{Anfechtungen und Leiden}. J.’s Anfechtungen und Leiden sind um so schwerer, da er allein stand II, 1226.; J. hatte viele A. II, 1228. 1252., schwere, von innen und außen II, 1318. ff.; durch seine A. ward er auf seinen hohen Beruf vorbereitet II, 1659. Er ward angefochten durch die egyptische falsche Lehre II, 1225. f., von seinem Fleisch, das ihn zur Unzucht reizte II, 1227. f.; seine Anfechtung und Versuchung zur Unzucht war schwer II, 1252. ff. 1259. 1264. 1268.; angefochten von Ungeduld II, 1317. f., da ihn der Schenke vergaß II, 1314., von Zweifel an der Wahrheit seiner Träume II, 1428. f. J. mußte \gesperrt{Unglück} leiden wie wir II, 1077. f.; ist in der Trübsal ein Erbe seines Vaters II, 1226.; sein Leben war erbärmlich nach dem Fleisch II, 1096.; seine Marter war größer denn der Märtyrer II, 1323. mußte nicht um Sünde willen leiden II, 1249. f.; mußte die Undankbarkeit der Welt erfahren II, 1314. f. J. war nicht stumpf und gefühllos gegen sein Leiden II, 1225. f.; Gott tröstete ihn II, 1157. 1286. ff.; Christus war mit ihm im Gefängniß II, 1295. f.; J. hielt sich ans Wort II, 1095. f. 1157. 1283. 1286. 1292. f.; wie er sich getröstet II, 1317. 1320. f. 1325.; sah in der Versuchung zur Unzucht auf Gottes Gebot II, 1266.; überwand mit unverzagtem Muth seine Marter II, 1379. f. 1384.; sein Kreuz hatte einen herrlichen Ausgang II, 1085. f. — S. unten 12. 13. 14.

12. \gesperrt{Joseph und seine Brüder}. II, 1020. ff. 1045. ff. 1069. ff. 1087. ff. 1115. ff. 1137. ff. 1450. ff. 1470. ff. 1485. ff. 1505. ff. 1589. ff. 1603. ff. 1613. ff. 1624. ff. 1637. ff. 1654. ff. 1670. ff. 1688. ff. 1711. ff. 1738. ff. 1784. ff. 2076. ff. 2082. ff. Die Sünde der Brüder J.’s soll man groß machen uns zum Trost II, 1101. Ihr \gesperrt{Haß}: sie haßten ihn, weil ihn Jakob lieber hatte als die Anderen II, 1032. ff., weil er vor ihren Vater brachte, was sie Uebeles thaten II, 1030. f., um der Erstgeburt willen II, 1110., welche Joseph gebührte II, 1025. 1057. 1111.; ihr Haß war unbillig und grundlos II, 1040. 1088., größer als Kains II, 1088., sie waren Bruder-, Vater-, Mutter-, Großvatermörder II,
\quellseite{558}
1088. 1101.; ihr Haß gemehrt durch Josephs Träume II, 1045. f. 1059.; waren nicht alle gleich im Haß entbrannt II, 1041. 1058.; zogen vielleicht aus Haß gegen ihn gen Sichem auf die Weide II, 1070. Sie \gesperrt{verachteten} ihn II, 1032., nannten ihn darum „Träumer“ II, 1108. f.; schnauzten ihn an II, 1041. Warum Joseph ihnen seine \gesperrt{Träume} erzählt II, 1045. f.; sie verachteten die Träume II, 1064. und um derselben willen auch Jakob und J. II, 1065. Ihre \gesperrt{blutigen Anschläge} wider J.: Alter J.’s und seiner Brüder zu der Zeit II, 1100. f.; Simeon und Levi Rädelsführer II, 1100. 1117. f.; Juda und Ruben willigten nicht gern ein II, 1100.; Ruben widersetzt sich vergeblich II, 1101. 1115. 1138.; die anderen Brüder waren noch zu jung II, 1100.; waren bei ihren Anschlägen vom Teufel besessen II, 1100.; gingen listig zu Werke II, 1099.; sind in ihrer List ein Bild der Mörder Christi II, 1100.; meinten, ihre Anschläge würden verborgen bleiben II, 1113.; ob sie den bunten Rock zerschnitten II, 1121.; ihre Sünde größer als Kains II, 1111. f.; mit der Einen Sünde haben sie viele andere begangen II, 1102.; wie Gott ihre blutigen Anschläge gestraft II, 1102. ff. Wie sich J. verhalten, da er in die \gesperrt{Grube} geworfen ward II, 1118. f.; seine Brüder waren unerbittlich II, 1119., ihre Grausamkeit größer als die der Juden gegen Christum II, 1119.; wollten ihn in der Grube Hungers sterben lassen II, 1120.; in die Grube geworfen zu werden, war für J. schlimmer als der Tod II, 1119.; da war er ohne Tröster II, 1121.; wie J.’s Brüder für diese That gestraft worden II, 1116. Warum Juda räth, J. zu \gesperrt{verkaufen} II, 1126. f.; ward seinen Vettern verkauft II, 1123. f.; die Midianiter und Ismaeliter waren Eine Gesellschaft II, 1130. ff.; warum er nicht den Nachbarn verkauft ward II, 1132. f.; was für eine Ursache des Verkaufs angegeben ward II, 1131. ff.; bei diesem Verkauf verbanden sich Juden und Heiden II, 1133. f.; wie oft Joseph verkauft worden II, 1130. f.; wie J. beim Verkauf zu Muthe war, übergeht Mose mit Stillschweigen II, 1129. f.; er hat wohl kläglich geweint und gebeten II, 1134.; der Verkauf war für Joseph schlimmer als der leibliche Tod II, 1126. f., daher die Sünde seiner Brüder schwerer, als wenn sie ihn getödtet hätten II, 1128. f.; ihre That war schrecklich II, 1498.; wie schwer es ihm war, daß er ein Knecht werden sollte II, 1157. f., da er von seinem Vater weggerissen wurde II, 1225., an Hebron vorbeigeführt wurde II, 1156.; die Absicht, die That vor Jakob zu verhehlen, ist gut II, 1140.; Ruben verhehlt die Sünde mit II, 1139. ff., hätte sie offenbaren sollen II, 1140. ff.; der Verkauf gab groß Aergerniß II, 1089.; Gottes Regierung ist wunderbar, da er J. verkaufen läßt II, 1155. f. Die \gesperrt{Brüder} waren nicht in gleichem Grade schuldig II, 1154.; haben sich selbst am meisten geschadet II, 1154. f.; in den 22 Jahren bis zum Zuge nach Egypten schlief ihr Gewissen II, 1141. f. 1144. f.; waren verzweifelte Buben und unbußfertig II, 1528. ff.; erkannten die Sünde erst nicht, thaten endlich Buße II, 1108.; waren tugendstolz II, 1631. f. 1634.; konnten in diesen 22 Jahren nichts Gutes thun, nicht beten II, 1141. f.; ihre Sünde in Egypten gestraft II, 1155.; sie hatten ein bös Gewissen II, 1520. ff. 1590. ff.; kamen in Egypten zur Erkenntniß ihrer Sünde II, 1459. ff. J. ist voll Barmherzigkeit und Liebe gegen sie II, 1458. f.; übt an ihnen den vollkommenen Nutzen des Gesetzes II, 1494.; sein Spiel mit ihnen war Gott angenehm II, 1629.; stellte sich, als sei er ein Weissager II, 1621. f.; J.’s und Simeons und Levis Gefühle, da er sie küßte II, 1718. ff.

13. J. \gesperrt{in Potiphars Hause}. II, 1220. ff. 1231. ff. 1250. ff. 1267. ff. J. hat daselbst viel gute Werke gethan II, 1229.; sich selbst Gott geopfert II, 1230.; das Gesinde gelehrt und gestraft II, 1229., in der Lehre von der Gottseligkeit unterrichtet II, 1596. f.; Viele bekehrt II, 1272.; Potiphars Weib eine Predigt gehalten II, 1257. ff.; viel erlitten und ertragen II, 1233.; verlor Ehre und guten Namen, mußte ein Heuchler und Bösewicht sein II, 1270. ff.; es war für ihn ein schweres Kreuz, daß er des Ehebruchs beschuldigt ward II, 1279. f.

14. J. \gesperrt{im Gefängniß}. II, 1282. ff. 1292. ff. 1306. ff. J.’s Werke im Gefängniß II, 1287. f.; vertheidigte seinen guten Namen II, 1304. f.; suchte auf rechte Weise seine Freiheit II, 1300. f.; betete II, 1299. f.; lehrte und tröstete II, 1286. f. 1297. f. 1302.

15. J. \gesperrt{an Pharaos Hofe}. II, 1322. ff. 1338. ff. 1357. ff. 1373. ff. 1398. ff. 1421. ff. 1434. ff. 1725. ff. 1805. ff. 1818. ff. J. sollte Großes in Egypten ausrichten II, 1156.; war der rechte Christus, ein Heiland in Egypten II, 1331. 1398.; war daselbst ein leiblicher und geistlicher Segen II, 1330. f.; Egypten stand unter ihm leiblich und geistlich in großem Flor II, 1730.; betete II, 1436.; übte sich im 4., 6. und 7. Gebot II, 1676.; verachtete das Irdische II, 1676.; führte das Regiment um Gottes willen II, 1405. f.; war der Oberste im Krieg und Frieden II, 1374. f.; hatte bei den Egyptern Ansehen und Gunst II, 1776. 1820. f.; hat den Egyptern gepredigt II, 1356. 1365. 1372. 1436. f., ihnen die rechte Lehre vorgetragen II, 1351. 1729. 1822. ff.; beschwerte sie nicht mit der Beschneidung und Satzungen der Väter I, 1014. 1023.; II, 1351. 1427. 1433. 1822. 1831.; schaffte nicht alle egyptischen Gebräuche ab II, 1616.; konnte nicht alles Verkehrte abthun II, 1618. 1637. f.; ging ihnen mit gutem Beispiel voran II, 1783.; bekehrte Viele II, 1331., nicht Alle II, 1436. f. 1773. 1823.; hatte auch seine Mißgönner II, 1726. f.; da war eine schöne Kirche II, 1831. 1833.; viele Egypter selig II, 1828.; die Egypter fielen später wieder ab II, 1831. f.

16. J.’s \gesperrt{Ehestand}. II, 1413. ff. 1426. ff.

17. J.’s \gesperrt{Tod u. Begräbniß}. II, 2072. ff. Warum er in Canaan begraben sein wollte II, 2089. ff.

Vgl. \gesperrt{Jakob} I, 19. 20.; II.

Josephus. Lügt oft zu Ehren seines Volks, sonst ein trefflicher Mann II, 889.; schreibt, man habe Stücke der Arche gefunden I, 562.

Josia, König. Vertrieb die Camarim II, 1610.

Jubal. Bedeutung des Namens I, 387.

Juda. I. 1. \gesperrt{Jakobs Sohn}. II, 1158. ff. 1165. f. 1959. ff. Wann geboren II, 548. ff.; ob er mit Bewilligung seines Vaters ehelich geworden II, 1161. f.; heirathete, 12 Jahre alt II, 1165. ff.; blieb nach seines Weibes Tode Wittwer II, 1185.; hatte die pharisäische Meinung, ein Mörder sei nur, wer selbst todtschlägt II, 1127. f.; war fromm, ehrlich, ein Exempel guter Sitten II, 1575.; gehorsam II, 1162.; hielt im Hause Zucht II, 1575.; verwaltete wohl das Priesteramt unter den Cananitern II, 1202.; hatte Ansehen und Gewalt bei den Cananitern II, 1203. 1205., besonderes Ansehen bei seinem Vater II, 1564. f. 1575. f.; König und Priester in Jakobs Hause II, 1771.; sündigte II, 1183. 1185. 1187. 1191. ff. 1210. ff., war auch leichtfertig II, 1194.; bekannte seine Sünde, that Buße II, 1210. ff.; Jakobs letzter Segen über ihn hat sich theilweise in David erfüllt II, 1961. ff. — S. \gesperrt{Jakob} I, 19.; II; \gesperrt{Joseph} 12.

2. \gesperrt{Juda u. Ger}. J. gibt Ger den Namen II, 1163.; gibt ihm ein Weib II, 1162.

3. \gesperrt{Juda u. Thamar}. II, 1177. 1185. ff. 1207. ff. Wann die Blutschande begangen worden II, 1166.; warum sie erzählt wird II, 1167. ff.; J. hat sich später nicht mehr zu T. gehalten II, 1564. f.

II. \gesperrt{Der Stamm}. Hat seinen Ursprung von einer Heidin II, 1172. ff.

III. \gesperrt{Das Reich}. Hatte im babylonischen Gefängniß die Hoffnung der Rückkehr II, 2051.

Judas Ischarioth. Soll vom Stamm Simeon gewesen sein II, 1041. 1044. 1117. 1717. 2055.; was seiner Buße fehlte II, 1467.

Juden. 1. \gesperrt{Das Volk der Juden} vor Christi Geburt. War durch die Beschneidung von allen anderen Völkern abgesondert I, 1027.; ein großes, wunderbares Volk I, 741. ff. 744. f.; II, 1377.; von allen anderen Reichen verschieden II, 2047.; trotz aller Plagen ein gesegnetes Volk I, 745.; berühmt I, 745. f.; mäßig I, 1155.; stolz darauf, daß es Gottes Volk sei II, 876.; wird im A. T. mit einem Opfer verglichen I, 958.; Teufel und Welt tobten dawider II, 2036.; Gott schützte es II, 1965. 2013. \gesperrt{Geschichtliches}: war 430 Jahre in der Fremde I, 769. 959. f.; das Ertränken der Knäblein in Egypten war eine Strafe für die blutigen Anschläge wider Joseph II, 1102. ff.; Raub an den Egyptern keine Sünde II, 631. f.; war unter den Richtern nicht Ein Volk II, 2027. f.; hatten Schulen in den Städten I, 1117.; den Edomitern nie unterworfen II, 823. f. \gesperrt{Religion}: ihr Beruf war, sich beschneiden zu lassen I, 1075.; wie sie in Egypten und der Wüste getröstet wurden I, 955. ff.; im jüdischen Volk war die rechte Kirche II, 2036.; wo sie opfern sollten, waren ungehorsam I, 779. f.; geriethen in Irrthum I, 1046.; wollten oft ohne Glauben, durch selbsterwählte Gottesdienste Gott versöhnen II, 2014.; sehr abgöttisch I, 1592. f.; II, 650. f. 776. ff. 960. 1616. f. 1825. 1914. ff. 1918. f. 1954.; schändliches Volk durch seine Sünden, aufrührerisch gegen die Propheten II, 995. f.; glaubten den Propheten und Christo nicht I, 495. f.; rechte Nachfolger der Brüder Josephs II, 1100. 1111. \gesperrt{Sitten u. Gebräuche}: das weibliche Geschlecht war dienstbar und verachtet II, 1183. f.; durften nicht Kleider aus mancherlei Garn tragen II, 1040.; nahmen den Hut nicht im Tempel ab II, 821.; Schüler saßen zu der Lehrer Füßen II, 1970.; Steuern und Abgaben II, 1362. — S. \gesperrt{Baalsdienst, Juda} III, \gesperrt{Kälberdienst, Moloch}.

2. \gesperrt{Die} J. \gesperrt{nach Christi Geburt}. Sind mit Unrecht \gesperrt{stolz} auf die Beschneidung I, 1008. 1100.; sind hoffährtig und ruhmredig II, 547. wegen ihrer fleischlichen Abstammung II, 34. 37. ff. 995. f. 1940., als ob sie Gottes Volk seien II, 618. Sind nicht die rechte Kirche II, 39. f. Sind jetzt \gesperrt{kein Volk} I, 750. 1053. 1099. f., Versuche und Hoffnungen, ihr Reich wieder herzustellen, vergeblich I, 1100.; sind nicht der rechte Same Abrahams I, 750. ff. 1396. f. 1669., trachten, Abraham ungleich, nach irdischen Gütern I, 917. f.; können ihren Wucher nicht mit Abrahams Beispiel beschönigen I, 829. f.; sind nicht mehr Gottes Volk, sondern verstoßen I, 754. 871. ff. 876. f. 961. ff. 1053. f. 1056. 1099. f. 1102.; haben Leib und Seele durch ihr Pochen auf die fleischliche Geburt verloren I, 1032.; sind von Gott in verkehrten Sinn gegeben II, 995. ff., vom Teufel und seiner Bosheit besessen I, 1152.; haben \gesperrt{keine Erkenntniß} in geistlichen Dingen I, 1553.; verwerfen die Lehre von der Dreieinigkeit I, 1149. f. 1278. ff.; II, 213. ff., haben den wahren Gott nicht I, 1551.; unterscheiden nicht zwischen Gesetz und Evangelium I, 1030.; wissen vom rechten Verstand der heil. Schrift nichts I, 998. f.; verstehen die Weissagungen von Christo und seinem Reich nicht I, 1097.; II, 1347.; hängen an dem Bunde der Beschneidung und verwerfen den Bund der Verheißung I, 1116.; verstehen nichts von der Beschneidung I, 1039. f.; bitten noch um den Messias II, 319.; haben selbsterwählte Gottesdienste II, 777.; an ihnen ist alle Ermahnung vergeblich I, 1033., wollen keine wahre Buße thun II, 320.; wie Verachtung des Worts an ihnen gestraft wird II, 319. f.; die es mit ihnen halten, haben keinen Segen davon I, 751. \gesperrt{Zu Luthers Zeit}: Lage der Juden I, 750. ff. 961. f. 1109.; waren Verräther Deutschlands I, 750.; hatten bei den Türken Aufenthalt I, 750.; trieben Wucher I, 751., davon sie selbst keinen Segen hatten I, 752.; \gesperrt{Gebräuche}: aßen nicht von der Spannader II, 801.; gaben ihren Kindern Namen von Blumen etc. I, 269. — S. \gesperrt{Rabbinen}.

3. \gesperrt{Juden und Heiden oder Unbeschnittene}. Die Juden, nicht die Heiden, hatten die Verheißung vom Samen und es war ihnen vertrauet, was Gott geredet hatte II, 984. f.; die J. waren nicht allein Gottes Volk I, 876. f.; J. und H. wollte Christus in Eine Kirche versammeln II, 1174.; J. und H. sind verwandt II, 909. 1172. ff. 1418. 1426.; die J. verachten die H. I, 1434. 1564.; II, 833. f. 888. f. 975. 995.; die J. hassen das Römische Reich und die Christenheit und beten dawider II, 984.; sie hassen die H. und dichten ihnen Laster an II, 1007. 1174. 1414.; meinen, sie dürfen die H. erwürgen II, 889.; sie möchten über die H. herrschen und sie umbringen II, 1400.; der J. Fall und Zerstreuung diente zur Bekehrung der H. II, 1758.

Judith. Heißt Lobsängerin II, 989. — S. \gesperrt{Esau} 3.

Jugend. Soll fröhlich sein II, 1800.; ist leicht verführt II, 1247.; man soll sie über die häuslichen, bürgerlichen Werke unterrichten I, 1179., zu Höflichkeit und Ehrerbietung anhalten I, 1635., an Hams Exempel fleißig erinnern I, 639. — S. \gesperrt{Anfechtungen} 2. 6.

Julius, Pabst. Hieß früher Julianus II, 988.; scheußliches Ungeheuer II, 1660. — S. \gesperrt{Bund}.

Jünger. Nach Christi Auferstehung schwer zu trösten II, 1721. f. — S. \gesperrt{Apostel}.

Jungfrauen. Rebekka ein Vorbild für J. I, 1706. ff.

Jungfrauenstand. Jungfrauen werden in der Schrift nicht wie Eheleute gelobt II, 571.; wer den J. mehr liebt, soll doch den Ehestand nicht verachten I, 1164. f.

Jüngling. Ihm ist keusch bleiben schwer II, 1255. f.

Jupiter. Darstellung J.’s in der Bildersprache I, 1476. — S. \gesperrt{Dionysius} 1.

Juristen. Ihr Amt II, 889. f.; J. bedürfte man nicht, wenn Alles nach Recht geschähe II, 889.; eine ihrer Künste ist, daß sie mehrdeutige Wörter brauchen II, 584. 1927.; junge J. verwirren Alles in Hader I, 847.

Justinus, Kaiser. Ehedem Kuhhirt II, 1786.

Juvenal. Ausspruch über das böse Gewissen II, 712.; spottet über der Egypter Aberglauben II, 1829.

\section{K}

Kabbalisten. Lehren: In lege non esse prius aut posterius II, 1024.

Kades. Lage I, 1006. 1306.

Kahath. I, 959.

Kain. 1. \gesperrt{Kain}. I, 295. ff. 307. ff. 317. ff. 320. ff. 330. ff. 334. ff. 347. ff. 361. ff. 379. ff. Bedeutung des Namens und Anlaß dazu I, 296. ff. 383.; II, 75.;
\quellseite{559}
hatte wohl eine Schwester zum Weibe I, 377. 381. ff.; war hoffährtig und darum verworfen I, 310. ff. Ward von Adam, nicht von Gott, wegen des Brudermordes zur Rede gestellt I, 334. 346. 350.; ihm blieb noch Raum zur Buße I, 360. 368.; verzweifelte I, 363.; von dem Zeichen, das er bekam I, 374. f.; wie er gestorben I, 370. 373. 389.; ist ein Vater aller, die die Heiligen tödten I, 343.

2. K. \gesperrt{und Abel}. I, 295. ff. 307. ff. 317. ff. 330. ff. 334. ff. 347. ff. 361. ff. Sind ein Bild der heuchlerischen und rechten Kirche I, 303.; K. opfert in Selbstgerechtigkeit, A. im Glauben I, 307. f. 315. f.; dem A., aber nicht dem K., ist sein Opfer ein Zeichen der Gerechtigkeit I, 1045.; K. wird verworfen, daß man nicht denke, es gebe zwei Kirchen I, 1050. f.

3. K.’s \gesperrt{und der Frommen Geschlecht}. I, 383. f. 408. ff.; II, 379. ff. 394. ff. K. hat eine neue Kirche mit besonderen Gottesdiensten aufgerichtet I, 356. f. 384. ff.; K.’s. Geschlecht war unstät und flüchtig I, 367.; K. und seine Kirche bedrängte die rechte Kirche I, 391. f. 408. 466.; Etliche aus K.’s Geschlecht hielten sich wohl zur rechten Kirche, wurden selig I, 356. 368. 372. 382. f. 388. 1031. 1050. f.; II, 12. 1828.

Kainszeichen. Worin es besteht I, 374. — S. \gesperrt{Kain} 1.

Kälberdienst. Warum man grade das Bild der Kälber gewählt II, 672. f.; die Propheten eiferten dagegen, die Juden vertheidigten ihn II, 425. f.

Kalne. Hieß später Seleucia I, 672. — S. \gesperrt{Seleucia}.

Kämpfer. Haben ihre Kunstgriffe II, 791. f.

Kayser, Leonh. Märtyrer I, 352.

Kebsweib. Bedeutung des Worts I, 1750. f.

Kedar. Bedeutung des Namens II, 13.; ist das steinigte Arabien I, 659.; II, 17.

Kedor Laomor. König in Elam I, 878.

Kelch. Bedeutet Noth und Gefahr I, 617.; Satans K. ist süß, aber man muß ihn wieder von sich geben I, 617.

Kelchentziehung. S. \gesperrt{Abendmahl}.

Kemuel. Vater der Syrer I, 1601.

Ketura. Bedeutung des Namens I, 1750.; ob K. und Hagar dieselbe Person I, 1418. f. 1744. f.; wer K. war I, 1750. f. — S. \gesperrt{Abraham} 27.

Ketzer. Erheben sich stets in der Kirche II, 99. aus Hochmuth II, 245.; verachten die rechte Kirche II, 245.; sind hoffährtig und sicher, daß ihr Ding Gott gefällt II, 82., reich an Worten II, 715.; aus frommen Christen werden oft K. I, 676., aus Abtrünnigen, Verfolger I, 191. f.; sind Nachfolger Satans, suchen die reine Lehre unter einem heil. Schein zu unterdrücken I, 181. 187.; finden Beifall I, 1283. — S. \gesperrt{Irrlehrer}.

Keuschheit. 1. \gesperrt{Im Allgemeinen}. Dreierlei K.: der Jungfrauen, Wittwen, Eheleute I, 1160. f.; besteht nicht in bloßer Ehelosigkeit II, 530.; was die Mönche für K. erklären I, 430. f.; Papisten verstehen von der rechten K. nichts II, 489. f.; im Papstthum findet man keine solche K., wie bei den Patriarchen I, 434.; ihre K., keine K. II, 528. ff., gehört ins gemeine Haus II, 540.

2. \gesperrt{Die jungfräuliche} K. Besondere Gabe I, 165.; II, 539. — S. \gesperrt{Ehestand} 9, 15.

Kinder. 1. \gesperrt{Im Allgemeinen}. Eine Gabe Gottes I, 1079.; können glauben I, 1041. f.; alle beschnittenen oder getauften K. selig I, 1049.; werden auf eine Weise erhalten, von der wir nichts wissen I, 1199. Erziehung der K. durch die Sünde beschwerlich I, 125.; gerathen gemeiniglich nach den Vätern II, 1936.; großer Leute Kinder gerathen selten wohl II, 1368. f. Dürfen auf ein Ehegemahl denken I, 1682. „Kinder“ heißen in der Schrift auch Kindeskinder II, 307. — S. \gesperrt{Beschneidung} 4, 7, \gesperrt{Bethlehem, Eltern, Moloch, Sodom, Taufe}.

2. K. \gesperrt{Gottes}. Wie wir K. G. werden I, 1033. 1433. f.; II, 35. 813.; Stolz und Vermessenheit folgt, wo die Lehre verlischt, daß wir aus Gnaden K. G. werden I, 1453. f. Wer die K. G. waren, von denen \gesperrt{Mose} redet I, 446. ff.; sie verachteten die Jungfrauen ihres Geschlechts I, 470. f.; worin ihre Sünde eigentlich bestand I, 471. f.; zeugten Tyrannen I, 473.

3. K. \gesperrt{der Welt}. Deuten die Schrift zu ihren Gunsten II, 826. f.

4. K. \gesperrt{der Welt} und K. \gesperrt{des Lichts}. Jene wissen ihre Gaben viel herrlicher zu machen als diese II, 826. f.; jene sind viel eifriger und rühriger in ihren Künsten, als diese in himmlischen Dingen II, 827.

Kindertaufe. S. \gesperrt{Taufe}.

Kindisch werden. Was das sei II, 1862.

Kindheit. S. \gesperrt{Alter}.

Kirche. 1. \gesperrt{Natur und Wesen} der K. Eine rechte Beschreibung der K. II, 436. f.; was und wo sie sei II, 430. ff. 437.; ist eine Pforte des Himmels II, 436., ein Reich der Gnade, außer ihr ist nur Drohung I, 1194.; ist das Himmelreich II, 429. f. Der \gesperrt{Titel}, daß man die K. sei, gilt vor Gott nicht I, 458. f. 496. 1405.; der \gesperrt{Haufe} macht nicht die K. I, 554. 1431. ff.; die K. sind, die das Wort hören und behalten I, 554. f. 1407. f. 1431. ff.; II, 42.; wir Christen sind die K. I, 963.; die K. ist nur, wo Christus regiert I, 1102., wo Gott redet I, 1446., bei den Getauften und Gläubigen I, 1044. 1102. 1402. 1404., auch mitten unter den Feinden Gottes II, 424. f. und Papisten II, 430. ff.; wer die wahre K. sei, wird am jüngsten Tage offenbar werden I, 459.; wo die K. ist, da ist Gott und der Dienst der Engel II, 1731. Ehre und Würde der K. II, 435. ff.

2. \gesperrt{Eine heilige, allgemeine} K. Es ist nur \gesperrt{Eine} K. und Ein Weg dazu I, 1053., Eine aus vielen Völkern I, 1095.; alle bekommen darin gleiche Gaben I, 1074.; der Teufel erweckt immer Secten I, 938. Die K. ist \gesperrt{heilig} um des Worts willen II, 391.; Unkraut ist auch immer unter den guten Samen gemengt II, 651.; in der K. kommen greuliche Sündenfälle vor II, 976.; da finden sich Böse und Heuchler II, 384. f. 652. 1128. 1617. f. 1938. f., aber das thut ihr keinen Abbruch an ihrem Wesen II, 1730.; sie ist \gesperrt{rein}, wenn die Lehre rein ist II, 651. f.; diese ist nicht immer gleich rein II, 709.; das ist ein großer Schade, wenn die Predigt und Aemter in der K. verderbt sind II, 1730. Die K. ist \gesperrt{allgemein}, schon im A. T. aus Juden und Heiden gesammelt I, 646. f. 1044.; II, 909. 1830.; nicht mehr an gewisse Personen gebunden I, 709., auch nicht an einen gewissen Ort I, 1065.; II, 429.; die K. besteht \gesperrt{zu allen Zeiten} vom Anfang bis zum Ende der Welt I, 708. ff. 1103. 1341.; sie geht nicht unter II, 342.; ist noch vorhanden, so lange Gott die Welt noch segnet I, 894.; der Heilige Geist findet allezeit Etliche, die er geistlich trunken macht II, 2002. — S. \gesperrt{Heiden, Himmelreich}.

3. \gesperrt{Feinde und Wohlthäter} der K. Die K. hat immer \gesperrt{Feinde} I, 747. 755.; der Teufel ist ihr Feind I, 393. 637. 695. 938.; II, 407. f. 979., der Unheil anrichtet II, 327., sie verunreinigt II, 438. f.; die Gottlosen verfolgen und plagen sie I, 936.; Tyrannen und irdisch Gesinnte bedrängen sie I, 391. f.; die Heiligsten verfolgen sie I, 310.; Abtrünnige werden die größesten Feinde I, 689. die K. bittet für die Feinde I, 494.; Gott sieht den Anschlägen der Feinde zu, lacht ihrer II, 1113. f.; straft sie I, 394. 644. 747. 1462. f. Die \gesperrt{Wohlthäter} der K. segnet Gott I, 747.

4. \gesperrt{Die falsche} K. Ist nur dem Namen nach eine K. II, 119., derselbe wird oft mißbraucht I, 1403. f.; sie maßt sich den Namen der K. an II, 99.; hat von Ham ihren Ursprung I, 685.; mißbraucht den Namen Gottes I, 1402. f.; will aus ihrem eigenen Kopf prophezeien I, 390.; rühmt sich ihres Glücks als eines Beweises, daß sie Gottes Volk sei II, 72. f.; will durch Werke Gott gefallen I, 317. \gesperrt{Kennzeichen} der falschen K.: Lüge und Mord II, 76. f.; Hoffahrt II, 118.; Verachtung der geistlichen, Liebe der irdischen Güter II, 118. ff. Esau ein \gesperrt{Bild} der falschen K. II, 118. — S. \gesperrt{Pabst} 1, \gesperrt{Pabstthum, Papisten, Türke} 2.

5. \gesperrt{Die wahre und die falsche} K. Es gibt zweierlei K., die wahre und die falsche I, 308. f. 1406.; II, 119.; so war es gleich im Anfang I, 380., wird es bis ans Ende sein II, 648. f.; die rechte K. ist Adams, die falsche Kains I, 356.; Bilder der falschen und rechten K. sind Kain und Abel I, 303., Ismael und Isaak, Esau und Jakob II, 67.; Kain ist der Anfang der falschen, Abel der wahren K. I, 308. f. Bei der rechten ist Gott, bei der falschen schweigt er II, 73. f. 77.; die rechte hat das Wort, die falsche nur leibliche Gaben II, 39. ff.; die falsche lästert und verfolgt das Wort, die rechte behält es und leidet Verfolgungen I, 1406.; die falsche mißbraucht Gottes Namen, die rechte hält ihn heilig I, 1406. f.; der falschen geht es wohl, der rechten übel I, 309. f.; II, 50. ff. 71. ff.; die rechte muß unstät und flüchtig sein, die falsche nicht I, 360. f.; der Zorn der falschen gegen die wahre ist teuflisch I, 318. f.; zwischen beiden ist steter Kampf I, 466. 632.; II, 34. ff. 118. um die Erstgeburt II, 1506.; die falsche gibt sich für die rechte aus, die rechte muß für falsch und ketzerisch gelten I, 309. 311. 388. 688. f. 933. f. 1504.; II, 308. 675. f.; die falsche will herrschen um der ersten Geburt willen II, 35.; die rechte wird von den Ketzern verachtet II, 245.; die falsche thut der rechten Gewalt und verdammt sie I, 310. f. 690. 984. f., kehrt sie um und verstört sie I, 380., entschuldigt ihre Tyrannei I, 392.; die wahre muß der falschen ihr Blut zu saufen geben I, 393. Gott verwirft die falsche I, 552.; Gott straft sie und rettet die wahre I, 380. 551. ff.; II, 13. — S. \gesperrt{Pabst} 1, \gesperrt{Pabstthum, Papisten, Türke} 3.

6. \gesperrt{Kreuz und Trost} der wahren K. Gott regiert und führt sie wunderbar II, 50. f.; sie trägt das Bild des Fluches I, 1002.; II, 304. f.; ist verborgen unter der schrecklichen Decke der Schmach und des Kreuzes II, 802. ff.; muß Kreuz tragen, hat weder Ehre noch Gut I, 478. f.; hat immer einerlei Gestalt: Rotten und Aergernisse I, 711.; wird verachtet I, 810.; hat niemals Ruhe I, 716. f.; von Anfang an verfolgt I, 394.; II, 456. f.; muß Verfolgung leiden I, 1408.; II, 1136.; Satan erregt viel Noth und Unglück I, 1135. f.; muß arm sein und vom Raube leben I, 1461. f.; II, 632. f. 647.; ist eine kleine verachtete Heerde I, 709.; II, 806.; das Allerschwerste ist, wenn man neue Lehre sucht und die Lehreinigkeit stört I, 1283. f.; II, 245. Gott drückt die K. mit dem Kreuz I, 1080., damit sie demüthig bleibe I, 823. Wie sie sich trösten soll im Leiden I, 309.; II, 52. 72. ff. 805., in Verfolgung I, 311. f. 350. 643., in Verachtung I, 497., bei der irdischen Macht der falschen K. I, 649. f., wider die Anschläge der Feinde I, 703. ff., über ihre geringe Zahl I, 899.; sie hat ewige Verheißungen der Gnade I, 372. f.; sie seufzt zu Gott um Rache wider die Gottlosen I, 507. f.; erkennt nach der Trübsal Gottes Liebe I, 1003.; II, 808. f.

7. \gesperrt{Stiftung u. Erhaltung} der K. Die K. ist vor dem Hausregiment und Polizei eingesetzt I, 126. ff.; hat die göttliche Majestät gegenwärtig II, 369.; hat im Wort gewisse Zeugnisse, daß Gott sich ihrer annehme I, 962.; Gott nimmt sich ihrer an, schützt und vertheidigt sie I, 231. 716. 728. f. 755. 776. 963.; II, 406. ff. 656. f. 1016. 1095. 1114.; die K. gewinnt endlich wider die Gottlosen I, 699.; Gott läßt sie nicht immer unter Tyrannen liegen I, 1357., erweckt fromme Fürsten, die sie pflegen und schützen II, 186. ff. 229. f.

8. \gesperrt{Kennzeichen} der wahren K. Es hat zu allen Zeiten Kennzeichen gegeben, wo die K. und Gottes Volk zu finden sei I, 1044. ff.; jetzt sind es andere als im A. T. I, 1056.; es sind die Gnadenmittel I, 709.; II, 112. 386. 806. Der Pabst gibt falsche Kennzeichen I, 1047. f.; irdisches Glück ist kein Kennzeichen I, 495. f. 1400. ff., vielmehr Kreuz I, 310.

9. K. u. \gesperrt{Wort} Gottes. Die K. kann nicht sein, wo man nicht Gottes Wort stets übt I, 669.; sie ist die Tochter, nicht die Mutter des Worts I, 554., Christi Schülerin I, 863. 867. f. 1317.; in ihr soll das W. regieren I, 766.; in ihr gilt nichts, denn was Gott thut II, 436. f.; man soll darin nichts annehmen oder anfangen ohne das W. II, 1760.; da soll nur Gottes W. gelehrt werden II, 2010.

10. \gesperrt{Verhalten} zur K. Die K. war zu allen Zeiten eine Zuflucht der Elenden und Armen I, 1135., so soll man ihr auch die Mittel geben, sie zu versorgen I, 1133. f.; vor ihr und ihrem Namen soll man Ehrfurcht haben und sich scheuen I, 1407. f.; es soll nicht Jeder, wie die Rotten thun, sich eine eigene K. machen I, 1095.

11. K. u. \gesperrt{Welt}. Die Welt hat nichts Köstlicheres denn die K., obwohl dieselbe sehr verachtet ist I, 1372.; die K. ist allezeit eine Mauer wider den Zorn Gottes über die W. I, 494., die längst zerstört wäre I, 893. f.; von der K. hat die W. leiblichen Segen I, 768.; II, 1235., um ihretwillen erhält Gott die W. II, 895. f.; doch muß die K. der Knechtschaft unterworfen sein II, 552. ff., die Gottlosen berauben sie I, 776.; die W. ist Lamechs Art, der K. feind I, 393.

12. K. u. \gesperrt{weltliche Reiche}. In der K. achte man darauf, wer da redet, in weltlichen Reichen darauf, was da geredet wird I, 1266. 1595.; in der K. gilt die Regel der Juristen nicht, daß wer einmal eine Uebelthat begangen, zu dem ist sich zu vermuthen, daß er allezeit übel thut I, 796.; die K. hebt die anderen Stände nicht auf I, 1465. 1477. f.; um der K. willen thut Gott den weltlichen R. Gutes I, 1409. und erhält sie II, 1607. 1705. ff.; das Gebet der K. erhält sie I, 1409.; wo die K. verfällt, verfällt auch Zucht und Ordnung I, 1192.; die K. hat schlechten Dank I, 1409. f.; II, 1706. f.; die weltlichen R. sollen Frieden schaffen, daß die K. lehren möge II, 2000. f. — S. \gesperrt{Obrigkeit} 12.

13. \gesperrt{Die lutherische u. die päbstische} K. Die luth. K. that recht, sich von der päbstischen zu trennen I, 1265.; sie ist die rechte K. I, 1405.; II, 42. 112.; ob vor ihrer Trennung Alle verloren waren I, 1281.; worin sie sich von der päbst. K. unterscheidet II, 42. f.; verwirft gute Werke nicht I, 1347.; wie sie mit den päbst. Kirchengütern verfahren II, 926. f.

14. \gesperrt{Geschichtliches}. 1 Mos. 12, 1. ff. ist eine kurze Geschichte der K. aller Zeiten I, 754. f.; Anfang und Mehrung der K. ist lieblich zu betrachten I, 729.; die K. bekam in der Zeit von Noah auf Abraham einen großen
\quellseite{560}
Stoß I, 728.; zur Zeit des Untergangs Sodoms war eine K. wie nachher nicht wieder bis auf Christum I, 1200.; die K. zu Rom war recht heilig, mit trefflichen Märtyrern geziert I, 448. — S. \gesperrt{Abraham} 19, \gesperrt{Enos}, \gesperrt{Ham} 2, \gesperrt{Jakob} 11, \gesperrt{Kain} 3, \gesperrt{Nimrod}, \gesperrt{Sodom}, \gesperrt{Väter} I; II, 11.

15. \gesperrt{Kirchen} u. \gesperrt{Schulen}. Wir sollen sie erhalten II, 453. ff. 455. f.; sie bauen und erhalten, ein herrlich Werk II, 449. ff.; was uns dazu treiben soll II, 451.; die Welt ist denselben feind II, 455.

16. \gesperrt{Kirchen, die Gebäude}. Zweck derselben I, 778.; welche Werke darin zu verrichten II, 222.; sind Herbergen für verwundete Gewissen II, 1723.; reine Lehre ihr rechter Schmuck I, 839.; II, 222.; wann sie Gotteshäuser sind II, 433. f. 939.; sind eher Kuhställe, wenn nicht darin gepredigt wird II, 917.; darin komme man ehrlich und züchtig zusammen II, 914.; waren im Pabstthum Theater des Teufels II, 917.; was für eine K. Isaak zu Bersaba gehabt II, 222.

Vgl. \gesperrt{Christ}, \gesperrt{Gastfreiheit}, \gesperrt{Himmelreich}, \gesperrt{Vereinigung}, \gesperrt{Volk} 2.

Kirchenväter. Sind wohl in Ehren zu halten, doch nicht den Aposteln und ihren Schriften gleichzustellen I, 148.; aus ihren Meinungen mache man keine Glaubensartikel II, 1547.; nicht ihnen, sondern dem Wort folge man II, 79., man nehme nur an, was mit der Schrift stimmt I, 1290. 1316. 1382.; weiche um ihretwillen nicht von der Schrift ab I, 149.; es ist angenehm und lieblich, wenn sie mit uns in der Lehre stimmen, wenn sie gleich oft nicht bequemlich und eigentlich reden I, 1153.; führen zu Zeiten schwache Beweise I, 1149. ff.; folgen gemeiniglich den Rabbinen I, 364.; sind nicht ohne Gefahr und führen von dem rechten Verstand der Schrift ab I, 953.; man lese sie mit Bedacht und lasse sich durch ihr Ansehen nicht betrügen I, 285.; II, 956.; daß sie geirrt, dient uns zum Trost I, 148. 1382.

Kirchhöfe. Erst später, bei den Christen I, 1625.

Kiriath Arba. Bedeutung des Namens I, 1609.; Veranlassung dazu I, 1610.; ist Hebron oder Mamre I, 1609.

Kittim. Völkername I, 661.

Kleidung. 1. \gesperrt{Im Allgemeinen}. Eine Erinnerung an den Fall I, 270. f.; ehrliche und ziemliche ist zu loben I, 272.; man darf sich schmücken, doch meide man Uebermaß I, 1715.; die Morgenländer hatten weißseidene Kleider, der Römer Tracht war Purpur II, 1001. f. 1037. 1039. 1307. 1987. f., die Türken trugen weißseidene II, 1307., ebenso die Egypter II, 1399.; die Juden trugen Pelze II, 1039.; man treibt mit der K. nicht bloß Uebermaß, sondern Unsinnigkeit I, 272. — S. \gesperrt{Stand der Unschuld}.

2. \gesperrt{Einzelne Gegenstände}. \gesperrt{Armbänder} trugen die Juden, auch die Deutschen zu Luthers Zeit I, 1713. f.; wer und was für \gesperrt{Barette} sie getragen I, 1713.; \gesperrt{Ohrringe} trugen die Mohren, Cleopatra I, 1712., auch die Juden I, 1713.; \gesperrt{Studentenkappen} II, 1188.

Kleinglaube. Ist uns eine Schande II, 379.

Kloster. Das K. der St. Agnes zu Rom rissen die Cardinäle an sich II, 105.; Klöster und Stifte sollte man zerstören II, 1896.

Klostergelübde. Gefallen Gott nicht II, 448. f.; sind Abgötterei II, 449.; verwerflich II, 447.; wie die Mönche sie mit der Schrift zu bestätigen suchen II, 445. ff.

Knaben und Mädchen. S. \gesperrt{Geschlechter}.

Knechte. S. \gesperrt{Gesinde}.

Kniebeugung. K. geschieht vor Gott und Fürsten II, 1046.; der Pabst fordert K. II, 1400. f.; eine schöne Sitte im Gottesdienst II, 1840. ff.

Köln. Die Theologen zu K. hetzten die Fürsten wider die Evangelischen II, 2001.

Kometen. Sind Schreckzeichen I, 607.

Könige und Fürsten. Ihr höchstes Amt ist, die Kirche schützen und erhalten I, 1356. ff.; II, 229. f.; hießen zu Abrahams Zeiten Heilande I, 879.; können Verachtung nicht leiden, gerathen darüber in Tyrannei I, 819. f.; die Namen aller frommen K. und F. kann man auf einen Ring stechen I, 789.; sind oft gut, aber ihre Räthe und Diener sind böse I, 1458. f.; um des Königs willen segnet oder straft Gott ein ganzes Volk I, 1324. f.; großer Schade, wenn fromme K. und F. abscheiden I, 1325.; fromme leicht verachtet II, 1294. f.; ihre Regalien soll man ihnen nicht nehmen II, 84. — S. \gesperrt{Höfe}.

Königreiche. Gott gibt sie und setzt ihre Zeit I, 741. f., theilt sie auch unter die Gottlosen aus I, 1409., regiert sie I, 741. f. — S. \gesperrt{Evangelium} 2, \gesperrt{Kirche} 12, \gesperrt{Regierung}, \gesperrt{Regimente}.

Königsthal. I, 897.

Korah. Kam selbst um, seine Nachkommen dichteten viele schöne Psalmen II, 1952.

Koran. S. \gesperrt{Alkoran}.

Kräfte. 1. \gesperrt{Die leiblichen}. Werden durch Schrecken und Bestürzung gemehrt II, 789. f.

2. \gesperrt{Die natürlichen}. Welches die natürlichen K. und Gaben im Menschen waren I, 202.; die Sophisten lehren, sie seien unversehrt I, 577. ff.; von dieser Lehre sind auch Spuren bei den Vätern I, 577.; daraus folgen andere irrige Lehren I, 580. f.; sind nicht unversehrt im Menschen I, 173. f. 201. ff. 204. 584. f., viel weniger im Teufel I, 174. 202.; worin das Verderben besteht I, 202. 210. f.; sie sind aufs Arge gerichtet I, 578. ff., der Mensch kann nichts Gutes I, 731.; der Mensch erkennt von Natur Gott nicht I, 578. 580. f., ist zu Faulheit und Laster geneigt I, 584.; wie sich das Verderben in der Jugend zeigt I, 585.; wie unversehrt sie sind, zeigt die Keuschheit der Mönche I, 831.; der Mensch kann durch sie sich des Teufels nicht erwehren I, 334. — S. \gesperrt{Mensch} 2, \gesperrt{Wille} II, 2.

Krankheiten. 1. \gesperrt{Im Allgemeinen}. Immer neue stellen sich um der Sünde willen ein I, 253. f.; der Schmerz ist am 3. Tage am heftigsten II, 885.

2. \gesperrt{Einzelne} K. Der englische Schweiß I, 253.; die französische Krankheit I, 253.; der Stein II, 885.

Kreuz. 1. \gesperrt{Urheber} des K.’s. Es kommt nicht von Ohngefähr, sondern aus Gottes Rath zu unserer Seligkeit II, 1086.; kommt selbst von Freunden, Hausgenossen und Kindern II, 1088. f.

2. \gesperrt{Wer} K. \gesperrt{tragen muß}. Ein Christianus muß auch ein Crucianus sein II, 467.; die Heiligsten sind demselben unterworfen II, 860. f., der Christen Leben ist verborgen unter der Decke des K.’s II, 802. ff. 1130.

3. K. ist ein \gesperrt{Zeichen der Gnade} II, 1083. ff. 1158. 1282. ff. 1466. 1490., der \gesperrt{Kindschaft} II, 1326., eine väterliche Züchtigung II, 1468., ein väterliches Spiel II, 1123. 1453. f. 1577.; Gott hat im K. Wohlgefallen an uns II, 1146.; das Leben unter dem K. ist das köstlichste II, 1147.

4. K. \gesperrt{ist nöthig}. II, 846. f. 920. f. 1492. f., um der Erbsünde willen II, 1463.

5. \gesperrt{Zweck und Nutzen} des K.’s. Gott schickt K. zu unserem Besten I, 1327., aus seinem Rath zu unserer Seligkeit II, 1086.; besser ist’s unter dem K. sein mit der Hoffnung ewiger Erlösung, als ohne K. in die ewige Verdammniß rennen II, 920. \gesperrt{Zweck}: K. soll zur Prüfung dienen II, 730. 734. f. 1158. 1282. ff., zur Besserung II, 1326. f. 1455. ff. 1459. ff. 1493. 1994. 2003. f., zur Erhaltung in der Demuth I, 929., uns in der rechten Mitte zu halten II, 761., zur Förderung in Geduld und Demuth I, 883.; soll zum Gebet treiben II, 903. \gesperrt{Nutzen}: K. ist sehr segensreich I, 822. ff.; II, 302. ff. 306. 1084. ff.; dient uns und Andern zum Heil II, 1079.; dient, daß wir Gottes Gnade und Barmherzigkeit, dazu das Geheimniß des Glaubens und der rechten Hoffnung verstehen lernen II, 1082.; treibt zum Gebet II, 1215.

6. \gesperrt{K. und Trost}. Sind bei einander I, 1080.; folgen immer auf einander I, 400. 921. 929. 1289; II, 191. f. 209. 234. 528. f. 732. f. 739. 846. f. 948. 1749., so bei Abraham I, 1340. ff. 1357. 1370.; das K. geht vorher I, 420. f. 938. f. 958. 993.; dieser Wechsel ist nöthig II, 846. 1668. und nützlich II, 761. Welches der rechte Trost ist II, 50. ff.; der rechte Trost ist Gottes Wille und Verheißung I, 1412. f.; Gott tröstet auch durch leibliche Güter I, 829. — S. \gesperrt{Kirche} 6.

7. \gesperrt{Verhalten im} K. K. thut dem Fleische wehe I, 958. f.; es flieht vor dem K. I, 996.; K. erbittert das Herz gegen Gott I, 988.; wir sind undankbar und fliehen vor dem K. II, 920. f. Wie Gläubige sich darin verhalten sollen II, 1151.; wir sollen nicht fliehen I, 979. f. 1374., nicht verzagen II, 755. f. 1225. f. 1587. f., sondern am Wort und Gebet halten II, 952. 1146. 1578. 1713. f., gehorchen und uns darein schicken II, 1136., geduldig sein I, 825. 828. f. 1130. f.; II, 1086. f. 1157. f. 1420. f., uns dazu reizen mit den Exempeln der Väter I, 693. f.; II, 855. 860. f. 900. f. 1078. ff.; womit wir uns im K. trösten und stärken sollen II, 567. f. 920. f., mit Gottes Wort II, 1082. f., Gottes Fürsorge II, 706. f., seiner wunderbaren Führung II, 354. ff., dem Wechsel von K. und Trost II, 739. f., daß Gott es nicht zu schwer werden läßt II, 739. f. 846., mit den Exempeln Anderer I, 967. 974.; II, 1078. ff. 1135. Gläubige rühmen sich ihres K.’s I, 784.

8. \gesperrt{Ende des} K.’s. Es ist ein fröhliches II, 1083.

9. \gesperrt{K. und Herrlichkeit}. K. geht der H. voraus II, 964.

10. \gesperrt{Kreuz Christi}. Wo es gestanden haben soll II, 428.

Vgl. \gesperrt{Anfechtungen}, \gesperrt{Ehestand} 13, \gesperrt{Landplagen}, \gesperrt{Leiden}, \gesperrt{Noth}, \gesperrt{Trübsal}, \gesperrt{Verachtete}, \gesperrt{Verfolgungen}, \gesperrt{Widerwärtigkeiten}, \gesperrt{Züchtigung}.

Kreuz der Grammatiker. Was das sei I, 719.

Kreuzigung des Fleisches. Worin sie besteht I, 1423. 1505.; II, 354.; ist nöthig II, 354., schwer I, 1523. — S. \gesperrt{Abraham} 17, \gesperrt{Beschneidung} 8, \gesperrt{Fleisch} 2, \gesperrt{Mönche} 2, \gesperrt{Selbstverleugnung}.

Kriege. Sind Strafen der Verachtung göttlichen Worts I, 882.; entstehen, weil die Könige mit dem Ihren nicht zufrieden sind I, 878.; den Feind darin auf Befehl der Obrigkeit tödten, ein gut Werk I, 763.; das Allertraurigste im K. ist, daß die Ehen getrennt werden I, 892. f.; in alten Zeiten tödtete man die Einwohner eroberter Städte II, 886.; auf K. muß man sich im Frieden rüsten II, 1375.; K. zu Abrahams Zeit I, 876. ff. 886. ff. — S. \gesperrt{Feldherren}, \gesperrt{Kriegsleute}, \gesperrt{Kriegslist}, \gesperrt{Schlacht}, \gesperrt{Sieg}.

Kriegsleute. Gott ist ein rechter Kriegsmann, unter ihm ist gut streiten II, 928. f.; Gott gibt und nimmt den K. den Muth II, 1378. f.; wenige recht tapfer II, 1375. f.; Ismael ein frommer Kriegsmann II, 13. — S. \gesperrt{Feldherren}.

Kriegslist. Ob sie unrecht sei II, 311.

Krone. Ursprung I, 1713.; mehrfache K. I, 1713.

Krystallhimmel. Meinung davon I, 34.

Kühe. Bedeuten Völker II, 1348.

Kunstausdrücke. Sind recht und nützlich I, 57. f., nicht Jedem verständlich I, 1359. — S. \gesperrt{Redeweisen}.

Künste, gute. Man studire sie I, 525., treibe sie II, 827.

Küssen. Die Hände küssen I, 1635.

\section{L}

Laban. 1. \gesperrt{Laban}. Bedeutung des Namens II, 481.; Nahors Enkel, Bethuels Sohn II, 475.; älterer Bruder Bethuels I, 1720.; seine Habe II, 593.; sein Gerücht bei den Nachbarn nicht gut II, 494.; lud wohl nur lose Fähnichen auf Leas Hochzeit II, 496.; war gastfrei gegen Elieser I, 1720. ff.; war zu Jakobs Zeit: ein hoffährtiger Heuchler und Werkheiliger II, 481. ff. 582. ff. 665. ff. 708. ff. 714. ff. 720. und sollte Nabal heißen II, 482. 581.; ein Götzendiener und Wahrsager II, 584. 643. 648. ff.; ein Schalk und im Geiz ersoffen II, 494. 580. ff., worin sein Geiz sich zeigte II, 497. 508. f. 593. ff. 599. 605. 621. f. 624. f. 679., sein Geiz machte ihn zu einem Unmenschen II, 627. ff. 682. f. 693. ff., sein Name sollte sprichwörtlich werden, Geiz anzuzeigen II, 612.; ein Tyrann II, 701.; war argwöhnisch II, 621.; falsch II, 583. 597., ging nur mit Ränken um II, 581. f.; hochmüthig II, 662.; undankbar II, 696. ff.; hatte ein böses Gewissen II, 712. 718. f.; nicht zu bessern II, 660. 662. ff.; ist ein Exempel, das jedermann verfluchen und fliehen soll II, 630., Bild eines Heuchlers II, 680. 720. 723., Götzendieners II, 720., der gottlosen Welt II, 713., des Pabsts II, 680. — S. \gesperrt{Jakob} 24.

2. \gesperrt{Labans Kinder}. Geizig und neidisch II, 610. ff.

3. \gesperrt{Laban und seine Töchter}. L. ein tyrannischer Vater II, 502. f.; hielt seine T. übel II, 626. ff.; sie führen gerechte Klage über ihn II, 625. 630. 639. f.; trösten sich in ihrem Elende II, 630. f.

Lachen. Verschiedene Arten II, 165.; Zeichen der Freude I, 1103.

Lais. Stadt, hieß später Dan II, 2022.

Lamech. 1. Bedeutung des \gesperrt{Namens}. I, 385.

2. \gesperrt{L., der Patriarch}. I, 428. ff. Nennt seinen Sohn Noah I, 428. ff.; starb fünf Jahre vor der Sintfluth I, 531.

4. \gesperrt{L., Kains Nachkomme}. I, 385. ff. Ein Mann von vortrefflichem Verstand, gewaltig und berühmt I, 391.; stolz und rühmt sich seiner Sünde I, 392.; Feind der Kirche I, 391.; ein Mann etwa wie Julius II. oder Clemens VII. I, 393.; der erste Polygamist I, 391.; warum er zwei Weiber nahm I, 386.; suchte bei Adams Tode die Herrschaft an sich zu ziehen I, 391.; ob er Kain getödtet I, 370. 389. 391.; seine Strafe soll die Sintfluth gewesen sein I, 371.; ist ein Exempel der Welt I, 393. Seine \gesperrt{Weiber} waren wohl gottesfürchtig, über seinen Todtschlag betrübt I, 390.

Länder. Sind glücklich, wenn darin Gottes Wort wohnt II, 1350.; in den besten L. wohnen die ärgsten Leute I, 854. — S. \gesperrt{Fluch}, \gesperrt{Volk} 1.

Landplagen. Gottlose meinen, daran sei Gottes Wort Schuld II, 132. ff.; sind Strafen für Verachtung des Worts II, 132. 134. f.; treffen die Gläubigen zur Prüfung und Uebung ihres Glaubens II, 131. f. 135.

Landstreicher. Die Obrigkeit soll gegen sie einschreiten II, 1602. f.

\quellseite{561}
Langmuth Gottes. Gott ist langmüthig I, 195.; II, 1677.; läßt aber darum die Strafe nicht immer nach I, 353. 507. f. 540.; Satan mißbraucht sie, die Menschen in Sicherheit zu stürzen I, 195.

Lasa. Lage I, 678.

Lästrygonen. I, 596.

Lazarus. Bedeutung des Namens I, 937.

Lea. 1. \gesperrt{Lea}. Sehr arm II, 541.; ehrlich, fromm, stark, aber nicht schön II, 491. f.; worin ihr blödes Gesicht bestand II, 491. f.; eine gottselige Hausmutter II, 540.; keusch und fromm, suchte nicht fleischliche Wohllust, sondern Kinder um des verheißenen Samens willen II, 539. ff. 544. 549. 556. 561. 563. f. 570.; erkannte ihre Fruchtbarkeit als eine Gabe Gottes II, 521. ff. 571.; fleißig II, 1239.; freundlich gegen die Wittwen der Sichemiten II, 935.; angefochten durch Verachtung, getröstet durch ihre Fruchtbarkeit II, 519. ff. 527. f. — S. \gesperrt{Jakob} 18, \gesperrt{Laban} 3.

2. \gesperrt{L. und Rahel}. II, 518. ff. 528. ff. 551. 565. ff. L. das Weib der Linken, R. der Rechten II, 969. f.; R. heißt vorzugsweise Jakobs Weib II, 1765.; L. verachtet, R. geehrt II, 519. ff. 524.; R. bevorzugt, bei ihr wohnte Jakob II, 525.; stritten, weil jede die Mutter des verheißenen Samens werden wollte II, 563. f.; bedeuten nach papistischer Allegorie das beschauliche und das wirkliche Leben II, 558. ff.

Leabim. Sollen die Libyer sein I, 676.

Leben. I. \gesperrt{L. Gottes}. Gott ist ein lebendiger Gott, daher auch wir ewig leben werden I, 1005.; II, 1759.; ist ein Gott des Lebens, nicht des Verderbens I, 1016., gibt ewiges Leben I, 1059., ein Gott der Lebendigen I, 1059.

II. \gesperrt{L. der Menschen}. 1. \gesperrt{Das zeitliche} L. Gott hat Lust dazu I, 601. f.; damit ist kein ander Gut zu vergleichen I, 269.; ist ein eitler Traum II, 1054. ff., ein steter Gang zum \gesperrt{Tode} I, 240.; nach der Sintfluth mehr Gefahren als vorher I, 253. f.; wird immer kürzer I, 803.; jedes Menschen L. ist voll Veränderungen I, 1604., voll Noth und Unglück I, 423.; Zweck des L.’s ist Glaube und Seligkeit II, 950.; wir sollen es zu erhalten suchen, aber auch den Tod nach Gottes Willen gern leiden II, 1301.; zu Erhaltung des L.’s brauche man die natürlichen Mittel II, 1306.

2. \gesperrt{Das bürgerliche} L. mit seinen Sitten, Ordnungen und Werken mißfällt Gott nicht I, 1464.

3. \gesperrt{Das geistliche} L. Ein gewisses Zeichen desselben ist, wenn man Gottes Zorn fühlt II, 1462.

4. \gesperrt{Göttliches} L. Ein göttliches L. führen heißt, Gottes Sache öffentlich führen I, 418. ff. 500. f. — S. \gesperrt{Wandel}.

5. \gesperrt{Heiliges} L. Welches L. ein heiliges sei II, 1583.; das heiligste L. ist unter dem Kreuz II, 1158. f.

6. \gesperrt{Gottseliges} L. Was es sei II, 1641.; die Welt versteht nicht, was es sei II, 480.

7. \gesperrt{Selbsterwähltes} L. Gefällt Gott nicht II, 1625.

8. \gesperrt{Speculatives oder beschauliches und wirkliches} L. Päbstische Meinungen und Allegorien davon II, 558.; wonach man beides unterscheiden muß II, 956. Was die Mönche \gesperrt{beschauliches} L. genannt II, 954.; ihr besch. L. ist Teufelsbetrug I, 1253. f.; II, 558. f.; dasselbe gefällt der Vernunft I, 1252.; führt zu Verzweiflung und Verdammniß I, 1252.; man hüte sich davor I, 1252. f.; II, 954. f.; ein Altvater warnt davor I, 1251.; worin das rechte b. L. bestehe I, 1252. f.; II, 560. 954. f.; ein Exempel desselben ist das cananäische Weib II, 955. f. Was die Mönche \gesperrt{wirkliches} L. genannt haben II, 954.; das rechte w. L. ist die Liebe II, 561. 955. f. — S. \gesperrt{Speculationen}.

9. \gesperrt{L. der Christen}. Ist ein steter Kampf des Geistes wider die Sünde I, 329.; besteht in Tödtung des Fleisches I, 1495.; steht auf Glauben und in der Hoffnung I, 642.; ist ein lustig Spiel Gottes mit den Christen II, 1577. 1586. f. 1629.; ist eine Pilgrimschaft II, 1801. f., ein Zug durchs Rothe Meer I, 42., ein steter Gang zum Tode, doch bleibt darin die Hoffnung des ewigen Lebens I, 240. 242., ein Schlaf II, 1578.; wie wir darin Gott dienen sollen I, 427. Das christliche L. ist keine schläfrige müßige Geschicklichkeit II, 1996. f.; ist, daß man geistlich trunken und täglich gewaschen wird II, 2005. f. — S. \gesperrt{Christ}, \gesperrt{Wandel}.

10. \gesperrt{L. der Gottlosen}. Ein Schlaf, ihr Glück ein Traum II, 1579. — S. \gesperrt{Gottlose}.

11. \gesperrt{Ewiges} L. Es gibt ein ewiges Leben I, 964. ff. 1064., weil unser Gott ein lebendiger Gott ist I, 1005. 1059.; die Schrift lehrt es I, 406.; wir haben die Verheißung davon reichlicher als Abraham I, 1066.; da wird uns wiedererstattet, was wir im Paradiese verloren I, 82.; da werden die sündlichen Gebrechen aufgehoben I, 136. 1088.; wie unsere Angesichter da sein werden I, 289.; da keine leiblichen Werke mehr I, 79. f.; ist den Menschen vorbehalten und bestellt I, 421., den Heiligen bereitet I, 421.; Abraham lebt I, 1060.; die Hoffnung des ew. L.’s sollte uns muthig machen I, 1066. f.; nach dem ew. L. sollen wir Verlangen haben I, 423. f.; unser Fleisch frägt nichts darnach I, 424.; die Vernunft kanns nicht glauben II, 1962. f. — S. \gesperrt{Himmel} 4, \gesperrt{Mensch} 1, \gesperrt{Seligkeit}, \gesperrt{Tod} 3, \gesperrt{Unsterblichkeit}.

Leben der Väter. Das Buch „L. d. V.“ bringt Exempel der Gefühllosigkeit der Mönche II, 1608. f.

Legenden der Heiligen im Pabstthum. Sind verderblich, lehren nichts vom Glauben I, 1315.; mit Bedacht zu lesen I, 1121.

Lehre. 1. \gesperrt{Im Allgemeinen}. Ist dreierlei: von Verheißung, Werken, Ceremonien II, 223. f.; auf die L., nicht auf die Kirchgebäude soll man sehen I, 839.; in der L. soll man sich auf das Wort, nicht auf Menschen verlassen I, 79.; ist nicht immer gleich rein in der Kirche I, 709. — S. \gesperrt{Gewißheit}, \gesperrt{Reformation}, \gesperrt{Religion}.

2. \gesperrt{Die reine} L. Wirkung derselben I, 1133.; II, 911.; wo sie ausgebreitet wird, da mengt auch der Teufel sein Unkraut darein II, 1417.; wie sie sich im A.\,T. ausgebreitet II, 1015. f.; Verhalten dagegen seitens des gemeinen Haufens I, 726. f., der Widersacher I, 637., der Papisten I, 1646.; sie wird verfolgt I, 840.; II, 779. f. 1249.; Prediger sollen darüber wachen II, 1424., Alle zur Erhaltung beitragen II, 648.; daran soll man halten I, 538. f. 1283. f.; man mache sich kein Gewissen aus dem Aergerniß, das man an der reinen L. nimmt I, 793.; man schweige nicht, wenn sie gelästert wird II, 681. 691.; beim Kampf darum halte man am Wort und kümmere sich nicht um die Schmähungen der Feinde I, 936. f.; man führe im Kampfe gute Beweise I, 1149. ff.; um Erhaltung derselben soll man bitten I, 534.

3. \gesperrt{Die falsche} L. L., die nicht mit Gottes Wort stimmt, ist vom Teufel I, 695. f.; daraus folgen allerlei Sünden I, 1551.; der Kampf dawider ist schwerer denn gegen äußerliche Gewalt II, 1983. f. — S. \gesperrt{Irrlehren}, \gesperrt{Irrthum}.

4. \gesperrt{Die reine und die falsche} L. Selig, wer die reine L. hat und nicht erst in falscher gesteckt hat II, 1991.; die falsche hängt Einem noch lange an II, 1990. f.; die Menschen geben reichlicher für die falsche, denn für die rechte L. II, 1954. ff.

Lehreinigkeit. Das ist das Betrübteste, wenn sie gestört wird II, 245.

Lehren. L. und vermahnen muß verschieden getrieben werden I, 1149.; Einfältige soll man lehren I, 18.

Lehrer in der Kirche. Sollen nicht zweideutige Worte brauchen I, 1407.; II, 1373., wider Ehrgeiz und Vermessenheit als ihre schädlichsten Feinde bitten I, 928.

Leib. Man soll ihn zähmen, nicht tödten II, 1879.; Pflege desselben geboten I, 859. 1710. f., darin die Mittelstraße zu halten I, 1716.; wie und wozu man sein warten soll I, 846.

Leibeigene. Gott hat L. ebenso lieb wie Freie I, 995. f.

Leibeigenschaft. Ist bürgerlicher Tod, oft noch schlimmer als der leibliche Tod II, 1126. f.; bestand unter den Patriarchen II, 1762., bei den Türken II, 1762.

Leichenbegängniß. Man soll es beibehalten II, 971. — S. \gesperrt{Begräbniß}.

Leichtgläubige. Werden leicht betrogen II, 273. f.

Leiden. 1. \gesperrt{Leiden}. Das schwerste L. ist Traurigkeit des Herzens II, 1684.; L. sind keine Genugthuung für die Sünde II, 1454. f. 1461.; alle L. der Heiligen gefallen Gott II, 772., sind vor ihm angeschrieben II, 1769.; dadurch will Gott geehrt sein II, 1284.; Leben im L. das heiligste II, 1158. f. 1635.; Christen leiden nicht um Sünde, sondern um Wohlthat willen II, 1270. 1273. f., um Gerechtigkeit willen II, 1249. 1269.; sie sind um Gottes willen ein Spott und Fluch der Leute II, 676.; darin soll sich der Glaube üben II, 1151.; thun weh, man sei geduldig II, 207. 1148. ff. 1157., denn Gott wird die Feinde strafen II, 680. f. 965.; dadurch sollen Christen in der Demuth erhalten werden I, 823.; wie man sich trösten soll II, 676.; das L. der Patriarchen ist uns faßlicher als das L. Christi II, 1078. — S. \gesperrt{Abraham} 17, \gesperrt{Isaak} 11, \gesperrt{Jakob} 15, \gesperrt{Joseph} 11, \gesperrt{Sarah}.

2. \gesperrt{L. und Freude}. Auf L. folgt F. I, 883. 922.; II, 1623. 1712. f.

3. \gesperrt{L. und Herrlichkeit}. Auf L. folgt H. II, 1330. f. 1431. ff.

Vgl. \gesperrt{Kreuz}. \gesperrt{L. Christi} s. \gesperrt{Christus} 2.

Lenden. Obwohl Sitz der Unzucht, doch nach der Schrift ein heiliges Ding I, 1656. ff.

Leo, Pabst. L.’s Meinung von der Unsterblichkeit I, 1243.

Levi. 1. \gesperrt{Jakobs Sohn}. Wann geboren II, 548. ff.; warum L. genannt II, 525. — S. \gesperrt{Jakob} I, 19.; II; \gesperrt{Joseph} 12, \gesperrt{Simeon} 2.

2. \gesperrt{Der Stamm}. Hatte nichts Eigenes II, 1110.; woher sie ihren Unterhalt hatten II, 451. f.; wurden von den Andern schlecht ernährt II, 1953. f.; aus L. sind große Propheten und Regenten gekommen II,; 1951. — S. \gesperrt{Simeon} 3.

Leviratsehe. Was sie war II, 1178. f.; schwer und verdrießlich I, 974. f.; II, 1179. ff.; Zweck II, 1181.; von Gott geboten schon zu Jakobs Zeit II, 1190.; damals aufs strengste gehalten II, 1189. f.; wohl nur von Frommen beobachtet II, 1180. f.; Strafe der Uebertreter II, 1180.; ob auch die Heiden dies Gesetz hatten, ist ungewiß II, 1189.

Leviten. Richteten falschen Gottesdienst an, um mehr Ehre und Güter zu erlangen II, 1954. — S. \gesperrt{Levi} 2.

Licht. Das am ersten Tage geschaffene L. war wohl beweglich, ähnlich wie die Sonne I, 23. f. 51.; ob es nach Erschaffung der Sonne und des Mondes verschwunden I, 48. L. bedeutet Gutes II, 76. Allegorie des Pabstes vom großen und kleinen Licht I, 612.; allegorische Deutung beider Lichter auf Christum und die Kirche I, 616. f.

Lichtenberger. L.’s Weissagungen sind vom Teufel II, 1050.; sagte den bayrischen Krieg vorher II, 1054.

Liebe. 1. \gesperrt{Im Allgemeinen}. Ein Fünktlein L. ist nach der Sophisten Lehre genug, das ewige Leben zu erlangen I, 1334.

2. \gesperrt{L. Gottes zu den Menschen}. Gott liebt seine Gläubigen I, 1342.; diese L. muß durch Erfahrung gelernt werden II, 809.; worin sich G.’s L. zeigt I, 588. f. 590. 592. 604. 1548. f. 1690.

3. \gesperrt{L. zu Gott}. Gott sollen wir mehr lieben als Weib, Kind oder eigenes Leben I, 1412. 1420. f. 1422. f.; Beweggrund dazu I, 1690.; ob der Mensch durch seine natürlichen Kräfte Gott über alle Dinge lieben könne I, 580. — S. \gesperrt{Glaube} 10.

4. \gesperrt{L. zu den Menschen}. Die natürliche L. ist unbeständig I, 798. L. ist geduldig und thut auch, was man um Geld nicht thun würde II, 491.; die rechte L. ist nicht argwöhnisch I, 799.; das höchste Werk der L. ist in der wahren Religion unterweisen I, 805.; lieben soll man alle Menschen, gute und böse I, 797.; L. ist mit Recht größer gegen ganze Brüder als gegen Fremde II, 1717.; L. zu den Brüdern heißt Philadelphia I, 1142.; die brüderliche L. ist viel größer als die allgemeine I, 1146., sie übertrifft die Freundschaft des Pylades und Orestes I, 1147.; beiderlei L. ist allein bei den Christen I, 1146. f.; L. offenbart sich in der Fürbitte für die Feinde I, 1198. ff., erweist auch den Geringeren Dienst und Ehrerbietung II, 824. — S. \gesperrt{Brautliebe}, \gesperrt{Eltern} 2, \gesperrt{Glaube} 10, \gesperrt{Neigungen}.

5. \gesperrt{Eheliche und unkeusche} L. Die reine eheliche L. thut und leidet viel II, 491., die unkeusche schlägt in Haß um II, 1268. f.

6. \gesperrt{L. und menschliche Gesetze}. L. ist Ziel und Meisterin aller menschlichen Gesetze I, 843. ff. 846. f.

Lieblosigkeit. Ist Sünde II, 944. — S. \gesperrt{Gefühllosigkeit}.

Lied. Angeführt: Mitten wir im Leben sind I, 1514.

Limbus patrum. Lehre der Papisten I, 1762.; II, 1546.; was vom l. p. zu halten I, 1762. f.; II, 2066. 2071.

Lindigkeit. Was L. ist I, 848.

Linsengericht Jakobs. Wie zubereitet II, 121.

List der Weiber. II, 1272. — S. \gesperrt{Kriegslist}.

Livius. Seine Auslegung des Traumes Hannibals II, 1049. f.

Lob Gottes. Der angenehmste Gottesdienst I, 1000. 1005. f.

Lot. 1. \gesperrt{Lot}. I, 1240. ff. 1254. ff. 1288. ff. \gesperrt{Erscheinung} I, 1206. ff.; hielt die drei Engel für heilige Männer I, 1236. \gesperrt{Glaube} L.’s: er erkannte die Lehre von der heiligen Dreieinigkeit I, 1214.; die Schrift rühmt seinen Glauben I, 739. \gesperrt{Tugenden} L.’s: war gastfrei I, 1213. f. 1219. f.; liebreich I, 1254. f. 1264. ff. \gesperrt{Werke} L.’s I, 858. ff.; errettet durch sein Gebet Zoar I, 1271.; seine Bitte ist ein rechtes Mustergebet I, 1267.; ob er daran wohlgethan, daß er seine Töchter den Sodomitern preisgeben wollte I, 1229. ff. 1232. \gesperrt{Sünden}: sein Zaudern bei dem Ausgang aus Sodom war Sünde I, 1260. f.; s. unten 4. Wie Lot über seinen Fall \gesperrt{Buße} that I, 1298. \gesperrt{Anfechtungen} L.’s: die Sodomiter plagten seine Seele durch ihre gottlosen Werke I, 693.; hatte die Seinen um sich, die ihn trösten konnten II, 1226. \gesperrt{Schwiegersöhne} L.’s I, 1243. f.; hatten
\quellseite{562}
sich wohl zu seiner Kirche gehalten I, 1243. Sein \gesperrt{Gesinde} kam wohl mit in Sodom um I, 1241. f.

2. \gesperrt{L.’s Weib}. I, 1273. ff. Ist uns ein Exempel I, 1281. War ein gläubiges, heiliges Weib I, 1280. f., sündigte aus Schwachheit I, 1282. f.; ihr Tod I, 1280. ff.; die Salzsäule ist wohl zugleich mit Zoar zerstört worden I, 1280. f.; ist wohl selig geworden I, 1282. f.; zu loben ist, daß sie Lot ins Elend folgte I, 773.

3. \gesperrt{L. und die Sodomiter}. I, 1222. ff.; um seinetwillen widerfährt den S. Gutes I, 893.

4. \gesperrt{L. und seine Töchter}. I, 772. 1240. ff. 1288. ff. 1298. ff. L.’s Töchter waren heilige Jungfrauen I, 1255.; von der Blutschande L.’s mit ihnen soll man gelinde urtheilen I, 1294.; L. handelte in der Bestürzung I, 1292. 1294.; II, 500. und auch die Töchter I, 1295. f.; ihr närrisches Vornehmen geschah nicht aus fleischlicher Wohllust I, 1292. 1295. f.; ihr Exempel soll uns zur Warnung dienen I, 1296. f. — S. \gesperrt{Abraham} 28.

Löwen. Die Sophisten zu L. sind lose leichtfertige Streicher II, 1931.; wissen keinen Unterschied zwischen Gottes Wort und Menschenwort II, 2010.; verstehen die Schrift nicht II, 2030.; kämpfen wider die erkannte Wahrheit II, 1966. f.; sagen, die Schrift lehre nicht, daß der Heil. Geist durch den Glauben gegeben werde II, 1997., noch, daß man durch den Glauben Vergebung und Seligkeit erlange II, 2059., leugnen damit, daß Christus ein Prophet sei II, 2062.; lehren falsch vom Gebet II, 2014.; die groben Esel vertheidigen das Fegfeuer II, 1900.; hetzen die Fürsten gegen die Evangelischen II, 2001. f.

Lucia, St. Ihre Antwort, da man sie zu Götzendienst und Unzucht zwingen wollte II, 505.; hielt den Tod für einen Scherz II, 1989.

Lucretia. Beging keinen Ehebruch II, 505. — S. \gesperrt{Tarquinius}.

Lud. Völkername I, 680.

Lüder. Kommt von Lothar oder Luther II, 1415.

Ludim. Sollen die Lydier sein I, 676.

Ludwig von Frankreich. S. \gesperrt{Bund}.

Lüge. Eintheilung in Scherz-, Noth- und Schadenlüge I, 788. f.; II, 171. ff.; Beispiel einer schädlichen Lüge II, 875.; gibt nur eine Art Lüge, Scherz- und Nothlügen sind keine L. I, 788. 1317. f.; eine L. gebiert die andere I, 335.

Lügner. Müssen klug und listig sein und ein gut Gedächtniß haben I, 334.; II, 272.; verrathen sich selbst I, 334. f.; II, 1121. f.

Lus. Lage II, 736. 936.; hieß später Bethel II, 935. f.; ob L. und Jerusalem dieselbe Stadt sei II, 936. — S. \gesperrt{Bethel}.

Lüste. Was L. \gesperrt{des Fleisches} sind I, 224.; fleischliche L. sind des Teufels Odem II, 1265.; diese Brunst des Fleisches ist nicht zu heilen II, 1257. f. — S. \gesperrt{Unzucht}.

Luther. I. \gesperrt{Hans} L. Verwarf das Mönchsleben II, 1890.

II. \gesperrt{Martin} L. 1. \gesperrt{L. selbst}. Klagt über sich selbst I, 1338.; nennt sich den Allergeringsten I, 1., einen von den Eselsfüßen und Fleischlichen I, 1516. 1523. \gesperrt{Aussprüche} Luthers über sich selbst in Beziehung auf: Abrahams und Melchisedeks Geschichte, wie er sie ehedem gelesen I, 909.; Absolution, wenn ihm versagt I, 1426.; Allegorien I, 284. f. 610.; II, 474. 557. f. 561.; seine Anfechtungen von Unglauben II, 146., bösem Gewissen II, 1521., über die Versehung II, 181., von Werkgerechtigkeit I, 1433.; II, 463., wegen der Verachtung des Evangeliums II, 804.; seine Anfeindungen I, 936. f. 1186. 1465.; II, 245. 1273. f. 1309.; Angefochtene, wie er sie getröstet II, 179.; seine Anschläge I, 1329.; II, 1482.; Antinomer, beriefen sich auf L.’s Beispiel I, 1428. f.; seine Arbeiten und Geschäfte I, 920. 967. 1301.; seinen Beruf zum Lehramt I, 1102.; II, 267. f. 278. f.; seine Bibelübersetzung II, 937. 1836., rechtfertigt seine Uebersetzung von 1 Mos. 47, 31. II, 1839. 1844.; seine Blindheit im Pabstthum II, 1889.; Bosheit der Menschen gegen die erkannte Wahrheit wundert ihn II, 1966, 1968.; Bratwurst, erschrak davor II, 1598.; seine Bücher I, 1289.; seinen Bund mit Gott II, 1047. 1053.; Concilium, zu Schmalkalden abgeschlagen I, 1329., seine Erkenntniß, kommt zum Verständniß von Röm. 1, 17. II, 320. f.; Ehestand, hielt ihn als Knabe für sündlich I, 164. f., warum er so fleißig davon handelt I, 1679., hielt ihn als Mönch für verdammten Stand II, 166.; Ehelosigkeit aufzuheben II, 364.; Ehrgeiz focht ihn nicht besonders an I, 924.; Einschlafen und Erwachen wollte er beobachten II, 216. 2071.; Eliesers Geschenke, der Rebekka gegeben, wundert sich, daß die Schrift so viel Worte davon macht I, 1711. f.; Erscheinungen, begehrt sie nicht, wie er sich dabei verhalten würde I, 1118. f.; 1527. f.; II, 782. f. 1047. 1309.; seine Fehler und Thorheiten II, 272. f. 275. 685.; Festungen, ist ihnen feind II, 932.; sein Fleisch II, 1322. f. 1383.; seine Frömmigkeit II, 772.; seine Gaben und Kenntnisse II, 1022. 1224.; seine Galgenreue I, 313.; sein Gebet I, 1547. 1581. f. 1702. 1704.; II, 60. 756. 1142. 1144. 1487. f. 1643. 1644. f. 1652.; äußerliche Gefahren, verachtete er II, 1983. f.; seine Gesundheit I, 1338.; II, 174. f.; seinen Glauben und Kleinglauben II, 286. 411. f. 434. 469. 1722. f. 2057. 2061. 2064. f.; Gott: lieber mit Gott in der Hölle, als ohne Gott im Himmel II, 569.; was Gott wohl gedacht, wenn L. ihm Vorschriften machen wollte II, 1289. f.; seinen Gram, worüber II, 1090.; Hausstand, als Mönch verachtete er ihn II, 1075.; Hofdiener, vereitelten oft seine Absicht I, 1467. f.; seine Hoffnung von der Wirksamkeit des Gebets der Kirche gegen Pabst und Türken I, 1112. f.; Irrlehren des Pabstthums, hingen ihm an II, 1990. f.; Irrthum, bekennt er II, 845. 856.; Isaaks Opferung, hätte nicht zusehen können I, 1511. f.; Ismaels und Hagars Geschichte, bewegt ihn sehr I, 1423.; seine Lehre I, 1428. f. 1573., vom geoffenbarten und verborgenen Gott II, 184. f.; Leibeigenschaft, der Tod wäre ihm lieber II, 1127.; Lügen, ist dazu nicht geschickt II, 272.; s. \gesperrt{Majestät}; Melanchthon I, 1286.; sein Messelesen II, 30.; Milchverbrauch, berechnete er II, 2049.; sein Mönchsleben I, 1576.; II, 318. f. 320. 988. 1107. 1115. 1301. 1569. 1879. f. 1890. 1957. 2004.; Mütter, die sich betrübten, weil ihre Söhne ins Kloster gingen, wie er sie gestraft II, 1679.; Pabsts Lehre, wornach er sie beurtheilt hat I, 696., Drohungen, er verachtet sie I, 1407.; II, 1985.; Papisten, hielten ihm den Namen der Kirche entgegen I, 790.; Plinius, ist ihm feind I, 1344.; Rabbinen und ihre Schriften, erkennt sie I, 370., stimmt nicht damit II, 698., widerlegte sie II, 997.; Rottengeister, sie fochten ihn an I, 130. f.; II, 1309.; seine Schriftforschung und -auslegung I, 1712. f.; II, 703. 1380. 1487. f.; seine sündliche Schwachheit I, 636.; frommen Schwiegersohn, schätzt ihn hoch II, 694.; die Siebenzig gelten bei ihm nicht viel II, 1767.; Sintfluth, Geschichte derselben bewegt ihn sehr I, 1206. 1209. f.; Sodoms Untergang ebenfalls I, 1206. 1209. f.; hebräische Sprache versteht er nicht vollkommen I, 322. 597.; II, 546. 678. 1998.; s. \gesperrt{Staupitz}; jüngster Tag, wartet mit Begier darauf II, 1713. 2059.; sein Testament I, 422. f.; seinen Tod I, 409. f. 1061. f.; II, 1150. f.; Träume, begehrt er nicht, kann sie auch nicht auslegen II, 1047. 1050. f. 1053. 1057. 1309.; Trost I, 1073.; II, 1411.; seine Unwissenheit II, 1932. f.; Vergleiche: L. vergleicht sich mit Isaak II, 128. f. 204., Jakob II, 270. f. 283. 684., Jakobs Söhnen II, 1034. 1142., Joseph II, 1433. 1881., Juda II, 1960., Lea II, 495. f., Lea und Rahel II, 630., Noah I, 537. f., den Patriarchen II, 300. Ruben II, 1928.; sein Verschieben des Betens, Abendmahlsgenusses etc. I, 1740.; seine Versuchung, ein Zeichen zu begehren I, 1527.; Verzug der Gnade, seine Anfechtung davon I, 927.; Vincentius, Summa angelica, er schaffte sich dies Buch an, ärgerte sich dann I, 815. f.; Weib, Leib desselben, seine Verwunderung darüber I, 247.; Welt, betrübt sich über ihre Bosheit I, 457. f., hofft, Gott werde sie strafen I, 509., meint, sie müsse vom Teufel besessen sein I, 1341.; seine Werke und Handlungen II, 463. f.; seine Widersacher I, 936. f.; Wohlthaten Gottes, klagt, daß er sie leicht vergißt I, 1338.; seine Zuhörer nicht aufmerksam II, 940. f. — S. \gesperrt{Prophezeiungen} 2.

2. \gesperrt{Luthers Schriften}. Seine Bücher von Manchen nicht gelesen, um die Lehre nicht kennen zu lernen I, 817.; werden verdreht werden II, 184. f.; von seiner großen Auslegung der Genesis I, 1. 1480. f.; von seinem Buche de servo arbitrio II, 176.; erwähnt sein Buch wider die Juden II, 1966.

3. \gesperrt{Luthers Zeit}. Aussprüche L.’s über seine Zeit in Beziehung auf: Abfall zu den Türken II, 51. f.; Adel: stolz I, 312.; II, 1371.; gegenseitige Beraubung II, 105., Beschäftigung adeliger Matronen II, 1213.; Beamte, Steuerbeamte etc. II, 1359. 1728. 1783.; Begräbnißgebräuche II, 1760.; Begrüßungen und Anreden II, 478. 821. 832. 1604. f.; Betrug im Handel und Wandel II, 1820.; Biersteuer II, 1820.; Bischöfe, tyrannisch I, 811.; Branntwein, wie gebraucht I, 712.; Deutschen: grob, bäurisch I, 409. 1629. f. 1709., Trunkenheit derselben II, 1623. (s. \gesperrt{Völlerei}, unten), Verhalten derselben gegen die Türken II, 931.; Deutschland I, 989. f. 1307.; Donat II, 58. f.; Ehestand: weltliche Rechte II, 856., verachtet I, 164. f.; II, 187., gemieden I, 144.; II, 1198., Heirathen nach Geld I, 1707., man wollte keine Kinder I, 1748.; II, 532. 579., besondere Schlafkammern der Eheleute I, 1719.; Eidgebräuche I, 1657. 1673.; Epicurer I, 1764.; Erbrecht II, 1925.; Ernte II, 192. 1185.; Evangelischen: karg, bedrückten einander I, 1138. f. 1193., man warf ihnen Abfall von der Kirche vor I, 555.; Feiertage I, 1598.; Familiennamen II, 2076.; Fleisch durch Frost vor Verfaulen geschützt II, 2073.; Frömmigkeit mangelt I, 1478.; Fürsten: tyrannisch I, 811., räth den Fürsten seiner Zeit, sich vom Pabst Frieden zu erkaufen II, 112. ff. 117.; Fürstenhöfe bargen viele Drohnen II, 1240.; Gärten, mit Dornen verzäunt II, 2075.; Gastfreiheit, gefährlich I, 1214.; Geiz, unmäßig I, 1197.; II, 586. ff.; Geographie Palästinas II, 2011. f.; Gesinde, untreu I, 1722. f. 1726.; II, 696. 1231. 1795. 1820.; s. \gesperrt{Gleichkauf}; Gottesdienste, Gebräuche dabei II, 1840. 1843.; Gottlosigkeit I, 453. 506.; Götzenbilder II, 650. f.; Grade der Doctoren etc., wie ertheilt II, 260. 1839.; Handelsstädte Deutschlands, in üblem Ruf I, 1314.; Hirten, sehr boshaft II, 697. 699. f.; Hochzeitsgebräuche II, 260. 499.; Hoffschranzen I, 1461. 1467. f.; Hurerei, allgemein I, 433., galt für keine Sünde I, 1225.; Indifferentismus II, 1830.; s. \gesperrt{Juden}; Jugend, frech I, 1636. f.; II, 1879.; Jünglinge, mit 18 Jahren noch nicht reif, Kinder zu zeugen II, 1167.; Käufe, wie bestätigt II, 117.; Kirche, Noth und kärgliche Versorgung derselben II, 60. 449. f. 632. ff., wovon erhalten II, 926. 1816. f., beraubt I, 1461. 1816.; II, 1733. ff.; Kirchgebäude, herrlich I, 839.; Kleider: Luxus I, 272. 1708. 1713. 1715.; II, 1039. f. 1809. f., aus Pelzen II, 1039., Schutz derselben wider Motten II, 293. f., Kopfbedeckung II, 1188. 1196.; Krankheiten I, 253. f.; Landstreicherplage I, 1139. 1146. 1720. f.; II, 1602. f.; Laster I, 253.; Lebensalter I, 1312.; II, 1800.; Lehenertheilung II, 260.; Lehre, die falsche nahm zu I, 253., die reine verfolgt II, 1249.; Mönche und Dompfaffen I, 165.; II, 223. f.; Nepotismus II, 1792.; Noth mehrte sich I, 785.; Obrigkeit, Klage darüber II, 1392.; neue Offenbarungen I, 1525. f.; Ordinationsgebräuche II, 261. 1345.; Pabst, hatte Glück II, 1113. (s. unten \gesperrt{Verfolgungen}); Pachtzins II, 1818.; Papisten, unbarmherzig I, 1138., achteten scharf auf den Wandel der Evangelischen II, 1029.; Praktiken der Bürger und Bauern II, 1099.; Prediger, hatten Ornat II, 281., verfolgt, beraubt I, 1358. 1396. 1459., litten Noth I, 1139. 1218.; II, 1235. f. 1242. 1735. 1816. 1956.; Predigtamt, verachtet I, 1632. ff. 1636.; II, 1367. (s. unten \gesperrt{Wort Gottes}); Reformation, wie beurtheilt II, 245. 276.; Regenten, Klage über sie II, 186. 1343.; s. \gesperrt{Rom}; Rottengeister II, 1347.; Sachsen, Zufluchtsstätte Vertriebener I, 1138.; Schildwachen I, 282.; Sicherheit II, 1811.; Sodomie I, 1222.; Stände, manche verachtet II, 1781. 1786.; Steuern II, 1362. 1810. 1817. 1819. f.; Sünden, schmückte man mit dem Namen der Heiligkeit I, 509.; Taufgebräuche II, 513.; Testamente, 7 Zeugen I, 1476.; Trauungsceremonien II, 1344. f.; Türke II, 1113.; Unbußfertigkeit II, 336.; Undankbarkeit I, 1341.; Unterthanen, bedrückt II, 85. f.; Verfolgungen durch Pabst und Türken I, 264. 311. 507. 960. 1135. 1138. f. 1255.; II, 53. f. 1698. 1712. 1954.; Völlerei I, 88. 256. f. 712. f. 858. 1219. 1384.; II, 1620. 1809. f.; Wittenberg: faul und fräßig Volk II, 1240., Klage über Fleischer etc. daselbst II, 1293. f.; Wort Gottes: schnell ausgebreitet I, 1711., viele Feinde bekehrt II, 817., sollte Schuld an Theurung sein I, 784. f.; II, 132. f., verachtet I, 409. 457. 462. 501. 527. 882. 1105. 1255. 1335. ff.; II, 321. 804. 1033. 1806. 1811., man könne daheim lesen II, 90., ging zu, wie zur Zeit der Sintfluth I. 462. f.; Wucher I, 1197. 1218. 1636. f.; II, 880. 887. 1362. f. 1806. 1810. 1817. 1820. — S. \gesperrt{Absolution}, \gesperrt{Köln}, \gesperrt{Löwen}, \gesperrt{Paris}, \gesperrt{Prophezeiungen}, \gesperrt{Reformation}, \gesperrt{Sacramentirer}, \gesperrt{Wein}.

Lyra. Ein Kenner der hebräischen Sprache I, 998.; geht dem Texte fleißig nach, weshalb ihn Luther allen andern Auslegern vorzieht I, 626.; seine Regel vom zwiefachen Verstande der Schrift I, 950.; folgt den Rabbinen I, 1002., hält viel von ihren Lügen I, 998.; bleibt fleißig bei der Historie, lenkt aber bisweilen zu unbequemen Allegorien I, 114.; Beispiele seiner Allegorien I, 283. 1004. \gesperrt{Meinungen und Aussprüche} über: 1 Mos. 46, 8. ff. II, 1764. f.; 1 Mos. 48, 6. II, 1855.; 1 Mos. 49, 11. 12. II, 1977. f.; 1 Mos. 49, 17. II, 2022.; Bethel II, 904.; Dreieinigkeit I, 274.; Egypter, Feindschaft wider die Juden II, 1615.; Gelübde II, 453. f.; Glauben, fides formata I, 945.; Himmelsleiter II, 394.; Hiob II, 1015.; Isaaks Tod II, 1067.; Jakobs Kampf II, 774. 780., Liebe zu Joseph II, 1064. f., Jakobs und Ismaels Namen II, 953., Jakob und Lea II, 499.; Josephs Ehrentitel II, 1402., bunter Rock II, 1037., zweiter Traum II, 1059. f. 1065., warum er seine Brüder herzutreten heißt II, 1682.; Josephs Brüder II, 1727.; Judas Wegzug nach Odollam II, 1160. f.;
\quellseite{563}
Leas Anfechtungen und Trost II, 527. f.; Leben, beschauliches und wirkliches II, 953. 956.; Lots Blutschande I, 1294. 1298., Erbieten, den Sodomitern seine Töchter zu geben I, 1220.; Materie I, 10.; Melchisedeks Brod und Wein I, 900. ff.; Messias, sollte Gott sein II, 938.; zweierlei Neid II, 534.; Pharao II, 1726., seinen Bäcker und Schenken II, 1294., Träume II, 1348.; Redefigur anticipatio und recapitulatio II, 1510.; Simeon und Levi II, 1943.; Simeon, den Stamm II, 1717.; Teraphim II, 649.; Träume II, 1308.; Tyrannen vor der Sintfluth II, 479.

\section{M}

Macedonien. Ableitung des Worts I, 661.

Macht Gottes. Man kann sie nicht durch Speculationen erkennen, muß sie erfahren II, 1383.

Macrobius. Schreibt über Träume II, 1051.

Madai. Sind die Meder I, 658.

Mädchen. Sollen nicht wild und Fensterguckerinnen, sondern häuslich sein, nicht ohne Erlaubniß, noch ohne Gefährten ausgehen II, 863. — S. \gesperrt{Geschlecht}.

Magd. S. \gesperrt{Beruf} I, \gesperrt{Gesinde}.

Magog. S. \gesperrt{Gog}.

Mahalath. II, 989. ff. Bedeutung des Namens II, 989. — S. \gesperrt{Esau} 3.

Mahanaim. Ortsname II, 733.; Lage II, 763. 846.

Mahlzeiten. Gebräuche verschiedener Völker II, 1619. f.

Mahomet. Lehrt, Jeder könne in seiner Religion selig werden II, 777.; seine Verheißungen sind teuflisch I, 757. — S. \gesperrt{Alkoran}.

Mahujael. Bedeutung des Namens I, 385.

Mailand. Gute Lage, daher viel Krieg um M. I, 878.

Majestät Gottes. Luther umgeht Fragen, die uns hinanziehen an den Thron der göttlichen M. I, 486. f. 1253.; man enthalte sich der Gedanken von der bloßen M. I, 520. f. 1084.; Disputationen darüber, die allergefährlichsten I, 486. f. 489. f. 1084. f., führen zu Verzweiflung I, 521. 1252.; wir sollen bei den Zeichen bleiben, darin sich uns Gott offenbart hat I, 490. f.; am jüngsten Tage werden wir die göttliche M. selbst sehen I, 491. f.

Mameluke. Ein Verfolger seines Ordens I, 689.

Mamre. Personen- und Stadtname II, 980. \gesperrt{Wer} M. war, sein Wohnort I, 875.; war gottesfürchtig I, 875., leutselig und fromm I, 886. 888., hielt sich zu Abrahams Kirche II, 1069. Was der \gesperrt{Hain} M. war I, 1141.; spätere Namen desselben I, 1609. — Siehe \gesperrt{Abraham} 29.

Manasse. 1. \gesperrt{Josephs Sohn}. Bedeutung des Namens II, 801.; Anlaß dazu II, 1427. f. 1432. — S. \gesperrt{Jakob} 21.

2. \gesperrt{Der Stamm}. S. \gesperrt{Ephraim} 2.

Mangel. Warum dabei geduldig sein I, 257.

Manichäer. Lehre von der Sünde I, 451., daß Bäume weinen II, 556.; verlachen die Patriarchen II, 556.

Mann. Seine Arbeit mühsam I, 257. f. — S. \gesperrt{Geschlecht}.

Mardachai. Hat Viele bekehrt II, 1773.

Margaretha. Schutzheilige der Gebärenden II, 1919.

Maria. 1. \gesperrt{Die Jungfrau} M. Im Ehestande geboren I, 235.; ihre Ehre ist, daß sie Jungfrau bleiben und Gottes Sohn gebären sollte I, 235.; ihr Glaube ein großes Wunder II, 417.; Papisten legen 1 Mos. 3, 15. von ihr aus I, 225. 234. f.; galt im Pabstthum für die Helferin in allen Nöthen II, 1919.; ihr zu Ehren ein 10. Chor im Himmel erdichtet I, 288. — S. \gesperrt{Christus} 2.

2. \gesperrt{M. Magdalena}. Warum sie 7 böse Geister hatte I, 824.

Martin, St. In welchem Falle er auch den Teufel absolviren wollte I, 336.; II, 1515.

Märtyrer. Die \gesperrt{christlichen} M. suchte man durch Unzucht zu fällen II, 1259. f.; überwanden den Tod nicht ohne großen Kampf II, 288.; wodurch sie überwanden I, 1238.; die Hoffnung des ewigen Lebens machte ihnen ihr Leiden zum Kinderspiel I, 1067.; ihr Blut schreit zu Gott I, 352. f. Wer ein M. \gesperrt{des Teufels} ist I, 263.

Maß. Größe desselben bei den Juden I, 1155.

Materia. Was Verschiedene so genannt I, 7. 10. f.

Matthias, König von Ungarn. Gelehrt, freundlich, fromm, doch oft tyrannisch I, 822.; kam aus dem Gefängniß zum Königreich II, 1786. f.

Maulpferde. Züchtung derselben ist wider Gottes Ordnung und die Natur II, 997. 1007.; sind unfruchtbar II, 1007. — S. \gesperrt{Ana}.

Maus. Entstehung, Creatur Gottes, hat schöne Gestalt I, 63.

Maximilian, Kaiser. Mit Geschäften überladen I, 260.; hielt sich einen Pfaffen zum Schreiber II, 1789. ff.; wie er, wenn er an einem Gerichte vorüberzog, dasselbe grüßte II, 1922. — S. \gesperrt{Bund}.

Mechtilde. Von Unglauben angefochten II, 783.

Meer. 1. \gesperrt{Im Allgemeinen}. Gott hat dem M. ein Ziel gesetzt II, 930. f.

2. \gesperrt{Mittelländisches} M. Erst seit der Sintfluth I, 120.

3. \gesperrt{Rothes} M. Vor der Sintfluth ein fruchtbares Land I, 120., bedeutet die Taufe I, 612. ff.

4. \gesperrt{Todtes} M. Heißt auch Salzmeer, früher das Thal Siddim I, 879.; ein schrecklicher Pfuhl I, 852. 879.

Mehetabeel. Bedeutung des Namens II, 1017.; ein vornehmes und treffliches Weib II, 1017.

Mehrung. S. \gesperrt{Fortpflanzung}.

Meinung, gute. Gilt nicht ohne Gottes Gebot II, 1757.; daher kam aller Aberglaube II, 1626.

Melancholiker. Sind oft in Gedanken versunken II, 290. f.; merken oft nicht, was vorgeht II, 500. f. 1471.

Melchisedek. Mag wohl Sem gewesen sein I, 898.; warum nicht Sem, sondern König zu Salem genannt I, 900.; hat seinen Namen vom Amt I, 899. 905.; war König und Priester I, 899., ein armer König I, 899., von Nachbarkönigen verachtet I, 899.; lehrte Vergebung der Sünden durch den zukünftigen Samen I, 899.; erhielt von Lot und Anderen den Zehnten I, 912.; hat Brod und Wein bei Abrahams Sieg nicht geopfert I, 900. ff.; größer denn Abraham I, 912. f.; mit M. ist der Pabst nicht zu vergleichen I, 900.; Figur Christi I, 898. 900. f. 905.; seine Historie ist in päbstischen Liedern besungen worden I, 900. — S. \gesperrt{Abraham} 30.

Mensch. 1. \gesperrt{Im Allgemeinen}. Gott ist der Poet, die M. sind seine Gedichte II, 1640.; ist ein Wunderwerk II, 1246., die allerherrlichste Creatur I, 102. \gesperrt{Ursprung} des M.: davon lehrt nur die Schrift, nicht die Philosophie I, 152. 156. 680. f. 683. f.; Schriftlehre davon ist der Vernunft lächerlich I, 102. f. 156.; jeder M. ist gleich im Anfang mitgeschaffen worden I, 92. Aehnlichkeiten und Unterschiede zwischen \gesperrt{M. und Thieren} I, 68. 81. 98. f. 101. f. 127. 421.; II, 486. f. \gesperrt{Vor dem Fall} hatte er ein leibliches und ein unsterbliches Leben I, 69.; Gott hatte an ihm seine besondere Lust und Freude I, 83. \gesperrt{Nach dem Fall} hat er ein natürliches und geistliches Leben I, 105.; Gott gedenkt noch jetzt das Beste über ihn I, 83. \gesperrt{Ziel des} M.: wozu er geschaffen, ist höchst wichtig zu wissen I, 159.; wozu er geschaffen ist I, 160., zu Gottes Dienst und Erkenntniß I, 97. f., zu einem höheren Leben denn die Thiere I, 68. f., zu einem ewigen Leben I, 55. f. 98. f. 402. f. 421. 678.; wornach man M. beurtheilen soll I, 797. ff.; die Gemeinschaft mit M. fliehe man nicht II, 1642.

2. \gesperrt{Der natürliche} M. Definition desselben I, 579.; ist Fleisch, was das heißt I, 460., ein Sünder I, 404., einfach gottlos und in der allertiefsten Finsterniß I, 484., außer der Gnade, unter dem Fluch I, 1427. f. 1568. 1571. 1942.; ist Gott feind, hat keine Erkenntniß Gottes I, 485.; seine Art und Natur ist irren und fehlen I, 1595.; hält das Schändlichste für Gottesdienst I, 455.; kann geistliche Dinge nicht verstehen I, 449.; versteht nicht die überschwängliche Gnade und Barmherzigkeit Gottes I, 1109.; hört lieber Gottes Verheißungen denn die Drohungen I, 1185. f.; deutet auch böse Zeichen zu seinen Gunsten II, 75. f.; je herrlicher der Schrifttext, desto schändlicher verfälscht er ihn I, 450.; ist hart und fühllos gegen das Evangelium I, 1105. f.; verachtet die Ordnungen Gottes, hält die Ordnungen des Teufels hoch I, 1593.; ist unbeständig und leichtgläubig I, 797. f. 867., abgöttisch und abergläubisch II, 1955., trotzig und verzagt I, 1438., hoffährtig I, 1687. f., selbstgerecht I, 1449., irdisch gesinnt II, 121. 1789., mörderisch I, 333. f.; II, 43., sein eigener ärgster Feind I, 440.; auf allerlei Sünden beflissen II, 1942.; kann sich wohl vor groben Sünden hüten I, 1552.; trotz der Sünde ist ihm die Hoffnung des ewigen Lebens gelassen I, 405.; Hagar ist ein Bild desselben I, 990. — S. \gesperrt{Geburt} 2, \gesperrt{Kräfte} 2, \gesperrt{Philosophie}, \gesperrt{Vernunft}, \gesperrt{Welt} 2, \gesperrt{Wille} II, 2.

Menschenblut. Das Gebot, M. nicht zu vergießen, war nöthig I, 595. 597. f.

Menschenkinder. Wen die Schrift so nennt I, 703.

Menschensatzungen. Haben einen großen Schein II, 246.; darin wohnt Gott nicht II, 437.; Ceremonien, von Menschen erdacht, denen man Vergebung der Sünden zuschreibt, sind Abgötterei und Mißbrauch des göttlichen Namens I, 280.; andächtige Gebete und Exempel der Heiligen helfen dabei nichts I, 280., werden nicht gut, wenn sie auch von den Vätern und Concilien kommen I, 1376.; führen von den rechten christlichen Werken ab I, 1163. f.; machen vermessen, daß man sich für heilig hält I, 1165.; entstehen, wenn man Gottes Wort verläßt II, 1760. Die \gesperrt{papistischen} M. sind wider Gottes Wort und zu fliehen II, 431.; sind von unseren kirchlichen Ordnungen verschieden II, 436. Man kann M. dulden, wenn sie nicht für nöthig erklärt werden I, 1419. — S. \gesperrt{Gottesdienst} 2, \gesperrt{Pabst} 1.

Menschenwitz. Schützt nicht vor Blindheit und Gottlosigkeit I, 499.

Menschenworte und -werke. Oft werden die Gedanken und Affecten der Gottseligen auf Gott selbst gedeutet I, 459. f.; es reut Gott, wenn es die reut, die das Predigtamt führen I, 486.; wann der Menschen Worte und Werke die Worte und Werke Gottes sind I, 165. f. 320. f. 346. f. 453. f. 460. 531. f. 735. 1183. 1247. f. 1395.; II, 59. 949.; Gott sieht, wenn die sehen, die im Predigtamt sind I, 486. — S. \gesperrt{Bekümmerniß}, \gesperrt{Betrübniß}, \gesperrt{Reue}.

Menschwerdung Gottes. Die Lehre davon kann man nicht genug treiben II, 403. f.; es ist etwas Großes, daß Gott seinen Sohn für uns gegeben hat II, 1130.; dieselbe war im A.\,T. angedeutet I, 106. 748. ff.; II, 395. ff.; Gott ward ein Mensch, daß wir uns nicht vor ihm fürchten sollten I, 1449.; ist eine hohe Ehre für das menschliche Geschlecht I, 1571. f.; ob Satan daraus Ursache zu seinem Fall nahm II, 400. ff. — S. \gesperrt{Christus} 1.

Meritum congrui und condigni. Was Beides sei I, 580.; II, 1458. 1468.

Mesach. Völkername I, 659.

Mesopotamien. Woher der Name I, 772.; fast allezeit mit Babel unter Einem Könige I, 772.; gute Weide, reich an Vieh I, 772.; II, 695.; hatte nicht viel Wasser, daher man die Brunnen mit Fleiß verwahrte II, 476.

Messe. Entstehung der Lüge von Opfermessen II, 1547. f.; schrecklich, daß der Pabst aus der Messe ein Opfer gemacht hat II, 921.; das Meßopfer ist ein Levitisches Opfer I, 901., überflüssig I, 904., eine Lästerung des Opfers Christi I, 904. f.; ein Greuel, der viel Geld eingebracht hat I, 1664.; II, 1548., daher die Papisten so fest daran halten II, 1549.; ein genügender Grund, sich vom Pabstthum zu trennen I, 904.; worauf die Papisten das Meßopfer gründen I, 900. 903. ff. 1242.

Messen. Die Frankfurter M. II, 1810.

Messias. Mußte Gott und Mensch sein I, 748. f. 1352.; II, 937. f., was die Juden leugnen II, 937. f.; sollte nicht ein irdischer König sein I, 917. f.; falsche irdische Hoffnungen der Juden von ihm I, 753. f.; die Patriarchen wußten Zeit und Stunde seiner Zukunft nicht I, 297.; ist gekommen I, 753. 1100.; die Juden erwarten ihn noch II, 319. — S. \gesperrt{Christus}, \gesperrt{Schiloh}.

Metapher. Redefigur II, 1909.; Beispiele: I, 357.; II, 715. 1261. f. 1909. 2018.

Metellus. Ermahnte zum Ehestande I, 1733. f.

Meteore. Meinungen darüber I, 606. f.

Methusael. Bedeutung des Namens I, 385.

Methusalah. Wann gestorben I, 532.; sein Testament an Noah I, 531.; Trauerzeit bei seinem Tode I, 532.

Metonymie. Beispiele: I, 224. 229. f. 267.

Micha. Erhielt Viele bei reiner Lehre II, 2047.

Michael. Ableitung des Namens I, 1013.; welcher Engel so geheißen haben soll I, 1198.

Midas. Fabel von ihm II, 598.

Midian. 1. \gesperrt{Abrahams Nachkomme}. Hoffährtig, hatte neben der Lehre Abrahams auch abgöttische Gottesdienste II, 1016.; von Hadad geschlagen II, 1016. f.

2. \gesperrt{Das Land}. Woher der Name, Lage I, 1751.

Midianiter. Unzüchtig, dienten dem Priapus I, 1662. f.

Milchverbrauch. II, 2049.

Miplezeth. Götzendiener II, 1610.

Mispat. Der Born M. d.\,i. des Gerichts I, 881.

Mißgeburten. Ursachen, Beispiele II, 601. f.

Mißtrauen. Nöthig II, 798.; rechte Beschaffenheit II, 799.

Mittel. \gesperrt{Gott} bedient sich der natürlichen M. I, 525. f. 560. f. \gesperrt{Wir} sollen die natürlichen M. brauchen I, 525. f. 890. f.; II, 1302. 1306. 1312. f. 1443. 1448. 1774. ff. und uns nicht auf die Vorsehung berufen II, 340. ff. 346. ff. 444. f. 746. ff. 764. ff. 1561. ff., auch wenn wir Gottes Gebot und Verheißung haben II, 466. ff. 637. f. 746. ff. 1774. ff.; man soll sie brauchen, das Andere Gott befehlen II, 1580. f., sich auf Gott verlassen II, 2045.; Verachtung der Mittel ist Sünde II, 1562.; wo sie nicht zu haben sind, sollen wir leiden und von Gott Hülfe erwarten I, 526. — S. \gesperrt{Gefahr}, \gesperrt{Verheißungen} 3, \gesperrt{Versuchen}, \gesperrt{Vertrauen}.

Mittheilung der Eigenschaften. Was man so nennt II, 775.

Mizpah. Bedeutung des Namens II, 716.; berühmte Stätte II, 716.

Mizraim. Ist Egypten I, 665.

\quellseite{564}
Moab. Bedeutung des Namens I, 1298.; warum Lots Tochter ihren Sohn so genannt I, 1298. ff.

Moabiter. Herrlich gesegnet I, 1298. f.; zu Kirchenämtern in Israel nicht zugelassen, doch viele selig I, 368. 1298. f.; waren Israel tributpflichtig, fielen ab II, 692.

Moloch. Wie gestaltet II, 1552. Entstehung des Molochsdienstes I, 1663.; wie die Kinder geopfert wurden II, 1552. f.; ihm dienten die Juden II, 1551., worauf sie sich dabei beriefen I, 568.; die von den Juden geopferten Kinder sind selig geworden II, 1553.

Mönche. 1. \gesperrt{Religion und Lehre der} M. Heißen Kemarim II, 1609. f.; laufen nicht zu Christo, sondern von Christo ins Kloster II, 1678. f.; leben ohne Wort in einer gemachten Religion I, 760.; ihre Gottesdienste selbsterwählt II, 1108.; was sie für den Weg zum Himmel halten I, 760.; verwerfen Reichthum I, 829. f., suchen ihre angebliche Armuth zu vertheidigen I, 835. f.; was sie für Keuschheit erklären I, 830. f., haben närrische Gedanken vom Ehestande II, 166. f., ärgern sich am Leben der Patriarchen I, 1315.; II, 1566.; heben den Gehorsam gegen die Eltern auf I, 1597.

2. \gesperrt{Leben der Mönche}. Ihr \gesperrt{Leben} gefällt Gott nicht, sondern dem Teufel II, 773.; ist ohne Gottes Wort II, 124. f. 643. 1731. f.; Henoch hat kein Mönchsleben geführt I, 418. f.; sie folgen ihren Stiftern nicht nach I, 759.; in ihrem Leben wird nichts Gutes ausgerichtet I, 1263.; es ist mit dem Leben der Patriarchen nicht zu vergleichen II, 127. ff., ist gar nichts, verglichen mit dem Hausstande I, 979. f.; ist nicht so schwer wie der Ehestand I, 1392. f.; man straft es billig I, 126. 760.; es ist Gottlosigkeit I, 1419., eine Versuchung Gottes I, 525., eine Zerrüttung aller Stände II, 1606. f. 1610. f., Verachtung des rechten Berufs I, 1071., eine Verachtung aller von Christo eingesetzten Werke II, 1611., eine Heuchelei II, 847. f.; ihr einsam „engelisch“ Leben ist teuflisch II, 943. Ihre \gesperrt{Geistlichkeit und Heiligkeit}: ihre Geistlichkeit ist Wollust und Wohlleben II, 128. f. 1075. 1199.; ihre Heiligkeit ist eine falsche I, 1315., selbsterwählt und Dreck I, 1392.; sie können in den Klöstern nicht vor dem Teufel fliehen I, 1374.; sind Baaliten II, 1611.; halten es für Heiligkeit, wenn sie über Tisch nicht reden I, 1159. f. \gesperrt{Selbstverleugnung, Tödtung des Fleisches}: da ist keine rechte Tödtung des Fleisches I, 1392. f. 1417. 1423.; tödten ihr Fleisch nicht wie Abraham I, 1505.; ihre Tödtung ist, wenn sie sehen, daß einem Andern besser Brod gegeben wird I, 1414.; ihre Kasteiungen sind nichts gegen die Tödtung des Fleisches bei Abraham I, 1414. f., oder Isaaks Anfechtungen in der Theurung II, 135. f. Ihre \gesperrt{Armuth} ist Thorheit und Undankbarkeit I, 836.; sie trachten nach irdischen Gütern II, 297.; finden in Klöstern mehr, denn sie verlassen I, 738. Wie \gesperrt{keusch} sie gewesen I, 831. ff. Ihr \gesperrt{Gehorsam} ist selbsterwählt und schändlich I, 1523., Teufelsgehorsam I, 764., voll Verachtung Gottes und der wahren Religion I, 759.; sie lassen den äußerlichen Gehorsam dem innerlichen vorhergehen I, 760. Sie dringen mit unbilliger Strenge auf ihre \gesperrt{Regeln} I, 845. f.; II, 1878. ff.; wollen Jedes Schuhe über einerlei Leisten schlagen I, 1074.; folgen der arithmetischen statt der geometrischen Proportion II, 511. f. Sie bleiben nicht beim \gesperrt{Wort} und gerathen in Irrthum I, 17., verachten die gemeinen, häuslichen Werke und sündigen damit I, 527. 1477. 1596. Ihre \gesperrt{Werke}: bei ihnen sind keine \gesperrt{gute Werke} I, 1163. 1311.; sie thun Alles ohne Gottes Befehl I, 569.; durch ihre scheußlichen Werke zermartern und verderben sie den Leib II, 1878. f.; verkaufen ihre Werke, Form eines Contracts II, 2004.; wollen Vergebung und Seligkeit verdienen II, 447.; ihr \gesperrt{Gebet} ist leichtsinnig und eine Lästerung Christi I, 1364. f., sie brauchen nach des Pabsts Decret nicht zu beten, sondern nur die Psalmen zu lesen und dabei nur auf die Worte zu achten I, 1365.; können nicht recht beten II, 2015.; verkaufen ihre Gebete I, 1363. 1365. Ihrer \gesperrt{Gefühllosigkeit} gegen Eltern und Verwandte rühmen sie sich II, 943., doch ist sie verwerflich II, 945. Ihre \gesperrt{Heiligen} rühmen sie über alle Maßen I, 733. f. Sie thun keine rechte \gesperrt{Buße} II, 320.; ihre \gesperrt{Werkheiligkeit} streitet wider den Segen, den Abrahams Same bringen sollte I, 1565. f. Fürchten sich vor \gesperrt{Tod} und Hölle II, 2004.; weisen einander in Todesnöthen auf ihre Patrone II, 2005. — S. \gesperrt{Mönchsleben}.

Vgl. \gesperrt{Barfüßermönch}.

Mönchsleben, Möncherei. Die Möncherei entstammt der Blindheit der menschlichen Vernunft I, 1093.; II, 705.; sie hat vor der Vernunft einen großen Schein II, 122. ff.; man hielt das M. vor der Reformation für den allerheiligsten Stand I, 1263., ist aber kein göttlicher Stand I, 1374.; man verglich es mit der Taufe I, 761.; durch die Möncherei sind die Menschen von dem Menschen der Sünde greulich verführt worden II, 1731. f.; es ist ein greulicher Irrthum, daß man wider den Willen der Eltern ins Kloster laufen soll II, 1679. f. Was von den Stiftern der Möncherei zu halten I, 759.; was für M. Augustin angerichtet I, 847. — S. \gesperrt{Mönche} 2.

Mond. I, 49. ff. — S. \gesperrt{Zeichen}.

Monica. Gebet für Augustin erhört I, 1111.; II, 574. 1329. f.; lachte ihn aus, wenn er in der Schule gestraft worden II, 1490. f.; stiftete Frieden II, 350.

Mord. M. und Todtschlag ist insgemein verboten I, 597.; ist eine greuliche, entsetzliche Sünde I, 340. ff.; den entsetzlichsten M. begehen, die die Heiligen tödten I, 343.; zieht schreckliche Strafe nach sich I, 355. f.; soll mit dem Tode gestraft werden I, 598. ff. — S. \gesperrt{Todtschlag}.

Mörder. Verschiedene Klassen derselben I, 597. — S. \gesperrt{Todtschläger}.

More. Anlaß zum Namen I, 776.; Lage I, 780.; wohl der Berg Garizim I, 779.

Morgen. „Land gegen M.“ I, 694.; II, 473. f.

Morijah. 1. \gesperrt{Der Berg}. Des Namens Ableitung I, 1492. f., allegorische Erklärung I, 1495.; Lage I, 1491. 1494. f. 1506.; ist nicht Bethel II, 903. f.; Abraham nennt ihn „der HErr siehet“ I, 1541. f.; war berühmt I, 1494. 1543.; da frug man den HErrn um Rath I, 1544.; Ereignisse daselbst II, 438. 452.; Schicksale des Orts II, 438. f. 903.

2. \gesperrt{Das Land} I, 1491. ff.

Mose. 1. \gesperrt{Mose}. Hatte einen starken Glauben II, 286., der mit dem Fleische kämpfen mußte II, 288. 742. f. 1503.; erkannte die Wohlthat des Wortes Gottes I, 863.; war voll Heil. Geistes und Glaubens und eines unverzagten Herzens II, 1380. f.; ein unverzagter Kriegsmann II, 1376. 1389. f.; hat das Leiden Christi gesehen II, 1048.; durfte nicht in das Gelobte Land, ob ihm wohl die Sünde vergeben war II, 1748.

2. \gesperrt{M.’s Schriften}. M. folgt der chronologischen Reihenfolge nicht I, 920. f.; unterscheidet sich von heidnischen Schreibern durch Kürze I, 542., zeigt mit wenig Worten viel an I, 340. ff. 408. f.; warum er Manches wiederholt I, 541. f. 547. f.; II, 992. f.; gebraucht bei Erzählung früherer Geschichten die zu seiner Zeit üblichen Ortsnamen I, 876.; warum er so manche unfläthige Dinge beschreibt II, 539. ff.; hat viel Ceremonien und Gesetze von den Vätern hergenommen I, 1220. f.; II, 143. 922. 1178. 1208.; ob er Alles selbst geschrieben II, 1010.; seine Schriften sind für die Kirche bis ans Ende II, 403. Das 1. \gesperrt{Buch Mose}: Urtheil frecher Juristen darüber I, 294. 968., der Papisten I, 434. 968. 1172.; II, 1167. f.; bei den Juden durfte es niemand vor dem 30. Jahre lesen I, 4.; die Kirche hat sich stets viel daran gelegen sein lassen I, 716.; ein gewaltiges und treffliches Buch I, 59.; das schönste Buch in der Schrift und darin wieder die Geschichte Josephs die schönste II, 1022.; wir sollen bedenken, daß der Heil. Geist der Meister desselben ist II, 566. \gesperrt{Das 1. Capitel} desselben: handelt von den wichtigsten, dunkelsten Sachen I, 4.; noch Keiner hat es eigentlich und richtig genug ausgelegt I, 4.

3. \gesperrt{M.’s und der Väter Gesetz}. Unterschied II, 144. f.

Münster. 1. Dr. \gesperrt{Sebaldus} M. Gestorben I, 1480.; lebt doch I, 1488.

2. \gesperrt{Sebastian} M. Seine Erklärung: des Namens Naphthali II, 546., von 1 Mos. 46, 8. II, 1767.

Münzer. Ein aufrührerischer Geist I, 610.; ehrgeizig I, 924.; verkehrte Alles in Allegorien I, 610.; reizte die Bauern zum Aufruhr II, 515. 533.; wollte die Patriarchen nachahmen I, 955., mit seinen Bauern Abrahams Exempel nach kriegen I, 889. f. 1233.; meinte, Luther habe wohl das Evangelium angefangen, aber es nicht hinausgeführt II, 245.; rühmte sich, Gott rede in sichtbarer Gestalt mit ihm I, 1119.; schrie immer: Geist, Geist I, 1668.; rühmte sich seiner Gefühllosigkeit II, 1500. 1572. f., die er auch in die Kirche einzuführen suchte II, 1608.; was er „entgroben“ nannte II, 954.

Musik. Von den Kainiten erfunden I, 387.

Muth. Gott gibt und nimmt ihn II, 1378. f.

Muttermale. Woher sie kommen II, 601. f.

Myrrhe. Ein Harz II, 1124.; schützt vor Verwesung II, 2073.; Christus ist die Myrrhe der Gläubigen II, 2073.

\section{N}

Nachkommen. Man soll täglich für sie bitten II, 1942., daß die Lehre rein erhalten werde I, 534. f.; sollen am Segen nicht verzweifeln, wenn auch ihre Vorfahren sich Gottes Zorn zugezogen I, 1299. — S. \gesperrt{Abfall}.

Nächster. Christen trauen Keinem, sind auch Keinem feind I, 801.; Christen sollen alles Böse verborgen halten und zum Besten kehren II, 348. f., sollen ihn vor Gefahr warnen II, 350. f.

Nacht. S. \gesperrt{Tag} 2.

Nachtschatten. Frucht und Gebrauch des N.’s I, 565. f.

Nackendgehen. War vor dem Fall ein Vorzug des Menschen, seither eine Schande I, 170. ff. 202. ff., das lehrt nur die Schrift I, 206.; ist schändlich, weil wir den Glauben verloren haben I, 204. ff.

Nadab und Abihu. Nicht verdammt I, 1282. f.

Naema. Bedeutung des Namens I, 387.

Naeman. Durch Elisa bekehrt II, 1830.

Nahor. 1. \gesperrt{Nahor}. Aelter als Abraham I, 1599. f.; regierte eine große und berühmte Kirche I, 1601.; II, 1015. — S. \gesperrt{Abraham} 31.

2. N.’s \gesperrt{Geschlecht}. I, 1599. ff. Erbte die Lehre von Nahor II, 1015.; hatte zu Jakobs Zeit noch Gottes Wort II, 383. f. 385. f.; hatte treffliche Lehrer, viele aus N.’s G. selig II, 385.; ob Hiob und Bileam dazu gehört II, 383. 385.

Nährstand. S. \gesperrt{Hausstand}.

Nahrung. Sie zu erlangen und bewahren, ist erst seit dem Fall schwer I, 125. f.

Namen. 1. \gesperrt{Gottes} N. „Gott der Höchste“ I, 906.; Hagar nennt ihn: „Der mich siehet“ I, 1001.; „elschaddai“ I, 1013. f. 1015. f. Gottes N. kann recht und falsch gebraucht werden I, 1406. f., wird oft mißbraucht I, 1403.; wird auch Engeln und Menschen beigelegt I, 14. f. — S. \gesperrt{Jehovah}, \gesperrt{Zebaoth}.

2. „\gesperrt{Der} N. \gesperrt{des HErrn}“ ist Christus I, 399.

3. \gesperrt{N. der Menschen}. Die N. der M. sind in der Schrift nicht bedeutungslos, Beispiele: I, 383. 398. 416. 664.; II, 13.; im Hebräischen ist ihnen oft der N. Gottes angehängt II, 937.; N. nimmt man gern von Königen und Heiligen II, 1415.; in den N. ist große Unordnung in der hebräischen Sprache II, 1016.; Veränderung der N. ist sehr gewöhnlich II, 988. ff. 1715. In der Taufe bekommen wir einen neuen N. II, 950. f.

4. \gesperrt{Guter} N. Ist bald geraubt II, 1723.; man soll ihn wahren II, 1303. f.

Naphthali. 1. \gesperrt{Jakobs Sohn}. Wann geboren II, 548. ff.; Bedeutung des Namens II, 546. f.; Jakobs Segen über N. II, 2034. ff. und Erfüllung des Segens II, 2037. f. — S. \gesperrt{Jakob} I, 19.; II; \gesperrt{Joseph} 12.

2. \gesperrt{Der Stamm}. Wohnsitz II, 2018. 2034. f.; erlitt viel II, 2037. f.

Narrheit. „Eine N. begehen“ II, 868. f.

Nasiräer. Waren theils für kurze Zeit, theils lebenslänglich N. II, 2053.; hatten sonderliche Ceremonien II, 2052.

Natur. 1. Die ganze N. u. \gesperrt{Welt}. Wird alt II, 857. f.

2. \gesperrt{Knechtische Naturen}. Müssen mit Strenge behandelt werden I, 989.

3. N. u. \gesperrt{Gnade}. Der Heilige Geist zerbricht die N. nicht in den Heiligen, tödtet die Affecten nicht, sondern heilt sie I, 1411. f.; II, 1500. f. 1558. f. 1606. ff. 1672. f. 1679. f., doch ist die Heilung keine vollkommene II, 1503.; die Heiligen haben auch noch menschliche Neigungen, ja, heftigere als Andere I, 1508.; je heiliger der Mensch, desto zarter die Neigungen I, 1613.; II, 1606. f.

Naturgesetze. Binden Gott nicht I, 33. 37.

Nebajoth. Sohn Ismaels, Wohnort II, 17.

Nebathea. Theil Arabiens, woher der Name II, 17.

Nebucadnezar. Erkannte den wahren Gott I, 1023.

Neid. Ist zweierlei II, 534. f.; eine Tugend des Teufels I, 1457.; gebiert Unrecht I, 1459.; gerechte Strafe des Geizes II, 612.; durch Anderer N. werden die Frommen geplagt I, 1459., das gibt Gelegenheit zu Geduld und guten Werken I, 1459.

Neigungen. Die natürliche N. zu Eltern etc. ist keine Sünde II, 944. f., findet sich auch bei den Thieren II, 947. f.; herzliche N. der Menschen zu einander gefallen Gott wohl II, 1778. — S. \gesperrt{Affecten}, \gesperrt{Brautliebe}, \gesperrt{Natur} 3, \gesperrt{Väter} II, 5.

Nepotismus zu Luthers Zeit II, 1792.

Nicäa. Daselbst 318 Bischöfe versammelt I, 888.

Nicolaus, St. Patron der Schiffsleute I, 895.

Niedrigkeit. Göttliche Herrlichkeit verbirgt sich in N. II, 482.

Nil. Vom N. hängt Egyptens Fruchtbarkeit ab II, 1701.

Nimrod. I, 668. ff. Bedeutung des Namens I, 668.; war entweder der erste Türke oder der erste Pabst I, 668. ff.; suchte weltliche und geistliche Gewalt I, 669.; sündigte wider Kirche und Polizei I, 687.; gründete Babel I, 877.; störte den seligen Zustand nach der Sintfluth I, 685.; richtete neue Gottesdienste an I, 670.
\quellseite{565}
711. 726.; er und sein Anhang richteten eine neue Kirche an I, 687. ff.; seine Religion hatte den größten Schein I, 731.; seine Rotte mehrte sich I, 722. f., verführte Tharah und sein Geschlecht I, 725.

Ninive. Von Assur erbaut I, 673. ff.; viertheilig I, 675.; wann und warum zerstört I, 676.; „große Stadt Gottes“ I, 675. 849.; II, 546. f. — S. \gesperrt{Assur}.

Niniviten. Verhalten zu Jonas Predigt I, 676.; ihre Buße bestand nicht im Fasten, sondern im Seufzen II, 1108.; nicht von Jona beschnitten I, 1014.

Noah. 1. \gesperrt{Im Allgemeinen.} I, 428. ff. 432. ff. 449. ff. 497. ff. 509. ff. 521. ff. 530. ff. 541. ff. 556. ff. 566. ff. 588. ff.

2. \gesperrt{Name} N.’s. Bedeutung desselben I, 429. ff. 443.

3. \gesperrt{Ehre und Bedeutung} N.’s. N. ist der höchste Prophet I, 464. f., der höchste unter den Vätern nach der Sintfluth I, 711.

4. \gesperrt{Offenbarungen.} N. empfing sie wohl erst durch Methusalah I, 531. f., darnach unmittelbar I, 532.; Gott hatte sonderliche Lust, mit N. zu reden I, 539.

5. \gesperrt{Bund Gottes mit} N. I, 518. ff. 598. ff.

6. \gesperrt{Glaube} N.’s \gesperrt{und Gerechtigkeit vor Gott.} I, 497. ff. N. war voll Heil. Geistes I, 639. 644. f.; hatte einen starken Glauben I, 466. f. 537.; glaubte die Drohung von der Sintfluth I, 466.; Gott stärkte seinen Glauben I, 518. f.; N. war allein gerecht vor Gott I, 499. 533. 537.

7. \gesperrt{Tugenden} N.’s. N. ist ein Exempel aller Tugenden I, 501.; war fromm I, 499., geduldig I, 501., gehorsam gegen Gott I, 466. 526. ff. 530. ff. 541. 566., keusch I, 432. ff., muthig I, 500., standhaft I, 467. 498. 538.

8. N.’s \gesperrt{Amt} u. \gesperrt{Werke.} N. hatte das Priesteramt I, 569.; war nach der Sintfluth Bischof, Regent und Hausvater I, 628. 637.; von N.’s Mühe im Kirchenamt I, 710. f.; von N.’s Predigen vor der Sintfluth I, 461. 463. f. 501.; wie er die Seinen getröstet I, 558. f. 573. f.; N.’s Opfer I, 567. ff. 572., Weissagung I, 696. f.

9. \gesperrt{Ehestand} N.’s. Warum er sich des Ehestandes enthalten wollte I, 433. ff. 443.; warum er in denselben trat I, 433. 435. f. 464. 501.; um desselben willen ward seine Predigt verachtet I, 463. f. 502.; hatte noch keine Kinder, da die 120jährige Frist anging I, 465.; aus seinem Ehestande entspringt eine neue Welt und Kirche I, 502.

10. \gesperrt{Sünden} N.’s. I, 627. ff. 711. Warum die Schrift den Fall N.’s erzählt I, 628. f. 637. f.

11. N.’s \gesperrt{Anfechtungen} u. \gesperrt{Leiden.} N. war ein Märtyrer über alle Märtyrer I, 442. f. 693.; ward auf mancherlei Weise angefochten I, 530.; sein Glaube ward angefochten und schwach I, 497. 556. ff.; fühlte wohl Verzweiflung I, 558. f.; es war eine große Noth, so lange in der Arche umhergetrieben zu werden I, 560.; ward geängstet und betrübt I, 467. 531. 541., angefeindet I, 498. ff., verachtet I, 478. f. 503. 537. 656. f. 693.; wie er geseufzt I, 506.; von dem Trost, den er empfing I, 537. 564. ff. 604.; hielt sich an die Verheißung I, 557.

12. \gesperrt{Tod} N.’s. I, 655.

13. \gesperrt{Zeit} u. \gesperrt{Zeitgenossen} N.’s. N.’s Zeit war die allerärgste I, 467.; N.’s Zeitgenossen I, 407. 639. ff. 710.

14. N. u. \gesperrt{seine Söhne.} I, 640. ff.

15. N. u. \gesperrt{die erste Welt.} N. hat die Welt geistlich und leiblich erhalten I, 502.; konnte nicht ohne Gefahr predigen I, 495.; seine Predigt verachtet I, 463. ff.; man hielt N.’s Predigt wohl Gottes Verheißungen entgegen I, 510.

Nod. Bedeutung des Namens, Lage I, 377. f.

Nonnen. Werden von Sarah übertroffen I, 1181.; verachten und verlassen den rechten Beruf I, 1071.; ihre Werke, Abgötterei I, 1628.; heißen kemaroth II, 1609. f. — S. \gesperrt{Mönche}.

Noth. Wird einem Kelche verglichen I, 617. Gott läßt seine Gläubigen in N. und Gefahr gerathen, damit sie nicht sicher werden II, 760. f., verläßt und vergißt sie aber nicht I, 557. 560.; II, 687. 901.; darin soll man weder Gott versuchen, noch verzagen II, 341. f., alle Sorge auf Gott werfen II, 1695., sich auf Gottes Hülfe verlassen I, 1319., am Wort halten I, 1342.; II, 1695. — S. \gesperrt{Kreuz}.

\section{O}

Obrigkeit. 1. \gesperrt{Gottes Ordnung} ist die O. I, 601.; II, 1340. f., die heidnische wie die christliche I, 1054., auch die gottlose I, 1409.; führt das Regiment auf Gottes Befehl I, 1456.; was sie thut, thut Gott I, 211. f.; gehorcht man ihr, so gehorcht man Gott I, 763.; ist geordnet zum Nutzen dieses zeitlichen Lebens II, 1342.

2. \gesperrt{Beruf dazu.} Ohne Beruf soll man sich nicht zu diesem Stande drängen II, 1341. f., um der Beschwerden willen ihn nicht fliehen II, 1408. f.

3. \gesperrt{Der Stand} der O. ist ein schwerer, mühevoller I, 259. f. 1371. f.; II, 1405. ff.; sie ist Knecht aller Knechte II, 552. f.; ein nöthiger Stand I, 1461. f.; II, 1397. 1404., nützlich II, 1397. 1404., Gott gefällig I, 1371. f.

4. \gesperrt{Amt und Werke} der O. Ihr \gesperrt{Amt} ist, Friede, Zucht, Ehrbarkeit zu erhalten II, 1394. f.; soll Heiland und Arzt sein I, 879. f., die Frommen schützen I, 821., II, 86., die Bösen strafen I, 821. 1259. 1319., wenn auch Gott die Sünde vergeben hat II, 1921. f. 1935.; säumt die O., so straft Gott die Uebelthäter II, 1922.; soll für der Unterthanen Nothdurft, in der Theurung für Speise sorgen II, 1359. f., gegen Landstreicher einschreiten II, 1602. f., auf die Ausländer Acht haben II, 1449. 1471. 1787. Ihre \gesperrt{Werke} gehören zu den allergrößten I, 528.; sind, im Glauben gethan, besser als der Mönche Werke I, 1320.; fromme O. will ihren Unterthanen nützen und dienen I, 1467.

5. \gesperrt{Amtsführung.} Die O. soll ihres Amtes warten I, 1461. f.; muß ein fleißiges Aufsehen haben und wachen II, 1241.; bedarf Diener und Gehülfen II, 1403.; soll demüthig sein I, 1471.; II, 1342. f. 1391. 1393., beten II, 278. 1391. f.; nur Gläubige haben den nöthigen Muth und Kraft I, 1319. f.; II, 1410.; schlechte Räthe hindern oft eine fromme O. I, 1467. f.; soll nicht auf Schmeichler und Hoffschranzen hören II, 1930. f.

6. \gesperrt{Sünden} der O. Sie sündigt oft I, 821.; daß eine O., die nicht vom Heil. Geist regiert wird, sündigt, ist kein Wunder I, 822.; sündigt damit, wenn sie nicht weiß, wo Unrecht geschieht I, 1466. ff.; solche Sünden unvermeidlich I, 1469. ff.; wie Gott die S. der O. straft I, 821.; ihre bösen Exempel sind schlimmer als die der Unterthanen II, 865. f.

7. \gesperrt{Gewalten und Rechte.} Irdische Gewalt ist an sich nichts Böses I, 669.; sie hat G. über Leben und Tod und alles, was geringer ist I, 599.; II, 644.; Vernunft und Philosophie halten die G. der O. für Tyrannei I, 600. f.; die G. der O. ist eine Wohlthat I, 599. f.; soll sich nicht durch das Gebot der Liebe vom Strafen abhalten lassen I, 800. Die O. darf nicht thun, was sie will, soll sich züchtig und gottselig halten II, 865.; Schoß, Renten und Zins werden ihr gegeben I, 259.; äußerlicher Schmuck ist recht und nöthig II, 1397. 1403. f. 1410.

8. \gesperrt{Schlechte und gute} O. \gesperrt{Schlechte} O. eine Strafe II, 1340. f. 1356.; gottlose O. sind gleichsam Gottes Mastsäue I, 476. f., sind Teufel und Tyrannen I, 1468.; O. wird tyrannisch, wenn nicht Gottes Wort sie regiert I, 789., durch des Teufels Treiben I, 882.; auch durch die böseste kann Gott das Regiment verwalten II, 2047. \gesperrt{Gute} O. gibt Gott, wenn er einem Volke günstig ist II, 1349. f. 1356.

9. \gesperrt{Feind} der O. ist der Teufel I, 637.; II, 1404. 1406. f.

10. \gesperrt{Schrift und} O. Die Schrift gebietet, die O. zu ehren, scheint aber selbst sie nicht zu ehren, sondern straft sie I, 475. f.

11. O. \gesperrt{und Unterthanen.} Weise O. und gehorsame U. sind eine große Gabe Gottes I, 1632.; wo O. und U. fromm sind, da ist ein rechtes Paradies I, 1476. f., ebenso, wo die O. fromm ist und die reine Lehre schützt I, 1631. f.; Beide werden erst zu Narren, wenn sie sollen gestraft werden II, 880.; U. müssen die Thorheit der O. büßen II, 880. ff.; die O. sollte nicht um der Strafe willen gefürchtet, sondern um Tugend und Weisheit willen geehrt werden II, 1422. f.; Christen sollen aller göttlichen Ordnung unterthänig sein II, 831.; man soll der O. gehorchen, doch nicht wider Gottes Wort II, 267. 645. 831. f. 943.; von diesem Gehorsam kann der Pabst nicht entbinden II, 642. f. 645.; man soll sie ehren II, 1404.; ob das Schmeicheln ist, wenn man sie „Gnädigste“ etc. nennt II, 832.; man soll für sie beten I, 821. 1468.; man soll sie strafen I, 1469.; II, 231. ff.; fromme O. zürnt nicht, wenn man sie straft I, 1469.; die Sünden einer frommen O. soll man zudecken I, 1468. f.; man soll ihr das begangene Unrecht anzeigen I, 1470. ff.

12. O. \gesperrt{und Kirche.} O. soll sich der Kirche annehmen II, 1734. f.; wo sie das herzlich thut, das ist eine große Gabe Gottes I, 1335.; O. dient der K., wenn sie Frieden erhält II, 1342. — S. \gesperrt{Kirche} 12.

13. \gesperrt{Christ und} O. Christen dürfen die Hülfe der O. anrufen I, 1468. f.; müssen auch das Weltregiment verwalten II, 458. ff.

Vgl. \gesperrt{Könige, Königreiche, Regiment, Stände}.

Ochsen. Erstgeborne O. lobt die Schrift II, 1948. f.; O., die man vor den Pflug spannte, durfte man nicht opfern II, 1948. f.; ihre Arbeit lobt die Schrift II, 1949.

Oecolampad. Ehrgeizig I, 924.

Odollam. Lage II, 1161.

Offenbarung. 1. \gesperrt{Die} O. \gesperrt{Gottes} ist die O. aller drei Personen I, 71.; Gott hat \gesperrt{stets} durch Zeichen offenbart, wo er zu finden sei I, 1044. ff.; Gott offenbart wohl dem Geringeren etwas und nicht den Großen I, 1410.; verschiedene \gesperrt{Weisen} der O. I, 13. f. 487. ff.; dreierlei Art: Gesicht, Traum, mündliches Wort I, 930. 1117. 1321.; II, 904. f.; O. durch Eingebung ist die höchste Art, dann durch Gesichte, dann durch Träume II, 1048.; O. durch das mündliche Wort ist die gewisseste Art I, 931. f.; O. durch Gesichte, Träume, Erscheinungen ist eine Gabe II, 1047. 1053.; zweierlei Art: durch das Predigtamt und durch innerliche O. I, 533.; Gott erscheint uns im N.\,T. im Wort und Sacrament I, 1118.; II, 177. ff. 1048. f. \gesperrt{Innerliche} O. sind selten, nur in sonderlichen Fällen I, 532. f.; Gott will nicht Jedem sonderliche O. geben, sondern weist uns an das Wort II, 560. f.; wir sollen auf keine besondere O. warten, wo das Predigtamt ist I, 532., sollen uns an Gottes O. halten I, 14. 16. ff.

2. \gesperrt{Neue} O. Der Heil. Geist lehrt nicht durch neue O., sondern durch das Predigtamt I, 674.; sind unnöthig I, 1527., verdächtig I, 1120. f.; darauf sollen wir nicht warten I, 1526.; man soll sie fliehen I, 1526.; müssen dem Glauben ähnlich sein I, 1526. f. 1531.; wer sich derselben gerühmt I, 1525. f.

3. O. \gesperrt{des Teufels.} Der Teufel gibt O. durch Gesichte und Träume II, 1298.

Vgl. \gesperrt{Erscheinungen, Gesichte, Träume}.

Og. Wann er gelebt I, 884.; woher der Name I, 1156.

Oel. Bedeutet Gnade und Vergebung I, 624.

Oelung, letzte. Hat keine Verheißung, wirkt keine Vergebung, ist Mißbrauch des Namens Gottes I, 280.

Onan. II, 1178. ff. Woher der Name II, 1163.

Opfer. I, 302. ff. Die Gewohnheit zu opfern, war vom Anfang der Welt I, 302.; O. sind nicht erst von Mose verordnet I, 539. 571. f.; haben sich von Adam forterhalten I, 966.; von Mose nur in gewisse Ordnung gebracht I, 302.; die von Mose beschriebene Weise hielt man schon früher ein II, 916.; die O. der ersten Eltern waren mit Predigt und Anrufung verbunden I, 303. und mit einem Opfermahl I, 1598. f.; im O. gilt nicht das Werk, sondern Gottes Gnade und Verheißung I, 306.; sie waren Zeichen der Gnade I, 303. f.; Gott offenbarte seine Gnade, indem er die O. mit Feuer vom Himmel anzündete I, 304. 307. 725.; sind ein Zeugniß, daß man einen Gott hat, den man ehrt und fürchtet I, 573.; wozu und auf wie lange verordnet II, 917. f. 921.; versöhnten Gott durch den Glauben an das zukünftige O. II, 1928.; durften nicht überall geschehen I, 1495.; die Schrift rühmt nicht, daß die Patriarchen Thiere geopfert hätten II, 916. f. O. brauchte man bei den Heiden bei \gesperrt{Bündnissen} I, 967.; wie diese O. angestellt wurden I, 967. Rechte O. sind Predigen und Hören, Danken und Beten II, 917. 922. — S. \gesperrt{Dankopfer}.

Ordination. Gott macht und ordinirt die Prediger II, 435. 437.; Ceremonien bei der O. II, 261. 1345.

Ordnung. Welche O. Gott zu halten pflegt, wenn er etwas thut I, 505. f.; Gottes O. verachtet der natürliche Mensch und hält die O. des Teufels hoch I, 1593.

Orestes. S. \gesperrt{Pylades}.

Origenes. Aergert sich an dem ehelichen Leben der Patriarchen II, 556.; rühmt Allegorien I, 284.; allegorisirt I, 284.; II, 557.; seine Allegorien sind fehlerhaft I, 612. 626.; II, 774. 811.; Allegorie vom Paradies I, 109.

Orimasda. Was O. sei I, 727. 1665.

Ornat. Trugen Prediger zu Luthers Zeit II, 281.

Osterrock. Wir wollten gern das alte Kleid behalten und den O. anziehen II, 1993. — S. \gesperrt{Himmel} 5.

Ostwind. War dürr und heiß II, 1346.

Ovid. Wurde allegorisirt II, 559.

\section{P}

Pabst. 1. \gesperrt{Pabst. Heißt} zu Rom Vater der Findlinge I, 293.; heißt nicht der Antichrist, sondern der Allerheiligste I, 479.; \gesperrt{will} trotz aller Bubenstücke der Statthalter Christi sein I, 505., Herr der Welt und Kirche I, 765. f., Herr über das Wort I, 766.; maßt sich den Namen der Kirche an II, 112. f., ist aber nicht
\quellseite{566}
die Kirche I, 1402. 1431. ff. 1454.; II, 112. f. 119. f. 431. f., sondern ein Nachkomme Kains und Hams I, 650., dem Nimrod ähnlich I, 668. ff.; hat Kains und Lamechs Gewissen I, 392. f.; ein Tyrann und den Sodomitern gleich I, 473. 478. f. 1147. f., ein Räuber, Dieb und Mörder I, 1621. 1623., ein Simonist I, 1621. ff.; II, 110. f., ein gottesfürchtiger und gelehrter Allegorienmeister I, 612.; ist der geweissagte Sündenknecht I, 766. 900. 1062. f. 1317. Seine \gesperrt{Religion} ist lauter Heuchelei I, 765.; treibt Pracht mit der geweihten Hostie, wie die Perser mit dem heiligen Feuer I, 727.; lebt und stirbt ohne Gott II, 1732.; gebietet schädliche Dinge I, 766. f., falsche Werke I, 867. \gesperrt{Pabst und Christus} lassen sich nicht vereinigen II, 431., er hat Christum gar begraben I, 1063.; erhebt sich über Christum, macht seine Decrete dem Worte Gottes gleich II, 2009.; will sich wie die Sodomiter nicht strafen lassen I, 1229.; glaubt den göttlichen Drohungen nicht I, 551. f., auch nicht die Auferstehung der Todten I, 1521.; übt unter dem Namen der Religion den höchsten Frevel I, 505.; will das Predigtamt aufheben I, 534.; streitet wider die erkannte Wahrheit II, 1965. ff.; macht aus den Sacramenten ein menschlich Werk I, 905.; mißbraucht Gottes Wort und Namen zu Gotteslästerung I, 905.; hat den rechten Weg zur Seligkeit fahren lassen, sich einen eigenen erwählt I, 1071.; seine Predigten sind nur Schreckpredigten I, 1205. Hat einen teuflischen \gesperrt{Haß} gegen die Kirche I, 343.; wartet nur auf Gelegenheit zu toben und zu wüthen I, 332.; verfolgt die Christen II, 53. f. 2002.; wüthet greulicher gegen sie denn Diocletian II, 1984.; macht die Welt voll Mordens I, 1275. Dringt hart auf \gesperrt{Gehorsam} I, 764.; hat durch den geforderten Gehorsam Viele in die Hölle gestürzt I, 765. 767.; wie man ihm auf seine Forderung gehorsam zu sein antworten soll I, 765.; maßt sich fälschlich an, vom Gehorsam gegen die Obrigkeit zu entbinden II, 642. f. 645.; erklärt die Beobachtung seiner Satzungen für nöthig und macht dem Gewissen Stricke I, 1316. f., hat dadurch die Welt mit Sünden erfüllt II, 2077.; fordert, daß man ihm die Füße küsse I, 766.; beraubt die ganze Welt II, 1956., bringt irdische Güter an sich I, 649., gebraucht dazu die Schlüssel II, 115.; seine Verheißungen sind teuflisch I, 757.; seine Drohungen sind nicht zu fürchten I, 1407.; wie allen Gottlosen gehts ihm wohl I, 667.; Gott wird ihn zerstreuen I, 702. ff.; ändert beim Regierungsantritt seinen Namen II, 988.; hat eine dreifache Krone I, 1713. Wer ihn als Meister anerkennt, hat kein Theil an Christo I, 868. Man könnte es wohl dulden, wenn er wollte aus Gnaden der Oberste sein I, 1419.; seine Satzungen könnte man dulden, wenn er sie nicht für nöthig erklärte I, 1419.; unter welchen Umständen man ihm auch den Fuß küssen könnte I, 766.; soll uns aus dem Wort beweisen, daß er Christi Statthalter ist, nicht durch Wunder I, 1546., wie er versucht I, 1545. Seine Decrete gehören auf das heimliche Gemach II, 2010.; man soll ihm zürnen und fluchen II, 1898. 2026. 2077. 2079. — S. \gesperrt{Simonie}.

2. \gesperrt{Pabst und Türke.} P. noch ärger als der Türke II, 53. 1898. 2062. f.; sie haben lange Gotteslästerung getrieben, Gott schien zu schlafen I, 1195. f. — S. \gesperrt{Antichrist}.

Pabstthum. \gesperrt{Entstand}, weil man durch Verachtung der Gnadenmittel in Irrthum gerieth I, 1046. 1595. f., durch den Ehrgeiz der Bischöfe und Trägheit der Obrigkeit I, 1461. f. Rühmt sich, es sei die \gesperrt{Kirche} I, 458. f. 507. 790.; II, 99.; verachtet die rechte Kirche I, 649., ist der evangelischen Kirche bitter feind II, 42. f.; verläßt sich auf Christi Verheißung von der Kirche I, 510., auf die Succession I, 1406., auf die große Menge, Schein und Personen I, 495. f. 554. 649.; II, 72. f.; ist nicht die rechte Kirche I, 643. 1402. f. 1406.; II, 39. 119. f. 437. f., eine Kainische Kirche I, 311. 318. ff., des Teufels Kirche I, 554., schlimmer als Sodom I, 1276., eine Versammlung der Teufel und allerschändlichsten Menschen II, 105., eine Ungeheuerlichkeit II, 1897., eine greuliche Verwüstung der Kirche II, 1890.; wie das Pabstthum im Bilde der Arche dargestellt wird I, 612. \gesperrt{Blindheit und Finsterniß} im P.: II, 1345.; da ist ein blindes Auge und ein hörendes Ohr I, 1633.; das ganze P. hat die Weisheit nicht, das Wort recht zu theilen I, 1211.; viele Lehrer wußten die 10 Gebote und Bitten des Vater Unsers nicht I, 1063.; der rechte Verstand der 10 Gebote war verloren II, 1895.; die ganze \gesperrt{Lehre} des P.’s ist die Frage warum? I, 1127., ist leichtfertig, ungewiß II, 1930.; hat nur rhetorische, keine dialektische Beweise für seine Lehre I, 1175. f.; lehrt vom Evangelio nichts Rechtes I, 219. 1386., sucht das Evangelium zu unterdrücken I, 459.; die geistliche Verheißung des N.\,T.’s ward verdunkelt und unterdrückt II, 1889. f. 1903. f.; die Lehre von Vergebung der Sünden war ganz aus den Augen gekommen I, 220.; man fürchtete sich mehr vor Christo denn dem Teufel II, 1591. 1990.; in Todesnöthen wies man zu Maria I, 220.; da war die Auslegung nicht Eines Psalms zu hören II, 1364.; Ehemänner beichteten als eine Sünde, daß sie mit ihren Weibern gescherzt II, 164. Der \gesperrt{Gottesdienst} im P. ist selbsterwählt II, 777., voll Abgötterei I, 778.; II, 913. 922. f. 1919.; da ist eitel Narrenwerk I, 864. f.; da zeigte man die Milch der heiligen Jungfrau II, 907. Da ist nichts denn \gesperrt{Menschensatzungen} I, 367.; II, 1977., die man an die Stelle der Verheißung setzte II, 1890. f. 1895.; da herrscht Unfug mit Wallfahrten, Messen und \gesperrt{andere Greuel} II, 1550. f.; das P. kann ohne Simonie nicht bestehen II, 104. f. 107. f.; da herrschte Hurerei und Unzucht I, 164.; II, 1264.; die Leute wurden jährlich einmal zur Beichte gezwungen II, 316.; da war kein Ende mit Geben für Klöster etc. II, 1235. Die \gesperrt{Seelen} sind im P. unstät und flüchtig I, 367., gedrückt und geplagt I, 219. 1428. ff. Noch Manche sind im P. \gesperrt{selig} geworden I, 220. 368.; II, 1334. Es hat die \gesperrt{Gnadenmittel} noch und ihre Verwaltung ist kräftig I, 962. 1405. \gesperrt{Man soll} das P. hassen I, 1417.; II, 108., davon ausgehen I, 1275.; es sollte zerstört werden I, 220. — S. \gesperrt{Papisten, Simonie}.

Palästina. Jede Stadt hatte einen König I, 1306.

Pallium. Preis desselben I, 1620. f.

Pannonien. Zweierlei P. II, 2033.

Papisten. \gesperrt{Sind} wie Noahs Rabe I, 621. f., wie Ham I, 632., wie die Hamiten I, 698. Maßen sich den Namen an: der \gesperrt{Kirche} und sorgen doch nur für den Bauch II, 674., des Volkes Gottes II, 618.; rühmen sich der fleischlichen Geburt II, 38. ff.; sehen nicht auf das Wort, sondern auf das äußerliche Werk I, 528. f., auf absonderliche Werke I, 1546.; haben nicht den rechten Gott I, 1551., keine Erkenntniß in geistlichen Dingen I, 1553.; verlassen sich auf ihr eigen Verdienst, weil sie den Artikel von der Schöpfung nicht verstehen I, 1441. f.; ihr Leben ist voll Geiz, Ehrsucht, Verachtung Gottes II, 459. \gesperrt{Verhalten gegen die reine Lehre:} hatten eine affectirte Unwissenheit der Lehre Luthers und verfolgten sie I, 817., verkehrten und schmähten sie I, 1151., verwarfen sie muthwillig I, 1277. f. 1646., verfolgten die erkannte
Wahrheit II, 1806. f., verstockten sich II, 1111. 1265.; vertheidigten falschen Gottesdienst wider ihr Gewissen I, 736.; rühmten ihre Canones und verachteten darüber Glaube, Liebe, Hoffnung I, 1310.; plagten Luther mit dem Argument: Meinst du, daß die Väter alle geirrt? I, 498. f. \gesperrt{Verhalten gegen die Evangelischen:} spürten den Sünden derselben nach und offenbarten sie I, 632.; wütheten gegen die Kirche I, 700. f.; verfolgten fromme und gelehrte Prediger I, 1010.; verachteten die Evangelischen als trunkene Narren II, 1995. f. \gesperrt{Ihre Strafe:} um ihretwillen ist Sodoms Untergang beschrieben I, 1274. f., aber sie fragen darnach nicht I, 1275. 1288.; ziehen die Verdammniß über sich I, 1277. f.; es wird sie ein Aergerniß treffen, dem nicht zu rathen ist I, 634.; Gott wird sie verstoßen I, 553., hat sie verworfen I, 1102.; II, 1751.; kein Papist wird deshalb selig, weil er Papist wäre II, 67.; wie sie \gesperrt{Gliedmaßen der Kirche} werden können I, 1433. Ihnen mußte Gesetz, nicht Evangelium gepredigt werden I, 1187. f.

Paradies. 1. \gesperrt{Im Allgemeinen.} I, 106. ff. 124. ff. P. heißt Garten I, 106. Größe und Lage des Paradieses I, 107. 109. 678. 851.; Meinung Etlicher von der Mitte des P. I, 378.; hatte eine Pforte nach Sonnenaufgang zu I, 282.; schöner als die übrige Erde I, 109. 111. f. 276.; darin steter Frühling I, 177.; ob wilde Thiere darin I, 522. f.; sollte ein gemeinschaftlicher Tempel der ganzen Welt sein I, 282. Das Paradies und die Hut durch die Engel blieb bis auf die Sintfluth I, 109. 282. f., durch welche es zerstört I, 107. 120. 378. Allegorische Erklärung des Paradieses I, 109. ff. — S. \gesperrt{Baum} 1, 2, \gesperrt{Eden, Schlange}.

2. \gesperrt{Das Wort „Paradies“} bedeutet auch das Wesen, darin Adam vor dem Fall gewesen I, 108. — S. \gesperrt{Stand der Unschuld}.

3. \gesperrt{Von dem} P., dahin der Schächer kam I, 1763.

Paris. Die Theologen zu P. verstanden die Schrift nicht II, 2030.; hetzten die Fürsten gegen die Evangelischen II, 2001.

Paronomasie. Dazu haben die Propheten Lust I, 727.; Beispiele: II, 166. 2011. 2021. 2030.

Pathrusim. Wer sie waren I, 676. f.

Patriarchen. Diesen Namen verdient Abraham vor Anderen I, 1054. f. — S. \gesperrt{Väter}.

Paulus, der Apostel. Er, nicht Einer aus den siebenzig Jüngern, ward aus Gnaden zum Apostel erwählt I, 732.; hatte etwas Eigenes I, 836.; war erst selbstgerecht I, 1434. f.; hatte einen Pfahl im Fleisch I, 823.; II, 1748.; ward angefochten II, 744. f.; hatte mit dem Fleisch zu kämpfen II, 288. 442., wider den Unglauben I, 1555.; schöpfte aus Mose I, 945. 952. f. 1559. 1571.; ist der beste Ausleger Moses I, 1029., der höchste Lehrer der heil. Schrift I, 953.; handelt fleißig von der Beschneidung und dem ganzen Gesetz I, 1088., von Gnade und Verdienst I, 1435.; an ihm erfüllt sich der Segen Jakobs über Benjamin II, 2053. f.

Peleg. I, 693.

Peripatetiker. Ihre Meinung von irdischen Gütern I, 854.

Periphrasis. Redefigur, ein Beispiel II, 1401.

Perser. Bei ihnen finden sich in dem heil. Feuer Spuren der Nimrodischen Ketzerei I, 727. — S. \gesperrt{Elam}.

Person. Gott sieht sie nicht an I, 553. f. 1434. 1441. 1446. f.

Pestilenz. Man darf davor fliehen I, 1480. f.; wie man sich darin trösten soll I, 1500.

Petrus. 1. \gesperrt{Der Apostel.} Ward auf dem Meer versucht, damit er nicht stolz würde II, 443.; wie ihn die Strafe der Sintfluth bewegt hat I, 535.; die Verheißungen Christi gelten von ihm nicht mehr, denn von den andern Aposteln I, 994.

2. P. \gesperrt{Lombardus.} Sucht in den Creaturen Fußstapfen der Dreieinigkeit I, 1151.; seine Meinung von den unbeschnittenen Kindern, die vor dem 8. Tage starben I, 1083.; geht kalt mit den 10 Geboten um II, 1896.

Pferd. Kriegerisch II, 1985. f.

Pharan. Wüste, stößt an Juda I, 880. 1455.

Pharao. Ist Königs-, nicht Personenname II, 185. — P., \gesperrt{Josephs Zeitgenosse:} ein Heide, wußte von dem Gott der Hebräer nichts II, 1330.; durch Joseph zur rechten Erkenntniß gebracht I, 1031.; II, 1366.; Joseph forderte von ihm die Beschneidung nicht II, 1367.; voll Glaubens und Heiligen Geistes II, 1733.; heiliger König II, 1737. f.; demüthig II, 1342. 1369. f.; freundlich II, 1728.; fromm und gottselig II, 1294. f. 1417. 1436. 1726. 1729. 1797.; weise II, 1294. f. 1375.; ernährte die Priester II, 1816. ff.; ward ein Pfleger der Kirche II, 1730. f. 1733. f.

Philadelphia. Heißt Bruderliebe I, 1142.

Philister. S. \gesperrt{Abraham} 32, \gesperrt{Isaak} 14.

Philosophen. Suchten großen Namen, verachteten Reichthum I, 829.; ihre Disputationen von Gott sind eitel Unwissenheit und Gotteslästerung I, 484. f.; die Meinung mancher P. vom Reichthum I, 830.

Philosophie. Ihre Lehre von der Endursache der Creaturen I, 607.; erkennt das natürliche Verderben nicht I, 798.; lehrt, was man thun soll, aber nicht, woher die Kraft nehmen I, 621.; kann das Gewissen nicht zufrieden stellen, gibt keine Hoffnung des ewigen Lebens I, 621.; Lehre von der Unsterblichkeit I, 406.; weiß von wahrer Freundschaft nichts I, 798.; Lehre von Geduld II, 1666., Feindesliebe I, 801.; lehrt, wer böse gewesen, sei allezeit für böse zu halten I, 801. — S. \gesperrt{Mensch} 2, \gesperrt{Vernunft}.

Phönicien. Brauch vor der Hochzeit I, 455.

Phrath. Ist der Euphrat I, 119.

Pighius. Lästerte Luther II, 1274.

Pison. Ist der Ganges I, 119. 123.

Plato. Hat wohl Fünklein der Wahrheit aus der Propheten Predigten gesammelt I, 6.; was er über Gott wußte II, 1674.; für welche drei Dinge er Gott gedankt I, 312. f.; Meinung: von den Himmelskörpern I, 57., von Materie und Idee I, 6., von der Schöpfung der Welt I, 6.

Plinius. Nennt die Augen den König der Glieder II, 492.; hält die Natur für eine Stiefmutter I, 1344.; verachtet die Unsterblichkeit II, 15.

Pniel. Bedeutung des Namens II, 799.; Lage II, 846.

Polizei. Geht damit um, der Sünde zu wehren I, 127.; vor dem Fall nicht nöthig I, 127.

Polycrates. Soll im Schlaf das Netz gezogen haben II, 1238.

Polygamie. Zweck ist: Mehrung des Volks I, 386.; von Lamech eingeführt I, 391.; zu Abrahams Zeit Sitte II, 239., auch zu Jakobs II, 532.; im A.\,T. gestattet, ja, geboten I, 974.; im Gesetz der Leviratsehe eingeschlossen II, 1179.; im A.\,T. gestattet, jetzt nicht I, 972. f.; läßt sich nicht durch Abrahams Beispiel rechtfertigen I, 974. 976.; von Jakobs Polygamie II, 510. ff. 528. f. 533. 539. ff.

Pommer, Dr. Fromm und gottselig II, 773.

Potiphar. II, 1222. ff. Bedeutung des Namens II, 1220.; nicht mit Potiphera zu verwechseln II, 1222. 1414. f.; was er war II, 1222. f. 1293.

\quellseite{567}
Potiphera. Bedeutung des Namens II, 1222. — S. \gesperrt{Potiphar}.

Prädestination. S. \gesperrt{Versehung}.

Prahler. Ihre Worte größer als die Thaten I, 1156.

Praesumptio pietatis. Was das sei I, 527.

Predigen. „Von dem Namen des HErrn p.“ I, 781. f.; wie Melchisedek das gethan I, 907.

Prediger. 1. \gesperrt{Ihr Amt.} Sie sind Gottes Mitarbeiter, Gott wirkt durch sie II, 1775., Christus handelt durch sie II, 1897.; Augen und Zähne II, 2006. f., Kriegsleute II, 408., Hader- und Gerichtsmänner I, 457.; ist ein gut Zeichen, wenn die Welt sie Hadermänner nennt I, 458.; ihr Amt ist „richten“ I, 457., lehren und beten I, 456., lehren, trösten II, 1722. f.; sollen auch von zeitlichen Dingen lehren, wie man sich darin verhalten soll II, 170. f.; sollen trösten und strafen II, 2007.; sollen die Sünde strafen I, 1259., das ist dann kein Haß II, 1141. f.; haben Macht den Himmel aufzuthun II, 1894. f.; sind größere Propheten denn Jakob II, 2057.; müssen über der reinen Lehre wachen II, 1424.; sollen Constantine, ein Fels in der Lehre sein, nicht ein Rohr II, 1930. f.; fromme, reine Lehrer, eine Gabe Gottes I, 1631. f.; II, 1364.; sind nöthig, Alle können nicht des Predigtamts warten II, 2006. f., sehr nützlich II, 231.; sind wenig gottesfürchtige Lehrer, dann auch wenig fromme Leute I, 534.; ihr Amt ist sorgen- und mühevoll I, 260. f.; essen ihr Brod im Schweiß ihres Angesichts I, 258. f.; haben nicht viel gute Tage I, 1339. f. 1341. f.; gottlose P. haben gemeiniglich mehr Glück, denn fromme I, 1632. — S. \gesperrt{Propheten} 1.

2. \gesperrt{Amtsführung.} Sie sollen ihr Amt führen, wenn sie auch selbst Sünder sind I, 636. f., ihres Berufs warten, den Erfolg Gott befehlen II, 748. 1619.; können nicht alles Verkehrte abschaffen II, 1637. f.; genug, wenn sie den geringsten Theil bekehren II, 1773. 1824.; sollen Tauben und auch Schlangen gleichen I, 624. f., erst beten und studiren, ehe sie lehren II, 395., Gott um Weisheit bitten II, 277. f., die Schrift studiren I, 1462.; ihnen soll die Rede aus dem Herzen, nicht aus dem Munde wachsen I, 1000.; sollen das Wort theilen und die Zuhörer unterscheiden I, 1208. f., geschickt sein, Gesetz und Evangelium zu unterscheiden I, 1430.; ein junger P. füllt die Hölle mit Seelen I, 847.; sollen nicht einzelne Werke der Heiligen zum Exempel, sondern Regeln vortragen I, 1233.; dürfen wohl die Sünde etwas härter strafen als Andere II, 232.; die reich werden wollen, dürfen die Sünde nicht strafen II, 633.; wie sie Angefochtenen begegnen sollen II, 784. 824.; sollen sich hüten, daß sie nicht durch ihren und der Ihren Wandel in Schande fallen II, 1030.; wie sie sich der Simonie schuldig machen können II, 111. f.; sollen ihren guten Namen vertheidigen II, 1304.; ob sie Gehalt nehmen dürfen I, 1623. f.; II, 1236. und vom Ersparten etwas kaufen I, 1624. f.; werden im Ueberfluß faul und geizig II, 1735. f.; wie sie sich trösten sollen, wenn ihre Predigt wenig Frucht schafft II, 1026. f., wenn sie wenig gute Tage haben I, 1339. f. 1341. f. 1358., wenn sie verachtet werden II, 634., wenns ihnen kümmerlich geht II, 1956. f. — S. \gesperrt{Anfechtung} 9, \gesperrt{Ehebrecher, Gesetz} 2, \gesperrt{Gesetzprediger, Sichere, Verbrecher, Verfolgte}.

3. \gesperrt{Verhalten gegen sie.} Es geht ihnen noch heute, wie ehedem Noah I, 497.; wenn sie die Sünde strafen, kommt Gewalt und Tyrannei über sie I, 470.; sie gelten für Betrüber des Volks I, 458.; werden verachtet II, 2007. f., geringer geachtet, denn Sauhirten und Büttel II, 634.; auf Verachtung derselben folgt Strafe I, 1632.; müssen darben, man gönnt ihnen nichts II, 632. ff. 637. Sie sind ihres Brodes werth I, 260. f.; man soll sie und die Ihren nähren I, 258. f.; II, 450. ff. 1815. ff. 1953.; obs genug ist, wenn sie Rock und Kropf haben I, 1461.

Vgl. \gesperrt{Lehrer}.

Predigtamt. 1. \gesperrt{Das Amt.} Kirchenamt und weltliches Regiment soll man unterscheiden I, 1462.; ist von Gott geordnet I, 1594., damit Gott mit den Menschen rede I, 532.; darin handelt Gott durch Menschen mit uns I, 869. 1448. f. 1456.; das Wort, das wir hören, ist Gottes Wort I, 1736. f.; II, 905.; wie tröstlich das ist II, 906.; der Prediger Worte und Handlungen sind Gottes II, 433. f. 949.; es soll strafen und trösten I, 628. 1208.; wo das öffentliche Predigtamt aufhört, kann weder Glaube, noch Gebet, noch rechtschaffener Gebrauch der Sacramente bestehen I, 534.; ist die höchste Wohlthat I, 456. f. 736.; darin wird Segen ertheilt, nicht nur gewünscht II, 298. ff.; darin wird das Blut Christi ausgesprengt II, 2006.; des Predigers Gottlosigkeit nimmt dem P. nicht die Kraft II, 119. f.; wird in der letzten Zeit fast ganz abkommen I, 534.; ist durch Noahs Tauben abgebildet I, 623.

2. \gesperrt{Verhalten dagegen.} Welt verachtet es I, 333.; II, 300. f. 1367. f. 1550., so ist’s in allen Ständen I, 467.; Pabst und Papisten heben es auf I, 533. f. 1593. f.; Gott hat es wider Pabst und Türken erhalten I, 963.; wo es ist, sollen wir auf keine bessere Offenbarung warten I, 532.; durch dasselbe, nicht durch neue Offenbarungen lehrt d. Heil. Geist I, 624.; man soll es hochrühmen II, 2059. f., in Ehren halten I, 1448. f. 1456.; so thaten die Patriarchen I, 1199. 1596.; dasselbe erhalten ist ein herrlich Werk II, 449. ff.; Gott straft, wenn man es nicht erhalten will II, 1820. — S. \gesperrt{Verachtung} 2.

Vgl. \gesperrt{Beruf} III, \gesperrt{Wort} 3, \gesperrt{Stände}.

Predigten. Worauf alle gerichtet sein sollten II, 222. f.; auf jede gute P. soll auch ein gut Vater Unser folgen II, 370.; erregen Haß und Feindschaft I, 501. — S. \gesperrt{Gebet} 13.

Priapus. Welche Völker diesem Götzen dienten I, 455. 1663. 1665.

Priester. Wer ein rechter P. ist I, 1521. f.; Abraham war ein P., aber nicht gesalbt und geschmiert I, 782. 1521.; ihr Amt ist lehren, anrufen, opfern I, 456. 837. f.; II, 370.; Kleidung im A.\,T. II, 1039.; manche heidnische P. waren verschnitten II, 1222.; hatten bei den Heiden großes Ansehen II, 1821. — S. \gesperrt{Segnen}.

Priesterthum. Das Levitische war nur Schatten und Vorbild I, 1590., sollte aufhören I, 1589. f., konnte den rechten Segen nicht geben I, 1591.; das P. ist den Eltern übergeben I, 302. f.

Propheten. 1. \gesperrt{Im Allgemeinen.} Ein P. ist, wer die Schrift versteht und Andere lehren kann I, 1366., wer das Sacrament austheilt oder empfängt II, 2060., jeder Christ I, 1367.; Noah ist der höchste P. I, 464. f.

2. \gesperrt{Die} P. \gesperrt{des} A.\,T.’s. Sind Seher I, 1349.; hatten nicht ohne Aufhören Erleuchtungen II, 1645.; viele sind nicht nach Jerusalem gekommen, daselbst zu lehren I, 1530.; haben Mose, die Schrift fleißig betrachtet I, 97. 106. 1055. 1559. 1584., durch sonderliche Erleuchtung viel aus Vergleichung der Worte und Dinge herausgelesen II, 1851. f., aus Mose geschöpft I, 749. 755. 806. f. 871. 961. 964. f. 1115. f. 1559. 1569. 1571. 1584.; II, 74. 430., Mose erklärt I, 1067.; wie sie Abrahams
Pilgrimschaft angesehen I, 1305.; eiferten wider den Götzendienst II, 959.; erkannten, daß man nicht durch Werke gerecht wird I, 1591. f.; klagen über ihre Schwachheit II, 2065.

3. P. \gesperrt{des Teufels.} Gibt es auch II, 1049. 1051.; die falschen P. heißen leichtfertig und unbeständig II, 1930.

4. \gesperrt{Der große Prophet.} S. \gesperrt{Christus} 3.

Prophezeiungen. 1. \gesperrt{Im Allgemeinen.} P., Drohungen und Verheißungen sind geistlich zu verstehen I, 648.; im A.\,T. sind herrliche P. auf Christum II, 2056. f.; die P. von Christo und seinem Reich werden nicht leichtlich verstanden II, 1347.; P. müssen geglaubt werden, scheinbar geschieht das Gegentheil I, 642. f.; warum die P., seien es Drohungen oder Verheißungen, nicht geglaubt werden I, 642.; ihre Erfüllung verzieht sich I, 643.; Erfüllung verschiedener P. I, 655. 994. 1588.; II, 296. f. 302. ff. 317. f. 323. f.; Noahs P. übertreffen alle andern I, 464. f. — S. \gesperrt{Drohungen, Verheißungen}.

2. \gesperrt{Luthers} P. \gesperrt{Deutschland} steht ein großes Unglück bevor I, 1255.; nach L.’s Tode werden die allergefährlichsten Zeiten folgen I, 457. f.; wodurch Gott gedrungen werde, Deutschland zu strafen I, 511. 641.; Gott wird die Verachtung des Worts an Deutschland nach L.’s Tode strafen I, 409. f. 864. 882. 896.; II, 881. f. 887. 1806. f.; Deutschland wird gestraft werden, weil Laster und Irrlehre zunehmen I, 253., weil man sündigt und sich nicht einreden lassen will II, 1043., weil Zucht und Ehrbarkeit gefallen ist I, 641., wegen des Geizes II, 586. f., wegen der Schinderei und des Wuchers II, 1806., wegen kärglicher Erhaltung des Predigtamts II, 1817. f.; Gott wird Deutschland durch allerlei Irrlehrer sein Wort entziehen I, 455. f., man wird von der reinen Lehre von der Rechtfertigung weichen I, 1668. 1670., die Nachkommen werden abfallen I, 1455.; II, 1827.; Luther wird in Frieden sterben, ehe der Jammer über Deutschland angeht I, 1758. An den \gesperrt{Ketzern} und Schwärmern wird man Wunder sehen II, 1827.; \gesperrt{Pabst und Türke} werden in Kürze ihre verdiente Strafe bekommen II, 55. 220.

Proportion. Arithmetische und geometrische I, 847. 1074.; II, 511. f.

Proselyten. Wen man so nannte II, 1706. — S. \gesperrt{Beschneidung} 2.

Psalter. S. \gesperrt{David, Korah}.

Ptolemäus. Seinen Leichnam sah Augustus II, 2073.

Put. Völkername I, 665.

Puteole. Ueberschwemmung daselbst I, 1277.

Pylades und Orestes. Ihre Freundschaft ist nichts gegen die Brüderschaft in der Kirche I, 1147.

Pyrrhus. Seinem Vater nicht gleich II, 1368.

Pythagoras. Meinung von Harmonie der Sphären I, 153.

\section{Q}

Quinctilian. Warnt vor mehrdeutigen Wörtern II, 1908.

\section{R}

Rabbinen. 1. \gesperrt{Im Allgemeinen.} Von ihnen darf man in der Schriftauslegung nichts Gutes und Reines erwarten I, 410. 812.; sie glauben nicht, können daher auch die Schrift nicht recht auslegen II, 2030.; verstehen die Bibel nicht II, 869.; sind Grammatiker, keine Theologen I, 362.; meinen, man könne die Schrift nur durch die Grammatik verstehen I, 451. f.; verstehen nichts von den Sachen, kleben an den Wörtern, geben ungeschickte Deutungen I, 322. 328. f. 998. ff. 1002. 1004. f.; II, 45.; verstehen die Sachen nicht, darum auch die Worte nicht I, 1381.; verstehen Vieles im Hebräischen nicht II, 677. 961. f. 1610.; können nichts Gutes lehren, haben die Wahrheit verloren I, 1359.; verachten das Licht des N.\,T.’s I, 369. f. 1004. f.; würden ohne das N.\,T. die hebräische Sprache nicht verstehen II, 1837.; nahmen heimlich viel aus L.’s Schriften II, 1834.; verdunkeln die Kraft der Worte I, 1698. f.; mengen deren Bedeutungen seltsam unter einander II, 677.; sind beflissen auf mancherlei Bedeutungen der Wörter II, 1908. f.; verstehen die Figuren nicht II, 1909.; kreuzigen mit ihrer Auslegung Mose I, 450. f.; fälschen fast überall den Sinn der Schrift I, 361. f. 429.; II, 789. 849., auch wissentlich II, 1966. ff.; ziehen das Geistliche auf das Leibliche I, 429.; verkehren die Stellen vom Sohne Gottes II, 853. f.; haben ungeschickte harte Deutungen II, 801. 1929.; haben aus der Schrift ein heimlich Gemach ihrer Unflätherei gemacht I, 1292. f. 1297. f.; II, 44. f.; rühmen gern ihre Vorfahren I, 813.; sind blutdürstig gegen die Heiden II, 1973. 1976. f. — Die Väter folgen gemeiniglich den Rabbinen I, 364. 452. 1005.; man soll es nicht II, 704.; sie hindern das Studium der Schrift II, 1927.; man lese sie mit Bedacht I, 361. f. 370. 381. f.; warum Luther sie zuweilen anführt I, 370.; sie machen ihm viel Arbeit I, 370.; er ist ihnen feind I, 362.

2. \gesperrt{Beispiele ihrer Schriftauslegung,} betreffend: 1 Mos. 38, 23. II, 1207.; 1 Mos. 49, 10. II, 1968. ff.; 1 Mos. 49, 11. 12. II, 1996.; 1 Mos. 49, 21. II, 2034. f.; 1 Mos. 49, 26. II, 2050. ff.; Abimelech I, 1331. 1361. f.; Abraham I, 722. 733. 805. 828. f. 1198. f., den Boten, der ihm Lots Gefangenschaft meldete I, 884.; Adam I, 85. 397. 410. 413.; Ahalibama II, 995. f.; Arche I, 515. f.; Beschneidung I, 1040. 1077.; Buhlteufel I, 446. f.; Canaan, Hams Sohn I, 650. f.; Caphthorim I, 676. f.; Casluhim I, 676. f.; Egypter I, 812.; unfruchtbaren Ehestand I, 974.; Hagar I, 995.; Haran und Abraham I, 722.; Henoch I, 425. f.; Isaak II, 44. f. 157. ff. 310., seine Opferung I, 1539.; Jakob II, 389. 466. 780.; Josephs Dolmetscher II, 1499. f.; Josephs Brüder vor Pharao II, 1784. f. 1789.; Juda II, 1160. f.; Kain I, 307. 370. 387. f. 389. 397.; Lamech, Kains Nachkommen I, 389. f. 397.; Lots Weib I, 1280.; Noah I, 433. 563. 564. f.; Pathrusim I, 676. f.; Pharaos Bäcker und Schenken II, 1294.; Philistim I, 677.; Potiphar II, 1222. 1414.; Rubens Blutschande II, 1932.; Sarah I, 802. 813. 1381. 1602. f.; Schorbor II, 1948., Sem I, 434.; Seth I, 397.; Teraphim II, 649.

Rabe. Noahs R., ob er wiedergekommen I, 563.; allegorische Deutung I, 618. ff.

Racha. Was R. heißt II, 1534.

Rache. Gott nimmt R. an den Verfolgern der Kirche I, 354. — S. \gesperrt{Unrecht}.

Rächen. Man räche sich nicht II, 874. 877. 1151.

Raema, Enkel Hams I, 666.

Raemses. Lage II, 1803.

Rahab. Ob Hure oder Wirthin II, 1205. f.; erlangte durch den Glauben Vergebung II, 1830.

Rahel. Worin ihre \gesperrt{Schönheit} bestand II, 492. f.; war fleißig II, 1239., fromm und gottselig II, 965., keusch und fromm, suchte nicht fleischliche Wohllust, sondern Kinder um des verheißenen Samens willen II, 536. ff. 539. ff. 544. 556. 561. 563. f. 570. 573., und freute sich aus Liebe zu der Frucht des Leibes über die Geburt
\quellseite{568}
Dans II, 545., Naphthalis II, 546. f.; war eine gottselige Hausmutter II, 540., freundlich gegen die Wittwen der Sichemiter II, 935., eine liebreiche Herrin II, 964. Ihre \gesperrt{Unfruchtbarkeit} war ihr ein schweres Kreuz II, 573. ff.; Leas Fruchtbarkeit erweckte bei ihr Haß und Neid II, 523. f. 528. ff., den sie aber bekämpfte II, 536.; ward stolz über Josephs Geburt II, 555.; war sehr arm II, 541. \gesperrt{Stahl} die Götzen Labans, was sie dazu bewogen II, 639. 643. f., ob sie daran unrecht gethan II, 639. ff. 648.; täuscht Laban durch eine Nothlüge II, 687. f. \gesperrt{Starb} unerwartet II, 965., vor Josephs Verkauf II, 1102., im Glauben II, 968. ff.; ihr Tod erregte große Traurigkeit II, 964.; sogleich nach demselben entstand Unordnung im Hause II, 973. — S. \gesperrt{Jakob} 18, \gesperrt{Laban} 3, \gesperrt{Lea} 2.

Raphael. Ableitung des Namens I, 1013.; Bedeutung desselben I, 1198.; welches Engels Name I, 1198.

Rath. Wird meist ausgeschlagen I, 978. \gesperrt{Gottes} R. und Wege II, 1689. ff. 1702. ff.; ist ein anderer als der unsere II, 1689.; unbegreiflich, man soll sich nicht daran stoßen II, 1689. ff.

Rathhaus. Im Morgenlande am Stadtthor I, 1213.

Rathschläge. Die allerbesten gerathen oft am übelsten II, 1479. f.; die verächtlichen führt Gott hinaus II, 1480.

Raub von Weibern und Jungfrauen erregt Krieg, Mord und Jammer II, 496. 863.; ist Sünde II, 873.

Rauh sein. Bedeutet Grausamkeit und Heuchelei II, 76.

Rebekka. I, 1706. ff. 1722. ff. \gesperrt{Gaben:} war schön von Angesicht I, 1707. f.; II, 156., mit schönen Gaben geziert I, 1707. f., von großem Geist und starkem Muth II, 262., listig und verständig II, 339. 348. f. 466. \gesperrt{Glaube:} hatte einen starken Glauben II, 283. ff., aus dem Wort und Erleuchtung des Heil. Geistes II, 285. 289. f.; wie sich derselbe offenbart II, 26. f. 269. 271. f. 289. 342. \gesperrt{Tugenden:} hatte viele Tugenden I, 1706. ff.; andere Jungfrauen waren mit R. nicht zu vergleichen I, 1708.; R. war behende I, 1710. 1719., dienstfertig I, 1709. f., ehrerbietig I, 1708. f., gastfrei I, 1708. 1717., geduldig II, 26. f., gottselig und gottesfürchtig II, 339., züchtig I, 1743., keine böse Schwiegermutter II, 243. R. hatte \gesperrt{Anfechtungen} II, 45. ff. 157. f. 196. 219.; wie sie sich dabei verhalten II, 52. f.; frug wohl Sem, da sie den HErrn um Rath frug II, 55. Durch \gesperrt{Wegsenden} Jakobs lehrte sie, man solle Gott nicht versuchen II, 346. f. \gesperrt{Wann sie gestorben} II, 348. 947. 1102. — S. \gesperrt{Eber, Isaak} 15, 16.

Recapitulatio. Redefigur II, 548. f. 572.; muß in der Schrift beobachtet werden II, 1510. f. — Siehe \gesperrt{Hysteron}.

Rechte. Um des Friedens willen gebe man daran nach I, 844. f.; so that Abraham I, 848.; das höchste Recht, das höchste Unrecht I, 1235. „R. des HErrn“ II, 144. „Weltliches \gesperrt{Recht}“ I, 814.; im weltlichen R. gilt die Regel: Wer einmal eine Uebelthat begangen, zu dem ist sich zu vermuthen, daß er allezeit übel thut I, 796. — S. \gesperrt{Gesetz} 4.

Rechtfertigung. \gesperrt{Wesen:} Die R. der Gläubigen ist wohl zu unterscheiden von ihrer Verklärung I, 1579. ff. 1587. f.; was die R. sei I, 1672.; kommt aus Gnaden durch den Glauben, nicht aus dem Gesetz durch Werke I, 942. ff. 949.; II, 446. f.; die Lehre von der R. durch Werke ist türkisch I, 734.; die Papisten verwerfen die Lehre von der R. durch den Glauben I, 1441. ff., nennen die Lutheraner Solarii I, 1573.; machten eine Schwenkung, lehrten: Glaube und Werke rechtfertigen I, 1573. ff. \gesperrt{Wirkungen:} die R. deckt die Sünde zu I, 1672.; Sünde bleibt noch in den Gerechtfertigten I, 1670. f.; verbietet gute Werke nicht I, 1151. \gesperrt{Wichtigkeit:} der Artikel von der R. durch den Glauben ist der vornehmste, Welt und Teufel sind ihm feind I, 941. 943.; II, 446.; durch denselben besteht die Kirche, behält Gott seine Ehre I, 1441.; fällt man von dieser Hauptlehre, so folgt Simonie II, 110. \gesperrt{Schriftgrund:} Luther hat diese Lehre nicht erdacht I, 949. 1573., sie liegt in der Verheißung von Abrahams Samen I, 1583. — S. \gesperrt{Gerechtigkeit, Glaube, Vergebung, Werke}.

Rechtspflege. Etliche Völker hingen den Dieb erst, dann frugen sie nach dem Recht II, 889.

Redemtio vexationis. II, 112. ff.

Redensart. „So Gott will“ I, 1168.

Redeweisen. Jede Kunst hat ihre besonderen R., so auch der Heil. Geist, daran muß man sich gewöhnen I, 57.; man soll sie nicht meistern I, 58. — S. \gesperrt{Kunstausdrücke}.

Redner. Wer ein R. ist I, 1005.

Reformation. Alle R. muß von der Lehre anfangen II, 910. f. 914. f., dann folgt die Reformation des Lebens und dann der Ceremonien II, 914. f.; die Papisten fangen bei den Ceremonien an II, 910. f. 915. R. \gesperrt{Luthers:} war eine große Wohlthat I, 1337., anfänglich schwach II, 1329.; von den Anschlägen der Feinde wider dieselbe und wie Gott sie verhindert II, 659.

Regenbogen. Es gibt ganz runde R. I, 606.; R. gäbe es ohne den Sündenfall nicht I, 94.; gabs nicht vor der Sintfluth I, 607. ff.; von den natürlichen Ursachen des R.’s I, 605. ff.; zeigen Gottes Zorn und Gnade an I, 608. f.; predigen Gottesfurcht und Glauben I, 607.; nur die Theologie, nicht die Philosophie kennt die causa finalis des R.’s I, 607.; bedeutet, daß keine allgemeine Fluth mehr kommen soll I, 605. f.; Bedeutung der Farben I, 437. 512. 609.

Regent. An ihm soll sein Ansehen und Gewalt und Ehrerbietung II, 1422. f.; muß beherzt sein II, 1412.; soll fromm sein in allen Geboten I, 1035., fleißig auf sein Amt sehen II, 1423.; die mit Recht und Gesetz regieren, heißen Hirten und Fürsten I, 475.; sollen in Wahrheit sein Gestrenge, Edle etc. II, 1792.

Regieramt. Eine Strafe der Erbsünde II, 1409. ff.; sehr beschwerlich II, 276. 1405. ff. 1791.; wer sich davor hüten kann, mag es wohl thun II, 1790.; wer dazu berufen ist, soll es verwalten II, 1790. f.; womit man sich dabei trösten soll II, 1411.; erfordert ein gottseliges, ritterliches Herz II, 1678.; will nicht faule, verzagte Leute haben II, 1790. f., auch nicht unbeständige, leichtfertige II, 1930.; darin ist Gebet nöthig II, 276. f.; geht nicht ohne Sünde ab I, 822. — S. \gesperrt{Obrigkeit, Regiment}.

Regierung, göttliche. Den Artikel von der R. Gottes soll man fleißig treiben II, 932. Gott regiert \gesperrt{Alles} II, 930. f., wie ein Hausvater I, 909. f.; regiert die Königreiche I, 741. f. 850., auch die geringen Stände II, 552., aller Menschen Herzen II, 815. 928. ff. Gott ist wunderbar in seinem Rath II, 50. 353. 1112. f. 1151. 1156., er erniedrigt die Erstgeburten und Großen, erwählt die Niedrigen I, 298. f.; II, 242. 552.; die Vernunft kann die R. Gottes nicht begreifen II, 1772.; warum wir sie nicht verstehen II, 731. f. 804. f.; ist unerträglich für das Fleisch II, 354. 1092. ff. 1291. f.; ist voll Aergernisse I, 784. 855., doch sollen wir uns nicht ärgern II, 1689. ff. Gott scheint schwach und närrisch: wo das Evangelium hinkommt, erhebt sich auch leibliche Noth I, 783. f.; Gott scheint oft zu schlafen I, 1195. ff.; schweigt oft zu bösen Anschlägen II, 1087. ff.; warum das Böse so oft die Oberhand gewinnt I, 726. ff.; II, 1113. f.; wenn der Teufel tobt, ist Gott auch gegenwärtig II, 1095. Gott läßt sich nicht hineinreden II, 1289. f.; die Sacramentschwärmer wollen Gott meistern II, 1290. f. Gottes R. leuchtet in den Historien II, 931. f. R. \gesperrt{und Führung der Heiligen:} Gott spielt mit ihnen wie ein Vater II, 11. 781. ff.; führt sie wunderlich I, 782. f. 890. ff. 922.; II, 1. 354. ff. 375. ff. 1090. 1288. f. 1565.; wir können es nicht fassen II, 1382. f.; Gott scheint seine Verheißung und Barmherzigkeit oft zu vergessen II, 1090. 1289. 1384. 1692.; er macht sie lebendig, tödtet sie aber erst I, 958.; II, 1068. f. 1150. f.; wozu Gott sie so führt II, 928. 1090.; bewahrt sie, daß sie nicht wider ihn sündigen I, 1331. f.; macht wieder gut, was sie verderben II, 685.; er führt Alles herrlich hinaus II, 1694.; Gott hat die Patriarchen geführt, wie er uns führt II, 1077. ff.; ein schöner Spiegel, wie Gott sie regiert, ist Josephs Spiel mit seinen Brüdern II, 1450. ff.; wir sollen dabei nicht murren II, 1693., sondern uns mit Gottes Wort trösten II, 1692. f. Gott regiert seine \gesperrt{Kirche} wunderbar II, 553. f. \gesperrt{Mittel:} Gott regiert Alles durch sein Wort I, 36.; bedient sich der Hülfe der Menschen und Engel I, 1245. 1249. f., der von ihm geordneten Mittel I, 1267. f. — S. \gesperrt{Königreiche, Regimente}.

Regimente. Haben nur zeitliche, die Kirche aber ewige Verheißungen I, 372.; Gott bestellt, erhält und beschützt sie I, 850.; sollen vornehmlich trachten, gemeinen Frieden zu erhalten I, 846. — S. \gesperrt{Königreiche, Regent, Regieramt}.

Reguel. Esaus Sohn, wann geboren II, 993.

Regulus, M. Attilius. Augustin lobt ihn über die Maßen I, 581. f.

Rehoboth. Theil Ninives I, 674.

Reich. 1. \gesperrt{Reich Gottes:} ist wohlthun, dienen I, 977.; wird Allen aus Gnaden gegeben I, 1418.; wer es zuerst sucht, dem fällt das Andere zu II, 88. 90. f. 98. f.

2. R. \gesperrt{Christi:} herrlicher denn die Reiche Moses und der Väter II, 1887.; ist geistlich II, 828., darin regiert das Wort, nicht äußerliche Gewalt, wie in andern Reichen II, 1973. 2000. 2008., welcher Unterschied zu merken II, 2000.; ein R. der Gnade II, 1490.; ist im Wort gewaltig wider Tod etc. II, 1974.; besteht im Gehorsam des Glaubens II, 1987. 2000.; darin gibts Starke und Schwache II, 152. f.; für die Herrlichkeit des R.’s Christi sollen wir Gott danken II, 1875. R. \gesperrt{der Gnade,} was es sei II, 1876.

3. R. \gesperrt{des Teufels:} ist Undankbarkeit I, 977.

4. \gesperrt{Reiche der Welt:} was man unter einem R. versteht I, 1099.; sind gefallen durch Verwirrung der Sprachen I, 704. — S. \gesperrt{Evangelium} 2, \gesperrt{Kirche} 12.

Reiche (divites). „Reich“ und „gottlos“ ist in der Schrift gleichbedeutend II, 85.; ist schwer, daß sie gerecht und gottselig seien II, 84.; stehen in Gefahr, daß sie versäumen den Armen zu helfen II, 85.; Hiob und David waren reich II, 85.

Reichthum. Ist an sich nicht böse II, 84., eine gute Creatur I, 836., Gabe Gottes II, 84. f., nichts gegen christliche Tugenden I, 1707.; gefährlich, weil wir Fleisch sind II, 83. ff.; nicht zu verwerfen, sondern zu gebrauchen I, 834.; Philosophen, Mönche und Wiedertäufer verwerfen ihn I, 829. f.; wie zu gebrauchen I, 835. f.; wer nicht Liebe damit übt, begeht Diebstahl II, 84. f. — S. \gesperrt{Abraham} 5, \gesperrt{Güter} 1, \gesperrt{Ueberfluß}.

Religion. Jede hat ihre sonderlichen Wahrzeichen und Ceremonien II, 1614.; darin soll man nichts nach eigenem Gutdünken vornehmen I, 1594. ff.; R. soll man sich nicht vorschreiben lassen II, 267. f.; wo es um die R. wohl steht, erhebt sich allerlei Noth I, 784.; bei der höchsten R. sind die höchsten Greuel I, 504. f.; welche Ordnung man in der Lehre der Religion halten soll I, 399. f.; die \gesperrt{väterliche} R. verlassen ist schwer I, 737.; die \gesperrt{päbstische} R. ist offenbar gottlos I, 737.

Rephaim. I, 879. f. — S. \gesperrt{Tyrannen}.

Resen. Ein Theil Ninives I, 674.

Reue. R. \gesperrt{Gottes:} bei Gott kann keine Veränderung seines Willens oder Rathes stattfinden I, 573.; von der R. Gottes über die Schöpfung der Menschen I, 480.; „Es reuete Gott“ I, 486.; Gottes R. streitet nicht mit seiner Weisheit I, 485. f. R. \gesperrt{der Menschen:} nur durch R. wird man Abrahams Same I, 1438.; durch R. verdient man die Seligkeit nicht I, 1438.; Lehre der Sophisten von der Reue II, 1458.; durch seine Lehre von der R. lehrt das Pabstthum Zweifel an Gottes Gnade II, 1900.

Reuel, Reuling. Ist das Hündlein „bös Gewissen“ I, 719. 1104. ff.; beißt uns nach der Uebertretung II, 1104. ff.

Rhetorik. S. \gesperrt{Dialektik} 2.

Richten. Heißt: Uebelthäter strafen, Unschuldige verschonen I, 1202.; wir sollen nicht richten d.\,h. nicht argwöhnisch sein I, 796.; wir sollen nicht alsbald richten, wenn jemand fällt I, 633. f. 636.; richten ist predigen, dadurch wir gestraft, ermahnt etc. werden I, 453. 500.

Richter. Zustand des jüdischen Volks unter den R. II, 2027. f.

Riphath. Völkername I, 660. f.

Ritterlich. Was ein gottselig ritterlich Herz sei II, 1678.

Rock Christi. Welcher Art er war II, 1038.

Rom. \gesperrt{Das alte} R.: in vier Stände getheilt I, 672.; die Weiber durften Gold tragen II, 523. 1214. f. \gesperrt{Zu Luthers Zeit:} konnte nicht ärger sein I, 698. 1225. f.; II, 1658.; nur Diebstahl und Bettelsack galt für Sünde I, 821.; ist das andere Babylon II, 2025.

Römer. Die alten R. hießen torquati II, 1002.; waren Epicurische Leute I, 854.; hatten neben den Rechtsverständigen auch Priester und Lehrer II, 1335. f.; ihre Könige waren sehr zur Unzucht geneigt I, 804.

Rorarius, M. Georg. I, 1.

Roßkäfer. Ihre Entstehung I, 63.

Röthe. Was R. bedeutet II, 76.

Rotten. Müssen kommen I, 711. — S. \gesperrt{Secten}.

Rottengeister. Versprechen viel von ihrer Lehre II, 912. — S. \gesperrt{Irrlehrer}.

Ruben. 1. \gesperrt{Jakobs Sohn.} Wann geboren II, 548. ff.; Anlaß zum Namen II, 523. f.; findet 8 oder 9 Jahre alt Dudaim II, 565.; etwas frömmer als Simeon und Levi II, 1498.; Jakob verkündigt Ruben den Fluch II, 1920. ff.; ist nicht gar verflucht II, 2055. — S. \gesperrt{Jakob} I, 19.; II; \gesperrt{Joseph} 12.

2. R. \gesperrt{und Bilha.} II, 972. ff. Die Juden leugnen fälschlich, daß sie Blutschande getrieben II, 975. f. 1932.; es war eine greuliche Sünde II, 973. f. 976. 1921. 1939.; geschah von Ruben aus Leichtsinn, Bilha war schön und unbeständig II, 973.; diese That gab Veranlassung zu Haß und Feindschaft II, 975., war eine Schmach für Bilhas Söhne II, 975.; dadurch verlor R. sein Ansehen als Erstgeborener II, 1101. 1109. 1138., schadete er seinen Nachkommen II, 974.; wie
\quellseite{569}
R. versuchte, sich beim Vater wieder einzuschmeicheln II, 1116. f.; die Sünde war R. vergeben, doch ward er dafür gestraft II, 1920. ff., Anderen zum Exempel II, 1922. 1934.

3. \gesperrt{Der Stamm} R. Schwach, verachtet, der geringste in Israel II, 974.; der schwächste Stamm, stolz, faul II, 1933.

Rückfall in Sünde ist gefährlich I, 276.

Ruhe Gottes. „Gott ruhte“ I, 94.

Ruth. Moabitin, Altmutter Christi I, 368. f.; erlangte durch den Glauben Vergebung II, 1830.

\section{S}

Saba. Die Königin von S. erkannte den rechten Gott I, 1052.

Sabbath. Ist für den Menschen, nicht für die Thiere geordnet I, 98.; Gott hat ihn geheiligt I, 96.; wozu geordnet I, 96. ff.; sein eigenstes Werk und Amt ist, daß Gottes Wort gepredigt und gehört werde I, 99.; Adam hat ihn nach dem Fall beobachtet I, 97. — S. \gesperrt{Stand der Unschuld}.

Sabbather. Lehre dieser Secte I, 873. f. 1009.; II, 1351.

Sacharja. Schöpfte aus Mose II, 1137. — S. \gesperrt{Propheten} 2.

Sachsen. Das Land hat seinen Namen von der Sachsenburg I, 1492.; die S. nennen ihre Edelleute „Gutmann“ I, 879.

Sackmann. II, 2031.

Sacramente. Was ein S. sei I, 1045.; sind Zeichen des göttlichen Willens zur Seligkeit I, 1049., Pfänder der Gnade I, 1105.; Beschneidung war ein S. I, 1049.; Eid ist gleichsam ein S. I, 1465.; die S. des N.\,T.’s unterscheiden nicht bloß von den Heiden, sondern machen zu Gottes Volk und selig I, 1050.; in den S. sind die Creaturen in das Wort gefaßt und haben daher ihre Kraft I, 279.; S. sind nichts ohne das Wort I, 1248. f.; sind aus Christi Seite geflossen, haben ihre Kraft aus dem Blute Christi II, 1992. f.; an S. darf man nichts ändern I, 844. 848.; gewöhnlich verachtet I, 1045.; der Pabst macht daraus menschliche Werke I, 905., lehrt, sie geben Gnade wegen des äußerlichen Werks II, 1904.; können ohne Glauben nicht Gnade wirken II, 1904. — S. \gesperrt{Erscheinungen, Gnadenmittel, Wort} 4, \gesperrt{Zeichen} 2.

Sacramentirer. Sind reich an Worten II, 715.; verachten die Sacramente, weil sie nicht darauf achten, daß Gottes Wort dabei ist I, 1248. f.; folgen ihrer Vernunft, wollen Gott meistern II, 1290. f.; wie sie beweisen, daß Christi Leib und Blut nicht im Abendmahl sei I, 1487. f.; halten ihre Lehre, Brod sei Brod, für Weisheit I, 196.; fanden Beifall I, 726.; werfen Luther Trennung der Einigkeit vor II, 1249.

Saeculum. Bedeutung des Worts I, 477. f.

Sadolet. Gelehrt, beredt, in Glaubenssachen unverständig I, 998.

Salah. Lebte mit Jakob 33 Jahr II, 89.

Salben. Geschah mit Oel oder Wein II, 960., um Personen, Kleider etc. vom weltlichen Gebrauch abzusondern II, 440. 960.

Salem. Bedeutung des Wortes II, 850.; obs ein Städtename war II, 851.; Jerusalem hieß auch Salem II, 851.

Salomo. 1. \gesperrt{Der König.} Hatte viele seiner Weiber unter dem Gesetz der Leviratsehe II, 1179.; bat Gott um Weisheit II, 276. f.; nahm viele seiner Sprüche aus Abrahams Historie I, 1156.; die Juden halten ihn für den verheißenen Samen Abrahams I, 753.

2. \gesperrt{Rabbi} S. Ist fein in der hebräischen Sprache II, 1995.; seine Meinungen über: „Hebräer“ I, 715., Jakobs Gesinnung bei der Nachricht, Joseph lebe II, 1745.

Samariter. S. \gesperrt{Sichemiter}.

Same. Ursache, daß er wächst und seine Art behält I, 44. f.; kräftige Knaben und Ehemänner verlieren S. II, 1928.

Sanctes Pagninus. Gelehrter, fleißiger Grammatiker I, 363.; wie er „Naphthali“ erklärt II, 546.

Sanctum saeculare. Was es sei II, 829. 1001.

Sanhedrin. Ursprung des S. II, 1969.

Saracenen. Sind eigentlich Hagarener I, 994., wollen nicht so heißen II, 13.; Ismaels abgefallene Nachkommen I, 1453.; warum sie sich S. nennen I, 1453.; ihr Gedicht von Isaaks Opferung I, 994.; rühmen sich der Verheißung Ismaels I, 1113.; Ismaels Name machte sie stolz I, 996.; rühmen sich Erben Abrahams um der Beschneidung willen I, 1008.; waren sehr gewaltig, vom Türken überwunden I, 994. — S. \gesperrt{Ismaeliter}.

Sarah. \gesperrt{Weß Tochter} sie war I, 720. f. 723. 1351. f. 1848.; Veränderung ihres \gesperrt{Namens} aus Sarai in Sarah, Bedeutung beider Namen I, 1094. f. 2042. Wird wie keine Andere in der Schrift \gesperrt{gerühmt} I, 1605. 1626.; übertrifft alle Nonnen I, 1161.; ist eine Königin und Hausmutter über alle anderen I, 1615., eine Aebtissin, herrlicher als die im Pabstthum I, 1628. \gesperrt{Glaube:} war ein gläubiges Weib I, 972. 1171.; Schrift rühmt ihren Glauben I, 739.; freute sich im Heiligen Geist über Isaaks Geburt I, 1378. f. \gesperrt{Tugenden:} war ein weises, gottseliges, mit vielen Tugenden geziertes Weib I, 970. f.; fromm I, 737. f.; an ihr kann man die höchsten und schwersten Tugenden einer heil. Hausmutter lernen I, 1160. ff. 1170. ff.; ehrerbietig I, 1170. 1389.; fleißig II, 1239.; geduldig I, 1174. 1603. ff.; gehorsam I, 773. 793. 1389.; keusch und züchtig I, 1169. f. 1173.; züchtig und bescheiden I, 972. \gesperrt{Anfechtungen} u. \gesperrt{Leiden:} hat viel und mancherlei Unglück muthig ertragen I, 1603. ff.; ihre Unfruchtbarkeit I, 723. f. 1080. 1378.; hatte wohl nicht geglaubt, als ihr Abraham sagte, sie werde nach göttlicher Verheißung noch gebären I, 1131. f.; lachte ungläubig I, 1168. f., welche Schwachheit ihr vergeben wurde I, 1169. 1174. \gesperrt{Schönheit} im Alter I, 802. f. 1312. Ehrte den \gesperrt{Ehestand} I, 972.; von dem Befehl, Isaak zu opfern, wußte S. wohl nichts I, 1490. f. Ist ein Exempel für Ehefrauen I, 970. ff. — S. \gesperrt{Abraham} 33, \gesperrt{Jisca}.

Satan. Ableitung des Worts II, 207.

Saul. Heißt: gefordert I, 385. Jakobs Enkel, in wiefern Simeons Sohn II, 1763. f.

Saumseligkeit. Leiblich und geistlich schädlich I, 1262. f. 1739. f.; II, 97.

Scepter. Gestalt des S.’s II, 1241.; Jupiters S. II, 1423.; Zeichen der Gewalt II, 1843.

Schädelstätte. Lage I, 378.

Schadenfreude. Dem Fleisch natürlich II, 738.

Schafe. Tragen an manchen Orten zweimal des Jahres II, 697.; Frühlinge besser als Spätlinge II, 604.

Scham. Die ehelichen Glieder I, 145. 205. f.

Schändung eines Mägdleins, eine große Schmach II, 859.

Scheba. Hams Urenkel, Wohnsitz I, 666.

Scheol. Was Sch. sei II, 1541. ff. — S. \gesperrt{Gehenna, Schooß}.

Schifffahrt. Von Menschen erfunden I, 66.; Javans Nachkommen darin geschickt I, 662.

Schiloh. Warum das Wort am Ende mit h geschrieben wird II, 1970. ff.; Bedeutung des Worts ib.; ist gekommen II, 1967. ff. — S. \gesperrt{Messias}.

Schlacht. Ausgang kann man vorher merken II, 929. 1378. f.

Schlaf. Was er sei, weiß man nicht II, 1546.; darin leben wir und leben doch auch nicht II, 216. f. 1545.; ein Bild des Zustandes der Seele nach dem Tode II, 216. f. 1545. f.; verschiedene Arten des S.’s I, 157.; eine Gabe Gottes I, 157.; man darf wohl etwas trinken, um S. zu bekommen I, 1716.

Schlange. Man bediente sich der S. zur Wahrsagerei II, 583. Im \gesperrt{Paradiese} war sie das schönste Thier I, 94. 228.; Gestalt I, 185. f. 230.; List I, 178.; war eine natürliche S., dadurch Eva versucht ward I, 226. f. 266. f.; sie redete aus Satans Eingeben I, 618.; Strafe, weil sie Eva verführt I, 186. 227. 229. Die \gesperrt{eherne} S.: woher die Kraft, gesund zu machen I, 278. f.; von den Juden angebetet I, 568.; Christi Allegorie davon I, 611.

Schlangensame. Des Teufels Same I, 266. ff.

Schlesien. Von den Sechsstädten in S. I, 675.

Schlüssel. Das mündliche Predigtamt I, 1066.

Scholastiker. Mit Bedacht zu lesen I, 1086.; warum man sie nicht recht versteht II, 849. f. Sie sollen nicht Theologen heißen II, 2079.; ihre Theologie ist Philosophie, bloßer Wahn I, 581.; II, 1927. 1997.; gehen an den 10 Geboten vorüber II, 1896.; verdunkeln die Verheißung II, 1999. Ihre \gesperrt{Lehre} betreffend: Ps. 4, 7. II, 1672.; ob Christus von sündlichem oder reinem Fleische geboren sei II, 1171.; Erbsünde II, 1464.; Genugthuung II, 1456. 1464.; Gottes Gnade I, 815.; Sphären in der Hölle II, 1546. ff.; Reue II, 1458. 1464.; Sünde II, 2003. ff.; Ungewißheit der Vergebung II, 2059. — S. \gesperrt{Sophisten}.

Schönheit. Worin sie bestehe II, 492. f.; Gabe Gottes I, 856. 1706. f.; geringer als Tugend II, 1251.; macht die Tugend um so angenehmer II, 1250.; nicht zu mißbrauchen I, 856.; man mag ein Weib um ihrer S. willen wählen II, 486.; Jakob liebte Rahel um ihrer S. willen II, 484. ff. — S. \gesperrt{Rahel, Sarah}.

Schooß. \gesperrt{Abrahams} S. ist die Verheißung, Abraham geschehen I, 1049. 1757., der Glaube Abrahams I, 1052.; dahin kamen alle Gläubigen nach dem Tode II, 2066. f.; warum er Abrahams S. heißt II, 2066. f.; ist im N.\,T. nicht mehr I, 1757. \gesperrt{Adams} S. ist die Verheißung, im Paradiese gegeben I, 1757. Im N.\,T. haben wir \gesperrt{Christi} Schooß I, 1757. — S. \gesperrt{Scheol}.

Schöpfer. Gott hat als Schöpfer Alles in seiner Hand II, 929. f.; hält das Meer inne II, 930. f.; des S.’s Werk ist: reich machen, Todte auferwecken, Betrübte aufrichten, Hoffärtige stürzen II, 928. Den \gesperrt{Artikel} von Gott dem S. soll man fleißig treiben II, 932.; wirkt großes Vertrauen II, 930.; wir sollen auf den S. vertrauen II, 933.; unser Vertrauen ist oft sehr gering II, 929. f.; ihn können wir nur dann anrufen, wenn wir wissen, daß wir einen gnädigen Gott haben II, 2014. 2029. — S. \gesperrt{Vertrauen}.

Schöpfung. I, 1. ff. 8. ff. 27. ff. 40. ff. 48. ff. 58. ff. 66. ff. 90. ff. 140. ff. Was jede Person in Gott bei der S. gethan I, 11. f.; das Mittel oder Instrument ist das Wort I, 19. 22.; Alles aus Nichts erschaffen I, 2. 12. 22.; alle Creaturen, die da waren und sein werden, sind alsbald mit geschaffen I, 92. f.; ob Gott auch nach dem Fall Neues geschaffen I, 93. f.; die sechs Schöpfungstage sind eigentlich zu verstehen I, 6. f. 45. 84. f. 147. ff.; beim 2. Tagewerk steht nicht, daß Gott es ansah und es sehr gut war I, 38. f.; Gott schmückt erst die Erde, darnach den Himmel I, 46.; die S. kann man nur aus dem Worte recht verstehen I, 152. 159., die Vernunft kann sie nicht begreifen I, 150.; wie weit man in seinen Gedanken über die S. gehen darf I, 22.; ob Gott die unfruchtbaren Bäume und Sträuche, Dornen und Disteln geschaffen habe I, 46.; ob die S. im Frühjahr oder Herbst geschah I, 45. f.

Schöpfungswort. Schafft und wirkt noch jetzt I, 92. f. 155.

Schrecken. Was der rechte S. ist I, 701.

Schreiben konnten vielleicht die Patriarchen I, 418.

Schrift, die heilige. Ist Gottes Brief an uns I, 1055. 1069.; darin kann ein Lamm gründen, ein Elephant schwimmen I, 2.; ist ein Brunnen, der Durst erweckt II, 2.; der theuerste Schatz I, 684. f.; allein die rechte Weisheit I, 152.; übertrifft alle Erscheinungen I, 1122. f.; lehrt von höheren Dingen, denn alle anderen Bücher I, 1554.; wo sie von geringen Dingen handelt, sind köstliche Edelsteine verborgen II, 543. f.; darin ist nichts so gering, es dient zur Besserung I, 793.; übertrifft mit ihren Historien alle weltlichen Historien I, 717. f. 861. 960.; II, 1925. 1964.; warum sie so geringe kindische Dinge beschreibt II, 566. f.; redet zuweilen von leiblichen Dingen ausführlich, von geistlichen kürzer I, 1712.; die ganze S. ist um Christi willen gegeben II, 1177. 1200.; ihre Historien sind rein und keusch II, 557.; von den scheinbar unfläthigen Historien der S. II, 1167. ff. 1200. ff.; in ihr stimmt das Vorige mit dem Letzten I, 654.; sie trügt nicht, weder im Vergangenen noch Zukünftigen I, 961.; kann nicht verführen I, 1047.; kein Wunder, wenn uns etliche Dinge darin \gesperrt{nicht verständlich} sind I, 1359.; was darin steht, ist uns zur Lehre geschrieben I, 1712., \gesperrt{dient} dazu, unser Leben zu bessern I, 794.; sie zerfällt in Drohungen und Verheißungen I, 1190.; hebt das Vertrauen auf Menschen auf, lehrt aber die Liebe I, 801.; handelt die Lehre von guten Sitten besser als Cicero und Aristoteles I, 802. Sie soll unser ganzes Leben regieren II, 211.; ihr sollen unsere Gedanken weichen I, 32.; sie soll in der Kirche \gesperrt{regieren} I, 766. Was sie lehrt, können wir getrost lehren, aber nicht mehr I, 1657.; sie soll allein Richterin und Meisterin bleiben unter allen Büchern I, 1290.; aus ihr können wir Alles richten I, 863. ff.; was damit nicht übereinkommt, ist vom Teufel I, 695.; was in den Canones nicht damit übereinkommt, ist zu verwerfen I, 1311.; wir müssen bei der S. Art zu reden \gesperrt{bleiben} I, 36.; nichts davon noch dazu thun I, 1064.; nichts nachlassen I, 844. 848.; nicht darüber hinaus forschen I, 1084.; die Papisten lehren, sie enthalte nicht alles, was zu glauben sei I, 1526.; die Papisten \gesperrt{verfälschen} sie I, 902.; aus Verfälschung der S. folgt Verderben I, 180. ff. Wir sollen sie \gesperrt{werth} halten I, 960.; die Menschen setzen Menschenwerke über Gottes Wort I, 1544.; sollen sie fleißig lesen I, 1055. 1069.; II, 2057., lebenslänglich studiren II, 79., Christum darin suchen I, 1047.; liest man sie als göttliche Weisheit, dann hat man auch den rechten Nutzen davon II, 2.; der Heil. Geist segnet das Lesen I, 1055.; wer sie ohne den Heil. Geist liest, der wandelt in Finsterniß II, 2030.; die Schrift soll man lesen, sich nicht auf menschliche Commentare legen I, 1281. — S. \gesperrt{Bibel, Christus} 6, \gesperrt{Geheimnisse, Obrigkeit} 10, \gesperrt{Widersprüche, Wort} 3.

Schriftauslegung. Die Schrift zu verstehen, dazu gehört der Heil. Geist I, 452. f.; der natürliche Mensch kann geistliche Dinge nicht verstehen I, 449.; je herrlicher der Text, desto schändlicher verfälscht ihn die Vernunft I,
\quellseite{570}
450.; Beispiele gemarterter Texte I, 389. 449. ff.; II, 1966. 1968. 1976. f. 2018. 2035. \gesperrt{Regel} für Schriftauslegung, betreffend: Bedeutung der Wörter II, 1486. ff. 1610.; den Buchstaben der Schrift I, 626. 903.; die Dreieinigkeit II, 915.; dunkele Stellen I, 719.; II, 695. 997. 1978.; Erscheinungen II, 775. f.; Freiheit in der Auslegung II, 1784.; Glaubensgewißheit II, 1380.; Gebote und Vermögen, sie zu halten I, 329.; der Heiligen Gebrechen und gute Werke II, 836. 843.; ihre Thaten I, 695.; Mehrzahl statt Einzahl I, 1381.; Mehrzahl von Gott gebraucht I, 70. f.; wie Menschen urtheilen, so redet die Schrift oft II, 654.; Nachdruck der Wörter II, 702. f.; Namen II, 18. 988.; Predigtamt: wo man füglich sagen kann, Gott habe durch Menschen geredet, soll man es dem P. zu Ehren thun I, 531.; Subject, davon gehandelt wird I, 322. f. 364. 451. f. 596.; „Tage“ II, 1192.; Uebersetzungen I, 321.; Verheißungen und Drohungen I, 648. 1083. 1106. f.; Verstand der Schrift: vierfachen I, 950. ff.; II, 1844. 1927., buchstäblichen und historischen I, 284. f. 950. ff.; II, 774., fleischlichen II, 1988., figürlichen I, 1150., heimlichen II, 1844. f.; Wort- und Sachkenntniß s. oben \gesperrt{Subject}; Worte Gottes, was die Schrift so nennt I, 1183. f.; Zeit und Zeitrechnung I, 95.; II, 1159.; Zusammenhang I, 370. f. — S. \gesperrt{Allegorien, Drohungen} 2, \gesperrt{Erscheinungen} 1, \gesperrt{Grammatik, Hysteron, Menschenworte, Prophezeiungen} 1, \gesperrt{Rabbinen, Recapitulatio, Verfälschungen, Verheißungen} 1, \gesperrt{Wiederholung}.

Schuldner. Bleiben auch die Christen II, 1496.; bleibt man Gott, so lange man dem Bruder nicht vergibt II, 1496. f.

Schulen. Hatten die Juden in den Städten I, 1117.

Schüler. Ein S., der dem Lehrer im Glauben zuhört, thut ein goldenes Werk II, 1743.; solchen, die sich aufs Predigtamt vorbereiten, soll man gern helfen II, 449. ff.; das thun die Reichen ungern II, 633.

Schulmeister. Ihr Beruf ist lehren und erziehen I, 978.; es ist genug, wenn Einer aus zehn Schülern fleißig lernt und sich ziehen läßt I, 978.; sollen die Schüler züchtigen II, 807., aus Liebe II, 1460. f.; suchen so der Schüler Bestes II, 1491.; das sollen die Schüler merken können II, 1491. f.

Schwaben. Galten für närrisch II, 851.

Schwachgläubige. Trösten sich leichter mit dem Tode der Heiligen als mit dem Tode Christi I, 1607. ff.; gefallen Gott auch II, 1571.; Gott hat mit ihnen Geduld II, 1332., verläßt sie nicht II, 927. f., will sie stärken, nicht verdammen I, 1160.; man soll sie nicht verachten I, 1555. f.; wir müssen bisweilen etwas schwach werden II, 1870.

Schwachheiten, sündliche. Kommen aus dem Fleisch II, 836.; Isaaks S. II, 150. ff.; Gott hat damit Geduld II, 151. f. 173. 517. 836. 1329., rechnet sie uns nicht zu II, 526. f., wendet sie zum Besten II, 274. f.; darin sollen wir uns mit dem Beispiel der heil. Väter trösten II, 151. ff. Ein Christ ist nicht verdammt, wenn er auch in einer Schwachheitssünde vom Tode ereilt wird I, 1282. f.

Schwangerschaft. Man soll sie nicht hindern I, 1748.; äußere Eindrücke haben großen Einfluß auf die Frucht II, 601. f.; besonders während derselben soll der Mann mit Vernunft bei dem Weibe wohnen II, 602. f. — S. \gesperrt{Empfängniß}.

Schwärmer. Fehlen oft in der Beweisführung I, 132.

Schwarmgeister. Sind unnütze Wäscher II, 1385.

Schweiß. Das ist verfluchtes Gut, das mit anderer Leute Schweiß erworben wird II, 699.; S. \gesperrt{des Angesichts} ist mancherlei I, 259. ff.

Schweiz. Verfassung der S. I, 675.

Schweizer. Halten auf Zucht, aber nicht am Wort II, 1827.

Schwenkfeldt. Verwarf Predigt u. Sacrament I, 1249.

Schwester. Bedeutung des Worts I, 1351.

Schwiegermütter. Hassen gemeiniglich die Schwiegertöchter II, 243. 993.

Schwiegertöchter. Leiden nicht gern die Herrschaft der Schwiegermütter II, 331.

Schwören. Was schwören heißt I, 1674. — S. \gesperrt{Eid}.

Scipio. Mit S. ist kein jüdischer oder christlicher König zu vergleichen II, 462.; ward im Glück nicht stolz II, 1387.; kein Held wie Abraham I, 889.

Scotus. Meinung: vom Feuer des jüngsten Tags I, 1540., von den natürlichen Kräften des Menschen I, 580. f.

Seba. I, 666.

Sebulon. 1. \gesperrt{Jakobs Sohn}. II, 551. 571. Segen Jakobs II, 2011. ff. — S. \gesperrt{Jakob} I, 19.; II; \gesperrt{Joseph} 12.

2. \gesperrt{Der Stamm}. Wohnsitz II, 2011. ff.; Hantierung II, 2017. f.; nach Assyrien geführt II, 2017.

Secten. Brechen ohne Gebot in die Kirche ein und verwirren Alles II, 278.; maßen sich den Namen der Kirche an II, 99.; immer neue erheben sich, die Kirche zu plagen I, 938.; II, 99.; es ist unmöglich, ihnen alle Wege zu versperren II, 938. — S. \gesperrt{Beschneidung} 8, \gesperrt{Rotten}.

Seele. Ursprung II, 1770.; eine wunderbare, uns unbekannte Creatur, großer Freude und Traurigkeit fähig II, 1670. f.; hat große Kraft, beherrscht den Leib II, 1670. f., macht ihn lebendig I, 596. — S. \gesperrt{Zustand}.

Segen. Der göttliche S. ist das Werk aller drei Personen II, 1866. 1868.; ist nicht ein Wunsch, sondern theilt mit, was er besagt II, 298. ff. 369. f.; solche thatsächliche S. haben wir im N. T. II, 298. f.; die S. im N. T. sind größer als im A. T. II, 300. f.; segnet Gott, so weist das darauf hin, daß der Mensch zuvor unter dem Fluch ist II, 1867.; göttlicher S. ist Vergebung und ewiges Leben I, 909.; S. sind auch die leiblichen Güter I, 1652.; II, 296.; in Gottes S. wird leiblicher und geistlicher S. angeboten II, 1867.; Gottes S. ist nicht ohne Kreuz II, 301. ff.; Gottes S. können die Menschen um der Erbsünde willen nicht vertragen I, 852. f.; allen S. gibt Gott um der Frommen willen I, 893. f.; zum geistlichen S. bedürfen wir auch des leiblichen II, 294. ff. 950.; der leibliche geht vor dem geistlichen her II, 1000.; der leibliche S. hat einen großen Schein, nicht so der geistliche II, 827.

Segnen. Verschiedene Bedeutungen II, 1798.; heißt mehren I, 65. 744. 1563.; geschieht mit Worten und mit der That I, 1563.; Gott segnet mit der That I, 1563.; segnen gehört zum priesterlichen Amt I, 1558. 1585. f. 1589. f.; der Größere pflegt den Kleineren zu segnen II, 1798.

Sehen. „Gott sahe“ I, 485. f.

Seir. Bedeutung dieses Namens Esaus II, 101. f. 1001. Auch Name eines Landes II, 101. f.; Lage desselben II, 987.

Sela. Anlaß zum Namen II, 1163.

Selbstgerechtigkeit. Macht sicher und halsstarrig II, 1633.; verdammt I, 1420. 1428.; ficht auch Christen an I, 1433.; haftet den Besten an II, 1660.; Gott allein kann sie ausrotten I, 1428.; sollte an Hagar und Ismael getödtet werden I, 1420. 1428. — S. \gesperrt{Werkgerechtigkeit, Werkheiligkeit}.

Selbstüberhebung. Aergster Feind der Lehrer der Kirche I, 928. — S. \gesperrt{Vermessenheit}.

Selbstverleugnung. Worin sie bestehe I, 1392. f.; Abraham ein Exempel I, 1497. — S. \gesperrt{Kreuzigung}.

Selbstvertrauen. Diese Sünde steckt in unserer Natur I, 923.; Gott ist dem S. feind I, 923. — S. \gesperrt{Vermessenheit}.

Seleucia. Hieß früher Kalne I, 672.

Seligkeit. Von der Freude der Seligen II, 1431. 1667. f. Wie man dazu kommt I, 1433. ff.; erlangt kein Heide als Heide II, 1829., Keiner, er werde denn durch Tod und Verzweiflung geführt I, 1430. f.; erlangt man nicht durch die fleischliche Geburt I, 1030. f., worauf sich doch die Menschen verlassen I, 1430. f.; nicht durch Reue verdient I, 1438.; wird erlangt nicht dadurch, daß man thut, was an Einem ist II, 1828., nicht durch Werke, sondern aus Gnaden I, 1440. f. 1447. f., durch die Wiedergeburt I, 1030. f., durch den Glauben an die Verheißung I, 1431. 1437., nicht aber von Jedem nach seinem Glauben II, 1830.; warum die Papisten die Lehre von der Seligkeit allein aus Gnaden verwerfen I, 1441. ff.; wie Isaak, Jakob und auch Manche aus Esaus Geschlecht selig geworden II, 65. f.; wie unter dem Pabst und Türken noch Manche selig geworden II, 67. Ein Christ soll seiner Seligkeit gewiß sein I, 1066.; II, 1353. f.; wir können derselben nicht gewiß sein, wenn wir sie nicht allein aus Gnaden erlangen I, 1441.; im Pabstthum lehrte man, man müsse an der Seligkeit zweifeln I, 1066.; II, 1353. ff. — S. \gesperrt{Gewißheit, Leben} 4, \gesperrt{Verklärung}.

Sem. 1. \gesperrt{Sem}. Bedeutung des Namens I, 652.; wann geboren I, 435.; ob er der Erstgeborene Noahs war I, 434. f. 502.; warum er dem Japheth vorgezogen wird I, 435.; heißt Vater aller Kinder Ebers I, 679.; war Vater aller Gemeinden II, 89.; regierte zu Jerusalem I, 1492.; wohnte zu Salem, wen er daselbst gelehrt I, 732.; predigte den Seinen von der Sintfluth I, 898.; genoß hohes Ansehen bei seinen Enkeln I, 711. f.; II, 55. f.; von den Nimroditen verachtet I, 722.; ein großer Märtyrer I, 693.; seine Zeitgenossen I, 710. ff.; II, 89. — S. \gesperrt{Abraham} 34, \gesperrt{Melchisedek, Noah} 14.

2. S.’s \gesperrt{Geschlecht}. I, 679. ff. Wohnsitz I, 657. 666.; durch Canaans Nachkommen verdrängt I, 677.

Seneca. Aussprüche: Man beachte nicht, wer da rede, sondern was geredet wird I, 1595., über Laster I, 821.; II, 1198.

Sententiarier. Handeln vom Ursprung der Seele II, 1770.

Sephar. Wohl der Ararat I, 680.

Septuaginta. S. \gesperrt{Siebenzig}.

Seraphim. Ihre Gestalt I, 288. f.

Seth. Bedeutung des Namens I, 396. ff. 414.; hatte eine gewisse Verheißung, Kain war unstät und flüchtig I, 367.; ist der Erste, auf den die Verheißung vom Weibessamen gerichtet wird I, 395.; mit ihm geht ein neues Geschlecht an I, 397.; war Adams Bilde ähnlich I, 413.; Gottes Bild ward in ihm durch das Wort aufgerichtet I, 414.; war Hoherpriester nach Henochs Aufnahme I, 422. f.; seine Zeitgenossen I, 407. — S. \gesperrt{Adam} 2.

Seufzen. Gott versteht es II, 574.; das unaussprechliche S. ist Gott liebliche Musik I, 1435. f., ein angenehmes Opfer II, 755.; findet sich auch noch in Anfechtungen II, 535. 573. f., widersteht dem Unglauben und Ungeduld I, 810.; dadurch erwecken sich die Gottseligen wider Verzweiflung II, 754.; scheinbar ist es oft schwach, doch unaussprechlich und stark II, 750. f. 1080. f. 1584. f. 1587. ff., nicht ohne große Frucht I, 1110. 1112.; II, 814., bringt den Feinden Verderben II, 1080. f.

Sforza, Herzog. Bauernsohn, Koch II, 1786.

Sibaris. Zerstörung von S. I, 1258. f.

Sichar. Schreibfehler für Sichem II, 1912.

Sichem. 1. \gesperrt{Hemors Sohn}. Ein verzogenes Kind II, 864. 866. f.; stolz II, 882.; unbußfertig II, 866. 871. 873. 882.; heuchelt II, 872.

2. S. u. \gesperrt{Dina}. II, 855. ff. 864. ff. 879. ff. Die Schändung Dinas war eine große Sünde II, 875. f.

3. \gesperrt{Die Stadt} S. Woher der Name II, 852.; hieß später Neapolis II, 1161.; Lage I, 775.; II, 851. 853. 1161.; von Jerobeam wiedererbaut II, 851. — S. \gesperrt{Sichar}.

Sichemiten. Stolz, närrisch I, 776. 779.; II, 851.; von Jakobs Söhnen erwürgt II, 885. ff.; ihre Ermordung unrecht II, 874.; bauten einen Tempel auf Garizim I, 779. — S. \gesperrt{Jakob} II, 2.

Sichere. Satan macht sie hart I, 1425.; Lots Eidame ein Bild der S. I, 1243.; S. soll man strafen, Verzagte trösten II, 1686. f.

Sicherheit. Satan stürzt die Menschen in S. I, 195.; folgt aus Verachtung des Worts II, 98.; Folgen der S.: Blindheit und Thorheit II, 882., greuliche Sünden II, 98., Unglück II, 333. f.; S. soll man meiden I, 376. 723.

Siddim. I, 879.

Sidon. Lage I, 678.; berühmte Handelsstadt I, 666.; gewaltig II, 2012.

Siebenzig, Septuaginta. Nicht geschickt zu ihrer großen Arbeit I, 321.; II, 1766. f.; Hieronymus klagt über ihre Uebersetzung II, 593. 1847. f.; Uebersetzung von 1 Mos. 47, 31. II, 1838. f.

Sieg. Im Kriege kommt S. von Gott II, 929.; der herrlichste ist, des Menschen Herz ändern II, 816. f.

Silberling. Werth eines S.’s I, 1360.; II, 1740.

Silo. Lage II, 851.; was sich da zugetragen II, 851.

Silpa. S. \gesperrt{Jakob} 18.

Simeon. 1. \gesperrt{Jakobs Sohn}. Wann geboren II, 548. ff.; hielt zu Thimnath Hochzeit II, 1164.; der ärgste unter Jakobs Söhnen II, 1505. f. — S. \gesperrt{Jakob} I, 19.; II; \gesperrt{Joseph} 12.

2. \gesperrt{Simeon und Levi}. Waren die Junker und schädlichsten Buben in Jakobs Hause II, 1101.; hoffährtiger denn die Anderen II, 1026.; maßten sich Erstgeburt und Herrschaft an II, 1045. f. 1058. 1109. ff.; Rädelsführer beim Mord der Sichemiten II, 888., ihre Strafe dafür II, 888.; Rädelsführer im Haß gegen Joseph II, 1041. 1044. 1058. 1100. 1943. f., darin „Brüder“ II, 1943. f.; Vorfahren der Mörder Christi II, 1117.; Jakob verfluchte sie II, 1925. ff., Andern zum Exempel II, 1936. ff.; der Fluch hat sich besonders in Christi Leiden erfüllt II, 1937., auch in der Verfolgung der Propheten etc. II, 1940. 1942.; sind nicht gar verflucht II, 2055.; ihre Hoffahrt soll man groß machen zum Preis der göttlichen Gnade II, 1027.; sind wohl selig geworden II, 1947.

3. \gesperrt{Der Stamm} S. Wohnsitz II, 2022.; wird von Mose in seinem Segen ausgelassen II, 1717. 1936.; der verächtlichste Stamm II, 1717. f. 1854.; bestand nur aus armen Schreibern und Schulmeistern II, 1110.; von S. und Levi sind auch treffliche Männer gekommen II, 1943. — S. \gesperrt{Judas Ischarioth}.

\quellseite{571}
Simonie. Was S. sei I, 1618. ff.; II, 106. ff.; was die Rechte und Canones dafür erklären II, 104.; woher der Name I, 1618. ff.; steckt in aller Menschen Herzen II, 108. ff.; folgt, wo man von der Lehre von der Rechtfertigung abfällt II, 110.; die Türken sind Simonisten II, 110., auch die Papisten I, 1620. f.; S. ist die Sonne im Pabstthum II, 104. f.; herrschte im Pabstthum II, 107.; Beispiele päbstischer S. I, 1620. f.; auch in Klöstern war S. I, 1623.; S. im Pabstthum ist eigentlich weltlicher Raub II, 105. f.; Pabst ist wirklicher S. schuldig I, 1622.; II, 110. f.; S. ward von Huß gestraft I, 105.; Abraham war kein Simonist II, 103. f.; wie sich ein Prediger der S. schuldig machen kann II, 111. f.; wenn er Gehalt nimmt, das ist keine S. I, 1624.; es wäre keine S., sich vom Pabst die Freiheit zu erkaufen, das Evangelium rein zu predigen II, 112. ff. — \gesperrt{Gaben} 2, \gesperrt{Pallium}.

Simson. Kein Held wie Abraham I, 890.; ein wunderbarer Kriegsheld und Heiland, vom Heiligen Geist getrieben II, 2022. f.; in ihm erfüllt sich Jakobs Segen über Dan II, 2021. ff.

Sinear. Ebene, da Babel erbaut war I, 672. 694.

Sinne. Wenn das Gemüth durch Einen S. beansprucht ist, ist seine Wirkung auf die anderen S. schwächer II, 1193.

Sintfluth. I, 541. ff. 556. ff. Warum die Beschreibung derselben so kurz I, 409.; Verkündigung derselben nicht geglaubt I, 517.; wann sie eingebrochen I, 543. f.; kam, weil auch die Kinder Gottes in Sünde und Laster geriethen I, 448.; warum die Thiere mit verderbt worden I, 503.; Wirkungen auf: die Gestalt der Erde I, 120. 378., das Paradies I, 378., den Erdboden I, 110. 495. 511. f. 1610., die Kräfte der Menschen I, 511. f., die Früchte I, 415.; nicht nur Lamechs, sondern eine allgemeine Strafe I, 371.; zeigt den überschwänglichen Zorn Gottes an I, 406. f. 437. 548.; soll die Gottlosen schrecken I, 1087.; wen dies Exempel nicht bewegt, hat ein steinern Herz I, 388. f.; zeigt an, daß der Frommen Geschlecht in Ewigkeit währe I, 437.; ob der Beischlaf während der S. verboten war I, 570. f.; allegorische Bedeutung der S. I, 611. 613.; es wird keine allgemeine S. mehr kommen I, 605. — S. \gesperrt{Paradies} 1.

Sirach. Ableitung des Namens II, 1979.

Sittigkeit. Gefällt Gott I, 1629.; man soll sich der S. befleißigen I, 1709., die Jugend dazu anhalten I, 1635.

Sixtus, Pabst. Seine Wunderwerke I, 1545.

Socrates. Lehrte die Unsterblichkeit I, 102.; seine Antwort auf die Frage: ob man ein Weib nehmen soll II, 362.; Erasmus macht ihn zu einem vollkommenen Christen I, 581.

Sodom. Untergang Sodoms I, 849. ff. 1181. ff. 1206. ff. 1222. ff. 1240. ff. 1254. ff. 1273. ff. Lage I, 1212. 1266. f.; Gegend S.’s war ein Garten Gottes I, 849. ff.; da war wohl das Paradies I, 851.; um der Sünde willen verflucht I, 594. f.; ging in S. wohl greulicher zu als vor der Sintfluth I, 1191. ff. 1199. f. 1225. ff.; welche Patriarchen zur Zeit des Untergangs S.’s gelebt I, 1200.; wie sich dieselben bemüht, S. zu bekehren I, 1200.; S. verwarf Gottes Wort I, 1340. f.; hatten wohl am Tage vor dem Untergang einen Festtag in S. I, 1218. 1224.; die Sünde wider die drei Engel war eine Sünde der ganzen Stadt I, 1256.; ob zu Ostern untergegangen I, 1220.; ein schrecklicher Schauspiel, als wir sonstwo lesen I, 1274. 1276.; wozu beschrieben I, 1087. 1273. ff. 1288.; von den Kindern zu S. will Luther nicht urtheilen I, 1199.

Sodomiter. Ihnen widerfuhr Gutes um Lots willen I, 893., durch Abrahams Glauben I, 896.; waren unbeschnitten I, 1075.; hatten Glück und gute Tage I, 1217. f.; fielen bald nach Abrahams Sieg wieder ab I, 896.; waren ohne Gottesfurcht und Nächstenliebe I, 1217.; waren nicht gastfrei, übten Schmach und Muthwillen an Fremdlingen I, 1216. f.; suchten die Sünden nicht zu verbergen I, 1219. 1223. ff. 1226. — S. \gesperrt{Abraham} 35, \gesperrt{Blindheit, Lot} 3.

Sodomiterei. Worin diese Sünde besteht I, 1223. 1226.; Onan beging eine Sodomitische Sünde II, 1181.; durch wen die S. nach Deutschland gebracht worden I, 1223.

Sodomsäpfel. Auswendig schön, inwendig voll Asche und Stanks I, 594. f. 759. 851.; ein Denkmal der früheren Sünde I, 851.

Solarii. Name der Lutheraner I, 1573.

Sonne. Ein wunderbares Werk Gottes I, 31. 49. f.; heißt ein großes Licht I, 49.; durch sie, sagen die Astronomen, leuchten Mond und Sterne I, 49.; Herr des Tages, ihr Licht zur Arbeit geschickt I, 51.; war vor dem Sündenfall heller und schöner als jetzt I, 95. — S. \gesperrt{Licht, Zeichen} 1.

Sophisten. Ihre Lehre: von den natürlichen Kräften des Menschen I, 577. ff., ein Fünklein Liebe sei genug, das ewige Leben zu erlangen I, 1334. — S. \gesperrt{Scholastiker}.

Sorbonne. Was die Theologen der S. vom Seligwerden lehrten II, 130.

Sorge. Soll man auf Gott werfen II, 1713.

Spannader. Die Juden essen nicht davon, Viele machen eine Delicatesse daraus II, 801.; allegorische Deutung derselben II, 812. ff.

Speculation. Dadurch Keiner selig I, 1251. f. — S. \gesperrt{Leben} 8.

Speise. Pflanzenspeise gesünder als Fleischspeise I, 44.; S. der Menschen vor und nach dem Fall I, 256.; ursprüngliche S. der Menschen und Thiere I, 88.; S. vor der Sintfluth I, 44. 524., da war Alles kräftiger I, 415.; wir sollen die S. mit Danksagung genießen I, 525. 868., derselben so brauchen, daß wir der Gesundheit nicht schaden I, 868.; wir schaden uns durch unsere Speise an Leib und Leben I, 415. 712. — S. \gesperrt{Fleisch} 1, \gesperrt{Thiere} 1.

Sphären. Was Verschiedene davon gelehrt I, 32. ff. 153.; die Annahme der S. ist nicht gewiß, doch nützlich I, 33. f.; ihre Ordnung, Art und Eigenschaft ist sehr lieblich und wunderbar I, 153. — S. \gesperrt{Himmel}.

Spiegel. S. der Wahrsager II, 1056.

Spott Gottes über Adam I, 272. f. 275.

Sprache. 1. \gesperrt{Im Allgemeinen}. Einerlei S. dient zur Einigkeit I, 684. 690.; Verschiedenheit der S. führt zu Feindschaft und Krieg I, 690. ff.; durch Verwirrung der S. fallen Königreiche I, 704.; Sprachen soll man studiren I, 525.; alle S. haben ihre Eigenthümlichkeiten II, 1529. f.

2. Die \gesperrt{deutsche} S. hat hebräische Wörter II, 1936.

3. Die \gesperrt{griechische} S. hat hebräische Wörter II, 1936.

4. Die \gesperrt{hebräische} S. Man hat kein vollkommenes Verständniß der hebr. S. II, 677. f. 789. 848. ff. 1035. 1358. 1373. 1472. 1509., hatte es auch nicht zur Apostel Zeit II, 1997. f.; zum richtigen Verständniß hat das N. T. viel geholfen II, 678. f. 1837.; soll die h. S. vollkommen aufgerichtet werden, so muß es durch die Christen geschehen II, 849. 1532., ohne deren Hülfe dieselbe noch viel dunkler wäre II, 961.; woher es kommt, daß sie uns unverständlich ist II, 849. f.; ihre Art zu reden ist uns fremd I, 1168.; II, 1529. f.; von dem Ausdruck: Man thut nicht also II, 507.; ist figurenreicher als andere S. II, 1440. 1532.; von den Puncten II, 1007. f. 1837. ff., Unordnung in den Eigennamen II, 1016., Numerus II, 1769. 1788., Conjugationen I, 1561., dem Hülfszeitwort „sein“ II, 1440.; Studenten der Theologie sollen die hebräische S. studiren I, 1563.; II, 912. Gott redet mit uns hebräisch II, 964. 1145. — S. \gesperrt{Rabbinen}.

5. Die \gesperrt{lateinische} S. ist gewandter als die griechische und hebräische II, 1487. f.; hat hebräische Wörter II, 1936.

Sprachgebrauch. Demselben müssen alle Regeln weichen II, 1533.; darauf zu achten II, 677. f.

Sprüchwörter. 1. \gesperrt{Nach Anfängen}: Aut regem, aut fatuum II, 1343.; Je länger, je ärger I, 575.; Laß gehen, oder: Mitte vadere II, 1406. 1619.; O formose puer II, 1012.; Pluribus intentus II, 1193.; Qui quae vult, dicit II, 1042.; Si fecisti II, 1524.

2. \gesperrt{Nach den Stichwörtern}: Abfall II, 629.; Achivi II, 881.; Alltägliches I, 1118. 1371.; Aequivocationen II, 1908.; Arbeiten — Genießen II, 203.; Argwöhnische I, 794.; Ars II, 312.; Arzt I, 847.; Aufgeschoben II, 377. f.; Ausleger II, 1053.; Bauch II, 591.; Beredsamkeit I, 1005.; Besser machen, als man kann I, 2.; Bier II, 1415.; Dieb, fauler II, 1231.; Einsamer I, 834.; Engel II, 629.; Esel II, 1505.; Faule II, 319.; Feinde II, 877.; Feuer I, 1354.; Fiscus II, 1817.; Frau, Hausfrau II, 1241.; Frieden I, 1076.; Frons II, 1241. 1423.; Fürsten I, 447.; II, 865.; Fürstenbriefe I, 1055. 1069.; II, 214.; Fürsten, Exempel II, 865.; Fürstengunst I, 1065.; Gabe I, 1263. 1710.; Gast II, 482.; Geduld II, 1314.; Geld II, 1787.; Gelegenheit I, 1262.; sich gemein machen I, 1442.; Gerechtigkeit — Glaube II, 1745.; Geschenke I, 1474.; II, 1581. f.; Gesetz II, 1862.; Gewissen II, 205. 1593. 1596. 1684.; Gewonnen — zerronnen II, 205.; Gleiche I, 690.; Glück I, 1218.; II, 1914.; Gott II, 1698.; Gut I, 1218.; II, 86., unrecht Gut II, 205. 590.; was gut ist I, 878.; Habgier II, 588.; Häkchen II, 81.; Herr, Hausvater I, 1723. f.; II, 552. 1240. f. 1423.; Hoffahrt I, 980.; Hühner im Korbe I, 1071.; Hunger I, 1390.; II, 1087.; Jagd I, 54.; Jugend I, 585.; Jurist I, 847.; Kaufen I, 1262.; Kinder großer Leute I, 313.; II, 1368.; Kirchen — Herzen I, 839.; Kirche II, 1002.; Klappern II, 257.; Knecht, fauler II, 1231.; Krankheit I, 894. f.; Kunst II, 451.; Land I, 852.; Land — Sitte II, 478. 1184.; Lauf eines Dinges I, 1192.; Laus I, 980.; List II, 877.; Löwen II, 162.; Lügen I, 335. 1533.; Lügner I, 334.; II, 1121.; Magd II, 1241.; Mameluke I, 689.; Mitten inne I, 1438.; II, 761.; Mönch II, 1408.; Nacht II, 772.; Name, guter II, 1723.; Niemand I, 1226.; Obrigkeit I, 846.; Pferde — Ochsen II, 385.; Phrygier II, 866.; Propheten II, 1645.; Rauch — Feuer II, 363.; Recht II, 1745. f., das höchste Recht I, 843. 1476.; Recht — recht ausführen II, 874. 877.; Reichthum II, 2038.; Römer I, 337.; Ruthe II, 1460.; Säen — ernten II, 203.; Sauer II, 1087. 2038.; Schaden I, 1314. 1462.; Schalk — Glück II, 507. 727.; Schuld, alte II, 1116. 1454.; Schwalbe I, 54.; Stein II, 1406.; Sünde I, 575.; II, 1121.; Tage, gute, I, 1218.; II, 1849.; Teufel I, 1424.; II, 103.; Theologe I, 847.; Türken II, 1986. f.; Undank I, 1337. 1340.; Undankbarkeit II, 1262.; Unglück I, 180. 1294.; II, 392. 1192.; Unrecht II, 207. 1147.; Unzufriedenheit I, 1071.; Urlaub II, 889.; Venus II, 1310.; Verborgen II, 1002.; Verbrennen II, 1772.; Verzweiflung II, 1408.; Volk I, 852.; II, 884.; Wege II, 2035.; Weib I, 1372. 1672.; Wieder schlagen II, 874.; Worte II, 720.; Zeit II, 830.; Zorn II, 340.

3. \gesperrt{Sprichwörtliche Redensarten}. Sich selbst in den Backen hauen I, 1576.; den Baum auf beiden Achseln tragen II, 122.; man muß es ihm besiebenen I, 1476.; laßt mir den Schäfer bezähmen II, 1727.; vom Himmel zur Erde bitten II, 1118.; einen Bock melken, ein Sieb unterhalten II, 107. f.; in die Ernte stoßen II, 2006.; wie ein Esel geschickt zur Sackpfeife I, 3.; wie die Faust aufs Auge I, 761. 1227.; Oel und Pech ins Feuer gießen II, 350.; im Trüben fischen II, 1319.; ich möchte grau werden II, 1799.; er ließe Holz auf sich hauen II, 2019.; ein gut Jahr haben II, 1411.; das Kind muß einen Namen haben II, 760.; mit Kolben lausen II, 691.; wie ein Kranz auf dem Kopf einer Sau I, 1713.; über einerlei Leisten strecken I, 1074.; sich das Maul wischen II, 1099.; nicht einen Nestelstift I, 915.; non propter te, sed propter tuum te II, 257.; für die Raben, daß dich die Raben fressen I, 622.; auf das Rad und auf die Scheibe setzen II, 2006.; in die Kappuse werfen II, 1238.; eine Schelle anhängen II, 712. 714.; Einem die Schlüssel vor die Füße werfen II, 278.; Stück um Stück II, 878.; um einen Taubendreck I, 801.; bis das Wasser über die Körbe geht II, 1337.

Staat. Man soll sich von Jugend auf gewöhnen, demselben zu dienen II, 1409. ff.

Städte. Die erste St. von Kain erbaut I, 384.; oft nach ihren Erbauern genannt II, 852. — S. \gesperrt{Handelsstädte, Hauptstadt}.

Stadtthor. Gerichtsstätte im Morgenlande I, 1213.

Stand der Unschuld. Ein \gesperrt{Mittelstand}, darin der Mensch sündigen konnte I, 128.; sehr herrlich I, 88. 94. 128. Anerschaffene \gesperrt{Gaben} des Menschen: Unsterblichkeit I, 135., Jugendfrische I, 112. f.; Erkenntniß I, 81. 128. f. \gesperrt{Leben} des Menschen im St. d. U.: ein \gesperrt{natürliches} Leben I, 105.; Essen, Trinken, Kinderzeugen I, 69.; Speise der Menschen und Thiere I, 46. 87. f.; Verdauung I, 112. 135.; Kleider I, 171.; Nackendgehen eine Ehre I, 170. ff.; Herrschaft des Mannes und Weibes I, 167. f.; Ehestand I, 86. 125. 128. 142. ff. 163. f. 206. 265. 291.; Zeugung und Geburt I, 125.; Erziehung der Kinder I, 88. 125.; Verehelichung der Kinder I, 162. 168. f.; Arbeit I, 79. 100. 124.; Herrschaft des Menschen I, 77. 80. 87. 146. Der \gesperrt{Gottesdienst} I, 18. 96. f. 131.; Bibel und Gegenstand der Predigt I, 128. 188. f.; Gesetz Gottes I, 130. ff. Adam und seine Nachkommen würden im Glauben gelebt haben I, 188. Wie der Mensch in das \gesperrt{ewige Leben} versetzt worden wäre I, 69. 135. 145.; davon Adams Schlaf im Paradies ein Vorbild I, 158. \gesperrt{Ueberbleibsel} dieses Standes I, 82. 130. 161. ff. 167. f. 250. 412. Wir sollen oft an das Gute denken, das wir verloren haben I, 130. — S. \gesperrt{Ebenbild, Paradies}.

Stände. 1. \gesperrt{Im Allgemeinen}. Sind von Gott geordnet, nicht zu verachten I, 1164. f.; sind verschieden, als, Reiche und Arme, daran stoßen sich Viele I, 995. f.; dieser Unterschied ist von Gott geordnet II, 1404., nöthig und nützlich zu diesem Leben II, 1371. f. 1390. f. 1394. f. 1403. ff. 1786. f., vor Gott gilt solcher Unterschied nicht I, 996. 1089.; Gottes Wort heiligt alle St. II, 125.; Christen sollen in solchem St. wandeln, den Gott geboten hat II, 1626.; Gott nimmt sich auch der geringen St.
\quellseite{572}
an II, 552.; in jedem soll man treu sein I, 1180., nicht müßig II, 1791. f., tugendhaft II, 1795.; darin bleibe man, bis Gott Einen herausnimmt II, 1795.; Geduld in allen St. wirkt Gottes Wort II, 126. f.; auf den St. sei man nicht stolz II, 1371. f.

2. \gesperrt{Die drei Hauptstände}. Man theilt das ganze Leben in drei St. I, 1180., in das Hausregiment, das weltliche Regiment, das Priesterthum II, 297. f.; das sind die rechten Hierarchien I, 1394.; das Leben in diesen St. ist das allergeistlichste II, 1607.; wer außerhalb derselben lebt, lebt in einem verdammlichen St. I, 1394. Sind von Gott geordnet I, 532. 1321. 1373. ff.; II, 1567.; nöthig II, 1567.; darin handelt Gott durch Menschen mit uns I, 1456.; geordnet zu Gottes Ehre und des Nächsten Nutzen II, 264.; Pflichten derselben I, 532. 1256. f.; II, 297. f. Darin findet sich Unwille und Unlust I, 261. f. 982. 984. 987. f. 1371. ff.; II, 24. 139.; wer darin gottselig leben will, muß Verfolgung leiden I, 1393.; Satan ist ein Feind derselben I, 1459. In allen St. sind gläubige Leute nöthig II, 542. f.; man bleibe darin bei Gottes Geboten II, 1567. f.; man bitte um Weisheit, sonst geräth Alles übel II, 277. ff.; man sei geduldig und sehe auf Gottes Wort I, 1371. ff. 1373. ff.; der Hülfe Gottes kann man darin gewiß sein I, 1321.; die Engel sind darin unsere Hüter I, 1686.; alle drei St. bedürfen immer Vergebung der Sünden II, 487. f.; wie sie sich fremder Sünden theilhaftig machen I, 1257. Sind vor Gott in den höchsten Ehren gehalten II, 1061. f.; es ist Satans Meisterstück, daß er die St. in Verachtung bringt I, 1456., wie im Pabstthum geschehen II, 1567. 1606. f. 1611. — S. \gesperrt{Hausstand, Prediger, Predigtamt, Obrigkeit}.

Vgl. \gesperrt{Beruf} I.

Standhaftigkeit. St. in der Lehre gegen die ganze Welt ist schwer I, 498. ff. 537.; wir sollen standhaft sein I, 538. f.

Starke. Ihr Trost wider den Tod I, 1607. ff.

Stärke. Definition dieser Tugend II, 1877. — S. \gesperrt{Tapferkeit}.

Staupitz. Wie er regiert II, 1406., Luther getröstet II, 180.

Stein. Bedeutet oft einen Obersten II, 2045. f. — S. \gesperrt{Krankheit} 2.

Sterne. Was sie sind I, 50.

Steuern. Schwere St. eine Strafe II, 1817.; bei verschiedenen Völkern II, 1819. f.

Stifte. Eitel Hurenhäuser II, 1658.

Stimme. Sehr verschieden II, 1192. f.; gibt den Menschen am allermeisten zu erkennen II, 270.

Stirnblatt. S. \gesperrt{Hoherpriester}.

Stoiker. Sind Heuchler II, 1612.; ihre Lehre von zeitlichen Gütern I, 830. 854. f.

Strafe. 1. \gesperrt{Im Allgemeinen}. St. erregt Murren u. Feindschaft I, 989.; St. von Eltern, Lehrern, Obrigkeit soll man mit Geduld ertragen I, 999. f.

2. \gesperrt{Göttliche} St. \gesperrt{Art und Natur}: \gesperrt{St. und Verheißung} läßt sich durch unsere Vernunft nicht vereinigen I, 509. ff.; aus der St. folgt nicht, daß Gott wankelmüthig sei I, 510.; bei der St. ist immer noch Gnade II, 1951. Die gewöhnlichen St. sind Pestilenz, Hunger, Krieg, wilde Thiere I, 517.; die größte, wenn Gott sein Wort wegnimmt I, 455.; neue St. folgen, wenn neue Sünden sich mehren I, 594. \gesperrt{Wen} die St. trifft: nicht nur den Sünder, sondern auch die Creaturen, seine Habe I, 224. f. 503., Land und Leute, wenn die Obrigkeit sündigt I, 822. 892. 1324. f.; die Frommen müssen oft die St. mit tragen II, 337. f. \gesperrt{Zweck und Nutzen}: St. ist eine Wohlthat II, 1103. f., soll uns demüthigen I, 1440., zur Buße führen I, 895, zur Prüfung des Glaubens und der Geduld dienen I, 822. f., nicht aber zur Genugthuung II, 1103. \gesperrt{Erfolg und Wirkung}: St. haftet viel fester denn Wohlthat I, 604.; wird schwerer, wenn man sie als St. erkennt I, 255.; ist oft vergeblich I, 989.; die Gottlosen zwingen Gott zu immer härteren St. I, 896. \gesperrt{Verhalten}: Gott nimmt sich in der St. unserer an I, 999.; Christus verwandelt sie in seliges Kreuz I, 244. f.; verschiedenes Verhalten der Frommen und Gottlosen II, 338.; unsere Natur murrt und verzweifelt I, 999.; die Welt trägt Leid über die St., nicht über die Sünde II, 336. f.; man schicke sich mit Geduld in die St. II, 338.; beklagen wir die Sünde, so wird die St. milder I, 990. \gesperrt{Verzug und Eile mit der} St.: es ist ein Gnadenzorn, wenn Gott mit der St. eilt I, 702.; II, 1043.; unselige Leute, bei denen Gott mit der St. verzieht II, 1042. ff.; Gott verzieht, ob die Gottlosen noch wollten Buße thun I, 966.; wenn die St. verzieht, scheint Gott zu schlafen I, 699. f.; die St. folgt endlich I, 509. 701. f. 811. 965. f.; II, 1043. 1103., wenn das Maß voll ist I, 700. f.; je länger er geschwiegen, desto größer nachher der Zorn I, 702. — S. \gesperrt{Fluch, Sünde} 4, \gesperrt{Zorn, Züchtigung}.

Stufenjahre. Bringen Veränderungen I, 1604.

Suah. II, 1162.

Suchoth. Woher der Name II, 844.; Lage II, 846.

Sueton. Erzählt Träume des Nero etc. II, 1310. f.

Sünde. I. \gesperrt{Im Allgemeinen}. 1. \gesperrt{Art und Natur}. Sünde ist im Herzen, nicht in den Augen, Ohren oder anderen Creaturen I, 833.; ist etwas Natürliches, sonst wäre keine Erlösung nöthig I, 203.; haftet uns so an, daß wir doch einmal gänzlich davon können gereinigt werden II, 1171.; ist, was ohne Gottes Wort geschieht I, 567.; besteht vor Allem im Ungehorsam gegen Gottes Wort I, 179. ff.; sich der Sünde rühmen, ist keine geringe S. I, 697.; Unwissenheit kann Sünde sein I, 1466. ff. \gesperrt{Arten der} S.: Die Hauptsünden sind nicht die groben, äußerlichen, wie Ehebruch, sondern Verachtung des Worts, Eltern, Obrigkeit I, 1199.; der Papisten Verfolgung des Evangeliums ist eine schwerere S. als Davids Ehebruch I, 1277. f.; die größte Sünde ist, wenn sich die Menschen nicht wollen von Gottes Wort strafen lassen I, 456. ff., die S. der ersten Tafel bleiben gemeiniglich unter dem Schein der Heiligkeit verborgen I, 503.; Gott klagt mehr über die Gewalt, die dem Menschen geschieht, als die ihm geschieht I, 511.; was himmelschreiende S. sind I, 1197., was sodomitische I, 1219.; S. haben wir, wenn wir es auch nicht wissen I, 1348., man kann aus Unwissenheit sündigen I, 814. ff.; S., in der Bestürzung gethan, sind nicht so schwer, wie die man mit Bedacht thut I, 1296.; S. der Begehung und Unterlassung II, 692.; Begehungssünden sind greulicher als Unterlassungssünden I, 1365.; S. thut, nicht nur wer sie begeht, sondern auch wer sie billigt I, 1224. 1256. \gesperrt{Stufenweiser Fortschritt}: S. folgt, wenn man vom Wort abweicht I, 18.; geschieht erst im Herzen I, 635.; man geräth immer tiefer in S. I, 210.; eine S. zieht die andere nach sich I, 214. 219. 336. ff. 480.; II, 1192. 1276. f., wenn man nicht Buße thut II, 884. \gesperrt{Erwachen und Entschuldigungen} der S.: S. geht süß ein, sauer aus II, 1718.; man fühlt sie nicht, wenn man sie begeht I, 199. f. 699.; ihre Art ist, daß sie eine Zeitlang liegt und ruht I, 701.; ruht wie ein eingeschlafenes wildes Thier, weil sie im Werk ist I, 325. f.; das Gesetz macht sie lebendig I, 201.; sie wacht bald auf, wird bald offenbar I, 326. f. 346.; dann ängstet sie uns II, 1103. f. 2077. ff. und wird gewaltig II, 2079., Satan sucht uns in Verzweiflung zu stürzen I, 200. f.; II, 1704. 1707. f.; erwacht bei Manchen zur Verdammniß, bei Anderen zur Buße II, 1107. Die S. will nicht offenbar werden I, 222.; II, 1503., will nicht S. sein I, 218.; II, 1524.; man sucht sie zu leugnen und bemänteln I, 206. f. 212. f. 215. f. 991.; man wirft die Schuld auf Gott I, 214. f. 217. f.; II, 1529.; verräth sich selbst I, 335.; klagt sich damit an, womit sie sich entschuldigen will I, 212. f. 334. f. 338. f. 992.

2. \gesperrt{Größe und Greuel} der S. Ist ein greulicher Gestank vor Gott I, 573., die allergreulichste Krankheit II, 1723.; ist so groß, wie das Wort, dawider sie geschieht I, 198.; ist größer oder kleiner, je nach der Person, an der man sündigt II, 875. f.; ist geringer, wenn man sich ihrer schämt, als wenn man davon rühmt I, 1219.; die Größe der S. lernt man aus der Strafe der ersten Welt, Sodoms und der Engel I, 415.; wie greulich sie ist, lehrt das Todte Meer I, 852.

3. \gesperrt{Urheber der} S. Gott wirkt die S. nicht I, 1481.; Satan treibt von einer S. zur anderen I, 1226. — S. \gesperrt{Teufel}.

4. \gesperrt{Folgen und Strafen der} S. Der Fall ist leicht geschehen, aber es hält schwer, wieder Verlangen nach Barmherzigkeit zu bekommen II, 1721. 2077. ff.; je weiter der Mensch von Gott gekommen, desto weiter sucht er zu entfliehen I, 211. f. 214.; der Sünder haßt die Strafe und darum auch Gott I, 217. Gott zürnt über die S. II, 2084. 2087., das glaubt freilich die Welt nicht I, 1205. f.; die S. zieht schwere Strafen nach sich I, 242. ff.; II, 1102. f. 1259.; S., welche die Obrigkeit nicht straft, zieht die Strafe schnell herbei I, 1256.; Strafe folgt, wenn die S. zur Gewohnheit werden I, 1222.; mit den S. mehren sich auch die Strafen I, 252. f. 255. 264. f. Gott kann aus der S. Gutes hervorbringen I, 1326. ff.; II, 1703. 2083. ff., aber die S. bleibt trotzdem böse I, 1328. ff., und wir sollen nicht daraufhin sündigen II, 2083. — S. \gesperrt{Strafe}.

5. \gesperrt{Hülfe}. Es gibt keine Hülfe, wenn uns Gott nicht mit seiner Barmherzigkeit zu Hülfe eilt I, 214.; Gott läßt uns fallen, um uns die Augen zu öffnen II, 1456.; Erkenntniß der S. ist der Anfang unseres Heils I, 482. f.; einst werden wir ganz davon frei werden II, 1175. f. Nur Gottes Sohn kann die S. überwinden I, 198., kein Werk oder Opfer kann sie tilgen I, 323.; durch den Glauben werden wir frei und Herren über sie I, 324. 329. f.

6. \gesperrt{Sünden der Christen}. Auch die besten Christen fallen I, 633. 789. f.; II, 683. 2003. ff.; Christen müssen wider die S. kämpfen II, 467.; thun wohl in der Bestürzung, was sie nicht sollten I, 1292. f.; wenn sie fallen, sollen wir nicht gleich richten I, 633., uns nicht ärgern I, 635. f.; wenn sie auch sündigen, gehen sie doch nicht verloren I, 1467.; ihre S. gewinnen einen fröhlichen Ausgang II, 272. ff.; ihnen muß auch das Böse, das sie thun, zum Besten dienen I, 1326. f.; Gott läßt sie fallen, um zu vielen guten Dingen Gelegenheit zu haben I, 1326. f., um sie zu demüthigen I, 1296. f. 1558.; wir sollen unsere S. beklagen, ehe die Strafe da und es zu spät ist I, 989. f.; vor Gott, aber nicht vor Menschen, soll man sich aller S. schuldig geben I, 1348.; wenn wir fallen, sollen wir nicht verzagen II, 2084. ff., sondern an Gottes wunderbare Führungen und Absichten denken I, 1329.; II, 1212., mit den Exempeln der Väter uns trösten II, 1155., wieder aufstehen II, 2079.; wir sollen bitten, daß wir nicht sündigen II, 1941. — S. \gesperrt{Schwachheiten}.

7. \gesperrt{Sünden Anderer}. Die S. eines Mitbürgers oder Glaubensbruders soll man vor der Obrigkeit nicht verbergen helfen I, 1258.; mit den S. Anderer habe man Geduld, rechtfertige sie aber nicht I, 986.; S. soll man strafen, auch an den großen Hansen II, 231., ff. man richte sich dabei nach Alter, Stand etc. II, 231. f.; Eltern, Lehrer, Obrigkeit, auch Privatpersonen sollen die S. strafen II, 1149.; will die Strafe nicht fruchten, so überlasse man die Sache Gott II, 1149. \gesperrt{Fremder} S. soll man sich nicht theilhaftig machen I, 1256.; wie man das thut I, 1208. 1257.; wer es thut, muß die Strafe mit leiden II, 886. f.; von verstockten Sündern soll man weichen I, 1264. f.

8. \gesperrt{Eine größere} S. \gesperrt{zu meiden}, darf man nicht eine kleinere begehen I, 1229. ff.

9. \gesperrt{Von unserer} S. \gesperrt{predigt} die ganze Creatur I, 256., unsere Kleider I, 270. ff., Speise I, 276., der Verlust des Paradieses I, 281. f.

Vgl. \gesperrt{Abraham} 16, \gesperrt{Bekenntniß} 1, \gesperrt{David, Ehestand} 8, \gesperrt{Erbsünde, Erkenntniß, Esau, Gewissen} 2, 3, \gesperrt{Isaak} 10, \gesperrt{Jakob} 14, \gesperrt{Joseph} 10, \gesperrt{Lüste, Noah} 10, \gesperrt{Obrigkeit} 6, \gesperrt{Sodomiterei, Sündenfall, Väter} I; II, 9, \gesperrt{Wachen}.

II. S. \gesperrt{wider den Heil. Geist}. Ward vor der Sintfluth begangen, worin sie bestand I, 479. f.; wer absichtlich und wissentlich sündigt, ist nicht weit davon II, 1704.

Sündenfall. I, 172. ff. 196. ff. 207. ff. 223. ff. 242. ff. 268. ff. Warum Gott die Versuchung gestattet, eine vorwitzige Frage I, 176. f.; warum Mose nicht deutlich sagt, der Teufel habe Eva versucht I, 177. f.; Eva fiel nicht, weil sie ein Weib war I, 469., hielt die Sünde für etwas Geringes I, 528.; durch vorwitziges Forschen nach den Gründen des Verbots kam Adam und Eva zu Fall I, 193. 1084. f.; II, 184. Der S. \gesperrt{bestand} nicht im Apfelbiß, sondern im Abfall von Gott I, 117. 183. 194. 196. 198. 208. 278.; dadurch vergriff sich Adam an der Person Christi und göttlichen Majestät I, 274. f. \gesperrt{Folgen} des Falles I, 75. ff. 89. 172. f. 272. an der Seele I, 75. 81. 197. 212. ff. 216. 221., am Leibe I, 122. 253. 414., im ehelichen Leben I, 128. 265., in der äußerlichen Natur I, 78. 94. f. 110. 120. f. 249. f. 254.; das Paradies ist verloren I, 275. ff. \gesperrt{Erinnerungen} an den Fall sind Speise und Kleider I, 270. ff. 276. ff. Nur \gesperrt{Gottes Sohn} kann diese Sünde überwinden I, 198. \gesperrt{Wann} der S. geschehen I, 85. 99. f. 171. 176.

Sünder. Sind alle Menschen, auch die besten I, 633. 789. f.; II, 683. 1632. f. 2061.; S. lassen sich nicht gern im Zaum halten II, 1042.; fliehen den Tod als eine Strafe I, 369.; Hohn und Schmach überfällt sie zuletzt I, 480.; S., deren Lüsten Gott nicht widersteht, sind unselige Leute II, 1042. ff.; Gott läßt sie endlich in ihres Herzens Dünkel I, 480.; kein S. soll wegen der Größe seiner Sünde verzagen II, 978. f., er hat noch den Trost ewigen Lebens I, 405.; ist ein Mensch gefallen, so soll man nicht alsbald richten I, 633., oder an ihm verzagen II, 342., man strafe ihn II, 690. f.

Sündigen. S. \gesperrt{Gnade} 5.

Sur. Lage der Wüste I, 991. 1306.

Susim. Bedeutung des Worts I, 880.

Syncope. Beispiel einer S. II, 678.

\quellseite{573}
Synekdoche. Eine in der Schrift gebräuchliche Redefigur II, 1065. ff. 1762.; Beispiele: I, 915.; II, 593. f. 623. 651. 673. 1065. f. 1438. f. 1617. f. 1913.

Syrer. Gute Schützen II, 2043.

\section{T}

Tag. 1. \gesperrt{Tag}. Zwiefache Bedeutung I, 51.; bei den Juden I, 24.

2. \gesperrt{Tag und Nacht}. Zweck der Abwechselung I, 51.

3. \gesperrt{Jüngster Tag}. Zeit desselben verborgen I, 721. 1600.; Zustände vor demselben I, 409. f. 448.; ob mehrtägige Finsternisse vorhergehen, die Weiber sieben Jahre vorher nicht mehr gebären I, 587.; Gott wird damit eilen I, 448.; woher Gott das Feuer nehmen wird I, 1540.; da werden sich die Gottlosen selbst verdammen I, 213. f.; Freude der Gläubigen II, 1671.; bitte Gott, daß du ihn mit Freuden erwartest I, 437. — S. \gesperrt{Zeit} 3.

Tänze. Nicht alle verwerflich, ihr Zweck I, 1682.

Tapferkeit. Verschieden von Tollkühnheit II, 174. — S. \gesperrt{Stärke}.

Tarquinius. Zwang Lucretia zur Unzucht II, 1279. — S. \gesperrt{Lucretia}.

Tauben. Eigenschaften I, 623.; Bedeutung in der Schrift I, 618.; \gesperrt{Noahs} Tauben I, 563. ff., allegorische Deutung derselben I, 623. ff.

Taufe. \gesperrt{Gott} tauft II, 437. 1842., Christus tauft II, 1892.; geht alle Völker an I, 1042., auch die Kinder I, 1041. f. Woher das \gesperrt{Wasser} seine Kraft hat I, 279.; wir werden in Christi Tod und Blut getauft und mit seinem Blut gewaschen II, 1992.; Schwarmgeister haben ohne Wasser getauft I, 197. Die T. ist eine rechte, wenn auch der Täufling sie durch List erlangt II, 310.; soll nur \gesperrt{einmal} vollzogen werden I, 1037., selbst wenn wir abgefallen waren, nicht zu wiederholen I, 1042. \gesperrt{Wirkungen}: sie ist ein Zeichen der Gerechtigkeit I, 1045.; außerhalb der T. ist keine Seligkeit I, 1044.; darin bekommt man einen neuen Namen vor Gott II, 950. f.; dadurch werden wir zur Hoffnung des Lebens wiedergebracht I, 240.; sie zerstört die Werke des Teufels, versetzt uns in das Reich des Lichts II, 1891.; macht gerecht I, 1020. f. 1024.; verheißt nicht bloß, gibt Vergebung der Sünden II, 369.; macht selig I, 612. ff.; Kinder empfangen darin den Heil. Geist, werden kraft der Verheißung selig II, 1904.; macht alle getauften Kinder, wenn auch nicht alle Erwachsenen, selig I, 1049. Sie ist \gesperrt{nöthig} zur Seligkeit, doch ist Gott nicht unbedingt daran gebunden I, 1250.; von Kindern, die ohne T. sterben I, 1040. f. 1091. 1250. 1761. f.; im ewigen Leben bedürfen wir der T. nicht mehr I, 1027. \gesperrt{Verachtung} der T. verdammt I, 1041. Allegorien von der T. I, 611. ff. \gesperrt{Tod und} T. verkehren und vergleichen sich mit einander in der Schrift I, 613. Die T. ist nicht weniger, sondern \gesperrt{herrlicher} als die Beschneidung I, 1042. 1065. f.; wer sich der T. recht verwundern könnte, würde des Teufels spotten II, 147. Vom \gesperrt{Trost} der T. II, 513. f.; Christi und unsere T. verglichen I, 1020. f. 1024.; im \gesperrt{Nothfall} darf jeder Christ taufen II, 1891., auch ein Weib II, 2057. Bei Verwaltung der T. soll man ehrerbietig sein II, 1842. f. — S. \gesperrt{Gnadenmittel}. T. \gesperrt{Christi} s. \gesperrt{Christus} 2.

Taufnamen. Legte man im Pabstthum ab II, 988.

Tauler. Ausspruch über Ungeduld im Kreuz II, 1326. f.

Tautologie. S. \gesperrt{Wiederholung}.

Tempel. Ist jeder Ort, da Gottes Wort gepredigt wird I, 839. Der T. \gesperrt{zu Jerusalem}: Lage I, 1491.; Zeichen daran vor Jerusalems Zerstörung II, 76.; Figur des Paradieses I, 282.

Teraphim. Meinungen darüber II, 649. f.

Terenz. Ausspruch: das höchste Recht etc. I, 846.; Meinungen über: irdische Güter I, 857., Verderben der menschlichen Natur I, 584.

Testament. Durch ein T. soll man Zank vorbeugen I, 1752. f.; II, 1906. Das \gesperrt{Neue} T. macht das Alte klar I, 266. ff. 273. f. 1004. — Siehe \gesperrt{Christus} 4, \gesperrt{Gläubige, Zeit} 2.

Tetragrammaton. S. \gesperrt{Jehovah}.

Tetzel. Wollte nicht mit Petro tauschen II, 568.

Teufel. Woher die Lehre vom T. zu nehmen I, 230. f. \gesperrt{Zustand}: wodurch er gefallen I, 1085.; II, 400. ff.; was der Alkoran von seinem Fall sagt II, 401.; was er verloren I, 174. f. 230.; sein böses Gewissen II, 1720.; hat keine Hoffnung der Erlösung I, 245.; in welchem Falle er auch selig werden könnte I, 219.; II, 1515. \gesperrt{Was er ist und vermag}: ist hoffährtig I, 1687. f., listig I, 184. 1177., heuchlerisch und boshaft II, 1271., lügt II, 1049., nur ein Geist II, 439., boshaft und vermessen I, 184.; ein Feind Gottes I, 689. 695., des Worts I, 194., aller Ordnungen Gottes, auch des Ehestandes II, 487. 1264. f., der Kirche I, 689. 695.; II, 438. f., der Christen II, 194. f., unser ärgster Feind I, 1689. im Leiblichen und Geistlichen II, 2013.; Glaube und Wort Gottes ist ihm eine Finsterniß II, 1056. f.; Gott spottete sein in der ersten Verheißung, die er nicht verstand I, 236. ff. 239.; sah Christi durch die Menschheit verdeckte Gottheit nicht II, 312.; er vermag nicht in ein gläubiges Herz zu sehen II, 1057.; hat wohl den Willen, nicht die Macht zu schaden I, 231.; II, 727. ff. 733. f.; Gott hindert ihn II, 2013.; Christus hindert ihn I, 231. ff., hat seinen Kopf, Tod und Sünde, zertreten I, 233. \gesperrt{Seine Wirkungen}: Was er \gesperrt{gethan}: hat aus Neid in der Schlange versteckt den Menschen verführt I, 99. f. 266. f., zu allen Zeiten versucht, unsere Natur zu entstellen und zu verwüsten II, 1175.; wüthete um der Verheißung willen vom Weibessamen gegen die Menschen im Allgemeinen und gegen Christi Vorfahren I, 236. ff.; tobte stets gegen die Kirche II, 407. f.; wüthete gegen Christum von der Krippe bis zum Kreuz I, 239. Was und wie er es \gesperrt{thut}: thut bisweilen Gutes, daß er desto größeren Schaden thue II, 1271.; hat seine eigenen Lehrer in allen Arten der Studien II, 1336.; hilft seinen Dienern bald, verzieht nicht wie Gott II, 8.; wüthet I, 245.; wirft uns überall Hindernisse in den Weg I, 1699. f.; richtet im Hausstand Unheil an II, 327. f.; ist Vater aller Kriege und Uneinigkeit I, 692.; reizt zur Unzucht II, 1227. f.; treibt von einer Sünde zur anderen I, 1226.; hat er erst den Kopf hinein, so folgt bald der ganze Leib nach II, 884.; die Sicheren macht er hart I, 1425.; wüthet jetzt gegen die Christen, weil er Christo nicht mehr schaden kann I, 239. 696. 1140. 1238.; streitet wider die guten Werke II, 702. f.; macht aus den allerbesten Werken der Christen Sünden II, 393.; macht unsere Last und Betrübniß immer größer, heißt deshalb Leviathan I, 1424. f.; II, 392. 967. f.; sucht uns durch anhaltende Anfechtungen zu ermüden II, 140. 220.; ist Gottes Affe, baut seine Capelle neben Gottes Kirche I, 1663. ff.; führt vom Worte ab I, 180. f.; verfälscht Gottes Wort I, 181. f., läßt uns, was ihm nicht paßt I, 131.; ist ein listiger Disputirer I, 1177.; wohnt auch in der Wüste I, 1374. \gesperrt{Christen} sollen der Teufelskirche feind sein I, 689., wider den Satan kämpfen II, 22. 202.; brauchen sich nicht vor ihm zu fürchten II, 569. — S. \gesperrt{Engel} 3, 4, \gesperrt{Feinde} 2, \gesperrt{Fragen, Hofteufel, Werke} IV.

Thamar. 1. \gesperrt{Judas Schwiegertochter}. II, 1177. ff.; Cananäerin II, 1177.; berühmt, Mutter Christi II, 1162.; fruchtbar II, 1180.; ein wunderbares Weib II, 1197.; durch die Blutschande eine große Sünderin II, 1191. ff.; sündigte nicht aus Geilheit II, 1189. ff. 1199. f. — S. \gesperrt{Juda} I, 3.

2. \gesperrt{Davids Tochter}. Trug einen weißen Kittel II, 1038. f.

Tharah. Von der Nimrodischen Rotte verführt I, 725.; durch Sem zurechtgebracht I, 725.; wann gestorben II, 1511. und wo I, 725. — S. \gesperrt{Nimrod}.

Tharsis. Soll Cilicien sein I, 661.

Theman. \gesperrt{Enkel} Esaus II, 994.; das \gesperrt{Land} T., ein Theil Idumäas II, 994.; die \gesperrt{Verständigen} zu T. waren Esaus Nachkommen I, 1032.

Themistocles. Konnte sich im Glück nicht beherrschen II, 1387.

Theologie. Darin übe man sich II, 646. ff.

Theseus und Hercules. Ihre Freundschaft ist nichts gegen die Brüderschaft in der Kirche I, 1147.

Theurung. Das Evangelium muß Schuld daran sein I, 784.; viele Kinder Gottes haben T. erfahren müssen II, 131. f. 1805.; Beispiele von T. I, 783. f.; Gott erhält die Seinen in der T. II, 1805.; Verheißungen der Gläubigen in der T. II, 135.; Isaaks Verhalten darin II, 131.; man soll in guten Jahren dafür sorgen II, 1360. f.

Thideal. König I, 878.

Thiere. 1. \gesperrt{Im Allgemeinen}. Unterschied zwischen: lebendigen T., T. und Vieh I, 521. f.; wurden gleich in großer Menge geschaffen I, 46.; waren alsbald da durch das Wort I, 67.; heißen Gottes Fußstapfen I, 83.; haben ein gewisses Vorgefühl für Unfälle I, 523. f., keine Erkenntniß ihres Schöpfers I, 81.; das Gesetz der 1. und 2. Tafel geht sie nichts an I, 502. f.; ob sie vor der Sintfluth so verderbt waren, daß sich die mancherlei Arten vermischt haben I, 502.; warum sie in der Sintfluth verderbt wurden I, 503.; der Unterschied zwischen reinen und unreinen T. findet sich schon bei Noah I, 1220.; ob wilde T. im Paradiese waren I, 522. f.; warum sie sich vor den Menschen fürchten I, 590. 592. f.; ob dies schon vor der Sintfluth der Fall war I, 589. f.; sie zu tödten, bekam der Mensch erst nach der Sintfluth das Recht I, 591.; vom Schlachten der T. I, 596. f. 601.; Gott bedient sich ihrer, die Sünde zu strafen I, 594. — S. \gesperrt{Herrschaft}.

2. \gesperrt{Die} T. \gesperrt{in der Arche}. I, 521. ff. Wie Noah die T. bezwingen konnte I, 523. f.; wozu das 7te reine Thier bestimmt war I, 540.

Thimna. II, 994.

Thimnath. Lage II, 1186.

Thiras. I, 657.; Bedeutung des Namens I, 659.

Thogarma. Völkername I, 661.

Thomas Aquinas. Meinung von Hölle und Fegfeuer II, 2067.

Thongruben. Was man da grub I, 881.

Thopheth. Thal II, 1550. ff.

Thor. „Zur Stadt T. aus und eingehen“ II, 884. f. — S. \gesperrt{Stadtthor}.

Thorheiten und Fehler sind bei Christen unvermeidlich II, 275.; sollen uns demüthigen II, 276.

Thubal. Völkername I, 659.

Thubalkain. Bedeutung des Namens I, 387.

Thun. „Thun, was an Einem ist“ II, 746. 764. 770. 819.

Thüringerland. Unfruchtbar um der Sünde willen I, 852.; ehedem fruchtbar II, 1794.

Thurmbau zu Babel. I, 684. ff. 694. ff. Worin die Sünde beim Bau bestand I, 685. ff.; wie sich die Frommen dabei verhalten I, 705.; daher stammt die Fabel von den Giganten, die Jupiter stürzen wollten I, 687.; Fabeln des gemeinen Pöbels davon I, 686. — S. \gesperrt{Verwirrung}.

Tobias. Sein Weib murrte im Anfang nicht II, 1326.

Töchter. Gerathen nach der Mutter II, 1415.

Tod. 1. \gesperrt{Im Allgemeinen}. T. und Zerstörung thut den Menschen und anderen Creaturen wehe II, 1713. f.; wie die \gesperrt{Vernunft} den Tod ansieht I, 406. \gesperrt{Urheber}: Gott liebt den T. nicht, der T. ist vom Teufel I, 602. \gesperrt{Ursache}: im Stande der Unschuld war kein Tod I, 425.; ist Strafe der Sünde I, 263. ff. 404. 425. 453.; ist Beweis, daß wir Sünder sind I, 404. f. \gesperrt{Schrecken} des T.’s: von Leib und Seele des Todes I, 1512. f.; der Schrecken des Todes ist der Tod der Seele und erschrecklicher als der T. des Leibes I, 1512. f.; der T. ist kein T., sondern ein Schlaf, wenn der Schrecken hinweg ist I, 1512. f. \gesperrt{Sieg}: der T. ist besiegt I, 427.; der T. mit seiner Macht ist verächtlich bei Gott, er spielt damit I, 1512. ff.; der Glaube überwindet den T. I, 1513. ff.; je stärker der Glaube, desto schwächer der T., und umgekehrt I, 424. f. T. \gesperrt{der Gläubigen und Ungläubigen} scheint gleich und ist doch verschieden I, 613. f.; kommt den Gottlosen unerwartet I, 544.; der T. der Heiligen ist sanft und stille I, 1754. f., ist ihnen ein Kinderspiel I, 1513. ff., ein Schlaf I, 424. f., ein seliger Wechsel I, 1756., eine Ursache zum Leben I, 614. f.; eine liebliche Beschreibung ihres T.’s ist „zu den Vätern fahren“ I, 964.; diese Lehre vom T. verstehen die Gottlosen nicht I, 1513. f. 1515. f.; der T. der Heiligen ist werth gehalten vor dem HErrn I, 348. f. 1754. f. \gesperrt{Verhalten} gegen den T.: der Sünder flieht den T. als eine Strafe I, 369.; Christen sollen sich nicht davor fürchten II, 1902. f., sich durch den gebenedeiten Samen und seinen T. dawider trösten I, 405. f. 1607. ff. und mit dem T. und Begräbniß der Patriarchen I, 1605. ff. 1609. 1765.; bei dem T. der Verwandten tröste dich mit dem ewigen Leben II, 969., bei ihrem T. in deiner Abwesenheit mit Abrahams Beispiel I, 1612.; im T. tröste dich II, 1150. f. damit, daß man nur durch den T. dem T. entgehen kann I, 615., mit der Auferstehung I, 428.; die Märtyrer haben den T. verspottet II, 1902. f.

2. T. \gesperrt{und Auferstehung} der Gläubigen ist wie Einschlafen und Erwachen II, 2071.

3. T. \gesperrt{und Leben} sind näher bei einander, denn wir sehen können II, 1901.

4. T. \gesperrt{Christi}. S. \gesperrt{Christus} 2.

Todesstrafe. Vor der Sintfluth nicht erlaubt I, 347. 598. 600., wäre ohne Gottes Verordnung unrecht I, 599.; warum nach der Sintfluth geordnet I, 598.; muß nach Gottes Verordnung verhängt werden I, 599.; was zu thun, wenn eine Schwangere zum Tode verurtheilt wird II, 1210. — S. \gesperrt{Todtschläger}.

Todte. Leben in Gottes Augen I, 963. f.; Trauer um sie ist recht I, 1612. ff.; II, 944. f. 1185.; Trauer soll ihr Maß haben I, 1615. f.; wie man sich verhalten soll, wenn sie in Todsünden dahinsterben I, 1616.; ihnen die Augen zuzudrücken, die letzte Wohlthat, kommt den Verwandten zu II, 1759. f.; man pflegt sie zu loben I, 1603.

Todtenbestattung. War ursprünglich Beerdigung, später
\quellseite{574}
Verbrennung I, 264.; Sarahs Begräbniß I, 1615. ff.; Todte soll man ehrlich begraben um der Hoffnung willen der Auferstehung I, 1625. 1756.; auch die Heiden hatten ehrliche Ceremonien dabei I, 1625. — S. \gesperrt{Begräbniß}, \gesperrt{Kirchhöfe}, \gesperrt{Väter} II, 12.

Todtschlag. Eine grausame Uebelthat I, 357. f.; warum verboten I, 600. ff.; auf Befehl der Obrigkeit vollbracht, ist ein gut Werk II, 763. — S. \gesperrt{Mord}.

Todtschläger. Haben ein böses Gewissen I, 351. 373. f.; verlieren die natürliche Lieblichkeit der Augen I, 374.; sollen wieder getödtet werden I, 600.; Strafe derselben zu Isaaks Zeit II, 345. f.; sollen sich in Simeon und Levi ein Exempel nehmen II, 1936. ff. — S. \gesperrt{Mörder}.

Traumdeuter. Irren oft II, 1311.

Träume. \gesperrt{Sind}, wenn im Schlafe dem Gemüthe Bilder vorkommen I, 930.; die \gesperrt{Schrift} verwirft bald die T., bald lobt sie sie II, 1047. f.; \gesperrt{es gibt} teuflische, natürliche und wahre oder prophetische oder göttliche II, 1049. ff. 1311. f.; von den \gesperrt{natürlichen} gilt: Somnia ne cures II, 1051.; \gesperrt{prophetische} sind nicht zu verachten II, 1051.; von den T., darin etwas Zukünftiges angezeigt wird II, 419. ff.; wie sich T. von anderen Erscheinungen unterscheiden II, 419.; sind Offenbarungen, eine Stufe niedriger als Gesichte I, 930.; wie man \gesperrt{melancholische und wahrhaftige} T. unterscheiden kann I, 420. ff.; manche kommen \gesperrt{von Gott}, andere nicht II, 1309.; der \gesperrt{Teufel} kann auch wahre T. verursachen I, 931.; II, 1049. ff. 1054. 1056. 1298. 1311.; T. zu unterscheiden, dazu gehört ein sonderlicher Verstand II, 1309.; woran man göttliche prophetische T. erkennen kann II, 1051. ff. 1056. 1308. ff.; die göttlichen machen einen tiefen Eindruck I, 931. 1361.; II, 660. 1052. 1335.; Gott gibt mit den T. zugleich Auslegung und Erfüllung II, 1052.; der Vater gibt, der Sohn erklärt, der Heil. Geist erfüllt die T. II, 1053.; man prüfe die T. nach Gottes Wort I, 1321. f.; II, 421. ff. 1049. 1057.; manche T. \gesperrt{betreffen} weltliche, andere kirchliche Angelegenheiten und die Seligkeit II, 419.; weltliche T. hatten Pharao und Nebucadnezar II, 420.; wann weltliche T. falsch und vergeblich sind II, 423.; Gottselige und Gottlose haben T., die in Erfüllung gehen II, 419. 1049. ff.; menschliche Auslegungen sind ungewiß II, 423. 1310. T. sind ein \gesperrt{Bild} unseres Lebens II, 1054. ff. — S. \gesperrt{Offenbarungen}.

Träumer. Sind nicht zu verachten II, 1057.

Traurigkeit. T. macht zuweilen schläfrig, zuweilen schlaflos II, 388. f.; tödtet oft II, 1669.; Gott ist ihr feind I, 1716.; ist ein Zeichen, daß Gott von uns gewichen I, 927. T. \gesperrt{des Heil. Geistes}, was das sei I, 454.; T. \gesperrt{des Herzens}, das schwerste Leiden II, 1684., läßt sich nicht ausreden I, 1103., da ist schwer helfen II, 1685. f., dawider soll Jeder helfen II, 1684. — S. \gesperrt{Gewissen} 2, 3.

Trauung. Ceremonien zu Luthers Zeit II, 1344. f.

Treue. S. \gesperrt{Barmherzigkeit} 3.

Trockenes. Definition des T. II, 1929.

Troglodyten. Rauben und stehlen I, 997.

Trost. Geht Erschrockenen schwer ein I, 604. f.; es ist ein großer T., wenn man des göttlichen Gebotes gewiß ist I, 1501. ff.; Christen bedürfen auch leiblichen T. I, 777. f. — S. \gesperrt{Anfechtung} 5, \gesperrt{Kreuz} 6.

Trösten. Die Elenden trösten ist Gottes eigenstes Werk II, 746.; er tröstet nur Trostes Bedürftige II, 1754. ff., auf mancherlei Weise I, 574.; trösten soll man sich mit den Exempeln der göttlichen Gnade I, 732. f.

Trübsal. Findet sich in allen Ständen II, 139.; Christen müssen durch viel T. in das Reich Gottes eingehen II, 236.; alle T. der Christen kommt vom Teufel I, 1238.; darin sehen wir Gott von hinten, darnach läßt er uns sein Antlitz sehen I, 1003. 1005.; T. macht grau und runzlicht II, 1799. f.; ist nöthig I, 245., nützlich I, 246. f. 784.; wie Gläubige die T. ansehen I, 739.; darin verliere man die Hoffnung nicht I, 1004.; wie man sich trösten soll I, 244. ff. 261. f. 263.; II, 139. f.; man soll sie durch den Glauben vergessen II, 1431. f. — S. \gesperrt{Kreuz}.

Trunken. Was das heißt II, 1623.

Trunkenheit, geistliche. Von der geistlichen T. der Gläubigen im Reich Christi II, 1980. ff.; ist eine T. des Glaubens II, 1983., von den Gaben des Heil. Geistes II, 2003., weit verschieden von der T. von Wein II, 1984. 1995.; ist nöthig II, 1983.; macht muthig II, 1985. 1994. f.; soll zur Zeit des Messias regieren II, 1984.; fand sich bei den Aposteln und Märtyrern II, 1981. ff. — S. \gesperrt{Freude} 2.

Tugenden. Was T. sei II, 1876.; Ehrerbietung, Zucht, Sittigkeit gefallen Gott wohl I, 1629. f.; T. erben nicht fort II, 1368. T. \gesperrt{der Heiden und Gottlosen und der Christen} muß man unterscheiden I, 582. f.; die Heiden hatten T. I, 484.; der Glanz und Schein der äußerlichen politischen T. gefällt den Leuten I, 581.; die T. der Heiden gefallen Gott nicht II, 461. f., sind nicht gut I, 582., sind Laster I, 585.; trotz aller T. ist der Mensch ohne den Heil. Geist einfach gottlos I, 484.; in wiefern die T. der Gottlosen noch T. sind I, 583. — S. \gesperrt{Abraham} 14, \gesperrt{Esau} 1, \gesperrt{Glaube} 10, \gesperrt{Isaak} 7, \gesperrt{Jakob} 7, \gesperrt{Joseph} 8, \gesperrt{Lea} 1, \gesperrt{Noah} 7, \gesperrt{Rahel}, \gesperrt{Rebekka}, \gesperrt{Sarah}, \gesperrt{Schönheit}, \gesperrt{Väter} II, 4, \gesperrt{Vaterlandsliebe}, \gesperrt{Wahrheit}.

Türken. 1. \gesperrt{Im Allgemeinen}. Sind Scythen, reden arabisch II, 14.; ehrliche, weise Leute II, 36. f.; halten strenge Zucht I, 989.; im Kriege grausam II, 886., tollkühn II, 176.; rühmen sich, daß sie nicht herrliche Gebäude, Festungen etc. bauen, und verachten deshalb Andere I, 658.; ihr Kaiser soll eine sechsfache Krone haben I, 1713.; treiben unbekannte, geringe Leute mit Knütteln vom Schloß des Sultans II, 1614. f.; schossen mit Pfeilen II, 2043.; Kleidung der Männer II, 1196., der Weiber II, 1188.; es geht ihnen wie allen Gottlosen wohl I, 667., haben großes Glück II, 50. ff.; siegreich, weil Deutschland nicht einig war I, 706.; erfolgreich, weil sie Zeit und Gelegenheit wahrnahmen II, 98.; haben die Saracenen überwunden I, 994.; hielten Leibeigene II, 1762.

2. \gesperrt{Religion} der T. Sie sind Simonisten, wollen durch Werke selig werden II, 110.; Fatalisten II, 176. 748. 1313.; haben Träume und Weissagungen vom Ende des Römischen Reichs II, 1053.; verlassen sich auf die fleischliche Geburt II, 36. ff.; meinen, ihr Glaube sei recht, weil sie so sieghaft sind I, 496.; überheben sich ihres Glücks, glauben nicht, daß Gott zürne I, 1196.; rühmen, wegen fleischlicher Vorzüge und Glücks, daß sie Gottes Volk seien I, 1400. f.; II, 36. ff. 72. f. 220. 618. Sind die rechte Kirche nicht, werden nicht selig I, 1432. f.; II, 38. f.; haben keine rechten Propheten II, 1053. f.; sind Kains und Hams Nachkommen I, 650.; haben nicht den rechten Gott I, 1551.; verspotten die Lehre von der Dreieinigkeit II, 213. ff.; haben keine Erkenntniß in geistlichen Dingen I, 1553.; haben selbsterwählte Gottesdienste II, 777.; keiner wird deshalb selig, weil er ein T. ist II, 67.; der T. ist das in der Offenbarung geweissagte Thier der Lästerung I, 1064.

3. \gesperrt{Verhalten} der T. gegen die Christen. Sie verachten die rechte Kirche I, 649.; verachten die Christen um der Taufe willen II, 1614.; halten sich für frömmer und gerechter denn die C. I, 894.; sind Feinde der C. II, 51.; wie sie über C., die in der Schlacht gefallen waren, gespottet II, 51.; machen den Himmel voll Märtyrer II, 53.; ihre Grausamkeit schmücken sie mit dem Namen der Religion I, 505.; auch ihnen hilft das Gebet der Kirche I, 1409.; Gott wird den T. zerstreuen I, 702. f.; sind eine Zuchtruthe für Deutschlands Sünde II, 880.

Vgl. \gesperrt{Antichrist}, \gesperrt{Pabst} 2.

Tyrannei. Folgt der Erbsünde I, 880.; schädlich II, 1950.; führt zu Sünden I, 820.; wie Könige hineingerathen I, 820.

Tyrannen. Sind in der Beschreibung Nimrods abgebildet I, 671.; fromme Fürsten werden Tyrannen II, 1406. Die T. \gesperrt{vor der Sintfluth} I, 473. ff., sind dem Pabst gleich I, 473. — S. \gesperrt{Gewaltige}, \gesperrt{Rephaim}.

Tyrus. Lage, wann erbaut, gewaltig und berühmt II, 2012. f.; berühmteste Handelsstadt Asiens I, 666. 678.

\section{U}

Uebel. Wahl zwischen zwei Uebeln I, 1230. f.

Ueberfluß. Gabe Gottes, von uns mißbraucht I, 856.

Uebersetzung. Man kann nicht überall den Nachdruck und die Figuren wiedergeben II, 100. Die \gesperrt{chaldäische} U. von 1 Mos. 49, 10. II, 1973. — S. \gesperrt{Siebenzig}, \gesperrt{Vulgata}.

Ulrich, St. Ernährte einen Bettler I, 1363. f.

Umgang. Den U. mit Menschen fliehe nicht I, 1502. f.

Unbegrenzte. Lehre vom U. I, 680. ff.

Unbußfertige. Gott straft sie trotz seiner Zusage und trotz ihrer Menge und Gewalt I, 553.

Unbußfertigkeit. Macht die Sünde groß II, 866. 872.

Undankbarkeit. Ein großes Laster II, 614. 1262. 1315.; ganz natürlich I, 978.; ist in des Teufels Reich I, 977.; wie undankbar wir Christen sind, verglichen mit den Patriarchen I, 431. f.; die Welt ist undankbar I, 979., wir dürfen nur Undank erwarten I, 977. 980.; über U. klagt man und wird unwillig I, 977., die Welt wird darüber zornig I, 979.; wie sich Christen über die U. der Welt trösten sollen II, 1237. f. 1243.; wir sollen U. ertragen lernen I, 979., uns dadurch am Wohlthun nicht hindern lassen I, 979. 1139. f.; den Undankbaren soll man Gutes thun I, 977. f.

Uneinigkeit. Ein babylonisches Unglück, die allerschädlichste Plage I, 706. f.

Unfriede. S. \gesperrt{Ungerechtigkeit}.

Unfruchtbarkeit. U. der Bäume und Kräuter ist erst nach dem Fall eingetreten I, 46. — S. \gesperrt{Ehestand} 12.

Ungeduld. Regt sich, wenn Gott mit seinen Verheißungen verzieht II, 377.; der U. soll man sich schämen II, 1233. f.; wir werden uns am jüngsten Tage über unsere U. strafen II, 1329.; nützt nichts II, 1238.; durch U. im Kreuz schaden wir uns selbst II, 1326. f.; Trost wider U. II, 378.

Ungehorsam. Allgemeine Klage über U. I, 1726.; Strafe desselben I, 1726. f. — S. \gesperrt{Eltern} 2.

Ungerechtigkeit und Unfriede küssen sich II, 1811.

Ungeziefer. Ob im Anfange geschaffen I, 65. 77. f. 94. 255.

Unglaube. Der Ursprung aller Sünden I, 180. f. 182. 199.; die allerschwerste Sünde I, 1552.; U. straft Gott Lügen I, 1550. 1552. ff.; kommt aus der Erbsünde I, 1551. f.; steckt in Allen I, 1550. 1552.; dawider haben auch die Christen zu kämpfen und sollen sich deshalb demüthigen I, 1554. ff.; II, 146. f. 201.; Strafe desselben I, 884.; die rechte Erlösung davon geschieht im Tode I, 828. f.

Ungläubige. Werden stolz, selbstgerecht, suchen sonderliche Werke II, 1641. — S. \gesperrt{Gottlose}, \gesperrt{Welt} 2, 3.

Unglück. \gesperrt{Christen} müssen viel U. leiden I, 360. f. 811. 1416.; daran stoßen sie sich I, 361. \gesperrt{Zweck und Nutzen}: Christen sollen dadurch bewähret werden I, 822. f., geprüft I, 822., in der Demuth erhalten I, 823., gebessert I, 824.; dient zu Gottes Ehre I, 824.; kommt zur Versuchung der Frommen und Strafe der Gottlosen II, 728.; Unglücksfälle über ganze Länder schreiben Weltweise natürlichen Ursachen zu, es sind aber göttliche Strafen I, 1185. \gesperrt{Gott} nimmt sich der Seinen im U. an I, 806. 818. 826. f.; ist bei ihnen II, 1573. \gesperrt{Verhalten}: im U. und Gefahr macht uns das Gewissen angst und bange II, 1215.; die Gottlosen sind verzagt I, 895.; darin fliehe man zu Gott II, 1599. f.; man fürchte sich nicht, sondern glaube der Verheißung II, 933. und sehe auf die Hoffnung des ewigen Lebens I, 649.; zur Geduld reizen Gottes Wohlthaten I, 1338. f.; dazu haben auch die Heiden ermahnt II, 1421.; U. soll man sich nicht selbst aufladen, wenn Gott es schickt, geduldig tragen II, 1300. — S. \gesperrt{Glück} 2.

Unrecht. Man soll es tragen, sich nicht selbst rächen II, 207. f., tragen, aber nicht billigen I, 986.

Unschuld. Derselben können wir uns vor Gott nicht rühmen II, 1666.

Unstät und flüchtig. Was das heißt I, 366. ff.; ist, wer keine göttliche Verheißung hat I, 367. f.; sind die Seelen im Pabstthum I, 367.

Unsterblichkeit. Beweise dafür aus den Fähigkeiten der Seele I, 53. 55., aus der Gestalt des Menschen I, 151. f. Die U. ist angezeigt durch die Ruhe Gottes I, 97. f., durch die Weise, wie der Mensch geschaffen worden I, 102., durch die Worte: der Mensch ward eine lebendige Seele I, 105., durch die Worte: zu seinem Volk versammelt werden II, 15. Wir sind unsterblich, sonst würde Gott nicht mit uns reden II, 218. f., sich nicht den Gott Abrahams nennen II, 404.; die U. ist in Henochs Geschichte gelehrt I, 403. 405. Die Vernunft kann die U. nicht fassen II, 16.; Heiden und Epicurer verlachen diese Lehre II, 15.; etliche Philosophen haben die U. gelehrt, andere sie verlacht I, 102., sie wußten nichts Gewisses davon I, 406. — S. \gesperrt{Auferstehung}, \gesperrt{Leben} 11.

Unterthanen. Sollen sich nach den Gesetzen ihres, nicht eines anderen Landes richten II, 1184. — S. \gesperrt{Obrigkeit} 11.

Unveränderlichkeit. Gott ist unveränderlich II, 1677., nicht wandelbar und zweideutig II, 1534.; ändert oder widerruft sein Urtheil nicht II, 292.

Unwissenheit. Ist zweierlei: des Rechts und der That I, 814.; U. des Rechts entschuldigt nicht I, 814.; was U. der That sei I, 814. f. Eintheilung der Theologen in unüberwindliche, faule, affectirte U. I, 815. ff., in grobe und unüberwindliche U. II, 504.; die unüberwindliche U. entschuldigt in geistlichen Dingen nicht I, 815. ff.; II, 504. U. der Obrigkeit darüber, wo Unrecht geschieht, ist Sünde I, 1466. ff. U. gebiert viele Meinungen II, 2035.

Unzucht. \gesperrt{Folgt}, wenn man die erste Tafel verwirft I, 445.; wann sie sich im Menschen zu regen pflegt II, 20. f. \gesperrt{Reizung} dazu: damit wird insonderheit die Jugend angefochten II, 1252. f.; je gesünder der Leib, desto größer diese Anfechtung II, 1253. f.; der Teufel reizt
\quellseite{575}
dazu II, 1265., besonders die Jugend II, 1227. f.; wenn er durchs Kreuz nichts ausrichten kann, sucht er die Christen durch U. zu fällen II, 1259. und bietet ihnen die Gelegenheit, welche die Gottlosen suchen II, 1260. \gesperrt{Strafen und Folgen}: U. bleibt nicht ungestraft II, 863.; durch U. wird der Glaube verletzt II, 1272.; aus U. erwächst allerlei Noth und Unglück und andere Sünden I, 444. f.; vor U. Rasende pflegen das weibliche Geschlecht zu verachten I, 1228. \gesperrt{Verhalten}: die Heiden sind der U. nachgegangen, die heil. Väter nicht I, 1080.; dawider kämpfen ist eine Marter II, 2124., schwerer als der Kampf mit Kreuz II, 1260.; wie man dawider kämpfen soll I, 1666. f.; II, 20. ff. 24. 1276.; ihr kann man nicht dadurch wehren, wenn man kein Weib ansieht I, 832. f. — S. \gesperrt{Lüste}.

Ur in Chaldäa. Ursprung und Bedeutung des Namens I, 727. 729.; ob Name eines Ortes oder Abgottes I, 725. ff.

Uriel. Ableitung des Namens I, 1013.

Urtheil. Man fälle es nicht vorschnell I, 348. 1196. f.; man höre beide Theile II, 1279.

Uz. Ableitung des Namens, Vater der Syrer I, 1601.

\section{V}

Valla, Laurentius. Meinung über die lateinische und griechische Sprache II, 1487.

Vater. Die Neigung zu Weib und Kind ist die allerstärkste I, 1411. 1417.; soll auch mit Gefahr seines Lebens die Ehre und Zucht seiner Töchter schützen I, 1230.; wie ein frommer V. zu seinen Töchtern und Schwiegersöhnen reden soll II, 709.; erzeigt seine Liebe durch Geschenke II, 709. f.; je verständiger und weiser Einer ist, desto mehr begehrt man seinen V. zu sehen II, 1797.

Väter, Patriarchen. I. \gesperrt{Die} V. \gesperrt{vor der Sintfluth und ihre Zeit}. Sie waren nicht gemeine Leute, sondern die allergrößten Helden nach Christo und Johannes, dem Täufer I, 407. \gesperrt{Glaube}: hofften und sehnten sich herzlich nach Wiedererstattung deß, was wir im Paradiese verloren haben I, 430. f.; hielten ihre geringen Offenbarungen werther denn wir unsere reicheren I, 432.; wie sie sich der Aufnahme Henochs getröstet I, 424. f. 428., um Abels und Henochs willen Verlangen nach dem ewigen Leben hatten I, 423. f. Sind treffliche Exempel der Keuschheit I, 434. Hatten auch \gesperrt{böse Neigungen} I, 426.; ihre Sünden werden uns zum Trost erzählt I, 628. f. 633. f. Was sie \gesperrt{erlitten}, wird am jüngsten Tage offenbar werden I, 407.; wurden verachtet, als sich die Menschen mehrten I, 416. Ihre \gesperrt{Sprüche und Werke} wurden mündlich fortgepflanzt, oder vielleicht von ihnen selbst aufgeschrieben I, 418. Sie \gesperrt{entschliefen} im Glauben der Auferstehung II, 2066. \gesperrt{Kirchliche Zustände}: ihre Zeit ist eine heilige und goldene zu nennen I, 416. f. 436. f.; die erste Welt war die beste und heiligste, dann ists immer ärger geworden I, 431. f.; viele graue Väter haben zusammen gelebt I, 407. f. 416. 467.; es gab auch einen Haufen der allerärgsten Leute I, 409., von den Sünden derselben I, 436. ff. 480. ff.; Hoffahrt war die Sünde der ersten Welt I, 438. 440. f.; es gab da wohl noch keine blutschänderische Ehen oder Sodomiterei I, 445. f.; es ward immer ärger zu Noahs Zeit I, 449.; um der Sünde willen verderbte Gott die Welt I, 436. f.; die Väter legten sich wie eine Mauer zwischen Gottes Zorn und die Welt I, 494. f.; ermahnten mit Einer Stimme zur Buße I, 461.; wie sich die Welt gegen ihre Predigt verhielt I, 453. f. 458. 461. ff. 495. ff.; die V. wurden von den Kainiten bedrängt I, 408.; sie lehrten ihre Kirche desto treulicher I, 408.; warnten vor den Kainiten I, 468.; wollten keine ehelichen Verbindungen mit den Kainiten I, 470. Wie sie ihre Töchter erzogen I, 469. Die \gesperrt{Menschen} waren mäßiger als jetzt I, 256. f. 415., gesünder I, 253.; ihr langes Leben war ein Stück des glücklichen Zustandes der ersten Welt I, 414.; Ursachen ihres langen Lebens I, 414. f. Die Männer \gesperrt{heiratheten} erst, wenn sie 80–100 Jahre alt waren II, 23. Die \gesperrt{Früchte} und Kräuter waren damals kräftiger als jetzt I, 415.

II. \gesperrt{Die} V. \gesperrt{nach der Sintfluth} u. \gesperrt{ihre Zeit}.

1. \gesperrt{Leben} u. \gesperrt{Historien}. Die Vernunft hält ihre Historien für lächerlich und verächtlich I, 2.; II, 122. ff.; fleischliche Menschen verstehen sie nicht II, 1323.; die Papisten verachten das Leben der V. als weltlich und fleischlich II, 459. 1096., doch ist das „heilige“ Leben der Mönche und Papisten damit nicht zu vergleichen II, 127. f. 458. f.; der V. Historien soll man rühmen II, 1567., fleißiger erwägen als andere II, 1570.; sie übertreffen alle anderen, da Gott mit den V. redet I, 860. ff. 901. 1054..; in ihren Historien achte man vornehmlich auf Gottes Wort I, 860. f. 1544.; II, 211. 390. 415. 417., das ist ihr rechter Schmuck II, 2., dadurch unterscheiden sie sich von anderen Historien I, 861.; II, 567.; dieses, die Verheißung von Christo, ist darin der wichtigste Schatz II, 1731.; die vornehmsten Stücke darin sind die Verheißungen, ihr Glaube, ihr Kreuz II, 906. 1097. f.; Gott hat auf dreierlei Weise mit ihnen geredet II, 904., sie wunderbar geführt II, 1452.; ihr Leben ist voll Wunderwerke II, 129. f.; ist eine große Wohlthat, daß ihre Historien in der heil. Schrift erhalten sind II, 1.; die V. werden uns zum Trost als gar geringe Menschen beschrieben II, 458.; ihr Leben sollen wir uns zum Exempel betrachten II, 1272.; man sehe auf ihren Glauben und Gottes Verheißung II, 1096.

2. \gesperrt{Ehre und Bedeutung}. Die V. waren große Lichter in der Welt II, 1705.; ihre rechte Würde ist, daß sie Gottes Wort hören und glauben und im Glauben und Leiden versucht werden II, 456. f.

3. \gesperrt{Glaube} der V. Die V. mußten zunehmen an Erkenntniß II, 79.; haben die Lehre von der Dreieinigkeit geglaubt I, 708. 1718.; II, 940. 1866. ff.; wußten, der Weibessame sollte einer Jungfrau Sohn sein II, 1971. f.; hatten eine große Erkenntniß von den zwei Naturen in Christo II, 2056.; glaubten an Christum als Gott den HErrn II, 1835.; glaubten die Lehre von den guten Engeln II, 725. f. und von der Auferstehung der Todten und dem ewigen Leben I, 405. 1755. f.; II, 1304. ff. 1836., verließen aber doch ihren irdischen Beruf und Güter nicht II, 1906.; waren weise, voll Heiligen Geistes II, 1771.; sie hatten einen großen, wunderbaren Glauben I, 1501. ff.; II, 371. 905. f. 1560. f. 1870. 1884., den wir betrachten sollen I, 814. und uns unsers schwachen G.’s schämen II, 148. f. 2057. 2062.; sie wurden auch schwach im G. II, 1445. ff. 1865.; sind durch Erkenntniß der Sünden gerettet worden II, 1519.; haben sich herzlich nach Christo gesehnt, an ihn geglaubt I, 1591.; II, 703. 1306., von ihm das ewige Leben gehofft II, 1304.; trösteten sich des Eides, den Gott Abraham geschworen I, 1549.; sie und wir haben einerlei Glauben vom ewigen Leben I, 1764.; sie glaubten an den zukünftigen, wir an den gekommenen Christus I, 1757. f. 2069.

4. \gesperrt{Tugenden} der V. Ihre Demuth und Geduld war groß II, 238.; waren gastfrei, barmherzig, wohlthätig I, 1142.; II, 130., hofften, Engel zu beherbergen I, 1145., hielten sich nicht werth, arme Brüder zu beherbergen I, 1146.; sind Exempel der Keuschheit I, 1173.; II, 489. f., obwohl die Papisten und Manichäer sie der Geilheit beschuldigen II, 556.

5. \gesperrt{Affecten} u. \gesperrt{Neigungen} der V. Hatten menschliche N. wie wir I, 1308. f.; II, 942. ff. 967. 1501., welche der Heil. Geist regierte I, 1614. f. und viel herzlicher waren als unsere I, 1613.

6. \gesperrt{Werke der} V. Von ihren Werken soll man behutsam urtheilen I, 955. 1232. 1294.; sie waren Menschen wie wir, was sie Gutes thaten, thaten sie durch den Heil. Geist II, 1170.; sie schienen im gemeinen Leben von den Heiden nicht verschieden II, 1568.; sie blieben in ihrem Beruf und suchten keine sonderlichen W. I, 1309.; II, 1098.; verachteten die natürlichen Mittel nicht II, 1774. f.; doch zuweilen waren sie von Andern gar verschieden und Wunderleute II, 1865. f.; die Papisten stoßen sich daran, daß der Heil. Geist von ihren geringen häuslichen Werken so viel erzählt II, 537. f.; warum ihre gemeinen W. beschrieben sind II, 1565. f. 1624. ff., es geschieht uns zum Trost II, 468. ff. 516. ff. 566. 568.; nicht alle ihre W. sind Exempel für uns I, 890. 1232. ff.; II, 506. 890.; wollen wir ihnen folgen, so müssen wir zusehen, ob wir auch gleichen Befehl haben I, 1496.; sind in ihren sonderlichen W. keine Exempel für uns I, 568. 954. 972. f.; II, 510. ff. 538. f.; wir sollen nicht ihre W. nachäffen und ihren Geist fahren lassen I, 1233. ff. 1668.

7. \gesperrt{Häusliches und eheliches Leben} der V. Sie zogen umher wie Nomaden II, 1186. f., wohnten in Hütten I, 1168.; ehrten den Ehestand um der Kinderzucht willen I, 1746. ff.; II, 1198.; sie eilten, Kinder zu zeugen I, 714.; die Männer heiratheten mit 40, die Weiber mit 30 Jahren II, 23. 238. 856. f.; die Wohnungen der Männer und Weiber waren getrennt I, 1719.; II, 499.; ihr eheliches Leben ein Spiegel und Uebung der Keuschheit und aller Tugenden I, 1173.; Unfruchtbarkeit galt für eine Schmach II, 577. f., Fruchtbarkeit für einen großen Segen II, 532.; es war wohl Sitte, daß unfruchtbare Weiber ihren Männern ihre Mägde zu Weibern gaben II, 532. und ihnen dieselben, wann sie wollten, nahmen II, 551. f. 555.; die V. gingen freundlich mit dem weiblichen Geschlecht um II, 935.; sie hatten schon die Leviratsehe II, 1178.; die väterliche Gewalt war groß II, 502. f.; lehrten ihre Kinder Gastfreiheit und Ehrerbietung I, 1184. f. 1214. 1708. f. 1717. 1720. f.; hatten die Leibeigenschaft II, 1762.

8. \gesperrt{Die Weiber} der V. Lebten nicht in Müßiggang und guten Tagen II, 1213.; was ihr Amt und Arbeit war II, 1213. f.

9. \gesperrt{Sünden und Schwachheiten} der V. Sie hatten alle die Erbsünde II, 1519.; hatten Schwachheiten und Gebrechen wie wir I, 1354. f. 1466. ff. 1471.; II, 140. 380. 942. und waren doch heilig II, 19. f.; hatten auch Sünden an sich I, 1325. f., wo man sie aber entschuldigen kann, soll man es thun I, 1234. f.; die heiligsten hatten mit der Brunst des Fleisches zu kämpfen II, 530.; Hurerei war bei Leibesstrafe verboten II, 869., wurde im priesterlichen Geschlecht härter als sonst gestraft II, 1209. Ihre Sünden offenbaren uns Gottes Liebe und Barmherzigkeit II, 1185.; werden uns zum Trost erzählt I, 628. 723. f. 354. f.; II, 151. ff. 441. ff. 805. 843. 942. 1155. 1191. 1212. f., auch zu unserer Warnung und Ruhm göttlicher Gnade II, 1167. ff. 1200. Von ihren S. soll man behutsam urtheilen I, 1294., ihre Schwachheiten nicht rechtfertigen II, 445., ihre S. uns zum Trost groß machen II, 1041. f.

10. \gesperrt{Anfechtungen und Leiden} der V. Gott hat sie regiert und geführt wie uns II, 1077. f.; sie waren rechte heilige Märtyrer II, 983.; waren immer mit Anfechtung und Betrübniß beladen II, 49.; ihre Noth und Anfechtung war größer als unsere I, 716.; II, 196.; ihre Leiden mit Christi Leiden verglichen II, 345.; brauchten Glück und Unglück in Gottesfurcht II, 1628. f.; in der A. sieht man ihren Glauben II, 1098.; wurden wohl in der A. müde, aber durch Gottes Wort wieder gestärkt II, 620.; ihre Exempel sollen uns zu Trost dienen und Geduld erwecken II, 374. 900. f. 967. f.

11. \gesperrt{Das öffentliche Predigtamt und kirchliche Zustände}. Es gab herrliche und große Männer, die ein sonderliches Verlangen nach Christo hatten I, 431. f.; Sem, Salah und Eber regierten die Kirche II, 55. f.; die V. hielten das Predigtamt in hohen Ehren I, 1199. Sie hatten Macht, Segen zu ertheilen, nicht bloß zu wünschen II, 299. f.; mit welchen Ceremonien sie den Segen ertheilt und empfangen haben II, 260. f. 268. f. 280. ff. 293. 1860.; sie hatten besondere priesterliche Kleider II, 280. f. 293. f.; was sie für Kirchen hatten II, 88.; die Altäre waren im Freien I, 348.; hatten Gottesdienste überall I, 1307., allezeit I, 669., wobei sie lehrten und beteten II, 958. 1203. Wovon sie gepredigt II, 1750. f.; predigten vom Reich Christi II, 1994., von dem Weibessamen II, 1971. f. 2056., von Christo als Gott dem HErrn II, 1835 f. Unterschied des Gesetzes Moses und der Väter I, 144.; sie suchten Viele zur Erkenntniß Gottes zu bringen I, 877., bekehrten Viele II, 908. 1829.; wie sich die reine Lehre von Sems Zeiten ausgebreitet II, 1015. f. Sie lebten in einem seligen Zustande, der durch Ham und Nimrod gestört wurde I, 684. f.; sie prophezeiten die Zerstreuung der Gottlosen, wie diese Drohung von diesen angenommen worden I, 696. ff.; Zustand der Kirche zu Abrahams Zeit I, 1368., zu Jakobs Zeit II, 456. f.

12. \gesperrt{Tod und Begräbniß}. Die V. sind gestorben und begraben worden wie andere Menschen I, 1605. ff.; sie entschliefen in der Hoffnung der Auferstehung II, 2066.; die Anverwandten drückten den Todten die Augen zu II, 1759. ff.; man begrub, wo man wollte II, 1634.; warum die V. in Palästina begraben sein wollten I, 1606.; II, 1836. 2064. 2066.; haben die Gräber in Ehren gehalten II, 971.; ihr Tod ermuntert zu Geduld I, 1609.

13. \gesperrt{Die Leute zur Zeit} der V. Waren mäßiger, kräftiger und schöner als jetzt I, 803. 888. 1312.; II, 156. f. 489. 856. 1063.; heiratheten zu Jakobs Zeit mit 12 Jahren II, 1665. ff.; wurden in kurzer Zeit arg I, 898. f.; von den Sünden dieser Zeit I, 436. 445.; wie man Todtschläger strafte II, 345. f.; Gott hat ihrer nicht verschont I, 436.; die Leute waren nicht so gar grob und wild, hatten ehrliche Sitten und gute Gesetze II, 869.; Kleidung II, 1187. f. 1196.; wie man Hurerei im priesterlichen Geschlecht strafte II, 1208. f.; Unfruchtbarkeit eine Schande II, 1197.; man trug Leid um ein gestorbenes Eheweib II, 1185.; Wittwen konnte man leicht verstoßen II, 1183.; die Tage der Schafschur waren Feiertage II, 1185.; Könige trieben Viehzucht II, 1213. f.; der Handel war Tauschhandel II, 1196. f. 1213. f.; höfliche Anreden und Begrüßungen II, 478. 737. 821. 832.; man wohnte in Hütten II, 1734.; wie man reiste II, 841. 1986.

Vgl. \gesperrt{Fußwaschen}, \gesperrt{Gläubige} 2.

\quellseite{576}
Vaterlandsliebe. Sonderliche Tugend I, 581. f.

Vater Unser. Wie man denken soll, wenn man es betet II, 64.; in jeder Bitte ist Himmel und Erde begriffen II, 62.; Ein V.\,U. kann die ganze Welt nicht bezahlen I, 1363. f.; „Vater“ enthält Danksagung I, 1268.; 1. Bitte ausgelegt I, 1366.; 2. Bitte ausgelegt I, 1366.; wenn wir wüßten, was die 4. Bitte enthält, würde sich Keiner unterstehen, sie zu sprechen I, 1112.; „täglich Brod“ II, 296. 1243.; 5. Bitte: wir bedürfen täglich Vergebung I, 1347.; „Schuld“ II, 1497. Die Papisten verstehen das V.\,U. nicht, können es nicht recht beten I, 1366.

Venus. Wie abgemalt I, 248. 1161.

Verachtete. Auf Erden ist nichts verachteter, denn die Christen II, 567.; es ist besser, verachtet und betrübt sein, wenn Gott lacht, als Glück haben, wenn Gott zürnt II, 526.; Gott nimmt sich der V. an II, 518. ff. 522.; wie sich V. trösten sollen II, 567. f.

Verachtung. 1. \gesperrt{Im Allgemeinen}. Zu Verachtung des Alltäglichen sind wir geneigt I, 1118.; V. der Eltern, des Predigtamts, der Obrigkeit ist ein gewisses Zeichen des Zorns und Unglücks I, 441. — S. \gesperrt{Eltern} 2.

2. \gesperrt{V. des Worts und Predigtamts}. Die V. \gesperrt{des Worts} darf uns nicht wundern I, 1307. f.; das Wort würde verachtet, wenn auch Könige und Engel predigten I, 1335. f.; der \gesperrt{angebliche Grund} ist, daß nichts Gutes aus dem Wort folgt I, 1336. f.; V. kommt vom Fleisch und Teufel I, 432., daher, weil man das Wort für Menschenwort hält II, 148. f., weil Gott durch Menschen redet, da wir doch ihm grade dafür danken sollten II, 905.; V. und Abfall tritt gewöhnlich nach 20 Jahren ein I, 1635. f.; je reicher und heller das Wort, desto größer die V. II, 1941. \gesperrt{Folgen} der V.: bei solcher V. ist weder Rath noch Hülfe zu hoffen I, 456. f.; daraus folgt: gottlos Wesen und Sicherheit II, 98. f., allerlei Sünden I, 470. 504.; Zucht und Ehrbarkeit liegt darnieder I, 1192. \gesperrt{Strafe}: V. erweckt Gottes Zorn I, 1182.; Gott straft sie schwer I, 869. 1291. 1681.; II, 134. 881. f., mit allerlei Unglück I, 863. ff., mit Theurung etc. II, 132. 1805., falscher Lehre II, 212. 1832., Wegnehmen seines Worts I, 460. 1333.; wenn Gott sein Wort wegnimmt, ist sein Zorn schrecklich I, 455., dann kann man es nicht leichtlich wieder finden II, 318., dann folgt Frevel I, 504. Wie die V. an den Juden gestraft worden II, 319. f. und im Pabstthum I, 1528. — S. \gesperrt{Evangelium} 1, \gesperrt{Krieg}, \gesperrt{Landplagen}.

Verbrecher. Wie man zu Tode verurtheilte V. trösten soll II, 317.; Gott vergibt ihnen die Sünde, doch sollen sie von der Obrigkeit gestraft werden, Anderen zum Exempel II, 1921. f.; nach dem Tode fühlen sie die Schmach der Strafe nicht mehr II, 1921.

Verbum nominale. Was das sei II, 1872.

Verdammniß. Ist zweierlei: V. des Zorns zum Tode und der Gnade zum Leben II, 1106. ff.

Verderbt sein. Was das sei I, 504.

Verdienst. S. \gesperrt{Gnade} 3.

Vereinigung der Gläubigen und der Kirche mit Christo II, 402. f. 1998. f.

Verfälschungen der Schrift. Die Papisten lesen 1 Mos. 3, 15. „dieselbe“ statt „derselbe“ I, 225. 234.

Verfolger. Kain ein Vater aller V., die die Heiligen tödten I, 343.; Gott wird sie strafen, wenn sie nicht um Vergebung bitten I, 354.

Verfolgte. Wie zu trösten I, 1140.

Verfolgungen. Christen müssen viel V. leiden I, 1136.; V. muß leiden, wer gottselig leben will I, 1393.; Gott vergißt der Christen darin nicht I, 352. ff.; darin soll man fleißig beten II, 60., an der Lehre festhalten II, 1698. f.; wie man sich in V. wegen der Lehre trösten soll II, 1693. f., wie, wenn man um des Glaubens willen das Leben verliert II, 53. f. 1699.; über V. geht die Kirche nicht unter II, 1698. f.

Verführte. Mit den Verführern gestraft I, 765.

Vergeben sollen wir einander, aber die Strafe nachlassen thut nicht allezeit noth II, 1747. f. 1920. ff. 1934.

Vergebung der Sünden. Gott will gern vergeben I, 460. ff.; sie macht, daß die Sünde nicht verdammt, zerstört aber die Sünde nicht in uns I, 1670.; II, 1457.; hebt die leibliche Strafe und Züchtigung nicht auf II, 974. 1748. 1920. ff. 1934., oder das böse Gewissen II, 1116.; V. erlangt man nicht, wenn man die Sünde leugnet oder entschuldigt I, 336., sondern durch den Glauben II, 2059.; ein Christ soll der V. gewiß sein II, 1353. f.; über V. d. S. sollten wir uns am allermeisten freuen II, 1988. f.; wer V. hat, soll Gottes Gnade rühmen, damit auch Andere selig werden II, 1711.; die Lehre von der V. ist dem natürlichen Menschen eine Thorheit II, 1525. — S. \gesperrt{Gewißheit}.

Vergessen. Ob \gesperrt{Gott} seiner Heiligen vergessen kann I, 557. 560.; manchmal scheint es wohl so I, 556. ff.; dann ist es schwer, zu beten I, 557. f. \gesperrt{Wir} vergessen leicht Böses und Gutes I, 271., auch die Sünde und bedürfen der Gedenkzeichen I, 283.

Vergraben. Geld und Anderes v. ist eine alte Gewohnheit II, 1597.

Verhältniß. S. \gesperrt{Proportion}.

Verheißungen. 1. \gesperrt{Art und Natur}. Kennzeichen der \gesperrt{göttlichen} und \gesperrt{teuflischen} V. I, 757.; V. des Teufels sind scheinbarlich, gehen mit Lust ein, machen sicher und frech I, 756.; Mahomets und des Pabsts V. sind teuflisch I, 757. V. sind zweierlei: die \gesperrt{passive}, die uns geschieht, die \gesperrt{active}, die wir im Glauben annehmen I, 1399. Die \gesperrt{leiblichen} und zeitlichen V. fassen in sich die \gesperrt{geistlichen} und ewigen I, 1097. f.; II, 371. 412. f. 1888., wenn sie nicht über unsere Werke gethan und wenn sie ohne Bedingung sind I, 1106. f.; die V. in der Schrift gehören nicht zu diesem elenden und armen Leben I, 916.; die leiblichen V. sind bedingt, die geistlichen nicht II, 1887. f.; die leiblichen V., die an die Bedingung unseres Gehorsams geknüpft sind, sind nur leiblich I, 1107.; die geistlichen V. begreifen auch die leiblichen in sich II, 294. f. Sind zweierlei: des \gesperrt{Gesetzes} und der \gesperrt{Gnade}, Unterschied beider I, 371. f.; die des Gesetzes sind bedingt, die der Gnade unbedingt II, 736. f.; V. des Gesetzes sind besser als gar keine I, 375.; die V. von Vergebung der Sünde und Versöhnung mit Gott ist die höchste, ein neues Licht I, 932.; die V. \gesperrt{von Christo} werden immer klarer im A.\,T. I, 1008. 1092. 1526. 1585.; wir haben im Grunde dieselben V. wie die \gesperrt{Väter} I, 1385. f., in Taufe und Abendmahl I, 1385.; worin wir Abraham in den V. voraus sind I, 1385. f. V. sind fest und unbeweglich II, 374.; was wir in der V. haben, ist uns gewisser, als was wir in Händen haben II, 1869. f. 1884. f. 1893. f.; Gott bestätigt sie mit einem Eid I, 1548.; werden wiederholt, wenn etwas Neues dazu kommen soll I, 1585.; wie sich V. und Drohungen mit einander reimen I, 551. f. Wenn wir eine V. haben, dann sind auch die Engel bei uns II, 200.

2. \gesperrt{Erfüllung}. Gott hält fester an seinen V., als wir sie im Glauben ergreifen II, 1145. f.; er hält sie gewiß I, 1342.; II, 1916., wenn er sie auch zu vergessen scheint II, 1091.; es ist Gottes Weise, daß er mit der Erfüllung verzieht II, 29. 375. ff.; warum er verzieht II, 11. 377. 1915. f.; die geistlichen V. werden sehr langsam erfüllt II, 1010. ff.; Gott erfüllt die V. oft anders, als die Menschen begehren I, 551. f.; oft scheint das Gegentheil vom Verheißenen zu geschehen II, 417. f. 1755. ff. 1914.; Gott erfüllt sie, wie er will II, 1755. ff., nicht ohne Mittel II, 444. f.; verzieht, um nachher desto reichlicher zu erfüllen II, 1917.; wie Gott die Verheißung von der Herrschaft des Menschen über die Erde durch die Sintfluth nicht aufgehoben I, 552. f., die V. von Jerusalem und Erhaltung des Tempels erfüllt I, 551., die V. von der Erhaltung der Kirche erfüllt I, 551.; die Erfüllung der V. genießen auch solche, denen die V. nicht gegeben sind I, 768.

3. \gesperrt{Verhalten}. \gesperrt{Satan} verkehrt die V. durch seine Dialektik I, 935. Dem \gesperrt{Fleisch} gefallen nicht die göttlichen, sondern die teuflischen V. I, 757. f.; es wird ungeduldig, wenn Gott mit der Erfüllung verzieht II, 377. f. Die \gesperrt{Juden} verstehen die V. des A.\,T.’s nicht I, 1097. Die \gesperrt{Gottlosen} reißen die göttlichen V. an sich I, 936. Die V. \gesperrt{fordern} Demuth, durch Hoffahrt verscherzt man sie II, 243.; die V. soll man groß machen II, 1894., die V. des Evangeliums groß achten II, 1098.; V. und Glaube sind von Natur bei einander I, 755.; der V. theilhaftig zu werden, muß man sie glauben I, 647. 1030. f. 1384. f.; auch die leiblichen V. fordern Glauben II, 2013., welcher den Glauben an die geistlichen V. in sich schließt II, 2013. f. 2015. 2029.; auf Gottes V. sollen wir uns verlassen I, 792.; die V. sollen wir auch im Unglück glauben II, 200. f. und uns nicht fürchten II, 933.; halten wir daran, so überwinden wir Gott II, 796. 812.; bei Gottes V. muß der Glaube im Kampfe stehen I, 756.; man soll ihnen glauben, wenn auch das Gegentheil zu geschehen scheint II, 353. 1962. f.; wenn Gott verzieht, sollen wir Geduld haben II, 28. 1916. f.; sind wir geduldig, bis Menschen ihr Versprechen erfüllen, warum nicht auch, wenn Gott verzieht II, 377. f.; wenn man auch eine V. hat, soll man doch nicht schlafen, sondern wacker sein und arbeiten II, 467. — S. \gesperrt{Mittel}.

Vgl. \gesperrt{Abraham} 8, \gesperrt{Christus} 4, \gesperrt{Drohungen} 2, \gesperrt{Isaak} 4, \gesperrt{Strafe} 2.

Verkehrte. Bei den V. ist Gott verkehrt II, 891.

Verklärung. Wir können jetzt von der künftigen V. nichts Gewisses sagen I, 37.; die künftige Klarheit ist heller als das Licht der Sonne I, 24.; V. wäre auch im Stande der Unschuld auf das leibliche und zeitliche Leben erfolgt I, 79. f. 105.

Verlassen. Es ist die schwerste Anfechtung, wenn man denkt, daß man von Gott verlassen sei I, 558.

Verleumdung. Eine große Sünde I, 794. f.

Verlöbnisse. 1. \gesperrt{Im Allgemeinen}. Erfordern Einwilligung der Eltern wie auch der zu Verlobenden I, 1730. 1740.; rechtes Verhalten der Eltern und Kinder, wenn diese sich verloben wollen I, 1641. ff.; Eltern sollen sorgen, daß ihre Kinder ehelich werden, doch ihre Gewalt nicht mißbrauchen I, 1647. 1651.; Kinder sollen den Rath und Willen der Eltern befolgen I, 1648. ff.; durch das V. wird die Braut schon das eheliche Weib des Bräutigams II, 493.

2. \gesperrt{Heimliche} V. V. ohne Wissen und Willen der Eltern streiten wider das 4. Gebot II, 1479., sind ein Unrecht an den Eltern I, 1650. f., schändlich und ein Teufelsbetrug I, 1728. ff., eine Todsünde II, 1104. ff., verwerflich II, 365. ff., nicht gültig I, 1641. ff. 1648. 1652., wider die kaiserlichen Rechte II, 367.; darauf folgen unglückliche Ehen I, 1651. f., Unglück und Sünde I, 1643. f.; II, 1104. ff. Die Canonisten lehren, sie seien gültig I, 1641. ff.; der Pabst vertheidigt sie II, 1479., hat durch Gestattung heimlicher V. den Kupplern und Hurenwirthen das Thor geöffnet I, 1648. — S. \gesperrt{Ehestand} 14.

Vermessenheit. 1. \gesperrt{Vermessenheit}. Eine besonders gefährliche Sünde I, 928.; die schwerste Anfechtung I, 924. f., welche die allerhöchsten Heiligen trifft I, 923. f. — S. \gesperrt{Lehrer}, \gesperrt{Selbstüberhebung}, \gesperrt{Selbstvertrauen}.

2. \gesperrt{V. und Verzweiflung}. Beides will Gott nicht II, 1168. f. 1212. 1481.; wir sollen die Mittelstraße gehen II, 1482. 2084. ff.

Vernunft. Ist \gesperrt{blind} um der Erbsünde willen I, 156.; geräth ohne Gottes Wort im Geistlichen auf ungereimte Gedanken I, 16.; ihre allerbesten \gesperrt{Gedanken von Gott} sind die allertiefste Finsterniß I, 484. f.; was bei Gott Ordnung ist, scheint ihr Unordnung I, 39.; sie kann den Zorn Gottes nicht glauben noch verstehen I, 509.; ärgert sich an den göttlichen Verheißungen I, 758.; kann Gott nicht finden I, 815. 1046. f., sucht ihn auf Irrwegen I, 1093.; kommt nicht weit mit der Betrachtung der \gesperrt{Creatur} I, 151.; hält die Welt für ewig I, 5. f. 150. f.; weiß nicht, daß die Welt durch Gottes Wort erhalten wird I, 43.; kennt die Lehre von den guten Engeln nicht II, 726.; weiß von des Menschen erstem Ursprung nichts I, 156. 680. f., kann sich keinen ersten Menschen denken I, 150. f., die Schriftlehre von des Menschen Ursprung ist ihr eine lächerliche Fabel I, 156.; versteht die Geburt des Menschen nicht als ein göttliches Werk II, 16.; faßt die Lehre von der Unsterblichkeit nicht II, 16.; ärgert sich an der \gesperrt{göttlichen Regierung} I, 855.; stößt sich an der Gottlosen Glück und der Frommen Unglück I, 807. f.; läßt sich nicht an Gottes Gebot genügen, sondern frägt warum? und wozu? I, 1260.; Glaube, Geduld und andere \gesperrt{christliche Tugenden} sind ihr verborgen II, 27.; urtheilt, man solle die Frommen lieben, die Bösen hassen I, 796. f.; begreift nicht, warum wir allen Menschen gleicherweise mißtrauen sollen I, 796.; verachtet die rechten \gesperrt{guten Werke}, hat Lust zu falschen I, 758. f.; schließt fälschlich: was Gott an Abraham gefällt, gefällt ihm auch an mir I, 779.; urtheilt falsch von irdischen Gütern I, 855.; verachtet die Historien der Patriarchen II, 2. 19. f. 122. ff.; kann viele \gesperrt{natürliche Dinge} nicht begreifen I, 193.; II, 216. f.; folgt man ihr, so verliert man Ehre und Ruhm der heil. Schrift und die Majestät des Schöpfers I, 158. Ist Gott \gesperrt{feind} I, 175. \gesperrt{Sie treibt zum Besten} aber nur politischer und bürgerlicher Weise I, 175.; der Juden und Sophisten Lehre hievon I, 484. \gesperrt{Gebrauch der} V.: in äußerlichen leiblichen Dingen sollen wir ihr folgen, nicht in geistlichen I, 1313. 1376.; sie versteht nicht mehr, denn was dem Leibe gut ist I, 484.; kann in leiblichen Nöthen rathen I, 786.; man soll sie in irdischen Dingen brauchen II, 746. f. 765. f. Man soll die V. \gesperrt{ertödten} II, 56., gefangen nehmen unter den Gehorsam des Glaubens I, 193. — S. \gesperrt{Mensch} 2.

Verräther. Wer der Obrigkeit ansagt, wo Unrecht begangen worden, ist kein V. I, 1470. ff.

Versammelt werden zu den Vätern nur Gottselige II, 2066.

Versehung. Was versehen ist, muß geschehen II, 747.; nach der V. sollen wir außerhalb des Wortes nicht forschen
\quellseite{577}
I, 1549. f.; II, 174. ff. 181. f.; Disputation davon ist thöricht II, 747., gefährlich I, 492. f.; man soll um der V. willen die von Gott geordneten Mittel nicht verachten, Gott will die V. durch Mittel hinausführen II, 340. f. 346. f. 444. f. 747. ff.; der Zweifel an der V. ist ein Werk Satans II, 180. f.; wie man diesen Zweifel meiden kann und soll II, 179. f. 182. f.; wie Manche die Lehre von der V. zu Sicherheit und Gleichgültigkeit mißbrauchen II, 174. ff. 444. f. — Siehe \gesperrt{Erwählung}, \gesperrt{Majestät}, \gesperrt{Mittel}.

Versöhnung. Gott wirkt sie unter Menschen durch seinen Geist und äußerliche Mittel II, 833.; Abbitte und Geschenke wirken V. unter Menschen II, 766. 828. 834.

Verstockung. Wann Gott mit V. straft II, 1065.

Versuchen. „\gesperrt{Gott versuchen}“ I, 791.; man soll Gott nicht v. durch falsches Vertrauen auf die Verheißung II, 895., dadurch, daß man sich in offenbare Gefahr begibt II, 383. 637. f., durch Verachtung und Versäumung der Mittel II, 346. ff. 444. f. 746. ff. 764. 1774.; Gott v. ist eine eben so schwere Sünde, wie ihm nicht glauben II, 347.

Versuchung. \gesperrt{Art und Natur}: der Ursprung aller V. ist, wenn die Vernunft ohne das Wort von Gott und seinem Wort urtheilen will I, 189.; \gesperrt{Gott} versucht die Menschen, z.\,B. den Abraham I, 1483. ff., wenn er sich anders stellt, als er sich in seiner Verheißung anbietet I, 1486. 1532. ff., durch Irrlehrer I, 1533. f.; wenn er uns versucht, so kämpft er mit uns II, 809.; \gesperrt{Satan} sucht uns vom Worte abzuführen oder dasselbe zu fälschen I, 179. f.; diese V. ist schwerer, als die V. zum Ehebruch I, 179. f. 182.; wenn Satan Gott lobt, will er dem Menschen die Kehle abschneiden I, 194.; gibt man Satan erst ein wenig nach und zweifelt an dem Wort, so wird er immer kühner I, 190. f., und man fällt in immer größere Sünden I, 191.; je weiter der Mensch vom Worte kommt, desto weiser und gelehrter läßt er sich dünken I, 196. f. \gesperrt{Zweck}: \gesperrt{Gott} versucht seine Heiligen, um Glaube und Liebe zu erwecken I, 177. 1534. f.; der \gesperrt{Satan} versucht die Menschen zum Bösen I, 1534. f., will uns vom rechten zum falschen Gott führen I, 187. \gesperrt{Verhalten}: man bitte Gott, daß er uns nicht in V. führe II, 1871.; man unterliegt, wenn man vorwitzig disputirt I, 192. f.; man halte fest an Wort und Verheißung I, 193. f. 1534. f.; so that Christus I, 182.; in der V. zu Sünde denke man an Gottes Zorn und Strafe II, 1263.; Gottes V. nimmt einen herrlichen \gesperrt{Ausgang} II, 1695. — S. \gesperrt{Anfechtung}.

Vertrauen. 1. \gesperrt{Auf Gott}. Gott hat V. auf ihn geboten I, 1500. f.; dadurch wird das erste Gebot erfüllt II, 8.; wir sollen ihm vertrauen, ob wir Rath wissen oder nicht II, 1481., grade dann, wenn Alles verloren scheint II, 743.; auf Gott den Schöpfer, nicht auf Wälle und Festungen sollen wir unser V. setzen II, 932. f.; Gott vertrauen soll man und doch des Berufes warten II, 1446. und die natürlichen Mittel nicht unbenutzt lassen II, 1443.; V. auf Gottes Güte läßt nicht zu Schanden werden I, 982. — S. \gesperrt{Mittel}, \gesperrt{Schöpfer}.

2. \gesperrt{Auf Menschen}. V. auf Menschen ist gefährlich und gottlos I, 797.; welches das rechte V. auf Menschen ist I, 798. f. — S. \gesperrt{Eheleute} 3.

Verwandte. Gemeiniglich untreues Gesinde II, 483. f.

Verwirrung der Sprachen. I, 690. ff. War die gewisseste und leichteste Weise, den Thurmbau zu hindern I, 704., das Werk Gottes, nicht der Engel I, 707., eine schreckliche Strafe I, 690. ff.; hat großen Schaden angerichtet: in allen Ständen I, 691. f. 706., in der Kirche I, 898.; an dieser Strafe leiden die Menschen noch I, 692. 706.; dieser Schade ist durch Christum im N.\,T. aufgehoben I, 691. f. — S. \gesperrt{Thurmbau}.

Verzagen sollen wir an uns selbst I, 438. — S. \gesperrt{Hoffnung}.

Verzug der Gnade. Was das sei I, 558. 926. f. — S. \gesperrt{Saumseligkeit}.

Verzweiflung. Eine Sünde wider das 1. Gebot I, 1452.; von Gott verdammt I, 1501.; damit werden schier Alle angefochten I, 1486.; was man in solcher Anfechtung thun soll I, 1486. f.; die Pein derer, die verzweifeln und nicht Buße thun, ist schrecklich I, 363.; V. verzehrt Leib und Seele II, 2087.; in der V. sprechen die Leute in abgebrochenen Sätzen II, 48.; darin ist dem Herzen nichts zu viel noch genug, daß es von Gottes Willen unterrichtet wird I, 539. — S. \gesperrt{Vermessenheit} 2.

Veste. Hat den Namen vom Ausbreiten I, 31.; ist zwischen die Wasser gestellt I, 32.; ward und besteht durch Gottes Wort I, 31. — S. \gesperrt{Firmament}.

Victor, Pabst. Bannte die orientalischen Kirchen I, 706.

Viehzucht. Erfordert großen Fleiß und Sorge II, 1565.; V. in Spanien II, 602.

Vincentius. Spottete der Marter II, 1989.; Luthers Urtheil über sein Buch Summa ang. I, 815.

Virgil. Handelt von Aeneas II, 1217. 1421.; Aussprüche über: den Ackerbau II, 1407., das Glück II, 1386. f. 1396.

Vögel. I, 59. f. Sind den Fischen nicht ungleich I, 59.; fruchtbar II, 1872.; wo sie zur Winterszeit hinkommen I, 63.

Volk. 1. \gesperrt{Im Allgemeinen}. Was ein „Volk“ ist I, 744. 750. 1053.; kein V. läßt sich mit dem jüdischen vergleichen I, 741. ff. 745. f.; es ist keins so rein, daß nicht etliche Gottlose darunter sein sollten II, 1616. 1823. f.; je heiliger das V., desto böser und ärger ist es II, 1941.; Vertreibung und Vertilgung eines V.’s ist eine Strafe der Sünde I, 660. f. — S. \gesperrt{Länder}.

2. \gesperrt{V. Gottes}. Ist nur, wer glaubt II, 2017.; gefällt auch in den geringsten Dingen Gott wohl II, 618.; der Glaube, daß man Gottes V. sei, ist nicht leicht II, 618.; wie man dessen gewiß werden kann II, 618. f.; erkennt man nicht am irdischen Glück I, 1400. f. — S. \gesperrt{Kirche}.

Völlerei. Ist Mißbrauch der Gaben Gottes I, 1715. f.; man soll V. meiden I, 1715. f. — S. \gesperrt{Gastmähler}.

Vorbilder im A.\,T. Die Gläubigen wußten, daß es nur V. waren II, 1852.

Vorfahren. Man soll gottseligen V. im Glauben, nicht bloß in Werken nachfolgen II, 1826.

Vormundschaft. Verdrießlich aber trefflich Werk II, 1182.

Vorsehung. Gottes V. für uns ist viel größer als alle unsere Sorgfalt und Bemühung I, 57.; woraus sie erhellt I, 30. f.; durch die göttliche V. sind die Heiden zum Evangelium berufen I, 654. — S. \gesperrt{Fürsorge}, \gesperrt{Gebet} 14.

Vorsorge. S. \gesperrt{Fürsorge}.

Vorwitz. Gott kann den V., da man nach den Ursachen der Gebote Gottes forscht, nicht leiden I, 1125. f.; solcher V. ist thöricht I, 1127., eine greuliche Sünde I, 1127. f.; dadurch haben wir das Paradies verloren I, 1128.; V. war Sauls Sünde und Untergang I, 1127.; darüber gerathen wir in Gotteslästerung I, 1128. f.; dadurch kann man seinen Beruf und Erwählung ungewiß machen I, 1126.; diese Sünde ist uns vom Satan angehängt I, 1085. f., er verleitet uns dazu I, 1126.; die Papisten lehren die vorwitzige Frage warum? I, 1127.; die rechte Antwort darauf ist: weil es Gott so gefällt und uns nützlich ist I, 1128. — S. \gesperrt{Fragen}.

Vulgata. Die gemeine lateinische Uebersetzung ist nicht von Hieronymus II, 1975.; übersetzt dunkel I, 515., fahrlässig I, 598., verkehrt I, 234. 250. 577. ff.; II, 1424. 1434. f. 1509. 1586. 1975.

\section{W}

Wachen wider die Sünde, nöthig I, 275. f.; II, 98. f.

Wachsen. S. \gesperrt{Erkenntniß} 5, \gesperrt{Lernen}.

Wahrhaftig. Gott ist w. II, 416. — S. \gesperrt{Wahrheit}.

Wahrheit. Ist unter allen Tugenden die schönste I, 581. Führt schlichte einfältige Worte I, 389.; muß uns lieber sein als Freunde I, 149; wider die erkannte W. streiten ist ein schrecklich Ding II, 1966. ff.

Wahrsagen. Dazu brauchte man Schlangen II, 583.

Wallfische. Allerlei große Fische im Meer I, 62.

Wandel. Das Urtheil, das man allein aus des Menschen Leben und Wandel hat, trügt I, 797.; worin ein \gesperrt{göttlicher} oder gottesfürchtiger W. besteht I, 418. f.; der W. der Christen sei aufrichtig vor Gott II, 259., vorsichtig um der Feinde willen II, 1029.

Wandeln. „Vor Gott wandeln“ I, 1011. 1034.

Wandern. Was das heißt I, 1352.

Warnen soll man Gefährdete II, 350. f.

Wasser. Ist über und unter der Erde I, 545. W. \gesperrt{über dem Himmel} I, 32. 36. ff., sind von unseren Wassern verschieden I, 40. f., brachen bei der Sintfluth herunter I, 544. ff.

Wegscheide. Wie die Hebräer sie nennen II, 1188. f.

Wehmütter. Sollen beherzt sein und die armen Mütter trösten II, 1217.

Wehrstand. Ist nöthig II, 458. f. — S. \gesperrt{Obrigkeit}.

Weiber. \gesperrt{Bestimmung}: sind nicht für öffentliche Aemter bestimmt I, 248.; II, 686. f.; sollen häuslich sein und die Hausarbeit besorgen I, 1160. f.; sind zur Kinderzucht geschaffen II, 485. f. 702.; begehren von Natur fruchtbar zu sein II, 1176.; W., die nicht schön sind, sind gemeiniglich am fruchtbarsten II, 485.; W., die ihre Zeit nicht haben, sind unfruchtbar I, 1169. \gesperrt{Tugenden und Schwachheiten}: ihre rechte Zierde sind gute Sitten II, 1251.; von der W. List II, 1272.; Zucht und Schamhaftigkeit der W. übertrifft alle Nonnenwerke I, 1161.; haben um der Neugierde und Geschwätzigkeit willen einen bösen Namen I, 1160. \gesperrt{Verachtung}: man soll sie nicht schmähen I, 469.; die Heiden nennen das W. ein Uebel oder Unglück I, 82., einen verstümmelten Mann oder Ungeheuer I, 85. f. \gesperrt{Gebrechen} haben die W., die aber durch das Gute übertroffen werden I, 1731. ff.; haben auch ihre Beschneidung I, 1078.; ihr Geschlecht ist nicht Ursache des Falls I, 469. \gesperrt{Ehre}: das W. ist ein Gebäude Gottes I, 160. ff.; II, 702., von keiner Ehre der menschlichen Natur ausgeschlossen I, 84., ein Miterbe der Gnade I, 84., auch zum ewigen Leben geschaffen I, 141. 143.; wurde im A.\,T. ohne Beschneidung selig I, 1078.; heißt Männin I, 167. f.; die Mutterehre ist hoch zu achten I, 247. f.; man sollte das W. vornehmlich lieben um der Kinderzucht willen II, 485.; W. durften deshalb in Rom Gold tragen II, 523. Die Erhaltung des weiblichen Geschlechts ist eine große \gesperrt{Wohlthat} I, 144. f. W. haben oft großer Königreiche Zerstörung \gesperrt{verursacht} I, 469. — S. \gesperrt{Geschlecht}, \gesperrt{Raub}.

Weibessame. Die \gesperrt{Verheißung} vom W. war besser, als wenn Adam hätte vom Baum des Lebens essen dürfen I, 280. f., ist eine Verheißung der Gnade I, 371. f.; wie sie immer klarer wurde I, 237. f. 645. 718. f.; II, 1971.; wie weit die Gläubigen des A.\,T.’s sie verstanden haben I, 236. f. Er zertritt der Schlange Kopf I, 349.; den Segen durch den W. erlangt man nur durch den Glauben I, 647.; sollte einer Jungfrau Sohn sein II, 1971. f., in Canaan geboren werden I, 719.; ist allein Christus, alle anderen Menschen sind Mannessame I, 237. — S. \gesperrt{Christus}.

Weichlinge. Worin ihre Sünde besteht II, 1181.

Weiden. Heißt in der Schrift oft regieren II, 2046.

Weihrauch. Ein Harz II, 1124.

Weihwasser. Hat keine Verheißung, ist Abgötterei und Mißbrauch des Namens Gottes I, 280.

Wein. Die alleredelste Frucht II, 1979.; macht fröhlich II, 1979. f.; verschiedene Wirkung bei verschiedenen Leuten I, 629.; daran gewöhnte man zu Luthers Zeit schon die Kinder I, 712.

Weisheit. 1. \gesperrt{Gottes} W. muß sich wie kein Schüler tadeln lassen I, 1737.; man kann sie nicht durch Speculation erkennen, muß sie erfahren II, 1383.

2. \gesperrt{Die persönliche} W. ist Christus II, 1295. f.

3. \gesperrt{Der Menschen} W. Die natürliche W. ist Thorheit II, 275. f.; unserer W. pflegt Gott zu spotten II, 1574.; der Welt W. macht Gott zu Thorheit II, 150.; die rechte W. ist ans Evangelium glauben I, 197.; die höchste W. ist Gott erkennen II, 1688.

Weissager. Sind nicht zu verachten II, 1057.

Weissagungen. S. \gesperrt{Prophezeiungen}.

Welt. 1. \gesperrt{Das Weltgebäude}. Verschiedene Meinungen vom Anfang der W. I, 5. f.; Vernunft hält sie für ewig I, 5., das ist ein gefährlicher Wahn I, 5.; Gott erhält sie durch sein Wort I, 42. f. 1480. f.; wird älter und schwächer I, 803. — S. \gesperrt{Natur}, \gesperrt{Schöpfung}.

2. \gesperrt{Die Ungläubigen}. Welt \gesperrt{sind} nicht Leute geringen Standes allein, sondern die Allerbesten I, 1338. Behält ihre Art allezeit I, 427., ist Fleisch I, 463.; ihre Bosheit ist groß II, 1136.; fällt vom Geist aufs Fleisch, sucht gute Tage, Ehre u.\,s.\,w. II, 81. f.; in ihr regiert Argwohn, Haß, Feindschaft I, 801.; ist undankbar I, 979.; II, 1314. ff.; schreibt alles Gute sich selbst zu zu ihrer eigenen Ehre II, 1634.; will sich nicht strafen lassen I, 458. 506.; verfolgt die, welche ihre Sünde strafen II, 1033.; glaubt nicht, daß Gott über die Sünde zürnt I, 1205. f.; je mehr geplagt, desto verstockter I, 254. f.; wird immer ärger, je näher das Ende kommt I, 264. 442. und wenn Gott Leute erweckt, die sie strafen I, 449.; wird immer sicherer und frecher I, 463., je näher das Ende I, 1243.; ist nicht werth, ein gut Werk zu thun II, 1135. Alles, was in der W. ist, ist nichts, Gottes Wort ist Alles II, 77. — S. \gesperrt{Kinder} 3, 4, \gesperrt{Mensch} 2.

3. \gesperrt{Christ und Welt}. \gesperrt{Christen} verachten die W. mit ihrem Reichthum und warten auf ein ander Leben I, 1401.; können den Umgang mit der W. nicht ganz meiden II, 385.; müssen wider die W. kämpfen II, 22.; betrüben sich über die Bosheit der W. I, 693. f., wenn sie Gottes Zorn über die Sünden der W. fühlen I, 492. f.; beten für die, so sie plagen I, 494.; sind wie eine Mauer gegen den Zorn Gottes über die W. I, 494.; gleichen dem Atlas und tragen das himmlische Gebäude I, 1289.; sind Anderen ein Segen und nützlich II, 483. 767. f. 894.; wo Gutes geschieht, geschiehts um ihretwillen II, 726.; um Eines Christen willen thut Gott oft einem ganzen Lande Gutes I, 1202.; II, 894. f. Die \gesperrt{Welt} versteht die Werke Gottes und der Christen
\quellseite{578}
nicht I, 1519. f.; die Welt frohlockt und spottet ihrer, während die Christen betrübt sind I, 493. f.; hält die Christen für närrisch I, 1241. f., für die Ursache alles Unglücks I, 786.; spürt ihren Sünden nach und offenbart sie I, 632. f.; ist nicht werth, daß die Christen ihr dienen II, 1280.; ist undankbar gegen sie II, 1242. f. 1280. f. — S. \gesperrt{Kain} 3, \gesperrt{Kirche} 11, 12.

Weltalter. Die 7 Alter der Welt I, 1088.

Weltflucht. Ist nicht christlich I, 979.

Weltkinder. S. \gesperrt{Kinder} 3, 4.

Werke. I. W. \gesperrt{Gottes. Sind} die W. aller drei Personen I, 71.; II, 1868.; sind entweder W. der Barmherzigkeit oder des Zorns I, 1190.; ein W. Gottes ist’s, daß er aus Nichts Alles macht II, 1430. f., aus dem Besten das Aergste und umgekehrt II, 1944., aus den ärgsten Sündern die besten Pfarrherren II, 1950. f., daß er aus bösen Sündern etwas Gutes erwählt II, 1952., daß er die Elenden tröstet II, 746.; seine W. \gesperrt{richtet} Gott \gesperrt{aus} durch die Engel I, 1540. f., durch Menschen I, 15., durch natürliche Mittel II, 1775.; sind so groß, daß wir sie \gesperrt{nicht fassen} noch glauben können II, 1712.; um sie recht zu erkennen, muß man auf Gottes Wort sehen I, 1373. f.; die Welt versteht sie nicht I, 1519.; fleischliche Menschen ärgern sich daran II, 11.; die W. Gottes, die nach der Ordnung der Natur geschehen, sind ebenso groß wie Wunderwerke, aber wir achten sie nicht, weil sie so gemein sind I, 1370. f.; die Herrlichkeit der W. Gottes erkennen wir nicht, weil so viel Verdruß, Plage und Aergerniß daran hängt I, 1371. ff.; auch die geringsten soll man nicht verachten II, 477.; Gottes W. soll man verkündigen, daß die Menschen Buße thun I, 1190. — S. \gesperrt{Gebet} 7, \gesperrt{Hülfe, Fürsorge, Regierung}.

II. W. \gesperrt{der Menschen. Werk und Wille.} Ein W. ist vollendet, wenn der Wille, es zu thun, vollendet ist II, 1115. 1945. f. 1947.

III. \gesperrt{Gute} W. 1. \gesperrt{Art und Natur} der guten W. Ueber die W. muß man nach Gottes Wort, nicht nach der Vernunft \gesperrt{urtheilen} I, 1163.; die W. sind darnach zu beurtheilen, ob sie \gesperrt{Gott geboten} hat I, 865. f. 1179. f.; man soll Alles thun nach Gottes Gebot I, 570., nichts ohne dasselbe vornehmen I, 567. ff. 1375. f.; ohne Gottes Gebot ist kein W. gut I, 1497.; II, 124. 211. f., man soll es verwerfen I, 1544. f.; gute W. sind nur, die Gott geboten hat I, 866. 1180., doch ist nicht immer ein ausdrückliches Gebot nöthig I, 868.; Gott muß das W. auch eben \gesperrt{dir} befohlen haben I, 762.; gute W. sind nur, die im \gesperrt{Glauben} und Gehorsam geschehen I, 1324.; W. sind nicht deshalb gut, weil vom Pabst oder der Kirche befohlen I, 867.; W., die nicht mit Gottes Wort übereinkommen, sind vom Teufel I, 695.; W., die Gott nicht geboten, sind Zauberei und Abgötterei I, 866. Gott sieht nicht auf die W., sondern aufs \gesperrt{Herz} II, 258.; aufs Herz kommts an, ob die W. gut oder böse sind I, 1473. 1493., das natürliche verderbt unsere W. I, 575. f.; W. sind Gott gefällig als Danksagungen, nicht wenn man etwas damit verdienen will II, 448. f. Gott sieht erst auf die \gesperrt{Person}, dann aufs W. I, 214. f. 315. f.; wenn die Person Gott gefällt, dann gefällt ihm auch das W. I, 214. f.; II, 461. und umgekehrt, wenn ihm die Person mißfällt, dann auch das W. I, 315. f.; Person muß gut sein, soll das W. gut sein I, 1018.; W. machen nicht gerecht, sondern die gerechte Person thut gerechte W. I, 1123.; damit fällt der Juden und Papisten Ruhm von ihren W. I, 316. W. \gesperrt{des Gesetzes} sind die allerhöchsten Tugenden, machen nicht gerecht I, 944. Die \gesperrt{gemeinen} bürgerlichen W. soll ein Christ nicht meiden I, 1477. f.; wer mit geringen W. umgeht, an dem hat Gott ebensowohl Gefallen, wie an dem, der hohe und große thut II, 518.; gemeine W. werden im Evangelio gerühmt I, 527.; auch die geringsten sind groß und heilig, wenn geboten I, 1497.; W. der Obrigkeit, der Eheleute sind die allergrößten I, 528.; die größten W. sind Beten und Danken II, 571. Durch den Schein der \gesperrt{selbsterwählten} W. sollen wir uns nicht berücken lassen, sondern aufs Wort sehen I, 867. f.; man hüte sich vor selbsterwählten W. I, 1181. 1376.; sind überflüssig II, 1627. f., verwerflich I, 867., schädlich II, 1628.; gefallen Gott nicht I, 1496.; II, 1625. f.; Gott zürnt ihnen I, 529.; II, 1628. Wer \gesperrt{gute} W. \gesperrt{thun kann:} nur ein Gerechter I, 316. f., wer sich ans Wort hält I, 1503.; W. \gesperrt{der Ungläubigen} sind nicht gut II, 542. f.; gefallen Gott nicht II, 1777., um des Unglaubens willen II, 472.; verdienen Belohnung in diesem Leben II, 472.; die W. der \gesperrt{Papisten} taugen nichts I, 774. 1311.; die W. der \gesperrt{Gläubigen}, wie gering und unfläthig sie sind, sind groß und herrlich II, 705. f., sind Wunderwerke I, 1520., sind alle vor Gott angeschrieben II, 1769., sind alle gut II, 706., gefallen alle Gott wohl I, 859. f. 953. f. 1710.; II, 772. 1577. 1777., auch ihr Schlafen und Wachen II, 1640. f., auch die geringen II, 847. 1742. Man soll alle W. mit Gebet \gesperrt{anfangen} II, 1340. — S. \gesperrt{Beruf} 5.

2. \gesperrt{Nothwendigkeit.} Wir sollen gute W. thun I, 1577.; II, 467.; sie sind nöthig II, 414.; die lutherische Kirche verbietet sie nicht I, 620.

3. \gesperrt{Feinde und Verächter} der guten W. Der \gesperrt{Teufel} streitet wider die guten W. II, 703. Die \gesperrt{Vernunft} verachtet die rechten, hat Lust zu falschen I, 758.; Vernunft und \gesperrt{Welt} versteht nicht, was rechte gute W. sind II, 705.; das menschliche \gesperrt{Herz} verachtet die gemeinen W. und wundert sich nur über sonderliche I, 1311.; die Welt beurtheilt die W. nach dem äußeren Schein II, 463. 468. f. Was der \gesperrt{Pabst} für gute W. erklärt hat I, 1671.; die \gesperrt{Papisten} wissen nicht, was gute W. sind I, 528. f.; II, 1570.; was die Papisten dafür halten I, 1367.; sie verachten die von Gott gebotenen W., thun der Weisheit Gottes Abbruch I, 529.; sie verachten die häuslichen und bürgerlichen W. I, 528. 1161. f.

4. \gesperrt{Wirkungen} guter W. Von Natur wollen wir durch W. etwas verdienen, so ist die Lehre, daß wir nichts verdienen, sehr nöthig I, 1444. 1449.; es ist thöricht, etwas damit verdienen wollen I, 1441. f. 1443. f.; gute W. taugen nichts in Gottes Gerichten I, 1347.; durch W. gerecht werden und W. thun, wenn man gerecht ist, ist zweierlei I, 1012.; W. machen nicht gerecht I, 620. 949. 1018. 1123. 1582., das erkennt schon die Vernunft I, 1018.; können nicht gerecht machen, denn sie sind besudelt I, 944.; nützen nichts zur Seligkeit I, 1577. Beruf und Erwählung sollen wir durch W. fest machen I, 1123. Gott belohnt die W. seiner Gläubigen I, 1581. ff.

Vgl. \gesperrt{Abraham} 15, \gesperrt{David, Ehestand} 8, \gesperrt{Esau, Glaube} 9, \gesperrt{Hausstand, Isaak} 8, \gesperrt{Jakob} 8, \gesperrt{Joseph} 9, \gesperrt{Mönche} 2, \gesperrt{Noah} 8, \gesperrt{Obrigkeit} 4, \gesperrt{Väter} I; II, 6.

IV. W. \gesperrt{des Teufels.} Was darunter zu verstehen II, 725.

Werkgerechtigkeit. S. \gesperrt{Werkheiligkeit}.

Werkheilige. Ihre Art und Verhalten I, 315.; sind unter dem Fluch I, 1568. f.; Kain war einer I, 315.; Esau ein Bild derselben II, 99.

Werkheiligkeit. Dabei vertraut man auf seine Werke I, 1575. ff.; ist Abgötterei I, 1576. 1583.; macht stolz und vermessen II, 1882. f.; ist durch den Segen, den Abrahams Same bringen sollte, ausgeschlossen I, 1565. f.; Gott macht sie zu Schanden II, 1883. — S. \gesperrt{Heiligkeit, Mönche, Selbstgerechtigkeit}.

Wetzel, Georg. Heuchler, abtrünniger Gotteslästerer, was er Luther vorwarf I, 1252. 1630. f.; verkehrte Luthers Lehre II, 1151.

Widersprüche, scheinbare, in der Schrift. Man soll nicht schnell die Schrift eines Irrthums zeihen I, 721. 725. \gesperrt{Beispiele} solcher W. und ihre Auflösung: von den 60 Jahren, die sich bei \gesperrt{Abraham} verlieren I, 721.; wie stimmt es, daß wir argwöhnisch sein und doch nicht richten sollen? I, 795. ff.; wie kann Arphachsad 2 Jahre nach der Sintfluth geboren sein? I, 713.; wie konnte in Gosen während der sieben unfruchtbaren Jahre Weide sein? II, 1793. f.; wie kann den weisen Gott etwas reuen? I, 485. ff.; wie kann die Schrift sagen: Gott sahe, da er doch alles vorausgesehen? I, 485. ff.; wie konnte Jakob innerhalb 7 Jahren 12 Kinder zeugen? II, 547. f.; wie stimmt das Register derer, die mit Jakob nach Egypten zogen, mit sich selbst? II, 1762. ff.; Sarah heißt Harans und auch Tharahs Tochter I, 1351. f.; wie reimt sich die Verwüstung der Erde durch die Sintfluth mit der Verheißung? I, 550. ff. \gesperrt{Schriftstellen:} 1 Mos. 2, 2. u. Joh. 5, 17. I, 91.; 1 Mos. 2, 24. u. Matth. 19, 4. f. I, 165. 1248.; II, 59.; 1 Mos. 14. u. 17. I, 365. ff.; 1 Mos. 6, 5. 6. u. 8, 21. I, 576. f.; 1 Mos. 8, 21. u. die Verderbung ganzer Städte und Länder I, 574.; 1 Mos. 11, 31. f. u. Apostg. 7, 2. 3. I, 724. f.; 1 Mos. 12, 4. u. Apostg. 7, 2. I, 769. f.; 1 Mos. 15, 13., 2 Mos. 12, 40. u. Gal. 3, 17. I, 959. f.; 1 Mos. 25, 1. 2. u. Röm. 4, 19. I, 1743. ff.; 1 Mos. 25, 23. u. 33, 1–3. II, 323. f.; 1 Mos. 48, 22., 34, 25. ff. u. Apostg. 7, 5. II, 1911. ff. — S. \gesperrt{Geheimnisse}.

Widerwärtigkeiten der Christen. Sind keine Zeichen des göttlichen Zorns I, 1239. f.; kommen alle vom Teufel I, 1238. f.; wir sollen darin nicht verzagen, uns an die Verheißung halten II, 210. 220. f.; wie sich ein Christ darin trösten soll I, 1239. f. 1343. f., wenn er Gut und Leben verlieren soll II, 153.; zu Geduld darin sollen uns Gottes Wohlthaten reizen I, 1338. ff.; unsere W. sind gering gegen die der Patriarchen I, 237. — S. \gesperrt{Kreuz}.

Wiederbringung aller Dinge. I, 122.

Wiedergeborene. In ihnen ist die Erbsünde noch mächtig II, 148.

Wiedergeburt. Nöthig zur Seligkeit II, 33. ff. — S. \gesperrt{Geburt} 2.

Wiederholung. In der Schrift ist keine vergebliche W. I, 541. ff. 547. f. 602. ff. 1132.; II, 1356. 1947. f.

Wiedertäufer. Verachteten das Wort, gaben einen großen Schein der Heiligkeit vor II, 1827.; warum sie die dunkeln Bücher der Schrift so gern haben I, 610.; allegorisiren gern I, 610.; halten den nicht für einen Christen, der etwas Eigenes hat I, 829.; lehrten: ein Christ dürfe die Hülfe der Obrigkeit nicht anrufen I, 1468., das Wasser in der Taufe könne die Seele nicht anrühren, das hielten sie für Weisheit I, 197.; wollen also leben, daß sie nichts überall sündigen I, 485.; irrten in der Deutung eines Zeichens am Himmel II, 423.; warum sie beim Volk Beifall fanden I, 726. — S. \gesperrt{Taufe}.

Wille. I. \gesperrt{Gottes Wille.} Wird unterschieden in den W. des Zeichens und Wohlgefallens I, 485.; von den verschiedenen Bedeutungen dieser Ausdrücke I, 485. f. 488. ff.; wir sollen nicht in allen Dingen so genau nach Gottes Rath und W. forschen I, 371.; nach den Ursachen des W.’s Gottes soll man nicht forschen I, 1082.; Gott will, daß alle Menschen selig werden I, 1434.; seinen gnädigen W. beweist und offenbart er in Wort und Sacrament II, 183.; des gnädigen W.’s Gottes sollen wir gewiß sein und nicht, wie die Papisten lehren, daran zweifeln I, 1550. f.

II. \gesperrt{Der Menschen} W. 1. \gesperrt{Im Allgemeinen.} Ist sehr unbeständig II, 1772. f.; kann nur durch Güte geändert werden II, 817.; ist der W. geändert, so folgt alles Andere nach II, 830.; ist die Königin, die über die Güter herrscht II, 830. — S. \gesperrt{Werke} II.

2. \gesperrt{Der freie} W. Um was es sich bei der Lehre vom freien W. handelt I, 483. f.; wir haben einen freien W. über die Dinge, die unter uns sind I, 103., aber nicht in den Dingen, die Gott angehen und über uns sind I, 103.; unser W. ist dem W. Gottes zum höchsten entgegen I, 175. 485.; ohne den Heiligen Geist ist der Mensch einfach gottlos und in Finsterniß I, 484., kann er sich Satans nicht erwehren I, 334., kann er nichts thun denn sündigen I, 481. ff. Die \gesperrt{irrige} Lehre hat ihren Ursprung in der Lehre vom Bilde Gottes I, 74.; die falsche Lehre vom freien W. folgt nicht aus dem Befehle an Kain, er solle der Sünde den W. nicht lassen I, 329. Der \gesperrt{Kirchenväter} Aussprüche über diese Lehre muß man mit Ueberlegung lesen I, 74. 480.; die \gesperrt{Juden} u. \gesperrt{Sophisten} lehren, die Vernunft treibe zum Besten I, 484.; die \gesperrt{Mönche} und heidnische \gesperrt{Philosophen} hielten den freien W. für gut, der nur durch irdische Güter verführt werde I, 830.; die \gesperrt{Papisten} lehren falsch vom freien W. I, 482.; wie wichtig die reine Lehre hiervon sei I, 482. f. — S. \gesperrt{Freiheit, Kräfte} 2, \gesperrt{Mensch} 2, \gesperrt{Vernunft}.

Wind. War nicht gleich im Anfang I, 11.

Wittenberg. Um W. müssen vormals die allerärgsten Leute gewohnt haben I, 595.; gab unter dem Pabstthum den Mönchen jährlich 1000 Gulden II, 109.; jährliche Ausgabe für Gerste II, 1869.

Wittwenstand. Wer ihn mehr liebt, soll doch den Ehestand nicht verachten I, 1164. f.

Wöchnerinnen. Am 3. Tage am schwächsten II, 885.

Wohlthaten. 1. \gesperrt{Im Allgemeinen.} Man vergißt W. leichter denn die Strafe I, 604.

2. W. \gesperrt{Gottes.} Sind viel größer als alles Unglück I, 1339. f. 1343. f.; sie sollen uns zur Erkenntniß der Güte Gottes führen I, 1098., zu Geduld in Widerwärtigkeiten reizen I, 1338. ff.; wir sollen sie erkennen und Gott loben I, 1001., uns dieselben besonders im Unglück vor Augen stellen I, 1339. f.; unser Fleisch verdunkelt sie durch den Unglauben I, 1339.; die Welt vergißt sie I, 1337. f.; die Gottlosen vergessen sie im Unglück, Christen denken daran I, 1343. f.; auch Christen vergessen sie wohl bei geringen Schäden I, 1338.; woher es kommt, daß die geistlichen W. Gottes so verachtet werden II, 148. f. — S. \gesperrt{Gaben}.

3. W. \gesperrt{der Menschen.} Man soll sie mit aufrichtigem Herzen erweisen I, 1216.; die Welt meint, man müsse W. um W. kaufen I, 979.; man thut die W. meist um Dankes willen I, 977.; von W. soll uns nicht abhalten, daß es so viele Unwürdige gibt I, 1140.; man soll W. erweisen vornehmlich, doch nicht ausschließlich, den Glaubensgenossen I, 1140. ff., den Dankbaren und den Undankbaren I, 978. f., den Feinden I, 977. f.; W. sollen behende und willig sein I, 1263. 1710.; man frage nicht
\quellseite{579}
erst, ob sie angenommen werden I, 1216.; seine eigenen W. soll man nicht rühmen I, 1143. 1148. 1156.; wir dürfen nur Undank erwarten I, 977. f.; bei Undank fahre man, wie Gott, mit W. fort I, 979.; W., einem Undankbaren erwiesen, sind nicht verloren I, 978.; bescheidene Leute muß man schier zwingen, W. anzunehmen I, 1216.

Wolken. Werden durch das Wort an ihrem Ort gehalten I, 545.

Wort. 1. \gesperrt{Das persönliche} W. Ist ewig, eine vom Vater unterschiedene Person I, 20. ff.; wie das ungeschaffene Wort sich vom geschaffenen unterscheidet I, 26. f.; es hat die unförmliche, aus Nichts geschaffene Materie unterschieden und geschmückt I, 11.

2. \gesperrt{Worte Gottes.} Gott hat andere W. und eine andere Grammatik als wir Menschen, seine Worte sind alsbald wesentliche Dinge I, 26.; ein Wort Gottes macht möglich, was unglaublich ist I, 60.; Beispiele I, 116. f. — S. \gesperrt{Creaturen}.

3. \gesperrt{Das Wort Gottes. Preis} des W.’s: es ist der rechte Schmuck einer Kirche I, 839., der allerhöchste Schatz I, 1376., mehr denn der Engel Werke II, 211.; wer das W. hat, hat Alles I, 555.; wo das W. ist, da ist die Kirche, der Geist, Christus, Alles I, 709.; ohne dasselbe kommt der Heil. Geist nicht II, 263., wo es im Schwange geht, da ist Gott I, 780.; ist eine große Wohlthat I, 862. f. 1338.; hat man das W. rein, so kann man auch alles Andere leicht beurtheilen II, 1057.; es ist das einige Licht in der Finsterniß II, 805.; ohne das Wort ist Alles Finsterniß II, 646. f. 1055. 1761., das ganze Leben ungewiß und unstät I, 367.; wer sich darauf verläßt, wird nicht zu Schanden I, 870.; es ist das Vornehmste, dann kommt unser Thun I, 1117.; den Werth des W.’s erkennt die Welt nicht II, 211. \gesperrt{Wirkungen} des W.’s: es macht alle Dinge herrlich II, 567. ff.; ohne das W. sind alle äußerlichen Dinge, auch die Sacramente, nichts nütze zur Seligkeit I, 1248. f.; es heiligt Alles I, 779.; II, 125. 468., macht jeden Ort zu einem heiligen Tempel I, 839.; heiligt uns II, 705.; ohne das W. gefällt Gott nichts II, 907.; wo es geoffenbart wird, breitet es seinen Glanz aus II, 987.; geht nicht ohne Frucht und Wirkung ab I, 1000. 1023. 1332. 1439.; II, 773.; ist wie ein Hammer und zweischneidig Schwert I, 1346.; II, 1687.; Etliche werden bekehrt, der andere Haufe verachtet es I, 453. f.; den Gottlosen gereicht es zum Aergerniß I, 115.; es bringt oft Früchte, wie bei Simeon und Levi II, 1942. f.; es lehrt uns erst recht Gottes Werke erkennen I, 1373. f.; es gibt den Heil. Geist II, 2009.; es wirkt den Glauben II, 289. f., erweckt und stärkt ihn II, 1325. f.; errettet von Tod und Teufel II, 1974.; macht lebendig und tröstet I, 1451.; II, 904.; wäscht uns II, 1993.; wirkt ein christliches Leben I, 1544. 1546.; II, 134.; macht zu guten Werken geschickt II, 125.; wirkt die rechten Werke I, 1544.; auf das W. folgt das Gebet II, 137.; wirkt in Anfechtung Geduld II, 126. f.; wer die Wohlthat des W.’s erkennt, wird selig I, 862. f.; sammelt immer einen Haufen und bringt auch die herzu, die da fremd sind, nämlich, die dazu versehen sind II, 909.; seine göttliche Kraft erfährt man, wenn man sich daran hält II, 54. \gesperrt{Verkündigung} des W.’s: man muß es recht theilen I, 1187. 1205. 1208. f., dazu gehört große Weisheit I, 1209. ff., die das ganze Pabstthum nicht hat I, 1211.; zu Zeiten ist das W. klarer, zu Zeiten dunkler offenbart II, 1334.; es wird im Frieden gelehrt, damit man es in Unruhe und Anfechtung brauchen kann II, 223.; es kommt nicht zu denen, die satt und überdrüssig sind II, 137.; wo es im Schwange geht, folgt bald leibliche Theurung etc. II, 131. 1805.; es ist nicht Schuld am leiblichen Unglück II, 132. ff. \gesperrt{Verhalten} gegen das W.: Gott weist uns an sein W. und will nicht Jedem sonderliche Offenbarungen geben II, 560. f.; man soll es brauchen, solange man es hat II, 318. ff., hören, gläubig annehmen I, 1182. 1522., demselben anhangen, die Augen gegen alles Andere verschließen II, 1291., nicht weiter forschen II, 218., daran hängen, es gehe uns, wie es wolle II, 357.; die Vernunft lasse man sich dabei nicht irre machen II, 285. f.; wer sich nicht daran hält, ist bald verführt II, 561.; demselben folgen ist eine große Tugend I, 869. f.; auf das W. soll man sich verlassen, nicht auf eigene Größe, Macht etc. I, 553.; auf das W., nicht auf Menschen soll man sich in der Lehre verlassen II, 79.; Gottes W. hoch halten ist keine Hoffart I, 866.; das W. soll man allen Dingen vorziehen I, 1421., über alle Werke der Menschen setzen I, 1544.; man fliehe, wo das W. nicht ist II, 2010. Wie sich die \gesperrt{Menschen} bei der Verkündigung des W.’s verhalten I, 453. f.; viele verachten es, etliche nehmen es an I, 1333. 1341. f.; die \gesperrt{Welt} wird des W.’s bald überdrüssig I, 727. 1304.; es ist zu allen Zeiten verachtet worden I, 1307. f.; es scheint den Menschen Thorheit II, 1064.; ward bald nach der Sintfluth verachtet I, 899.; die Gottlosen meinen, es sei an allem Unglück Schuld II, 132. ff.; je reichlicher es gepredigt wird, desto größer der Undank I, 864. f.; die Welt glaubt nicht, daß es große Dinge ausrichtet II, 2008. f.; der \gesperrt{Pabst} mißbraucht es zu Gotteslästerung I, 905.; \gesperrt{Satan} sucht uns von Gottes Wort abzuführen II, 416. f.; mißbraucht es zu Gotteslästerung I, 905.; wüthet dagegen I, 839. f. 1135. f. Wie sich die Welt vor dem jüngsten Tage dagegen verhalten wird I, 409. — S. \gesperrt{Drohungen, Evangelium, Gebet} 15, \gesperrt{Gesetz} 1, 2, \gesperrt{Gehorsam, Glaube} 4, \gesperrt{Kirche} 9, \gesperrt{Prophezeiungen, Predigtamt, Schöpfungswort, Schrift, Verheißungen}.

4. \gesperrt{Wort und Sacrament.} Wo sie fehlen, folgen greuliche Irrthümer I, 305. f. — S. \gesperrt{Gnadenmittel}.

Wörter. Ihre Bedeutung kann man oft nicht aus der Ableitung, sondern nur aus dem Sprachgebrauch lernen II, 1738. f., aus dem Zusammenhang II, 1910. 1925. f.; es gibt W., die durch Zufall mehrdeutig sind II, 1908. f., andere werden dazu gemacht II, 1909.

Wucher. Ist eine Sodomitische Sünde I, 1218. f., wider die Natur, bringt keinen Segen I, 751. f.

Wucherer. Behalten nicht, was sie zusammenscharren II, 205.

Wunder. Die Engel, aber nicht die Menschen, können solche aus eigenen Kräften thun I, 1244. 1247.; sind nicht größer als andere Werke Gottes, die nach der Ordnung der Natur geschehen I, 1370. f.; W., welche täglich geschehen, achten wir gering I, 153. ff. 546., das ist Undankbarkeit I, 154. f.; die täglichen W. sind größer, als die Christus gethan II, 1245. f.; auf W. sollen wir nicht achten, wenn Einer ohne das Wort kommt I, 1545. f.; wie man den päbstischen W. begegnen soll I, 1545. ff. — S. \gesperrt{Werke} I.

\section{Z}

Zahlen. Die Z. zeigen an, daß der Mensch eine besondere Creatur sei I, 53.; Bedeutung der Zahl 2 I, 39., der Zahl 7 I, 1088., der Zahl 8 I, 1087.

Zähne. Bedeuten in der Schrift oft Lehrer II, 2006.; \gesperrt{weiße} Z. bedeuten oft Narrheit, Milchkinder II, 1996.

Zank. Fällt auch zwischen Freunden und Eheleuten vor I, 1392.; Jakob und Rahel zankten II, 537. f.

Zauberer. Erwecken plötzliche Ungewitter II, 8.

Zebaoth. Warum Gott so heißt I, 90.

Zehn Gebote. S. \gesperrt{Gebote}.

Zehnten. Von wem Melchisedek den Z. empfangen I, 912. f.; wie die Epistel an die Hebräer von dem Z. handelt, den Abraham dem Melchisedek gab I, 912. f.

Zeichen. 1. \gesperrt{Im Allgemeinen.} Inwiefern Sonne, Mond und Sterne Z. sind I, 51. ff.; böse Z. deuten die natürlichen Menschen zu ihren Gunsten II, 75. f.; die Gottlosen irren bei Auslegung der Z. II, 423.

2. Z. \gesperrt{der Gnade.} Sind eine unaussprechliche Gabe Gottes I, 306.; sind höher zu achten als alle Wunderwerke I, 1093.; ihre Art ist, daß sie trösten, nicht schrecken I, 306. f. 604.; sie sind nöthig I, 1046. 1093.; Gott hat immer neben dem Wort auch sichtbare Z. der Gnade gegeben I, 303. f. 401. f. 603. f. 953.; von den Z. im A.\,T. I, 303. f. 1050. 1093.; wir haben im N.\,T. herrlichere Z. I, 1066., den Sohn Gottes und die Gnadenmittel I, 305. 1045. ff. 1093.; wir sollen keine Z. d. G. nach eigenem Gutdünken fordern I, 953. f.; wir sollen darüber halten, ihnen folgen I, 305. f. 1047. ff. 1093. f.; sie werden verachtet I, 1045. 1094. von der Vernunft I, 1047., von den Papisten I, 306.; Gott straft diese Verachtung I, 305. f.; sie stürzt ins Verderben I, 1046.; der Pabst will uns von den Z. d. G. abführen I, 1047.; der Satan sucht uns die rechten Z. aus den Augen zu rücken und falsche aufzurichten I, 1046. ff. 1094. — S. \gesperrt{Gnadenmittel}.

Zeit. 1. \gesperrt{Im Allgemeinen.} Die Eintheilung der Z. in Tage und Jahre ist Gottes gnädige Ordnung I, 53.; warum Kinder und unvernünftige Thiere von der Zeit nichts wissen I, 53. „Zeit“ bedeutet bei den Hebräern „Fest“ I, 58.

2. \gesperrt{Die} Z. \gesperrt{des} N.\,T.’s. Ist durch Noahs Tauben als eine Zeit der Gnade und Vergebung abgebildet I, 625.; hat einen großen Vorzug vor der Z. Abrahams I, 863. 869. 1118. f. 1122. f. 1385. f.; von der Ehre und Herrlichkeit des Segens im A.\,T., verglichen mit dem im N.\,T. II, 1887. ff.; die Herrlichkeit des N.\,T.’s kann man nicht genug rühmen II, 2057. 2062. f.; die Väter, welche im A.\,T. göttliche Erscheinungen hatten, sind nicht seliger denn wir im N.\,T. II, 145. ff.; Engel erscheinen jetzt nicht so häufig, weil der Sohn Gottes erschienen ist I, 1525.; wir haben viel herrlichere Erscheinungen denn Abraham I, 1527.; die Z. des N.\,T.’s ist reicher an Offenbarung denn die Z. der Patriarchen I, 432., hat herrlichere Weissagungen II, 2056. ff.; wir haben eine klarere Erkenntniß von der Auferstehung als die Väter im A.\,T. I, 1755. f.; wir haben viel reichlichere, herrlichere Verheißungen denn Abraham I, 1065. ff.; II, 443. und die Gläubigen im A.\,T. II, 2060. f.; im A.\,T. waren die leiblichen Verheißungen größer als im N.\,T. II, 1897.; die Verheißungen des A.\,T.’s enthielten die des N.\,T.’s II, 1896. f.; wir haben auch leibliche Verheißungen, aber die geistliche ist viel herrlicher II, 1888. f. 1892. f. 1896. f.; im N.\,T. muß man auch Gesetz treiben, aber sein eigentliches Amt ist das Evangelium II, 1890.; ein Weib kann taufen, ein Weib und Kind kann absolviren II, 2057. ff.; wir haben größeren Segen, denn sie im A.\,T. hatten II, 300. f. \gesperrt{Die} Z. \gesperrt{Zeit,} die Z. des N.\,T.’s, ist voll Gnade, aber auch voll Gotteslästerung, muß gestraft werden mit der Sintfluth des Feuers I, 437.; ist die letzte und ärgste Zeit der Welt, Gott wird ihrer nicht schonen I, 432.; ist die Grundsuppe der Welt, verachtet Christum I, 442., Gott muß sie strafen I, 444., ist sehr undankbar I, 441. f.; je näher der jüngste Tag, desto größer die Verachtung des Worts I, \unsicher{409.} f.; Strafe wird folgen I, 449. — S. \gesperrt{Segen, Verheißungen} 1.

3. \gesperrt{Die letzte Zeit.} Wird gar erschrecklich sein, wie die Z. Noahs I, 409. f.; die Sünden werden so greulich sein, daß die Strafe durch keine Fürbitte abgewendet werden kann I, \unsicher{1409.}; das Licht des Evangeliums wird fast ganz ausgelöscht sein I, \unsicher{538.} f.; die Farben des Regenbogens auf die Strafen der letzten Zeit angewandt I, 512. — S. \gesperrt{Tag} 6; oben 2.

Zeitalter. Das \gesperrt{goldene} Z. war vor der Sintfluth I, 436.; woher die Poeten etwas davon wissen I, 436.

Zeitliches. Zeitliche Dinge, die zum Haus- und Weltregiment gehören, soll man nicht verachten I, 169. f.; auch von zeitlichen Dingen soll man in der Kirche lehren II, 170.; unser Fleisch wüthet aus Liebe zu dem Z. I, 424.

Zeitrechnung. Studium der Z. ist nöthig und nützlich II, 845. 1158. f.; sie gibt Anlaß zu schönen Gedanken II, 983.; gefällt Gott, daß man sie kennt II, 1803. f.

Zertrennung. Eine schädliche Plage I, 704.

Zeugung der Menschen. Von dem \gesperrt{Werk} der Z. I, 380. ff. Ist ein wunderbares Werk I, \unsicher{108.} f.; II, 50. 1414. f., das edelste in der Natur I, \unsicher{238.}; eine Creatur Gottes I, 291.; II, 49. f.; eine Gabe Gottes trotz der unkeuschen Brunst I, 86.; eine Gabe, welche die Heiden nicht erkennen I, 589.; ein Segen, der auch nach dem Fall geblieben I, 412.; ein Segen trotz unserer Erbsünde I, 744., trotz der damit verbundenen Schmerzen I, 1078. f.; man hielt sie im A.\,T. für einen Segen I, 728. f. 1079. \gesperrt{Geschieht} durch Vermischung des Samens von Mann und Weib I, 1169. An diesem Werke läßt sich die Tiefe des \gesperrt{Falls} erkennen I, 143. f.; es ist durch denselben verderbt I, 228.; durch die Sünde ist des Teufels Gift daran gekommen I, 201.; geschieht jetzt mit unkeuscher Brunst I, 86.; II, 660. f.; durch die Sünde ist sie mit Schmerz und Scham verbunden I, 128.; durch die Sünde ist sie ein Werk, deß man sich schämt I, 205. 248., davon man kaum ohne Scham reden kann I, 144. f.; ist durch den Fall für die Weiber mit Beschwerden verbunden I, 245. \gesperrt{Wie} man sich bei den Beschwerden \gesperrt{trösten} soll I, 246. f. 1078. f. Ist \gesperrt{geheiligt} durch die Verheißung im Paradies I, 233.; Gott hat Gefallen daran I, 1748.; Mose verdammt an der ersten Welt nicht die Z., sondern den Mißbrauch I, 444.; der Heil. Geist scheut sich nicht, von diesem Werk zu reden I, 291. f. \gesperrt{Man soll} die Z. hochhalten und begehren II, 1176. — S. \gesperrt{Ehestand, Empfängniß, Fortpflanzung, Geburt} 1, \gesperrt{Stand der Unschuld}.

Zibeon. Meinung der Juden über ihn II, 995. f.

Ziegel. I, 694. f.

Ziegler, Bernh. Erklärt das Wort „Schiloh“ II, 1970.

Zilla. Bedeutung des Namens I, 386.

Zoar. Hieß zuvor Bela, warum der Name geändert worden I, 1273.; Lage I, 1274.; zu Abrahams Zeit bekriegt I, 878. f.; durch Lot gerettet I, 1271.; wie Sodom zerstört I, 1290. f.

Zorn. 1. \gesperrt{Zorn Gottes.} Die Vernunft kann denselben nicht glauben noch verstehen I, 509. f.; wie derselbe mit Gottes Verheißungen zu vereinigen ist I, 509. ff.; es ist ein Gnadenzorn, wenn die Strafe eilt I, 702.; er ist desto unerträglicher, je länger er geschwiegen hat I, 702.; den Z. Gottes, wie er sich in der Sintfluth offenbart hat, kann niemand ausdenken I, 548.; er ist dann
\quellseite{580}
sehr schrecklich, wenn Gott das Wort wegnimmt I, 455. 460.; dawider schützt nicht der Name der Kirche und der große Haufe I, 496.; es gibt viele Exempel des Z.’s I, 437. ff.; wir verachten sie I, 439.; man soll sie einschärfen I, 1193.; wozu diese Exempel dienen sollen I, 438. ff. 465. 1185. 1207. 1277.; Betrachtung des Z.’s Gottes macht kleinmüthig I, 497.; auch in den Exempeln des Z.’s leuchtet Gottes Gütigkeit I, 1210. — S. \gesperrt{Strafe} 2.

2. Z. \gesperrt{der Menschen.} Ist ein natürlicher Affect und an sich nicht Sünde II, 1501.; er steckt in uns, wenn wir auch nicht von Anderen gereizt werden II, 1501. f.; er macht kühn zu allerlei bösen Thaten II, 339. f.; es gibt einen heiligen Zorn wider die Sünde II, 691.

Zornigen gehe man aus dem Wege II, 340.

Zucht und Ehrbarkeit liegt darnieder, wo man das Wort verachtet, gegen Gott undankbar ist I, 1192.

Züchtigung, göttliche. Die rechte \gesperrt{Lehre} von der Z. kennt die Vernunft nicht, findet sich nur im Wort II, 1664. f., auch die Philosophen und Juristen kennen sie nicht II, 1666.; dieselbe soll man oft treiben II, 1664. \gesperrt{Art und Natur:} wenn Gott uns züchtigt, kämpft er mit uns II, 809., stellt er sich, als sei er nicht unser Vater, offenbart sich aber doch endlich als solchen II, 1657. 1661. 1666. f.; in der Z. verstehen wir seinen Rath und Willen nicht, sondern erst hinterher II, 1661. f. 1667.; in der Z. kämpfen wir mit Unglauben und Verzweiflung II, 1662.; sie thut weh II, 1663.; ein schönes Bild, wie Gott züchtigt, ist die Behandlung, die Joseph seinen Brüdern widerfahren ließ II, 1657. Z. ist \gesperrt{nöthig und heilsam} für die Gottseligen II, 807. f. 809. f. 1660. f. 1664. 1671. f.; sie soll \gesperrt{dienen} zu ihrer Läuterung II, 1658. 1660. und Demüthigung II, 1658. f. \gesperrt{Verhalten:} die Z. zu ertragen, dazu gehört Glaube II, 1662.; die Kunst, in Z. Muth zu haben und zu triumphiren, können nur die Christen II, 1664., und auch sie nicht vollkommen II, 1665.; in der Z. klagt und murrt das Fleisch II, 1662. f.; wider dies Murren soll man kämpfen II, 1663.; wir sollen uns nicht weigern, die Z. zu erdulden II, 809., sollen sie mit Geduld tragen I, 999. 1001.; II, 1044.; man halte sich ans Wort und Taufe II, 1663. ff. — S. \gesperrt{Kreuz, Strafe}.

Zufriedenheit. Wie man Z. erlangt I, 1654.

Zuhörer. Sind oft nicht andächtig II, 940. f.; der Teufel führt ihre Herzen vom Wort ab II, 1364.; Gott gibt ein hörend Ohr II, 1364.

Zustand der Seele nach dem Tode. Z. der gläubigen Seele I, 1758. ff.; II, 216. ff. 1542. ff., der gottlosen I, 1761. 1763.; der Z. der gläubigen Seele ist verschieden von dem Z., darin Christus lebt I, 1763.; von den fünferlei Behältnissen der Papisten I, 1761.; was von den verschiedenen Behältnissen zu halten, davon die Väter disputiren I, 1760. ff.; nach diesem Z. der Seele forsche man nicht allzu genau II, 2070. — S. \gesperrt{Fegfeuer, Gehenna, Hölle, Limbus, Scheol}.

Zütphen, Heinr. v. Ein Märtyrer I, 352.

Zweifel an Gottes Gnade. Eine schwere Sünde I, 1550. ff.; der Pabst lehrt Z. an Gottes Gnade II, 1900. 2059. ff.; eine schändliche Lehre II, 2061. 2064.; wer so lehrt, ist unter die Teufel zu rechnen II, 2063.

Zwingli. Vom Ehrgeiz getrieben I, 924.; seine Lehre von der Heiden Seligkeit II, 1828., wobei er sich fälschlich auf Luther berief II, 1829. 1831.

Zwinglianer. Werden immer tiefer fallen II, 1830.
